% !TEX root = main.tex
\documentclass[11pt, a4paper, openany]{book}

\usepackage[
  a4paper,
  left=3.5cm,
  right=2.5cm,
  top=3.5cm,
  bottom=3.5cm
]{geometry}

\usepackage{setspace}
\setstretch{1.1}

\raggedbottom

\usepackage[utf8]{inputenc}
\usepackage{graphicx}
\usepackage{tikz}                
\usetikzlibrary{quantikz2}
\tikzset{operator/.append style={rounded corners}}
\tikzset{
    dashed line/.style={path picture={
        \draw[densely dashed]
            (path picture bounding box.west) -- (path picture bounding box.east);
    }}
}
\tikzset{
  halfctrl/.style={
    circle, draw=black, line width=0.4pt,
    fill=white, minimum size=6pt, inner sep=0pt,
    path picture={
      \fill[black]
        (path picture bounding box.south west) --
        (path picture bounding box.south east) --
        (path picture bounding box.north east) -- cycle;
      \draw[black, line width=0.3pt]
        (path picture bounding box.south west) --
        (path picture bounding box.north east);
    }
  }
}
\newcommand{\dashedthrough}{|[dashed line,thickness, inner ysep=3pt]| }
\usepackage{enumerate}
\usepackage{booktabs}
\usepackage{amsmath}
\usepackage{amssymb}
\usepackage{mathtools}
\usepackage{amsthm}
\usepackage{algorithm}
\usepackage{algpseudocode}

\renewcommand{\topfraction}{0.9}
\renewcommand{\bottomfraction}{0.9}               
\renewcommand{\textfraction}{0.1}
\renewcommand{\floatpagefraction}{0.8}

% \DeclarePairedDelimiter\bra{\langle}{\rvert}
% \DeclarePairedDelimiter\ket{\lvert}{\rangle}
% \DeclarePairedDelimiterX\braket[2]{\langle}{\rangle}{#1\,\delimsize\vert\,\mathopen{}#2}

\newtheorem{theorem}{Theorem}
\newtheorem{lemma}[theorem]{Lemma}
\newtheorem{proposition}[theorem]{Proposition}
\newtheorem{definition}[theorem]{Definition}
\newtheorem{corollary}[theorem]{Corollary}
\theoremstyle{definition}
\newtheorem{remark}[theorem]{Remark}

\usepackage{todonotes}

\usepackage{titlesec}

\setcounter{secnumdepth}{1} 
\setcounter{tocdepth}{2}    

% 3. Redefine \subsection to look like a "run-in" paragraph
\titleformat{\subsection}[runin]{\normalfont\bfseries}{}{0pt}{}\titlespacing*{\subsection}{0pt}{12pt}{1em}

\usepackage{dsfont}
\newcommand{\idm}{\mathds{1}}

\DeclareMathOperator{\tr}{Tr}
\DeclareMathOperator{\id}{id}
\DeclareMathOperator{\erfc}{erfc}
\newcommand{\dd}{\mathop{}\!\mathrm{d}}
\newcommand{\ii}{\mathrm{i}}
\newcommand{\ee}{\mathrm{e}}

\usepackage[
    backend=biber,      % The modern backend, better than bibtex
    style=alphabetic,   % This is the key! It creates labels like [Knu84]
    sorting=anyt,        % Sorts by author, year, title
    maxalphanames=4
]{biblatex}

\usepackage[toc, page]{appendix}
\usepackage{xcolor}
  \definecolor{pinegreen}{HTML}{2D5A3D}                           
  \definecolor{bordeaux}{HTML}{7A2E39}       
  \definecolor{dustyplum}{HTML}{8E6F8C}
  \definecolor{deepplum}{HTML}{735874}
  \definecolor{aubergine}{HTML}{4F3B5C}                      
  \definecolor{slateblue}{HTML}{5C7794}                          
  \definecolor{sage}{HTML}{8B9F7E}                                  
  \definecolor{ochre}{HTML}{B89143}
  \definecolor{terracotta}{HTML}{B5654A}                       
  \definecolor{dustyteal}{HTML}{5F8B8E}                         
  \definecolor{mustard}{HTML}{C9A86A}
  
  \usepackage[                                                                                      
      colorlinks=true,
      linkcolor=deepplum,            
      citecolor=bordeaux,
      urlcolor=deepplum,                  
  ]{hyperref}  
\usepackage{bookmark}

\addbibresource{references.bib}

% \part{Preliminaries}

\todo{this can be skipped pretty sure:}
As with any topic, we need to first understand the basic notions and tools that build the foundation. After the foundation is set, we can see what is possible to our current knowledge. Both the title and the introduction suggest that we would like to understand how a quantum system evolves to its thermal equilibrium state, the Gibbs state. For this, it is crucial to define the quantum dynamics that are at play here. But also to introduce the properties of the Gibbs state we would like to prepare. As we will see, these properties will set the difficulty level of the equilibration process. The underlying physics is however only one side of the problem, the other is the method with which we would like to prepare this target state. The tool we would \textit{eventually} like to use is a quantum computer, and accordingly we will need a quantum algorithm that could run on such a machinery. In classical computation Markov chain sampling was the quintessential technique to prepare and sample from the Gibbs distributions, and thus it comes as no surprise that we will work with quantum Markov chains in this thesis. Understanding how the quantum side works is definitely one aim of this chapter, but we also want to make the differences between the classical and quantum pictures as clear as possible. This should complement later analysis of the parameter regimes in which we can hope for a quantum advantage.

\chapter{Introduction}
\newpage
\chapter{Quantum Dynamics}

\label{sec:quantum-formalism}
\section{Quantum Formalism and Notation}
In order to have a self-contained thesis, we ought to start with the necessary concepts in quantum information. The aim here is not to give a complete introduction, but rather to provide all the definitions that are used later on in the text. In case the reader needs more details and pedagogic examples, we would advise them to open up already well-established textbooks and lecture notes like \cite{bible, wolf2012lecture, breuer2002theory}.

\smallskip
\label{sec:quantum-states}
\subsection{Quantum states} The state of a finite-dimensional quantum system is described within a complex vector space known as the Hilbert space, $\mathcal{H}$. While pure states are represented by vectors in $\mathcal{H}$, the most general description of a quantum state is given by a density operator, $\rho$.

These bounded operators belong to the space of all linear operators on $\mathcal{H}$, which we denote by $\mathcal{B}(\mathcal{H})$. For a physical state, an operator $\rho \in \mathcal{B}(\mathcal{H})$ must satisfy the following properties:
\begin{itemize}
    \item \textit{Hermiticity.} $\rho^\dagger = \rho$
    \item \textit{Positive semi-definiteness.} $\rho \geq 0$
    \item \textit{Normalization.} $\tr[\rho] = 1$
\end{itemize}
The set of all operators that fulfil these conditions forms the space of physical states, which we will denote by $\mathcal{D}(\mathcal{H})$. This space is a convex subset of $\mathcal{B}(\mathcal{H})$, i.e. any convex combination (or mixture) of density matrices is also a density matrix, $\sum_i \lambda_i \rho_i = \rho \in \mathcal{D}(\mathcal{H}).$

The states in $\mathcal{D}$ can be characterized by their purity, a measure of mixedness defined as $\tr [\rho^2]$. At one end there are the \textit{pure states}, for which the purity takes its maximal value of $1$. In this case the state is given by a rank-one projector onto a state vector $\ket{\psi} \in \mathcal{H}$
$$
\rho = \ket{\psi}\bra{\psi}
$$
States that are not pure, are called \textit{mixed states}. At this end we find the maximally mixed state, $\rho = 1 / \dim(\mathcal{H})$, which has the minimum possible purity of $\idm / \dim(\mathcal{H})$.
We can decompose any mixed state into a convex sum of orthogonal pure states as well:
\begin{equation}
    \rho = \sum_i \lambda_i \ket{\psi_i} \bra{\psi_i},
\end{equation}
where $\{\ket{\psi_i}\}$ is an orthonormal basis on $\mathcal{H}$ and are the eigenvectors of $\rho$. This is a powerful decomposition, as it tells us that mixed states can be represented by an ensemble of pure states with some probabilities $\lambda_i$. This will be quite useful later on for our numerical computations, as it is much easier to simulate the evolution of pure states than mixed states.

The task of preparing a specific state, such as the thermal Gibbs state, on a quantum computer involves navigating this state space $\mathcal{D}(\mathcal{H})$ through a sequence of physical operations. Once a state is prepared, we can then measure an observable quantity of interest, such as the system's energy.

\textbf{TODO: maybe correlation and entanglement and distances}

\smallskip
\subsection{Quantum measurements} In order to extract valuable information about the physical system, we still need to apply measurements to the system. In classical theory, predicting the outcome of a measurement should be possible, granted that we know the initial condition of the system and have the right theory. But due to the inherently probabilistic nature of quantum mechanics, such a prediction is just not possible. What is possible, is to predict probabilities and thus also the resulting measurement statistics that we expect for our statistical experiments. It is useful to first separate two notions, one is \textit{what} we can measure, and the other is \textit{how} we can measure it. A physical observable $A \in \mathcal{B}(\mathcal{H})$ that we can measure is described by a Hermitian matrix. The expectation value of the measurement can be then only understood with respect to a state $\rho$, and we can write it as $\langle A \rangle = \tr [\rho A]$. 
We can give this expression slightly more depth if we label the possible outcomes of the measurement by $\{a_\alpha\}$ and the probability with which we measure them by $\{p_\alpha\}$. Consequently, the expectation value is $\langle A \rangle = \sum_\alpha p_\alpha a_\alpha$. 

Now it is just a question of what do we really understand by the probabilities $p_\alpha$, which is also the same question of \textit{how} do we do measurements. In principle, we can always associate an outcome to a positive-valued operator $0 < E_\alpha \leq \idm$ such that,
$$p_\alpha = \tr [\rho E_\alpha].$$ 
If these operators (called \textit{effects}) have the normalization property $\sum_\alpha E_\alpha = \idm$, then they define a \textit{positive operator-valued measure} or POVM. This is still quite a general way of understanding quantum measurements. The way we connect back to the physical observable $A$ and its measurement statistics, is to find the spectral decomposition, $A = \sum_\alpha a_\alpha P_\alpha$, in which case $\{P_\alpha\}$ forms a special kind of POVM, a \textit{projection valued measure} (PVM). The naming is due to the fact that now each $E_\alpha = P_\alpha$ satisfy $P_\alpha^\dagger P_\alpha =  P_\alpha$, and thus an orthogonal projection to the eigenspace of $A$ with eigenvalue $a_\alpha$. While the POVM formalism is perfectly suited for describing measurement statistics, it does not care about the quantum state that we would find post measurement. However, in the PVM case, we definitely know in which state the system is after the measurement, as the operator projects to one of the eigenspaces of $A$:
\begin{equation}
    \label{eq:post_meas_state}
    \rho_\alpha = \frac{P_\alpha\rho P_\alpha}{\sqrt{\tr{\left[P_\alpha \rho\right]}}} = \frac{P_\alpha\rho P_\alpha}{\sqrt{p_\alpha}}.
\end{equation}
Due to the projection operator's property, if we were to repeat this measurement we would get the same $a_\alpha$ result and the same state $\rho_\alpha$. Famously, this is called as the collapse of the wave function, as the projection singles out a component of the superposition. Clearly not all measurements are repeatable in quantum mechanics, so the projective measurement must really be a special case. For instance when we measure a photon, we destroy it in the process and there is no way to repeat that measurement again. The post measurement state in the most general case can still be written up, and its form is similar to \eqref{eq:post_meas_state}, but without the projective property of the operator,
\begin{equation}
    \label{eq:post_meas_state_general}
    \rho_\alpha' = \frac{M_\alpha \rho M_\alpha^\dagger}{\sqrt{\tr\left[M_\alpha \rho M_\alpha^\dagger\right]}} = \frac{M_\alpha \rho M_\alpha^\dagger}{\sqrt{p_\alpha}}\quad\text{with}\quad M_\alpha M_\alpha^\dagger = E_\alpha.
\end{equation}
The difference here is that $\rho_\alpha'$ is in general not an eigenstate of $A$ and if we were to measure the same operator $M_\alpha$ we would get a different post measurement state.

\todo{THIS NEEDS REWRITING FOR THE NEW STRUCTURE}
\section{Closed systems}
With the introduction of quantum measurements in the previous chapter, we have already stumbled upon the notion of how a quantum system $\rho$ can change from one moment in time to another. This process can be described by an \textit{open quantum system evolution}, where we, the observers or (if we take away the human part) the environment, act on a quantum system. From the point of view of the observed system, this is an irreversible process, i.e. there is no physical map that would reverse the dynamics for all possible initial states. Irreversibility arises because we only observe a subsystem and lose the information of the rest. But the dynamics of the whole \textit{closed} system we can describe by a reversible process. 

Before we describe closed and open system dynamics in more detail, this is also a good place to introduce two frameworks in which one can understand evolutions. Since in quantum mechanics the only physically measurable quantities are the expectation values of some observables, $\langle A \rangle (t)$, and a physical process as we have seen is divided into a state preparation and a measurement, we have the freedom to choose which side we would like to evolve. In the \textit{Schrödinger picture} we evolve the system $\rho(t)$ and the operators $A$ are considered to be static, while in the \textit{Heisenberg picture} we keep the state $\rho$ static and evolve the operators $A(t)$. Naturally these pictures are equivalent to each other and give the same physical expectation values,
\begin{equation}
    \label{eq:schrödinger-heis}
    \langle A \rangle (t) = \tr \left[A \rho(t) \right] = \tr \left[ A(t) \rho \right].
\end{equation}

A reversible evolution of a closed system is described by a unitary map $u(\rho)$
\begin{equation}
    \label{eq:unitary-map}
    \rho \mapsto \rho(t) = \mathcal{U}(\rho) = U \rho U^\dagger,
\end{equation}
in the Schrödinger picture, or as the dual map $\mathcal{U}^\dagger(A)$ in the Heisenberg picture. These maps are naturally related via \eqref{eq:schrödinger-heis}, and by using the cyclic property of the trace we can see that $A(t) = \mathcal{U}^\dagger(A) = U^\dagger A U$. From Stone's theorem, we know that if these unitaries build a homogeneous group and have a continuous time parameter, i.e. multiple time evolutions can always be written as one: $U(t) U(t') = U(t + t'),\,\forall t, t'\in\mathbb{R}$, then there exists a Hermitian $H$, such that
$$
U(t) = \exp (-\ii H t).
$$
If we let a closed system evolve naturally, it will do so with respect to the Hermitian $H$ we call the Hamiltonian of the system. In principle, we can also evolve a system coherently via a unitary given by some other Hermitians, which we will see when we come to the quantum circuits. It is important to note that if we act with a unitary on a state, it will preserve its norm, which means that in a closed system evolution we can be certain that our evolving quantum state remains a density matrix. The same for open system evolutions cannot be said in general, however as we will see, we will always use map that will preserve all density matrix properties, for which reason we will call the map \textit{physical.}
\section{Open systems} \label{sec:OQS}
While closed system evolutions can describe a lot of scenarios, it requires for us to know the \textit{whole} system. In many cases this is impossible to do, either because the number of degrees of freedom is intractable, or because the interaction with the environment is not even known. 
For this exact reason, there has been a lot of research on how to describe the evolution of just a subsystem, while discarding the information of the environment. Naturally, this will not be a reversible process any more, because the discarded information is generally irrecoverable by any physical operation on the subsystem alone. 

A physical meaningful evolution in the Schrödinger picture is given by a linear map $\mathcal{E}: \mathcal{D}(\mathcal{H}) \rightarrow \mathcal{D}(\mathcal{H}')$ that preserves all the density matrix properties. 
This means that $\mathcal{E}$ has to be \textit{trace preserving} (TP), $\tr[\mathcal{E}(A)] = \tr[A]$ and also \textit{completely positive} (CP). Positivity for a map means, that it would map positive operators to positive operators, i.e. $\mathcal{E}(A^\dagger A) \geq 0$. Imagine that we only act with $T$ on just part of the system and rest we leave alone, $\mathcal{E}\otimes \id_n$. We will also require this map to be positive, for all $n \in \mathbb{N}$, which is exactly what we call complete positivity. If a linear map fulfils both of these conditions we call it a CPTP map or a \textit{quantum channel}. 

We can also define here the dual map $\mathcal{E}^\dagger : \mathcal{B}(\mathcal{H'})\rightarrow \mathcal{B}(\mathcal{H})$ that describes the same evolution in the Heisenberg picture, which we can always find from $T$ by $\tr[\mathcal{E}(\rho) A] = \tr[\rho \mathcal{E}^\dagger(A)]$. The only difference here is that the trace preserving property manifests itself as the \textit{unitality} of the map, $\mathcal{E}^\dagger(\idm) = \idm$.

While the unitary map, discussed for the closed system earlier, is a CPTP map, we can now introduce more general maps too. A quantum channel can be brought into different forms and representations that will be useful in different settings. 

\smallskip
\subsection{Kraus representation.} A linear map $\mathcal{E}: \mathcal{B}(\mathcal{H}) \rightarrow \mathcal{B}(\mathcal{H}')$ is completely positive iff there exists $r$ operators $K_j: \mathcal{H}\rightarrow\mathcal{H}'$ such that
\begin{equation}
    \label{eq:kraus}
    \mathcal{E}(A) = \sum_{j=1}^r K_j A K_j^\dagger,
\end{equation}
and $\mathcal{E}$ is trace preserving iff $\sum_j K_j^\dagger K_j = \idm$. The minimum number of Kraus operators needed to describe a map is called the Kraus rank. 
Now we can also better understand quantum measurements, as they are a very important example of irreversible evolution. A quantum channel can be understood as a general measurement process where the classical outcome is discarded. There is a clear similarity, between the post-measurement state we found in \eqref{eq:post_meas_state_general} and the Kraus representation: we find that the Kraus operators $\{K_j\}$ are precisely the measurement operators $\{M_\alpha\}$. From \eqref{eq:kraus} we also see, that the resulting state after applying all the measurements, or $\mathcal{E}(\rho)$, is the ensemble average over all possible post-measurement states.

This operator-sum representation is useful when our aim is to find out what state we get after applying a quantum channel to an initial state. However, sometimes we are interested in the map itself and want to analyse its static properties. The next representation could provide some help with that. 

\smallskip
\todo{rewrite this a bit more rigorously, mention isometry}
\subsection{Choi-Jamiolkowski representation.} A channel that acts on $d = \dim(\mathcal{H})$ dimensional operators, can be expressed in terms of a $d^2\times d^2$ matrix. But before we turn maps into matrices, let us first turn matrices into vectors. \textit{Vectorization} of a $d\times d$ matrix $A$ is the $d^2 \times 1$ column vector, denoted by $\mathrm{vec}(A)$, formed by stacking the columns of $A$ on top of one another:
$$
A = 
\begin{pmatrix}
    a &\! b \\
    c &\! d
\end{pmatrix}
\quad \Rightarrow \quad 
\mathrm{vec}(A) \equiv \ket{A}\rangle = 
\begin{pmatrix}
    a \\
    c \\
    b \\
    d
\end{pmatrix}.
$$
The way linear maps $\mathcal{E}$ are related to operators $\tau$ is slightly more complicated. Define a bipartite system $\mathrm{AB}$ with the joined Hilbert space $\mathcal{H}_A \otimes \mathcal{H}_B$. Then prepare this system in the maximally entangled state $\ket{\Omega}$
\begin{equation}
    \ket{\Omega} = \frac{1}{\sqrt{d}} \sum_{j=1}^d \ket{j}_A \otimes \ket{j}_B \equiv \frac{1}{\sqrt{d}} \sum_{j=1}^d \ket{jj}_{AB}.
\end{equation}
The Choi matrix $\tau$ is then found by applying the map $\mathcal{E}$ only on one part of the maximally entangled system $\tau = (\mathcal{E} \otimes \id_d)(\ket{\Omega}\bra{\Omega})$, and it is in one-to-one correspondence with the linear map $\mathcal{E}$ via
\begin{equation}
    \tr[A \mathcal{E}(B)] = d \tr[\tau A \otimes B^T].
\end{equation}
After some entry reshuffling (see \textbf{Appendix}) we can find a key vectorization identity:
\begin{equation}
    \mathrm{vec} (AXB) = (B^T \otimes A) \ket{X}\rangle.
\end{equation}
For instance if applied to the Kraus representation of the channel \eqref{eq:kraus}
we find that
\begin{equation}
    \label{eq:vectorized_kraus}
    \ket{\mathcal{E}(A)}\rangle = \left(\sum_j \bar K_j \otimes K_j\right)\ket{A}\rangle =: \hat{\mathcal{E}}\ket{A}\rangle.
\end{equation}

\smallskip
\subsection{Lindbladian generator.} So far we have described open quantum system evolutions with discrete maps that would transform some initial state to a final state over a fixed time. However, we are often interested in what hides in gaps between now and then, i.e. in continuous evolutions. Just as a single parameter, time, parametrized the unitary evolution of closed systems, it can also parametrize a family of quantum channels $(\mathcal{P}_t)_{t\geq 0}$: $\mathcal{D}(\mathcal{H}) \rightarrow \mathcal{D}(\mathcal{H})$. We call such a family a \textit{quantum dynamical semigroup} if it satisfies the following properties for all $t, s \in \mathbb{R}_+$
$$
\mathcal{P}_t\mathcal{P}_s = \mathcal{P}_{t + s} \quad\text{and}\quad \mathcal{P}_0 = \id.
$$
Mathematically, these are the properties that describe a continuous, homogeneous and \textit{Markovian} process. Physically, we can understand it as a memoryless process, i.e. that a future state depends only on the present state and not on any of the past states. In a moment we will see what physical assumptions are needed for the open quantum system to have its evolution described by a physical quantum dynamical semigroup. But first let us look into from the mathematical side.

An important result from the theory of one-parameter semigroups is that any such continuous family of maps is uniquely determined by a time-independent operator 
known as the generator, denoted by $\mathcal{L}$. The generator is a linear map on the space of operators, $\mathcal{L}: \mathcal{B}(\mathcal{H})\rightarrow \mathcal{B}(\mathcal{H})$, and is formally defined as the derivative of the dynamical map at $t=0$:, 
\begin{equation}
\mathcal{L} = \lim_{t \to 0^+} \frac{\mathcal{P}_t - \mathrm{id}}{t}.
\end{equation}
The generator thus describes the infinitesimal change that drives the dynamics. The naming also becomes clear since if we have $\mathcal{L}$, we can generate the semigroup $(\mathcal{P}_t)_{t\geq 0}$ through the exponentiation of the superoperator
\begin{equation}
    \mathcal{P}_t = \exp (t \mathcal{L}).
\end{equation}
For our purposes of course $\mathcal{L}$ cannot take up any arbitrary form. The evolution it generates $\mathcal{P}_t$ must be a valid physical process for all times $t \geq 0$, i.e. each $\mathcal{P}_t$ must be a CPTP map. In this case we call the semigroup a \textit{quantum Markov semigroup} (QMS) and we can describe the continuous evolution of a subsystem within an environment starting from some initial state $\rho(0)$ by
\begin{equation}
    \label{eq:oqs-evo-with-L}
    \rho(t) = \exp(t \mathcal{L})(\rho(0)).
\end{equation}
Or conversely,
\begin{equation}
    \frac{\dd}{\dd t}\rho(t) = \mathcal{L}(\rho(t)).
\end{equation}

What is the general form of $\mathcal{L}$ that generates physical evolutions? This was answered by Gorini, Kossakowski, Sudarshan, and Lindblad (GKSL) by a fundamental theorem they proved in \cite{gorini1976completely,lindblad1976generators}. The way to get to their result, starts by pointing out that any generator $\mathcal{L}$ of a QMS can be written in the form:
\begin{equation}
    \label{eq:qms_generator}
\mathcal{L}(\rho) = \Psi(\rho) + K\rho + \rho K^\dagger,
\end{equation}
where $\Psi: \mathcal{B}(\mathcal{H})\rightarrow \mathcal{B}(\mathcal{H})$ is a completely positive map and $K$ is an operator in $\mathcal{B}(\mathcal{H})$. The separation between the completely positive part and the remaining part allows us to reveal some structure, if we decompose the operator $K$ into its Hermitian and anti-Hermitian parts
\begin{equation}
    K = -V - \ii H,
\end{equation}
with the Hermitian operators: $H = \frac{\ii}{2}(K - K^\dagger)$ and $V = \frac{\ii}{2}(K + K^\dagger)$. Substituting this into the expression \eqref{eq:qms_generator} yields
\begin{equation}
\mathcal{L}(\rho) = -i[H, \rho] + \Psi(\rho) - \{V, \rho\}.
\end{equation}
Up until now the decomposition is completely general. But this is where the GKSL theorem comes in and finds a central constraint for $\mathcal{L}$ to be a generator of a CPTP semigroup: the map $\Psi$ must have a Kraus form $\Psi(\rho) = \sum_k L_k \rho L_k^\dagger$, and for the dynamics to be trace-preserving, the operator $V$ must be given by $\frac{1}{2}\sum_k L_k^\dagger L_k$. Putting these pieces together, we arrive at the celebrated GKSL form of the Lindbladian generator
\begin{equation}
    \label{eq:lindblad-generator}
    \mathcal{L}(\rho) = -\ii [H, \rho] + \sum_{k}\gamma_k \left(L_k^\dagger \rho L_k - \frac{1}{2}\left\{L_k^\dagger L_k, \rho\right\}\right).
\end{equation}
We can also clearly separate the different dynamics in this picture. The first term is the $\textit{coherent part}$, generated by the Hermitian $H$. Note, that $H$ does not necessarily have to be the Hamiltonian of the system, it could also be just some generator that coherently drives the system (foreshadowing). The second term is the \textit{dissipative part}, which describes the various ways of decoherence and dissipation via the Lindblad jump operators $\{L_k\}$ with a dissipation or jump rate given by $\gamma_k \geq 0$.

\smallskip
\subsection{Microscopic picture.} While it is nice to arrive at a mathematically rigorous result, it is also insightful to derive the same GKSL form of the generator from the physical principles and assumptions. For this we need to work with the microscopic model of a system $S$, interacting with a large environment or reservoir, $E$. The combined system is described by the following Hamiltonian
\begin{equation}
    H = H_S + H_E + H_I,
\end{equation}
where $H_S$ is the system Hamiltonian, $H_E$ the environment Hamiltonian, and $H_I$ defines the interaction between them. The total system is assumed to be closed, i.e. the evolution is purely unitary and governed by the von Neumann equation for the total density matrix $\rho$,
\begin{equation}
    \label{eq:von-neumann}
    \frac{\dd}{\dd t}\rho(t) = -\ii [H(t), \rho(t)].
\end{equation}
Note, that this is the same evolution as \eqref{eq:unitary-map} but written in the differential form. But of course, we are interested in the effective equation of motion of only the subsystem $\rho_S$. In order to derive it, we first have to translate the von Neumann equation to the interaction picture. This is another picture we can work in, and it is equivalent to the previously mentioned Schrödinger and Heisenberg pictures. The difference is that here we split the time evolutions between states and observables, and say that the states are evolved by the unitary generated by only the interaction Hamiltonian $H_I$ while the observables by $H_S + H_E$. In this picture the von Neumann equation in the Schrödinger picture \eqref{eq:von-neumann} will simply turn into
\begin{equation}
    \label{eq:von-neumann-int}
    \frac{\dd}{\dd t}\rho(t) = -\ii [H_I(t), \rho(t)],
\end{equation}
or in the integral form
\begin{equation}
    \rho(t) = \rho(0) - \ii \int_0^t \dd s [H_I(s), \rho(s)].
\end{equation}
Inserting this back into the von Neumann equation \eqref{eq:von-neumann-int} we find
\begin{equation}
    \label{eq:starting-integral-von-neumann}
    \frac{\dd}{\dd t}\rho_S(t) = -\int_0^t \dd s \tr_E \left[H_I(t), \left[H_I(s), \rho(s)\right]\right].
\end{equation}
Here, $\tr_E[\cdot]$ denotes the partial trace over the Hilbert space $\mathcal{H}_E$, which leads us to the system state, $\rho_S = \tr_E[\rho]$. Without loss of generality we assumed above that the interaction does not generate any dynamics in the environment from the initial state,
\begin{equation}
    \tr_E\left[H_I(t), \rho(0)\right] = 0.
\end{equation}
To proceed we need to use our first approximation, called the \textit{Born-Markov approximation}, which describes systems that interact weakly with their environment. Thus, the state of the environment is negligibly affected by the interaction and can be approximately separated from the system state, 
\begin{equation}
    \label{eq:sys-env-separable}
    \rho(t) \approx \rho_S(t) \otimes \rho_E.
\end{equation}
Though it might sound trivial, it is crucial to separate the system from the environment, because otherwise any statement about the system's dynamics alone is ill-defined. In \eqref{eq:sys-env-separable} we also assume that the environment is in a fixed stationary state $\rho_E$, which is essential if we want to derive a time-independent Lindblad generator \eqref{eq:lindblad-generator}. But it raises the question on which stationary state should we choose for $\rho_E$, and if this reflects at all what happens in Nature. In the next chapter we will see that choosing the environment to be the Gibbs state is necessary to derive thermalizing dynamics for the system $S$.
The Born-Markov approximation also assumes a timescale separation, in which the excitations in the environment decay much faster than what the system $\rho_S$ could ever resolve. This allows the replacement of $\rho(s)$ with $\rho(t)$, but also lets us replace $s$ with $t-s$ and send the upper limit to infinity (since the integrand decays quickly above the timescale of the environment). Applying all these steps to \eqref{eq:starting-integral-von-neumann} gets us
\begin{equation}
    \label{eq:born-markov-me}
    \frac{\dd}{\dd t}\rho_S(t) = -\int_0^\infty \dd s \tr_E \left[H_I(t), \left[H_I(t-s), \rho_S(t)\otimes\rho_E\right]\right].
\end{equation}
However, there is a problem here: these last steps indeed bring us closer to find a time-independent generator, it also fails to generate completely positive maps. Works from \cite{davies1974markovian,dumcke1979proper} illuminated why, and also provided the rigorous approximations we need to derive the Lindblad generator. Since the techniques derived in these works are necessary for us to understand the algorithms in this thesis, we give here a bit more thorough review. 

Let us write the interaction Hamiltonian in the Schrödinger picture as a tensor product of Hermitian system operators $A^a$ and Hermitian environment operators $B^a$
\begin{equation}
    H_I = \sum_a A^a \otimes B^a.
\end{equation}
The Bohr frequencies of the Hamiltonian $H_S$ are given by 
\begin{equation}
    B_{H_S} = \{\nu = \lambda_i - \lambda_j : \lambda_i,\lambda_j \in \mathrm{Spec}(H_S)\},
\end{equation}
where $\mathrm{Spec}(H_S)$ is the spectrum of $H_S$. Then we can write the Bohr decomposition of an operator $A$ as
\begin{equation}
    \label{eq:bohr-decomp}
    A = \sum_{\lambda_i, \lambda_j \in \mathrm{Spec}(H_S)} \Pi_i A \Pi_j = \sum_{\nu \in B_{H_S}} A_\nu,
\end{equation}
with
\begin{equation}
    A_\nu := \sum_{\lambda_i - \lambda_j = \nu} \Pi_i A \Pi_j.
\end{equation}
If we now do a unitary rotation, we find the operator's form in the interaction picture, 
\begin{equation}
    A(t) = e^{\ii H_S t} A e^{-\ii H_S t} = \sum_\nu e^{\ii \nu t} A_\nu \quad\text{or}\quad [H_S, A_\nu] = \nu A_\nu.
\end{equation}
Generally the Bohr spectrum $B_{H_S}$ is degenerate, meaning that there could be always different parts of an operator $A$ that would lead to the same energy difference $\nu$. The interaction Hamiltonian in the interaction picture follows
\begin{equation}
    H_I(t) = \sum_{a, \nu} e^{\ii \nu t} A^a_\nu \otimes B^a(t),
\end{equation}
with $B(t) = e^{\ii H_E t} B e^{-\ii H_E t}$. This finally lets us reinterpret the Born-Markov master equation \eqref{eq:born-markov-me} as
\begin{equation}
    \label{eq:bohr-decomposed-born-markov-me}
    \begin{split}
        \frac{\dd}{\dd t}\rho_S(t) =& \sum_{\nu, \nu'}\sum_{a, b}e^{\ii(\nu - \nu')t}\Gamma_{ab}(\nu)\left(A^b_\nu\, \rho_S(t)\, (A^{a}_{\nu'})^\dagger - (A^a_{\nu'})^\dagger A_{\nu'}^b\, \rho_S(t)\right) \\
        &+ \text{h.c.}
    \end{split}
\end{equation}
Here, h.c. stands for the Hermitian conjugate part. $\Gamma_{ab}(\nu)$ obtained by carrying out $\int \dd s$ in \eqref{eq:born-markov-me}, or as it turns out the half Fourier transforming the reservoir correlation functions:
\begin{equation}
    \label{eq:half-fourier-reservoir-corr-fn}
    \begin{split}
        \Gamma_{ab}(\nu) &:= \int_{0}^{\infty} \dd s\, e^{\ii\nu s} \tr_E\left[(B^a)^\dagger(t)B^b(t-s)\rho_E\right] \\
        &\equiv \int_{0}^{\infty} \dd s\,e^{\ii\nu s}\langle (B^a)^\dagger(t)B^b(t-s) \rangle \\
        &= \int_{0}^{\infty} \dd s\,e^{\ii\nu s}\langle (B^a)^\dagger(s)B^b(0) \rangle.
    \end{split}
\end{equation}
In the last equation we used the assumption that the bath is in a stationary state, $[H_E, \rho_E] = 0$, which leads to the correlation functions time translation invariant, with which we see that $\Gamma_{ab}(\nu)$, the function the determines the jump rates in the system-environment interactions is time-independent. 

For the Born-Markov form we also assumed a timescale separation between the system and the environment, which manifests here by assuming that the reservoir correlation functions decay rapidly $\lim_{s\rightarrow\infty}\langle (B^a)^\dagger(s)B^b(0) \rangle = 0$, before the system could ever be affected by it. But this also elucidate why a large or infinite environment is also required: if $\mathrm{Spec}(H)$ is not continuous then there will be quasi-periodic remnants left from the Fourier transform. From the physics point of view, in the case of a smaller, finite environment, the information that flowed from the system to the environment over the past interactions, would never die off, and after some time (Poincarré recurrence time) it would reappear and affect the system. Hence, ruining the Markovianity of the dynamics.

But there are still problems left after the Born-Markov approximation: $(i)$ there is still time-dependency left in the generator, $(ii)$ by assuming the coefficient matrix \eqref{eq:half-fourier-reservoir-corr-fn} to be constant over all times, it can in general break the positivity of the whole generator in some time regions as well. Thus, it cannot yet be a generator of a QMS. The solution to both of these problems is the \textit{secular approximation}. In \eqref{eq:bohr-decomposed-born-markov-me}, the secular terms are those that have $\nu = \nu'$, and they correspond to resonant processes where the total energy of the system and the environment is conserved. The non-secular terms, with $\nu \neq \nu'$ are off-resonant transitions to virtual states for which the total energy is temporarily not conserved. In the master equation they are also accompanied by the oscillating prefactor $e^{\ii(\nu - \nu')t}$. This oscillation is happening on the system's timescale, $\tau_S \sim \nu^{-1} \sim |\nu - \nu'|^{-1}$. Another timescale that is relevant here, is the relaxation time $\tau_R$ over which the system state changes through the interaction with the reservoir. It is inversely proportional to the coupling strength or jump rates, i.e. quite long in the weak-coupling limit we are working with. The secular approximation then says that we are only interested in cases where
\begin{equation}
    \label{eq:secular-approx}
    \tau_S \gg \tau_R \quad\text{or}\quad {|\nu - \nu'|} \gg 1.
\end{equation}
This effectively means that we can drop all the off-diagonal terms with $\nu \neq \nu'$ in \eqref{eq:bohr-decomposed-born-markov-me} and find
\begin{equation}
    \label{eq:after-secular-me}
    \frac{\dd}{\dd t}\rho_S(t) = \sum_{\nu}\sum_{a, b}\Gamma_{ab}(\nu)\left(A^b_\nu\, \rho_S(t)\, (A^{a}_{\nu})^\dagger - (A^a_{\nu})^\dagger A_{\nu}^b\, \rho_S(t)\right) + \text{h.c.},
\end{equation}
which is finally a QMS generator! To bring it into the Lindblad form we need to decompose the coefficient matrix into its Hermitian and anti-Hermitian parts:
\begin{equation}
    \Gamma_{ab}(\nu) = \frac{1}{2}\gamma_{ab}(\nu) + \ii S_{ab}(\nu)
\end{equation} 
with
\begin{equation}
    \gamma_{ab}(\nu) = \Gamma_{ab}(\nu) + \Gamma_{ab}^\dagger(\nu)
    = \int_{-\infty}^{\infty}\dd s\, e^{\ii\nu s}\langle (B^a(s))^\dagger B^b(0)\rangle
\end{equation}
and
\begin{equation}
       S_{ab}(\nu) = \frac{1}{2\ii}\left(\Gamma_{ab}(\nu) - \Gamma_{ab}^\dagger(\nu)\right).
\end{equation}
$S_{ab}(\nu)$ is a Hermitian and $\gamma_{ab}(\nu)$ is called the \textit{Kossakowski matrix} and it contains all jump rate coefficients; it is positive since now the Fourier transform over the reservoir correlations is a full Fourier (Bochner's theorem).
By inserting these into \eqref{eq:after-secular-me}, we have finally derived the (non-diagonal) Lindblad generator from microscopic principles and more-or-less physical assumptions:
\begin{equation}
    \label{eq:lindblad-microscopic}
    \mathcal{L}(\rho_S) = -\ii[H_{LS}, \rho_S] +  \sum_{ab}\sum_\nu \gamma_{ab}(\nu)\left(A^b_\nu\, \rho_S (A^a_\nu)^\dagger - \frac{1}{2}\left\{(A^a_\nu)^\dagger A^b_\nu, \rho_S\right\}\right).
\end{equation}
Comparing it to the general Lindblad generator \eqref{eq:lindblad-generator} we find that the Lindblad jumps are $L^a_\nu = \sqrt{\gamma_a(\nu)} A^a_\nu$, and that there is a contributing coherent term of the form:
\begin{equation}
    H_{LS} = \sum_{ab}\sum_\nu S_{ab}(\nu)(A_\nu^a)^\dagger A_\nu^b.
\end{equation}
This is the Lamb shift Hamiltonian, and it originates from the interactions with the environment which effectively shift the system's energy levels (since $[H_S, H_{LS}] = 0$). As the system $S$ evolves coherently with respect to its Hamiltonian $H_S$, the Lamb shift will add to it and lead to an evolution with $H_\text{eff} = H_S + H_{LS}$. 

There is one more comment to be made about \eqref{eq:lindblad-microscopic}: in its current form it is non-diagonal in the indices $a, b$. However, it is always possible to diagonalize the Kossakowski matrix $\gamma_{ab}(\nu)$, since it is Hermitian. Thus, we can always find some unitary $U$ such that $\gamma_a(\nu) = U \gamma_{ab}(\nu) U^\dagger$. We can accordingly find the rotated jump operators for which we can work in this diagonal picture by
\begin{equation}
    \label{eq:diagonalize-jumps-in-generator}
    L^a_\nu = \sum_b U_{ab} A^b_\nu.
\end{equation} 
Such a diagonalization `trick' can be quite useful for efficient numerical analysis as we will later see. 

Up to this point, we have introduced the main character of Markovian open quantum system dynamics: the Lindbladian generator. We have derived it from mathematical and physical principles, and have found that they generate CPTP maps, i.e. as the subsystem evolves within an environment, it preserves its density matrix properties. Of course, we are not only interested in how the system evolves but also where it is headed. Generally, these evolutions inch the system towards some stationary state(s). But a priori it is not clear if the system ever reaches that equilibrium (within a practical amount of time), since it can stumble upon a metastable state, where the evolution can get stuck. In order to understand this side of the dynamics, we will introduce the theory of Markov chains in the next chapter.
\chapter{Detailed Balanced Markov Dynamics}
A Markov chain is a stochastic process that describes the evolution of a probability distribution over a set of states $\mathcal{X}$. The dynamics are governed by a stochastic map that has the Markov property, which dictates that the probability of transitioning to any future state depends only upon the current state. Once we get to know the fundamentals of Markov chains, we can ask questions about to which state or states does the Markov chain evolve to, and how fast it reaches its goal. For more details on this topic, we refer to the indispensable source from Levin and Peres, \cite{levin2017markov}.

\section{Classical Markov Semigroups}
It is instructive to start the discussion with the classical theory of Markov chains. Even if the notions have to be generalized later to the quantum case, introducing them first classically, sheds a lot of light on the path ahead.

In classical theory, a probability distribution over the states $x\in\mathcal{X}$ at some time $t$ is given by a row vector $\mu_t$, where $\mu_t(x)$ gives the probability of finding the system in state $x$. The evolution of the probability distribution is described by a stochastic transition matrix $P$ as
$$
\mu_{t+1} = \mu_t P.
$$
The stochasticity of $P$, i.e. having non-negative entries and each of its row sums to $1$, ensures that $\mu_t$ stays a well-defined probability distribution at all times during the evolution.

In this work, we will be interested in Markov chains that for long enough times $t\rightarrow\infty$ let the distribution $\mu_t$ converge to an equilibrium distribution, known as the \textit{stationary distribution} $\pi$. This distribution is a fixed point of the dynamics, 
\begin{equation}
    \pi = \pi P.
\end{equation}
The existence of a unique stationary distribution to which an initial distribution converges, is guaranteed if the chain is \textit{irreducible} and \textit{aperiodic} due to the Perron-Frobenius theorem. Irreducible means that it is possible to reach any state from another state, while aperiodic means that the chain does not get stuck in a cycle of states. If both of these conditions hold, we will also call the chain \textit{ergodic}. As an example, if the system has some symmetry and the transitions prescribed by $P$ are all within that symmetric subspace, then the Markov chain would never be able to discover the states in the whole space, and would never reach its global stationary state $\pi$.

In Markov chain applications, the goal can be simply stated as: construct a transition matrix $P$ that prepares a desired probability distribution $\pi$ from some initial distribution(s). Naturally, we would first ask, why is there even a need for this, couldn't we directly sample from the target distribution $\pi$ if we know its form already?
In general the answer is no, and there are cases where it is indeed easier to walk through a Markov chain defined by $P$ that converges to $\pi$. Such a case is exactly what we have at our hands in this thesis, finding and sampling from the Gibbs (or Boltzmann) distribution
\begin{equation}
    \label{eq:cl-gibbs}
    \pi_\beta(x) = \frac{e^{-\beta E(x)}}{Z}, \quad  Z = \tr [e^{-\beta E(x)}]
\end{equation}
With $E(x)$ the energy of the state $x$, and $Z$ the partition function, which is a sum over all the possible energies in a system... which is a Herculean task to compute for larger system sizes.

Thankfully, there is a universal way to construct Markov chains with a given stationary distribution $\pi$. (Though universal, it does not always mean an efficient Markov chain.) The main concept for these processes, is \textit{detailed balance}. A chain is said to satisfy the detailed balance condition if, at equilibrium $\mu_t = \pi$, the probability flow from state $x$ to state $y$ is equal to the flow from $y$ to $x$:
\begin{equation}
    \label{eq:db-classical}
    \pi(x) P(x, y) = \pi(y) P(y, x), \quad\forall x, y\in\mathcal{X}.
\end{equation}
These chains are also called \textit{reversible}. This refers to a time-reversal symmetry of the dynamics, like a musical palindrome, for which we could not distinguish if it is played forwards or backwards. From a thermodynamic point of view these are processes that produce zero entropy at the equilibrium.

Important to note that while the detailed balance condition \eqref{eq:db-classical} is a sufficient condition for $\pi$ to be a stationary distribution, it is not strictly necessary. A weaker version of the condition, the \textit{global balance condition}, does not require for each pair of states to have reversible probability flow, but rather the net probability flow for each state $\pi(x)$ should remain zero under $P$:
\begin{equation}
    \sum_{x\in\mathcal{X}} \pi(x) P(x, y) = \sum_{x\in\mathcal{X}} \pi(y) P(y, x) = \pi(y),\quad\forall y\in\mathcal{X}.
\end{equation}
This latter condition however, is much less utilized in applications, since in terms of states it is not a local condition any more.

There is also another, operational, way of understanding detailed balance, which will help with the quantum detailed balance generalization later on. For this, we will work in the space of real-valued functions that act on the state space, $L^2(\pi)$, equipped with the $\pi$\textit{-weighted inner product}, $\langle\cdot,\cdot\rangle_\pi$. For two functions $f: \mathcal{X}\rightarrow \mathbb{R}$ and $g: \mathcal{X}\rightarrow \mathbb{R}$, it is defined as:
\begin{equation}
    \langle f, g\rangle_\pi := \sum_{x\in\mathcal{X}} f(x) g(x) \pi(x),
\end{equation}
where $\pi$ is still the stationary distribution of the chain. Since $\pi(x) > 0$ for all $x$ in an irreducible chain, the inner product is positive definite. Basically, it weighs each state's contribution by how likely that state is to occur at equilibrium. An inner product induces the natural norm in this space, $\|f\|_{2} = \sqrt{\langle f, f \rangle_\pi}$.

An operator (our transition matrix $P$) is then said to be self-adjoint with respect to the $\pi$-weighted inner product  if for any two functions $f$ and $g$, the following holds:
\begin{equation}
    \label{eq:P-selfadjoint-cl}
    \langle f, Pg\rangle_\pi = \langle Pf, g\rangle_\pi.
\end{equation}
Both sides written out gives,
\begin{equation}
    \sum_{x, y\in\mathcal{X}} f(x) g(y) \pi(x) P(x, y) = \sum_{x, y\in\mathcal{X}} f(x) g(y) \pi(y) P(y, x),
\end{equation}
which is exactly the detailed balance condition \eqref{eq:db-classical}, since it holds for all $f, g$. After showing the converse, we find a fundamental equivalence that, even if not in its exact form but in concept, survives in the quantum regime: a Markov chain satisfies detailed balance if and only if its transition matrix $P$ is a self-adjoint with respect to the $\pi$-weighted inner product.  
A consequence that comes from \eqref{eq:P-selfadjoint-cl} is that all the eigenvalues $\lambda_i(P)$ of $P$ are real, which allows us to order all eigenvalues on the real line, 
$$
1 = \lambda_1(P) > \lambda_2(P) \geq ... \geq -1,
$$
and have a clean definition for the spectral gap $\gamma(P) := 1 - \lambda_2(P)$. Intuitively, a large gap $\gamma(P)$ means that all terms in the distribution that are not the stationary state decay quickly. Conversely, a vanishing gap implies very slow convergence. We will soon see how to bound the mixing time of a chain with the spectral gap, but let us first explain how to construct a Markov chain that prepares a desired target distribution.

\label{parag:MH}
\smallskip
\subsection{Metropolis-Hastings algorithm.} The generic recipe for constructing a transition matrix $P$ that satisfies the detailed balance condition for any target distribution $\pi$. It was pioneered and was run on the MANIAC computer as early as 1953 \cite{metropolis1953equation}, and the algorithm has been generalized by Hastings in \cite{hastings1970monte}. The recipe is as follows:
\begin{enumerate}
    \item Decide upon an initial state $x\in\mathcal{X}$ where the Markov chain starts from. 
    \item Propose a new state $y\in\mathcal{X}$ based on a \textit{proposal matrix} $Q(x, y)$, which is some irreducible Markov chain on the state space $\mathcal{X}$. This can be as simple as randomly picking a state in the space. 
    \item Accept the state with probability given by the \textit{acceptance matrix} $A(x,y)$.
\end{enumerate}
Therefore, the transition matrix $P$ for the resulting Markov chain is a combination of a proposal $Q$ and acceptance $A$.
\begin{equation}
    P(x,y) =
\begin{cases}
    Q(x,y)A(x,y), \quad &\text{if}\; y\neq x, \\
    1 - \sum_{z\neq x} Q(x, z)A(x, z), \quad &\text{if}\;y = x.
\end{cases} 
\end{equation}
The key here is to define the acceptance probability $A(x, y)$ such that the final chain $P$ satisfied detailed balance condition \eqref{eq:db-classical} with respect to $\pi$, 
\begin{equation}
    \pi(x) Q(x, y) A(x, y) = \pi(y) Q(y, x) A(y, x),
\end{equation}
for which the choice that maximizes the acceptance rates $A$ (thus allowing for more efficient exploration in the state space) is given by the Metropolis-Hastings form: 
\begin{equation}
    \label{eq:MH-acceptance}
    A(x,y) = \min\left(1, \frac{\pi(y)Q(y, x)}{\pi(x)Q(x,y)}\right).
\end{equation}
It is straightforward to show that this choice of $A(x,y)$ forces the detailed balance condition to hold for $P$. The special case, when the proposal matrix is symmetric, i.e. $Q(x, y) = Q(y, x)$, is the Metropolis algorithm that Hastings later generalized to \eqref{eq:MH-acceptance}.

Finally, it becomes clear why this algorithm is so powerful. The acceptance $A$ and thus the full transition $P$ depends only on the \textit{ratio} of the target probabilities. For sampling from the Gibbs distribution $\pi(x) = e^{-\beta E(x)} / Z$, we would never have to compute the intractable partition function $Z$. The sole information the Markov chain needs in this case, is to the energy difference between $x\rightarrow y$: $\Delta E = E(y) - E(x)$. 

Note, that the Metropolis-Hastings algorithm makes no guarantee about the speed of convergence. It heavily depends on the state space, the proposal  matrix $Q$, and its interplay with the target distribution $\pi$. Quantifying this efficiency is the subject of mixing time analysis, which will be crucial for us for later findings as well. 

We have now come to the point where we can turn our attention to recent developments in quantum Gibbs sampling. The moment we mentioned Gibbs sampling in the classical setting, we also learnt the recipe that allowed researchers for decades to sample Gibbs states. So much so, that the \hyperref[parag:MH]{Metropolis-Hastings algorithm} is now an indispensable word of the sampling jargon. 

\smallskip
\subsection{Mixing time.} The go-to quantity that defines how fast the Markov chain converges to the target distribution. First, we need to define a distance between two probability distributions $\mu$ and $\nu$ using the \textit{total variation distance}:
\begin{equation}
    \label{eq:tv-distance}
    \|\mu - \nu\|_{TV} = \frac{1}{2} \sum_{x\in\mathcal{X}} |\mu(x) - \nu(x)|.
\end{equation} 
The distance from stationarity at time $t$ is then defined as the worst-case distance, maximizing over all possible starting states $x$:
\begin{equation}
    \label{eq:convergence-thm}
    d(t) := \sup_{x\in\mathcal{X}} \|P^t(x, \cdot) - \pi\|_{TV} 
\end{equation}
With respect to this distance, the mixing is the amount of time needed to reach a given $\varepsilon>0$ closeness to the stationary state $\pi$:
\begin{equation}
    \label{eq:mixing-time-cl}
    t_\text{mix}(\varepsilon) := \inf\left\{t \geq 0:\; d(t) \leq \varepsilon \right\}
\end{equation}
The mixing time clearly tells us how efficiently does a Markov chain prepare the target distribution $\pi$. Often we are interested in how the mixing time scales with the system size $n$ (or the state space size $|\mathcal{X}| = 2^n$). We distinguish between these cases
\begin{alignat*}{3}
    \text{(i)}\;   &\textit{Slow mixing}  &\quad& \text{with} &\quad& t_\text{mix}(\varepsilon) = \mathcal{O}(\exp(n) \log(1/\varepsilon)), \\
    \text{(ii)}\; &\textit{Rapid mixing}  &\quad& \text{with} &\quad& t_\text{mix}(\varepsilon) = \mathcal{O}(\mathrm{poly}(n) \log(1/\varepsilon)),\\
    \text{(ii)}\; &\textit{Optimal mixing}  &\quad& \text{with} &\quad& t_\text{mix}(\varepsilon) = \mathcal{O}(n \log(n) \log(1/\varepsilon)).
\end{alignat*}
Sometimes there could be also other parameters that can have a huge influence on the mixing times, like the temperature, which always have to be taken into account for Gibbs sampling. 
Broadly we can say, that high temperature systems unless there are geometric bottlenecks, or if the system is at a critical point of phase transition. In this latter case the system would have non-local correlations, that a local Markov chain like the Metropolis-Hastings would fail to resolve.

In order to bound the mixing times as the next step, we will make use of the following.
While the total variation distance is an $L^1$ distance, it can be rewritten in terms of the $L^1(\pi)$ distance $\|\cdot\|_{\pi, 1}$ induced by the $\pi$-weighted inner product, 
\begin{equation}
    \label{eq:tv-1norm}
    \begin{split}
        2\|P^t(x,\cdot) - \pi \|_{TV} &= \sum_{y\in\mathcal{X}} |P^t(x, y) - \pi(y)| \\
        &= \sum_{y} \pi(y) \left|\frac{P^t(x, y)}{\pi(y)} - \mathbf{1}\right| \\
        &= \left\|\frac{P^t(x, \cdot)}{\pi(\cdot)} - \mathbf{1}\right\|_{\pi, 1}.
    \end{split}
\end{equation}
After using Jensen's inequality, we can upper-bound the TV-distance from the stationary state $\pi$ via the $L^2(\pi)$ distance to $\pi$:
\begin{equation}
    \label{eq:tv-2norm}
    \begin{aligned}
          2\|P^t(x,\cdot) - \pi \|_{TV} &= \left\|\frac{P^t(x, \cdot)}{\pi(\cdot)} - \mathbf{1}\right\|_{\pi, 1} \leq \left\| \frac{P^t(x, \cdot)}{\pi(\cdot)} - \mathbf{1}\right\|_{\pi, 2} \\
          \quad \xRightarrow{\sup} d(t) &\leq \frac{1}{2} \sup_{x\in\mathcal{X}}\left\| \frac{P^t(x, \cdot)}{\pi(\cdot)} - \mathbf{1}\right\|_{\pi, 2} 
    \end{aligned}
\end{equation}

\smallskip
\subsection{Spectral bounds.} Bounding the mixing time is a central challenge in the analysis of Markov chains. While simulation can provide empirical estimates, rigorous analytical bounds are required for guaranteeing efficiency. One of the most fundamental approaches, relates the convergence rate to the spectral properties of the transition matrix $P$. While finding those spectral properties can be computationally difficult, this method at least provides a universal way to bound mixing times for all kinds problems.

Determining the mixing time is more subtle than simply identifying the decay rate of the second-largest eigenmode. While the eigenvalue spectrum dictates the asymptotic speed at which the system forgets its initial state, it does not account for the distance the system must traverse. To capture both effects, we must move from the natural metric of the total variation distance, to a metric compatible with spectral analysis.

As we have shown in \eqref{eq:tv-1norm} and \eqref{eq:tv-2norm}, the TV distance, is basically an $L^1(\pi)$ norm, and it can be upper bounded by the $L^2(\pi)$ norm. Working in the $L^2(\pi)$ space with the $\pi$-weighted inner norm is suitable, because the spectral theorem becomes applicable. This states, that for a self-adjoint operator, like our detailed balanced $P$, has an orthonormal basis $\{v_i\}$ and real eigenvalues $\lambda_i(P)$. (If the chain is taken to be `lazy', i.e. $P(x, x)\geq 1/2$, then $\lambda_i(P) \geq 0$.) 

In this function space, probabilistic notions turn into geometric objects. Probability distributions $\mu_t$ are represented by relative density vector $h_t = \mu_t / \pi$, while the stationary distribution $\pi$ corresponds to the constant vector $\mathbf{1}$ that stays constant under the transition: $P \mathbf{1} = \mathbf{1}$. Statistical expectations are now projections onto the axis given by $\pi$, $\mathbb{E}_\pi [f] = \langle f, \mathbf{1} \rangle_\pi$. Then the variance is the squared $L^2(\pi)$ norm of the orthogonal terms, $\text{Var}_\pi(f) = \|f - \mathbf{1}\|_{\pi, 2} ^2$. To see this, take the deviation from the equilibrium to be the vector $g_t = h_t - \mathbf{1}$. This is orthogonal to the stationary state because
\begin{equation}
    \mathbb{E}_\pi[g_t] \equiv \langle g_t, \mathbf{1} \rangle_\pi = \sum_x \left(\frac{\mu_t(x)}{\pi(x)}- \mathbf{1}\right) \pi(x) = \sum_x \mu_t(x) - \sum_x \pi(x) = 0.
\end{equation}
Thus, the deviation vector $g_t$ is spanned only by the rest of the eigenvector over its whole evolution,
\begin{equation}
    g_0 = \sum_{i=2}^{|\mathcal{X}|} c_i v_i \rightarrow P^t g_0 = g_t = \sum_{i=2}^{|\mathcal{X}|} c_i \lambda^t_i(P) v_i.
\end{equation}
The sum is dominated by the second-largest eigenvalue $\lambda_2$ and allows us to bound the decay of the distance (or variance) at time $t$ as,
\begin{equation}
    \begin{split}
        \sqrt{\text{Var}_\pi(h_t)} &\equiv \left\|g_t\right\|_{\pi, 2}  = \left\|\sum_{i=2}^{|\mathcal{X}|} c_i \lambda_i^t(P) v_i\right\|_{\pi, 2}\leq (1-\gamma(P))^t\sum_{i=2}^{|\mathcal{X}|} c_i v_i \\
        &\leq e^{-\gamma(P) t} \|g_0\|_{\pi, 2}.
    \end{split}
\end{equation}
Here, we used that $\lambda_2(P) = 1 - \gamma(P)$. This inequality shows already the structure we expected for the mixing process to have: the decay part is an exponential governed by the \textit{relaxation time}, $t_\text{rel} := 1 / \gamma(P)$, that multiplies a term quantifying the initial distance from stationarity. 

To derive a rigorous mixing time bound, we must still evaluate this initial distance for the worst-case starting configuration. If the system starts in a specific state $x_0$, the initial point mass will have a peak at $\mu_0(x_0) = 1$. Then the $L^2(\pi)$ norm of the initial deviation is:
\begin{equation}
    \begin{split}
            \left\|g_0\right\|_{\pi, 2} &= \left\|\frac{\mu_0}{\pi} - \mathbf{1}\right\|_{\pi, 2} = \sqrt{\sum_x \pi(x) \left(\frac{\mu_0}{\pi} - \mathbf{1}\right)^2} = \sqrt{\frac{1}{\pi(x_0)} -1} \\
            &\leq \frac{1}{\sqrt{\pi(x_0)}}.
    \end{split}
\end{equation}
Evidently, the worst-case means that the initial distribution $\mu_0$ concentrated all its probability mass on a state $x_0$ that has the least contribution to the stationary state: $\pi(x_0) = \pi_{\min} := \min_x \pi(x)$. Consequently, the mixing time is the required for the exponential decay to suppress this large initial factor, the \textit{burn-in} factor, down to a desired $\varepsilon$ size,
\begin{equation}
    \frac{1}{\sqrt{\pi_{\min}}} e^{-\gamma(P) t} \leq \varepsilon
\end{equation}
Rearranging, simplifying and not forgetting that mixing time was defined with the TV distance, this inequality yields the standard spectral bound on the mixing time:
\begin{equation}
\label{eq:spectral-bound}
t_\text{mix}(\varepsilon) \leq \frac{1}{\gamma(P)} \ln \left(\frac{1}{2\varepsilon \sqrt{\pi_{\min}}}\right).
\end{equation}

While the spectral bound \eqref{eq:spectral-bound} is universally applicable, it is often loose for high-dimensional systems where $\pi_{\min}$ is exponentially small. This motivates the need for sharper tools, such as the \textit{logarithmic Sobolev inequalities}, that were pioneered by Leonard Gross in \cite{gross1975logarithmic}.

\smallskip
\subsection{Logarithmic Sobolev bounds.} The weakness of the $1 / \sqrt{\pi_{\min}}$ scaling in the spectral bound stemmed from the $L^2(\pi)$ norm. Essentially, we have been tracking the convergence of the variance, $\mathrm{Var}_\pi(h) = \|h - \mathbf{1}\|_{\pi, 2}^2 = \|g\|_{\pi, 2}^2$ with $h = \mu / \pi$. But it turns out that looking at the convergence of entropy instead can lead to stronger results on the mixing time. 

The \textit{relative entropy} (or \textit{Kullback-Leibler divergence}) between the evolving distribution $\mu$ and the stationary distribution $\pi$ is defined as:
\begin{equation}
    \label{eq:relative-entropy}
    \text{Ent}_\pi(h) = \sum_{x\in\mathcal{X}} \pi(x) h(x) \ln h(x).
\end{equation}
Before we bound the relative entropy, it is constructive to explain why and when this measure is better than the variance. 

Consider the near-equilibrium regime, where the relative density takes the form $h(x) = \mathbf{1} + \varepsilon g(x)$ with $\varepsilon \ll 1$, thus $\mathbb{E}_\pi[g] = 0$. By performing a Taylor expansion of the function $f(t) = t \ln t $ from \eqref{eq:relative-entropy} around $t=1$, one finds that the entropy and variance agree to leading order:
\begin{equation} 
    \text{Ent}_\pi(\mathbf{1} + \varepsilon g(x)) = \frac{1}{2} \text{Var}_\pi(\mathbf{1} + \varepsilon g(x)) + \mathcal{O}(\varepsilon^3).
\end{equation}
In other words, when the initial distribution is close to $\pi$, the variance and the relative entropy are locally equivalent and both decay exponentially. 

Where they differ significantly is in their treatment of far-from-equilibrium states. Thus, consider now a chain starting in the worst-case initial distribution $h_0$, i.e. a point mass at the least likely state,
$$
h_0(x_0) = \frac{1}{\pi_{\min}}, \quad h_0(x\neq x_0) = 0.
$$
In this case, the variance in $L^2(\pi)$ is given by
\begin{equation}
    \mathrm{Var}_\pi(h_0) = \mathbb{E}_\pi[h_0^2] - 1 = \frac{1}{\pi_{\min}} - 1. 
\end{equation}
This scales exponentially with the system size $n$, since $\pi_{\min}\sim e^{-n}$. On the other hand, the relative entropy is now 
\begin{equation}
    \label{eq:initial-entropy}
    \mathrm{Ent}_\pi(h_0) = \pi_{\min} \left(\frac{1}{\pi_{\min}} \ln \frac{1}{\pi_{\min}}\right) = \ln \frac{1}{\pi_{\min}},
\end{equation}
which scales only linearly with the system size $n$. Of course, this is not the end of the story. To exploit this better scaling, we must prove that the entropy also contracts exponentially, which requires establishing a Log-Sobolev inequality.

Just as the spectral gap $\gamma(P)$ controls the decay of the variance, we define a new parameter, the constant $\alpha$, to control the decay of the entropy. It has been shown \cite{bobkov2006modified}, that the Markov chain satisfies a modified Log-Sobolev inequality (MLSI) with constant $\alpha > 0$ if for all positive functions $h$:
\begin{equation}
    \alpha \text{Ent}_\pi(h) \leq \mathcal{E}(h, \ln h),
\end{equation}
where $\mathcal{E}(f, g) = \langle f, (\id-P)g \rangle_\pi$ is the Dirichlet form associated with the chain. This specific form is chosen because the time derivative of the entropy is exactly the negative of the Dirichlet form on the right-hand side: $\frac{d}{dt}\text{Ent}_\pi(h_t) = -\mathcal{E}(h_t, \ln h_t)$ (in the continuous time limit). Consequently, this inequality guarantees that the relative entropy contracts exponentially at a rate $\alpha$:
\begin{equation}
    \text{Ent}_\pi(h_t) \leq e^{-\alpha t} \text{Ent}_\pi(h_0).
\end{equation}
To translate this entropy decay into a bound on the mixing time, we relate the entropy back to the TV distance using \textit{Pinsker's inequality}, which states that $2\|\mu - \pi\|_{TV}^2 \leq \text{Ent}_\pi(\mu/\pi)$.
Together with the initial entropy calculated above \eqref{eq:initial-entropy}:
$$
2 d(t)^2 \leq \text{Ent}_\pi(h_t) \leq e^{-\alpha t} \ln \frac{1}{\pi_{\min}} \leq 2\varepsilon^2.
$$
Solving for $t$ yields the Log-Sobolev bound on the mixing time:
\begin{equation}
    \label{eq:log-sobolev-bound}
    t_\text{mix}(\varepsilon) \leq \frac{1}{\alpha} \ln \left( \frac{\ln(1/\pi_{\min})}{2\varepsilon^2} \right).
\end{equation}
Compared with the spectral bound \eqref{eq:spectral-bound}, we see a significant improvement for high-dimensional systems. Recall that for a system of size $n$, the inverse stationary probability scales as $1/\pi_{\min} \sim e^{-n}$. The spectral bound depends linearly on this factor inside the logarithm, resulting in a mixing time scaling of $\mathcal{O}(n)$. In contrast, the Log-Sobolev bound takes the logarithm of this factor, resulting in a scaling of $\mathcal{O}(\ln n)$.

Intuitively, this improvement arises because the entropy is less sensitive to the spiky nature of the initial point-mass distribution than the $L^2$ norm is. The latter requires the dynamics to smooth out the sharp peak of the initial state before the bound becomes `useful'. The entropy, being logarithmic, penalizes this initial peak much less, thus having a tighter bound on mixing time. 

Unfortunately, the stronger bound comes with a significant analytical cost. While the spectral gap $\gamma(P)$ is well-defined for any reversible chain (by $1-\lambda_2$), establishing a strictly positive Log-Sobolev constant $\alpha$ is much harder. It generally requires proving geometric properties of the state space, e.g. Bakry-Émery criterion which relates to the Ricci curvature of the generator \cite{bakry2006diffusions}; or product structures (non-interacting systems) \cite{diaconis1996logarithmic}, that was extended to near-product structures (weakly-interacting systems, e.g. for high-temperatures) by \cite{stroock1992logarithmic}. But these results are neither universally present, nor easy to compute. Interestingly, we can relate the two decay rates as $2\alpha \leq \gamma(P)$ \cite{diaconis1996logarithmic}, which means the rate of entropy decay is actually slower than of the variance decay. But this factor is less consequential, than the difference in handling the far equilibrium starting points. All in all, while the Log-Sobolev inequalities are standard for proving rapid mixing in high-dimensional systems, the spectral gap remains the more universal and accessible tool for general contexts.

\smallskip
\subsection{Continuous-time Markov chains.} The natural generalizations of Markov chains, where the transitions occur at random times $t\in[0, \infty)$. The discrete nature of the chains we previously introduced, did also come with some artefacts, that will not be present for the continuous case, e.g. periodicity issues. 

In the continuous limit, the transition matrix $P$ is replaced by the classical Markov semigroup $(P_t)_{t\geq 0}$ generated by $L$. No surprise here, indeed, this is the classical analogue of the generator we introduced for open quantum system evolutions in Section \ref{sec:OQS}. Only that the generator $L$ will have slightly different but analogue properties to the quantum one. The evolution of the probability distribution $\mu(t)$ is governed by the classical Master Equation:
\begin{equation}
    \frac{\dd}{\dd t} \mu(t) = \mu(t)L
\end{equation}
The solution to this differential equation is formally given by the matrix exponential $\mu(t) = \mu(t=0) e^{t L}$. Just as the Lindbladian preserves the trace and positivity of the density matrix, the classical generator $L$ preserves the normalization and positivity of the probability vector $\mu(t)$ at all times. This requires the off diagonal entries of $L$ to be non-negative, and the row sums to be zero, $\sum_y L(x, y) = 0$.

The concept of detailed balance \eqref{eq:db-classical} gets rephrased in terms of the generator. Since detailed balance should hold at all times, then the derivatives of the probability flows must match at $t=0$ as well. The generator is defined exactly as that derivative, $L = \frac{\dd}{\dd t}P_t |_{t=0}$, and thus the detailed balance condition takes the form:
\begin{equation}
    \pi(x)L(x, y) = \pi(y) L(y, x), \quad \forall x\neq y.
\end{equation}
Note, that the diagonal is already fixed by $L(x, x) = - \sum_{y\neq x} L(x, y)$. This also implies, just as detailed balance implies global balance, that $\pi L = 0$. Equivalently to the above, in the operational framework, $L$ must be self-adjoint with respect to the $\pi$-weighted inner product:
\begin{equation}
    \langle f, Lg\rangle_\pi = \langle Lf, g\rangle_\pi.
\end{equation}
This self-adjointness of $L$ ensures that the semigroup $P_t$ is also self-adjoint, thereby preserving reversibility at all times $t\geq0$.

The convergence to the stationary state is still determined by the spectral gap, only that now the spectrum of a reversible $L$ is non-positive, and the unique stationary state $\pi$ has the eigenvalue $\lambda_1(L) = 0$, such that the spectral gap is simply $\gamma(L) = -\lambda_2(L)$. Of course, the relation between the eigenvalues of $P$ and $L$ at some time $t$ is $\lambda_i(P) = e^{t \lambda_i(L)}$. Nevertheless, there is no change about the contractive nature of the deviation from the stationary state, only that now we would replace $\gamma(P)$ with $\gamma(L)$, and have
$$\|h_t - \mathbf{1}\|_{\pi, 2} \leq e^{\lambda_2(L) t}\|h_0 - \mathbf{1}\|_{\pi, 2}.$$
Thus, the spectral bounds on mixing time \eqref{eq:mixing-time-cl}, also apply for the continuous case.

Similarly, the Log-Sobolev inequality adapts as expected to the generator setting. The Dirichlet form becomes $\mathcal{E}(f, g) = \langle f, -Lg \rangle_\pi$, and the modified Log-Sobolev inequality is now,
$$
    \alpha \text{Ent}_\pi(h) \leq \langle h, -L (\ln h) \rangle_\pi.
$$
This recovers the mixing time bound \eqref{eq:log-sobolev-bound} for the continuous case as well. 

With this formalism established, we can finally tackle the quantum setting in the next section. We will see that quantum Gibbs sampling is essentially the task of designing a quantum generator (the Lindbladian) that satisfies the \textit{quantum detailed balance} with respect to the quantum Gibbs state. For the convergence analysis of quantum Markov semigroups, we will also describe the required non-commutative generalizations of the spectral and Log-Sobolev results.

\section{Quantum Markov Semigroups}

It is now time to introduce the core subject of this work: the evolution of open quantum systems towards the thermal equilibrium. Just as we described classical mixing using generators and weighted inner products, we will formulate the quantum problem using reversible quantum Markov semigroups.

In the quantum setting, the probability distributions are given by density matrices $\rho(t)$, which in general does not commute with other density matrices. Gibbs sampling then means that we
construct a Lindbladian generator $\mathcal{L}$ of the GKSL form \eqref{eq:lindblad-generator}, such that the system state $\rho(t)$ converges to the stationary density matrix, i.e. the quantum Gibbs state,
\begin{equation}
    \rho_\beta = \frac{e^{-\beta H}}{Z}, \quad Z = \tr[e^{-\beta H}].
\end{equation}
What makes this quantum compared to the similar classical definition of  $\pi_\beta$ in \eqref{eq:cl-gibbs}, is the fact that now we have non-commuting terms in $H$ contributing to the energy $E(x)$. If $H$ has exclusively commuting terms inside (e.g. classical Ising model), then it is possible to write it as a diagonal matrix in some product basis, like the computational basis ($\sigma_Z$ basis). The Gibbs state $\rho_\beta$ would also be diagonal. Showing, that in very special cases quantum Gibbs sampling reduces to the classical problem.

The evolution is governed by the quantum Master equation $\frac{\dd}{\dd t} \rho = \mathcal{L}(\rho)$ which generates the semigroup $(\mathcal{P}_t)_{t\geq 0}$, just as we described for open quantum systems \eqref{eq:oqs-evo-with-L}. For the system to globally thermalize, $\rho_\beta$ must be the unique fixed point, $\mathcal{L}(\rho_\beta) = 0$. Note, that now the operators $\mathcal{L}$ (and $\mathcal{P}$) have been upgraded to superoperators $\mathcal{L}: \mathcal{B}(\mathcal{H})\rightarrow \mathcal{B}(\mathcal{H})$ since they act on operators (density matrices) instead of vector distributions.

However, generalizing the concept of detailed balance to the quantum case is not straightforward. Classical detailed balance was defined by the reversibility condition $\pi(x) P(x, y) = \pi(y) P(y, x)$. In the quantum case, the transition probabilities are replaced by operators, and the non-commutativity of the Hamiltonian and the states, $[\rho, H] \neq 0$ in general, means there is no unique way to `weight' the inner product by the stationary state.

\smallskip
\label{parag:qdb}
\subsection{Quantum detailed balance.} To define reversibility for a quantum process, we must adopt the operational framework we introduced classically: a generator is detailed balanced if it is self-adjoint with respect to a specific $\sigma$-weighted inner product. The non-weighted inner product on the space of bounded operators $\mathcal{B(H)}$ is the \textit{Hilbert-Schmidt inner product}:
\begin{equation}
    \label{eq:HS}
    \langle A, B \rangle := \tr[A^\dagger B].
\end{equation}
In fact, the Hilbert space we get with this inner product, can be thought of as the infinite temperature limit, where the stationary state is the maximally mixed state, $\lim_{\beta\rightarrow 0}\rho_\beta = \id_\mathcal{H} / \dim \mathcal{H}$. Since we juggle several inner products in this text, let us fix the following notation: the adjoint of the superoperator $\Phi\;:\;\mathcal{B(H)}\rightarrow\mathcal{B(H)}$ is always with respect to the Hilbert-Schmidt inner product $\langle \cdot,\cdot\rangle$, and it is denoted by $\Phi^\dagger$, e.g. 
\begin{equation}
    \langle A, \Phi(B) \rangle = \langle \Phi^\dagger(A), B \rangle.
\end{equation}
Written out explicitly for a general case we get,
\begin{equation}
    \Phi(\cdot) = \sum_j \alpha_j A_j (\cdot) B_j \quad\Rightarrow\quad \Phi^\dagger(\cdot) = \sum_j \bar{\alpha}_j A^\dagger_j (\cdot) B^\dagger_j
\end{equation}

The moment we have to deal with finite temperatures, we have to weight \eqref{eq:HS}. Due to non-commutativity there is no unique way to distribute the multiplication by $\sigma$ within the trace. Thus, we find a family of weighted inner-products:
\begin{equation}
    \label{eq:weighted-HS}
    \langle A, B \rangle_{\sigma, s} := \tr[\sigma^s A^\dagger \sigma^{1-s} B], \quad s\in[0,1].
\end{equation}
While all of these choices reduce to the classical product when matrices commute, they correspond to fundamentally different physical processes in the quantum setting.
The expression above \eqref{eq:weighted-HS} defines a continuous family of inner products parametrized by $s$, but they effectively boil down to just two distinct cases. As it turns out \cite{fagnola2007generators}, the inner products in \eqref{eq:weighted-HS} lead to coinciding QMS for all $s\in[0,1] \backslash \{\frac{1}{2}\}$. Consequently, all of these inner products are equivalent to the \textit{Gelfand-Naimark-Segal (GNS) inner product} defined by $s\in\{0, 1\}$ in \cite{alicki1976detailed}, i.e. to the following,
\begin{equation}
    \label{eq:GNS-product}
    \langle A, B\rangle_{\sigma, 1} = \tr[\sigma A^\dagger B].
\end{equation}
For $s=\frac{1}{2}$ on the other hand, we get a distinct, symmetrically weighted inner product, which was introduced in \cite{goldstein1995kms} and was named after \textit{Kubo-Martin-Schwinger (KMS)}:
\begin{equation}
\langle A, B \rangle_{\sigma, \frac{1}{2}} := \tr[\sigma^{1/2} A^\dagger \sigma^{1/2} B].
\end{equation}
A generator $\mathcal{L}$ satisfies quantum detailed balance if it is self-adjoint with respect to either of these weighted inner products, and it is called GNS or KMS detailed balanced for $s = 0, 1$ or $s=\frac{1}{2}$ respectively.
\begin{equation}
    \langle A, \mathcal{L}(B) \rangle_{\sigma, s} = \langle \mathcal{L}(A), B \rangle_{\sigma, s} \quad\forall A, B \in \mathcal{B(H)}
\end{equation}
Equivalently, if detailed balance holds we can express the adjoint of the generator as
\begin{equation}
    \label{eq:general-qdb-adjoint}
    \mathcal{L^\dagger}(\cdot) = \sigma^{s-1} \mathcal{L}(\sigma^{1-s}\, \cdot\, \sigma^{s}) \sigma^{-s}.
\end{equation}
Let us check if this is the case, then the stationary state $\sigma$ is indeed stationary. Since $\mathcal{L}$ generates unital maps in the Heisenberg picture, i.e. $\mathcal{L}^\dagger(\idm) = 0$, we can write 
\begin{equation}
    \label{eq:sigma-is-fix}
    \begin{split}
        0 &= \langle A, \mathcal{L}^\dagger(\idm) \rangle_{\sigma, s} \\
        &= \langle A, \sigma^{s-1} \mathcal{L}(\sigma) \sigma^{-s} \rangle_{\sigma, s} \\
        &= \tr[A^\dagger \mathcal{L}(\sigma)] \\
        &= \langle A, \mathcal{L}(\sigma)\rangle, \qquad\forall A\in\mathcal{B(H)}.
    \end{split}
\end{equation}
Where we used \eqref{eq:weighted-HS} and \eqref{eq:general-qdb-adjoint}, and assumed that $\sigma$ is a full-rank density matrix (e.g. the Gibbs state at finite temperatures) and are invertible. 
This confirms, that if the generator and its adjoint has the relation $\eqref{eq:general-qdb-adjoint}$ then $\sigma$ is the stationary state of the generated QMS, $\mathcal{P}_t(\sigma) = \sigma$.

An alternate but equivalent definition of quantum detailed balance that we will use later is given by its \textit{discriminant} or parent Hamiltonian, 
\begin{equation}
    \label{eq:discriminant}
    \begin{split}
        \mathcal{D}(\sigma, \mathcal{L}) &:= \sigma^{-\frac{1}{4}}\mathcal{L}(\sigma^{\frac{1}{4}} \cdot \sigma^{\frac{1}{4}})\sigma^{-\frac{1}{4}} \\
        &= \sigma^{\frac{1}{4}}\mathcal{L}^\dagger(\sigma^{-\frac{1}{4}} \cdot \sigma^{-\frac{1}{4}})\sigma^{\frac{1}{4}} \equiv \mathcal{D}^\dagger(\sigma, \mathcal{L}).
    \end{split}
\end{equation}
For the last equation we used the definition of $\mathcal{L}^\dagger$ from \eqref{eq:general-qdb-adjoint}. Thus, conjugating a Lindbladian with its unique stationary state in this manner, can be thought of as a similarity transformation that makes the Lindbladian Hermitian. 
Then breaking detailed balance can be understood as having a non-Hermitian discriminant. Or equivalently, since every operator can be decomposed into a Hermitian and anti-Hermitian part, that we have a non-vanishing anti-Hermitian part.
\begin{equation}
    \begin{split}
        \mathcal{D}(\rho, \mathcal{L}) &= \mathcal{H}(\rho, \mathcal{L}) + \mathcal{A}(\rho, \mathcal{L}) \\
        \mathcal{D}(\rho, \mathcal{L})^\dagger &= \mathcal{H}(\rho, \mathcal{L}) - \mathcal{A}(\rho, \mathcal{L}).
    \end{split}
\end{equation}
\begin{definition} \label{def:approx-db}
    \textup{(Approximate detailed balance condition)}. The Lindbladian $\mathcal{L}$ with full-ranked fixed point $\rho$ satisfies the $\varepsilon-$approximate detailed balance condition if the anti-Hermitian part $\mathcal{A}$ is small
    $$
    \frac{1}{2}\left\|\mathcal{D}(\rho, \mathcal{L}) - \mathcal{D}(\rho, \mathcal{L})^\dagger \right\|_{2\rightarrow2} = \left\|\mathcal{A}(\rho, \mathcal{L})\right\|_{2\rightarrow2} \leq \varepsilon
    $$
\end{definition}
This is one of the main quantities we will track in our numerical analyses, because it gives a clear picture on how much we break detailed balance via implementation errors.

An advantage of the transformation \eqref{eq:discriminant}, is that it keeps the spectrum the same, which makes spectral analysis possible for Lindbladians through their parent Hamiltonians. This is often useful if the jump operators are local or quasi-local. In these cases we often find an already known gapped quantum lattice model, or at least have more tools to use to bound that spectral gap. They used these techniques for GNS detailed balanced Lindbladians with commuting system Hamiltonians in \cite{kastoryano2016quantum}, and for the recent KMS Lindbladians where the parent Hamiltonians are provably frustration-free in \cite{rouze2025efficient,bergamaschi2024quantum,vsmid2025polynomial}.

At the moment it might seem that the difference between GNS and KMS detailed balances is purely mathematical, and it is merely a choice of the parameter $s$. But for the thermalizing process there are clear physical consequences of choosing one or the other. To understand the distinction, we must introduce the \textit{modular operator} $\Delta_\sigma\;:\;\mathcal{B(H)} \rightarrow \mathcal{B(H)}$ defined as:
\begin{equation}
    \Delta_\sigma(A) = \sigma A \sigma^{-1}
\end{equation}
In the classical limit where all operators commute, $\Delta_\sigma$ reduces to the identity. Now it encodes the non-commutativity structure of the algebra with respect to the steady state. Since $\rho_\beta \propto e^{-\beta H}$, the action of the modular operator is intimately linked to the (imaginary) time evolution by the system Hamiltonian $H$. It was shown in \cite{alicki1976detailed} that a generator $\mathcal{L}$ satisfies the GNS detailed balance if and only if it satisfies KMS detailed balance, \textit{and} it commutes with the modular operator:
\begin{equation}
    [\mathcal{L}, \Delta_\sigma] = 0.
\end{equation}
This additional condition implies that the dissipative dynamics generates a semigroup that commutes with the free unitary evolution of the system. In other words, the dissipation does not generate coherences between energy levels and thus respects the energy conservation of the system. Notice, that this is actually in line with the standard weak-coupling limit in which we explicitly needed the secular approximation \eqref{eq:secular-approx} (dropping non-secular or coherent terms) to derive a physical generator. The KMS detailed balance however leads to more permissive yet still physical dynamics that couple populations (diagonal terms) to coherences (off-diagonal terms). The fact that it retains non-secular terms but also physical, since we started from a physical Lindbladian and restricted it to have KMS detailed balance, means that it cannot be derived in the standard weak-coupling limit with a heat bath. On what kind of first principles would lead to KMS detailed balanced thermalization, we can refer to works where they use the \textit{singular coupling limit} \cite{alicki2007quantum}, or the \textit{thermal quantum collision / repeated interaction limit} \cite{attal2007langevin}.

\smallskip
\label{subsec:davies}
\subsection{Davies generator.} Canonically, the first thermalizing Lindbladian has been written up in the standard weak-coupling limit (with Born-Markov and secular approximations) by Davies in \cite{davies1974markovian}. This has been achieved by coupling the system to an infinite free heat bath, and the resulting GNS detailed balanced generator was such a milestone that it got named after Davies.

We have already derived most of what we need for it in \eqref{eq:lindblad-microscopic} where we derived the Lindblad generator from first principles. There the main idea was to decompose the interaction operator into eigenoperators of the system Hamiltonian, $A_\nu$, such that $[H, A_\nu] = -\nu A_\nu$. These operators characterize the transitions between the energy levels with an energy separation of $\nu$. In the Davies generator the transition rates given by $\gamma_a(\nu)$ satisfy a specific condition, called the \textit{KMS condition} (not to be confused with KMS symmetry), which relates the absorption and emission via the Boltzmann factor for any $a$ indexing the set of jump operators $\{A^a\}$:
\begin{equation}
    \label{eq:KMS-condition}
    \frac{\gamma_a(-\nu)}{\gamma_a(\nu)} = e^{\beta \nu}
\end{equation}
To remind ourselves, the Kossakowski matrix $\gamma(\nu)$ is the full Fourier transform of the reservoir correlation functions \eqref{eq:half-fourier-reservoir-corr-fn}. If we assume the environment is in the Gibbs state $\rho_E = e^{-\beta H_E} / Z_E$, then the correlation functions of the reservoir indeed satisfy the KMS condition \eqref{eq:KMS-condition}. This is a necessary condition for any thermalizing dynamics, but not sufficient. In this text we will discuss two distinct ways to thermalize a system dissipatively. If we work in the standard weak-coupling limit and drop all non-secular terms, we end up with the GNS detailed balanced Lindbladian from Davies
\begin{equation}
    \label{eq:davies-generator}
    \mathcal{L}_\text{Davies}(\rho) = -\ii[H_{LS}, \rho] +  \sum_{a}\sum_\nu \gamma_{a}(\nu)\left(A^a_\nu\, \rho (A^a_\nu)^\dagger - \frac{1}{2}\left\{(A^a_\nu)^\dagger A^a_\nu, \rho\right\}\right).
\end{equation}
We can see that due to the secular approximation, all terms with $\nu \neq \nu'$ have been dropped, while keeping all terms with $\nu=\nu'$. A unitary rotation via \eqref{eq:diagonalize-jumps-in-generator} was also used such that after diagonalization we have $a=b$ compared to \eqref{eq:lindblad-microscopic}. Note, that the only other difference to the Lindbladian derived in the microscopic picture with the Bohr decomposition, is that the transition rates $\gamma_{a}(\nu)$ now fulfil \eqref{eq:KMS-condition}. 

In a few steps we can even see how the Davies generator satisfies the GNS detailed balance. From $[H, A_\nu] = \nu A_\nu$, the following relation will be quite useful
\begin{equation}
    \label{eq:sigma-A-commutation}
    \rho_\beta A_\nu = e^{-\beta \nu} A_\nu \rho_\beta.
\end{equation}
This helps to show that since the Lamb-Shift Hamiltonian $H_{LS}$ commutes with the Gibbs state, and so does the anti-commutator term $\mathcal{R} \propto \{A_\nu^\dagger A_\nu, \rho\}$ in the Lindbladian, they are both self-adjoint with respect to the GNS inner product \eqref{eq:GNS-product}. Note, that this is only so simple because we have solely secular terms in the Davies generator. As we will later see, the insight for the KMS detailed balanced Lindbladian was to keep and make the non-secular terms in the coherent and reflective part $\mathcal{R}$ comply with each other. 

The remaining transition part needs a bit more work to prove that it is self-adjoint. But in both the GNS and KMS cases, the main requirement is for the transitions to fulfill KMS condition \eqref{eq:KMS-condition}, and to have  
the Bohr frequencies of the system have $B_H = - B_H$. The following derivation should apply for all jumps $a$ of $\{\sqrt{\gamma^a(\nu)A_\nu^a}\}$.
\begin{equation}
    \label{eq:davies-gns-db}
    \begin{alignedat}{2}
        \langle X, \mathcal{T}(Y) \rangle_{\rho_\beta, 1} &= \sum_\nu \gamma(\nu) \tr\left[\rho_\beta X^\dagger A_\nu^\dagger Y A_\nu\right] \\
        &= \sum_\nu \gamma(\nu) e^{\beta \nu} \tr\left[\rho_\beta (A_\nu X A_\nu^\dagger)^\dagger Y\right]  &\quad& (\text{cyclic trace; \eqref{eq:sigma-A-commutation}})\\
        &= \sum_\nu \gamma(-\nu) \tr\left[\rho_\beta ((A_{-\nu})^\dagger X A_{-\nu})^\dagger Y\right] &\quad& (\text{KMS condition})\\
        &= \sum_\nu \gamma(\nu) \tr\left[\rho_\beta (A_{\nu}^\dagger X A_{\nu})^\dagger Y\right] &\quad& (-\nu \rightarrow \nu)\\
        &= \langle \mathcal{T}(X), Y \rangle_{\rho_\beta, 1}
    \end{alignedat}
\end{equation}
Hence, all parts of the Davies generator are self-adjoint with respect to the GNS inner product, which implies that the Gibbs state $\rho_\beta$ is the fixed point, $\mathcal{L}_\text{Davies}(\rho_\beta) = 0.$

It is natural to question the physical ubiquity of the Davies semigroup. The secular approximation, which enforces the GNS detailed balance for the generator by discarding non-secular terms, implies that the environment seems to discriminate perfectly between the system's energy levels. But it is better to think of it as the ideal case where the coupling is so weak, and the environment is so large that the relaxation time $\tau_R\sim 1 / \lambda_2(\mathcal{L})$ is orders of magnitude longer than the system's intrinsic time $\tau_S \sim 1/ \nu_{\min}$. In this asymptotically exact description, all the non-secular terms average to zero, hence their absence in the Davies generator. The other obstacle that arises, is the fact that for a simulation \textit{we} would have to be the ones (instead of Nature) that discriminate, or resolve the energy levels of the system. An implementation that constantly has to distinguish between Bohr frequencies cannot be efficiently done for larger system sizes. With growing system size the gap between the Bohr frequencies gets exponentially small. Accordingly, the amount if simulation time needed, due to the Heisenberg's uncertainty principle $\Delta E \Delta t \geq \hbar / 2$, would exponentially increase as well, leading to an intractable computation even if it was done in a quantum computer. \todo{Refer to QPE section.}

Before we move on to introduce and work with the KMS detailed balanced Lindbladian, let us first go over the notion of mixing time and the techniques to bound it in the quantum setting. 

\smallskip
\subsection{Mixing time.} Similar to the classical setting, we need to first quantify the distance from the equilibrium before and then establish rigorous bounds on the convergence speed. The natural generalization of the total variation distance \eqref{eq:tv-distance} is the \textit{trace distance}. For any two quantum states $\rho$ and $\sigma$, it is defined via the non-commutative $L^1$ norm called the Schatten 1-norm, or trace norm $\|\cdot\|_1$
\begin{equation}
    \|\rho - \sigma\|_1 = \tr\left[ \sqrt{(\rho-\sigma)^\dagger (\rho-\sigma)} \right].
\end{equation}
Operationally, it represents the maximum probability of distinguishing two states with an optimal measurement. The mixing time is defined analogously to the classical case \eqref{eq:mixing-time-cl}. We look for the time $t$ required for the evolved state $\rho(t) = \mathcal{P}_t(\rho_0)$ to be $\varepsilon$ close to the stationary Gibbs state $\rho_\beta$ for any initial state $\rho_0$,
\begin{equation} \label{eq:mixing-time-defi}
    t_\text{mix}(\varepsilon) := \inf \left\{ t \ge 0 : \sup_{\rho_0} \|\mathcal{P}_t(\rho_0) - \rho_\beta\|_1 \le \varepsilon \right\}.
\end{equation}
Intuitively, the mixing of a quantum Markov semigroup is a more complex process than its classical counterpart. While a classical chain merely has to redistribute probability mass among the states (population or diagonal elements), a quantum process must simultaneously accomplish two tasks: it must drive the populations to the Gibbs distribution \textit{and} it must destroy all off-diagonal quantum coherences between energy levels. We will soon see what role this new decoherence timescale will play.
On top of that, the non-commutative nature of the space also creates new ways to thermalize. In the classical case, if an energy barrier stopped the convergence, a large enough kick could get the evolution moving. In the quantum case, such an energy barrier could be tunnelled through (if it is not too wide). But some transitions could also be possibly blocked if the phases do not align with the target subspace. Thus, a priori it is not always clear, that quantum mixing is faster than classical. To find out more about it, a recent paper \cite{rakovszky2024bottlenecks} discusses bottlenecks in both regimes. 

\todo{Possibly mention classical energy transport for Davies.}

\smallskip
\subsection{Quantum spectral bounds.} To bound the mixing time, we once again turn to the spectral properties of the generator. But just as before, we have to change to the (non-commutative) $L^2$ space. There, we can work with an orthonormal basis due to the spectral theorem that still applies.

We consider the Hilbert space of bounded operators $\mathcal{B(H)}$ equipped with the weighted inner product $\langle \cdot, \cdot \rangle_{\sigma, s}$ for which the generator $\mathcal{L}$ satisfies detailed balance. This means that even in the quantum setting, the generator is self-adjoint to an inner-product, and accordingly has real eigenvalues. Together with the fact that $\mathcal{L}$ generates CPTP maps, we know that all the eigenvalues lie on the non-positive real line, $0 = \lambda_1(\mathcal{L}) > \lambda_2(\mathcal{L}) \geq \dots$, and the spectral gap is clearly defined as:
\begin{equation}
    \lambda(\mathcal{L}) := - \lambda_2(\mathcal{L}).
\end{equation}
The arguments follow similarly to the classical case. The generator's only fixed point is $\rho_\beta$, as shown in $\eqref{eq:sigma-is-fix}$, and the dynamics in the subspace orthogonal to this kernel are governed by the remaining eigenvalues and eigenstates. The variance of any observable (equivalently, the distance of any state from the equilibrium) contracts exponentially at a rate given by $\lambda(\mathcal{L})$.

The trace distance used for the mixing time definition, can be bounded by the weighted $L^2(\sigma)$ norm via the non-commutative Cauchy-Schwarz inequality. Without loss of generality, we can use the KMS inner product with $s=\frac{1}{2}$, and write $\|\cdot\|_{2, \text{KMS}} := \|\cdot\|_{2, \rho_\beta, 1/2}$ for the following derivations. The evolved state $\rho_t = \mathcal{P}_t(\rho_0)$ the distance from the stationary state $\rho_\beta$ takes the form:
\begin{equation}
    \label{eq:trdist-ub-kms-norm}
    \begin{split}
        \|\rho_t - \rho_\beta\|_1  &= \|\rho_\beta^{-\frac{1}{2}}(X_t - \idm)\rho_\beta^{-\frac{1}{2}} \|_1 \\
        &= \|X_t - \idm\|_{1, \text{KMS}} \\
        &\leq \|X_t - \idm\|_{2, \text{KMS}} \\
        &\leq e^{-\lambda(\mathcal{L}) t} \|X_0 - \idm\|_{2, \text{KMS}}.
    \end{split}
\end{equation}
Here, $X_t := \rho_\beta^{-\frac{1}{2}} \rho_t \rho_\beta^{-\frac{1}{2}}$ can be interpreted as the relative density of $\rho_t$ to the Gibbs state with respect to the KMS inner product. (Admittedly, there is no unique "density" any more, as for the GNS inner product we would get a different form.) The evolution of $X_t$ is governed by the KMS self-adjoint operator: $$X_t = \rho_\beta^{-\frac{1}{2}}\mathcal{P}_t\left(\rho_\beta^{\frac{1}{2}} X_0 \rho_\beta^{\frac{1}{2}}\right) \rho_\beta^{-\frac{1}{2}}.$$
For a rigorous upper-bound we have to take the worst-case scenario for the remaining norm in \eqref{eq:trdist-ub-kms-norm}, i.e. when the initial state $\rho_0$ has a small overlap with the Gibbs state $\rho_\beta.$ An example would be a high energy initial state, that tries to thermalize to a low temperature Gibbs state, i.e. where $\beta \gg 1$. After squaring and writing out the remaining norm, it simplifies to
\begin{equation}
    \begin{split}
        \|X_0 - \idm\|_{2, \text{KMS}}^2 &= \| X_0 \|_{2, \text{KMS}} - \idm \\
        &\leq \tr\left[ \rho_0 \rho_\beta^{-1} \rho_0\right] \\
        &= \bra{\psi} \rho_\beta^{-1} \ket{\psi}.
    \end{split}
\end{equation}
Here, we first used that $X_0 = \rho_\beta^{-\frac{1}{2}} \rho_0 \rho_\beta^{-\frac{1}{2}}$, then chose a pure state for $\rho_0 = \ket{\psi}\bra{\psi}$ to have an initial state whose probability mass concentrates on a state that is the least compatible with $\rho_\beta$. This allows us to bound the norm, just as in the classical case, with the smallest eigenvalue of $\rho_\beta$:
$$
\sup_{\psi} \langle \psi | \rho_\beta^{-1} | \psi \rangle = \|\rho_\beta^{-1}\|_\infty = \frac{1}{\lambda_{\min}(\rho_\beta)}.
$$
Substituting back to \eqref{eq:trdist-ub-kms-norm} we get the upper bound on the trace distance from the fixed point,
$$
\|\rho_t - \rho_\beta\|_1 \leq e^{-\lambda(\mathcal{L})t} \frac{1}{\sqrt{\lambda_{\min}(\rho_\beta)}}.
$$
Finally, solving it for time for which the distance falls below some $\varepsilon$ gives the spectral bound on the mixing time,
\begin{equation}
    \label{eq:quantum-spectral-bound}
        t_\text{mix}(\varepsilon) \leq \frac{1}{\lambda(\mathcal{L})} \ln \left( \frac{1}{\varepsilon \sqrt{\lambda_{\min}(\rho_\beta)}} \right).
\end{equation}
A neat result, since it coincides with the classical bound \eqref{eq:spectral-bound}, only that now the relaxation time due to the spectral gap of the quantum generator $\mathcal{L}$ and the burn-in factor due to the smallest eigenvalue of the quantum state $\rho_\beta$ have quite different behaviour to their classical analogues. Note, that the result is the same even if we were to use the GNS inner product in the derivation. 

Unfortunately, the problem of handling far-equilibrium initial states for this bound is also inherited from the classical result. But worry not, there has been a vast amount of research done to use the logarithmic Sobolev techniques also in the non-commuting quantum setting. \todo{CITE from rouzé}

\smallskip
\label{sec:prelim-qlsi}
\subsection{Quantum logarithmic Sobolev bounds.} The problem with the spectral bound \eqref{eq:quantum-spectral-bound} is that the factor $1 / \sqrt{\lambda_{\min}(\rho_\beta)}$ typically scales as $e^{\beta n / 2}$ for an $n$-particle system, making the bound loose for large systems. To get a tighter bound even in these cases, we once again turn to the (quantum) relative entropy of the system, and try to find that logarithmic factor that helped us out for the classical case as well, \eqref{eq:log-sobolev-bound}.

The quantum relative entropy between the states $\rho, \sigma \in \mathcal{D(H)}$ is defined as:
\todo{write it up with Ent from rouzé, its not $\text{Ent}(\rho)$}
\begin{equation} 
    \text{Ent}_{\sigma}(\rho) := D(\rho \| \sigma) = \tr[ \rho (\ln \rho - \ln \sigma)]
\end{equation}
In analogy to the classical case, we seek a constant $\alpha > 0$ such that the entropy relative to the Gibbs state $\rho_\beta$ decays exponentially, i.e.
$$
D(\rho_t \| \rho_\beta) \leq e^{-\alpha t} D(\rho_0 \| \rho_\beta).
$$
The differential condition of this decay is the \textit{Quantum Logarithmic Sobolev Inequality} (QLSI), which is satisfied for a quantum Markov semigroup generated by $\mathcal{L}$
\begin{equation}
    \alpha D(\rho \| \rho_\beta) \leq \mathcal{E}
\end{equation}
where $\mathcal{E}_L(\rho) := -\frac{d}{dt} D(\rho_t \| \sigma)|_{t=0} = -\tr[\mathcal{L}(\rho)(\ln \rho - \ln \sigma)]$ is the quantum entropy production (or Dirichlet form for the log-relative entropy).

$$\mathcal{E_L}(X, Y) := -\langle X, \mathcal{L}(Y)\rangle_{\sigma}$$

\chapter{Quantum Circuits}
\smallskip
\section{Basics}
The smallest non-trivial system in $\mathcal{D(H)}$ is the \emph{qubit}, the quantum analogue of a classical bit. Its Hilbert space is $\mathcal{H} = \mathbb{C}^2$ and a pure state takes the form
$$
\ket{\psi} = \alpha\ket{0} + \beta\ket{1}\quad\text{with}\quad |\alpha|^2 + |\beta|^2 = 1.
$$
An $n$-qubit register lives in the tensor product $\mathcal{H}^{\otimes n} = (\mathbb{C}^2)^{\otimes n}$ of dimensions $2^n$, with computational basis $\{\ket{x}\}_{x\in\{0,1\}^n}$ where $\ket{x} = \ket{x_1}\otimes\cdots\otimes\ket{x_n}$. Mixed states $\rho$ on this register are the density operators on $\mathcal{H}^{\otimes n}$, in the same sense as introduced \hyperref[sec:quantum-states]{previously}.

A \emph{quantum gates} acting on $n$ qubits is a unitary $U \in \mathcal{B}((\mathbb{C}^2)^{\otimes n})$. Any unitary on $n$ qubits can be approximated to arbitrary accuracy by composing gates from the universal set consisting of all single-qubit unitaries together with the two-qubit \textsc{CNOT} gate \cite{bible}. The single-qubit Pauli matrices that will time to time show up in the thesis are defined as
\begin{equation}
    \label{eq:paulis}
    X = \begin{pmatrix} 0 & 1 \\ 1 & 0 \end{pmatrix}, \quad
    Y = \begin{pmatrix} 0 & -\ii \\ \ii & 0 \end{pmatrix}, \quad
    Z = \begin{pmatrix} 1 & 0 \\ 0 & -1 \end{pmatrix}.
\end{equation}
They are simultaneously Hermitian and unitary, and satisfy $X^2 = Y^2 = Z^2 = \idm$. For $\sigma_i, \sigma_j \in \{X, Y, Z\}$ they fulfil the following (anti-)commutation relations
$$
[\sigma_i, \sigma_j] = 2\ii\,\epsilon_{ijk}\,\sigma_k, \qquad \{\sigma_i, \sigma_j\} = 2\,\delta_{ij}\,\idm
$$
with $\epsilon_{ijk}$ the totally antisymmetric Levi-Civita symbol. Together with $\idm$ they span $\mathcal{B}(\mathbb{C}^2)$, and tensor products of Paulis form an orthogonal basis of $\mathcal{B}((\mathbb{C}^2)^{\otimes n})$ under the Hilbert-Schmidt inner product \eqref{eq:HS}. \todo{we could reason that this is why we like to write Hams as Paulis and ref to our Ham in numerical section}

Each Pauli generates a one-parameter family of \emph{Pauli rotations} $R_\sigma(\theta) := \ee^{-\ii\theta\sigma/2}$ for $\sigma\in\{X,Y,Z\}$ (which can also be extended to multiple qubit rotations $R_{\sigma_j\sigma_k}$). One of them is of particular interest to us, since it can encode arbitrary amplitudes, namely 
\begin{equation}
    R_Y(\theta) = \cos(\tfrac{\theta}{2})\,\idm - \ii\sin(\tfrac{\theta}{2})\,Y.
\end{equation}
Then, for a well-chosen angle $\theta$ its action is $R_Y(2\arcsin\sqrt{p})\ket{0} = \sqrt{1-p}\ket{0} + \sqrt{p}\ket{1}$, i.e. it deterministically loads a probability $p\in[0, 1]$ into the amplitude of $\ket{1}$. We will also make use of the shorthand: $Y_\theta := R_Y(2\arcsin\sqrt{\theta})$ throughout, particularly for the Boltzmann and weak measurement ancillas used in \hyperref[sec:algorithm]{Sec.~\ref*{sec:algorithm}}. 

A \emph{quantum circuit} is a horizontal-time diagram in which each wire carries one qubit (or a bundled register), boxes denote gates applied at that time slice, and a filled dot $\bullet$ on a wire indicates a control: the gate on the target wire fires conditionally on the control being $\ket{1}$ (or controlled on $\ket{0}$ if it is an open dot $\circ$). A half-filled dot \tikz[baseline=-0.5ex]{\node[halfctrl]{};} on a bundled register denotes a \emph{uniform} (SELECT-style) control, under which the target gate is applied in superposition to every computational-basis state of the control register (i.e. all possible bitstrings), each contributing its own branch to the state. All circuit figures in this thesis follow this convention and were created by the quantikz package \cite{kay2018tutorial}.

There is another key element in the quantum circuit: the read-out, which is performed by a \emph{computational-basis measurement}. This is a projective measurement \eqref{eq:post_meas_state} associated with the spectral decomposition of the $Z$-operators on each qubit. On a pure state $\ket{\psi} \in (\mathbb{C}^2)^{\otimes n}$ this returns a bitstring $x\in\{0, 1\}^n$ with probability $|\langle x|\psi\rangle|^2$ and collapses the register onto $\ket{x}$. \emph{Mid-circuit} measurements are also available on several types of quantum hardware, where the outcome can be discarded (yielding a non-unitary channel of Kraus form \eqref{eq:kraus}) or post-selected on a desired result.

The three primitives introduced so far, e.g. universal gate set, fresh ancilla qubits prepared in $\ket{0}$, and measurements, already suffice to realise \emph{any} CPTP map $\Phi: \mathcal{D}(\mathcal{H}^{\otimes n}) \to \mathcal{D}(\mathcal{H}^{\otimes n})$. By Stinespring's dilation theorem (see \hyperref[sec:OQS]{Sec.~\ref*{sec:OQS}} and \cite{bible}), every quantum channel can be lifted to a unitary on a larger Hilbert space. More concretely, there exist an $m$-qubit ancilla register $A$ and a joint unitary $U$ on $n + m$ qubits such that, 
$$
\Phi(\rho) = \tr_A\!\bigl[\,U\,\bigl(\rho \otimes \ket{0}\bra{0}^{\otimes m}\bigr)\,U^\dagger\,\bigr],
$$
with the partial trace implemented either by physically discarding the ancilla or, equivalently, by measuring it in the computational basis and forgetting the outcome. The resulting average over unread outcomes reproduces the Kraus form \eqref{eq:kraus}. Since single-qubit unitaries together with \textsc{CNOT} decomposes $U$ exactly, a circuit implementing $\Phi$ to any desired accuracy requires only the three primitives above.

For any work about quantum algorithms and their realizations we cannot forget about the different timsecales in which a quantum and a classical computer operate. The gates above are physical processes on real hardware, and their wall-clock duration is a hard constant in any complexity statement. On state-of-the-art \emph{superconducting} processors, single- and two-qubit gates take $10-100\, \text{ns}$ (e.g. IBM Heron r3 pushing median two-qubit fidelity above $99.9\%$ \cite{abughanem2025ibm}). On a \emph{trapped-ion} hardware pursued by QVLS in Hannover, the most recent fully chip-integrated two-qubit register on $^{9}\text{Be}^{+}$ runs microwave-driven Mølmer–Sørensen ($R_{XX}$) gates of duration around $1\,\text{ms}$ with a composite process fidelity of $96.6\%$ \cite{pulido2024arbitrary}. By contrast, a logic gate on a classical CMOS hardware switches in $\sim 10\,\text{ps}$. The quantum-classical wall-clock gap is therefore three to eight orders of magnitude per gate, a constant overhead that any claim of ``quantum advantage" must overcome.

\begin{figure}[ht]
  \centering
  \begin{quantikz}[row sep=0.7cm, column sep=0.45cm]
  %
  % Ancilla qubit 0 (MSB)
    \lstick{$\ket{0}$}
      & \gate[4, nwires={2}]{\textsc{Prep}_f}
      & \qw
      & \qw
      & \qw
      & \ctrl{4}
      & \gate[4, nwires={2}]{\text{QFT}}
      &
      &
      & \meter{} \\
  %
  % Ancilla qubit 1
    \lstick{$\ket{0}$}
      &
      & \qw
      & \qw
      & \ctrl{3}
      & \qw
      &
      &
      &
      & \meter{} \\
  %
  % Vdots row (invisible wire, just for the ellipsis)
    \lstick{$\vdots$} \setwiretype{n}
      &
      & \vdots
      &
      &
      &
      &
      &
      &
      & \vdots \\
  %
  % Ancilla qubit r-1 (LSB)
    \lstick{$\ket{0}$}
      &
      & \ctrl{1}
      & \qw
      & \qw
      & \qw
      &
      &
      &
      & \meter{} \\
  %
  % System register (bundled)
    \lstick{$\ket{\psi}$}
      & \qwbundle{}
      & \gate{U}
      & \ \ldots\ \qw
      & \gate{U^{2^{r-2}}}
      & \gate{U^{2^{r-1}}}
      & \qw
      & \qw
      & \qw
      & \rstick{$\ket{\psi}$} \qw
  \end{quantikz}
  \caption{Phase estimation
  circuit $\mathcal{C}_f$ with $r$ ancilla qubits and a bundled system register $\ket{\psi}$. The
  state-preparation block $\textsc{Prep}_f$ loads the time-domain window
  $\ket{f} = \sum_{\bar t\in S_{t_0}} f(\bar t)\ket{\bar t}$ on the ancilla register. (i) The uniform window $f\equiv 1/\sqrt{N}$ is achieved by Hadamard gates on each qubit: this is the standard QPE. (ii) A Gaussian $f(\bar{t}) \propto \ee^{-\bar{t}^2 \sigma^2}$ leads to a more smeared phase estimate (albeit controlled) with the advantage of requiring less time evolution.}
  \label{circ:qpe}
\end{figure}

\smallskip
\section{Phase Estimation} One of the key subroutines classical and quantum Gibbs sampling algorithms rest on is the phase (or energy) estimation. This allows us to label the system with an energy register on which a transition-rate $\gamma(\omega)$, satisfying the KMS condition \eqref{eq:KMS-condition}, can act in superposition. It was first introduced as \emph{Quantum phase estimation} (QPE) by Kitaev in \cite{kitaev1997quantum,cleve1998quantum}, and it was quickly applied to Hamiltonian eigenvalue problems \cite{abrams1999quantum}. It has underpinned not only recent developments in Gibbs sampling, but also the seminal quantum Metropolis algorithm of \cite{temme2011quantum}. As QPE is an already established textbook material, we only highlight the similarities and differences of it to the phase estimation we use later on: a one-parameter generalization of QPE, that can be used to compute \emph{Operator Fourier Transforms} (OFT).

To make this shared structure visible, see the circuit in \hyperref[circ:qpe]{Fig.~\ref*{circ:qpe}}. Let the Hamiltonian $H$ have Bohr frequencies $\nu \in B_H$, fix the time unit $t_0>0$, and an ancilla register of $r$ qubits with time grid $S_{t_0} = \{0, t_0, 2t_0, \dots, (N{-}1)t_0\}$ and a dual frequency grid $S_{\omega_0}$ of spacing $\omega_0$ satisfying $t_0\,\omega_0 = 2\pi/N$. The three stages of the circuit are
\begin{enumerate}
    \item $\textbf{Prep}_f$ on the ancilla, preparing the filter state $\ket{f} := \sum_{\bar t \in S_{t_0}} f(\bar t)\,\ket{\bar t}$ for a chosen window $f : S_{t_0} \to \mathbb{C}$ with $\sum_{\bar t}|f(\bar t)|^2 = 1$.
    \item \textbf{Controlled Hamiltonian evolution,} $\sum_{\bar t}\ket{\bar t}\!\bra{\bar t}\otimes U^{\bar t/t_0}$ with $U = \ee^{+\ii H t_0}$ (plus sign, so the readout register labels $\nu$ directly in our later applications \eqref{eq:oft}). This is achieved by the cascade of evolutions $U, U^2, \ldots, U^{2^{r-1}}$.
    \item \textbf{Quantum Fourier Transform (QFT)}, $\ket{\bar t}\mapsto\frac{1}{\sqrt N}\sum_{\bar\omega}\ee^{-\ii\bar t\bar\omega}\ket{\bar\omega}$, maps the time register to the frequency register $\ket{\bar \omega}$ on $S_{\omega_0}$.
\end{enumerate}
On an eigenstate $H\ket{\psi_\nu} = \nu\ket{\psi_\nu}$ the three stages together act as
\begin{equation} \label{eq:qpe-unified}
    \mathcal{C}_f \,\bigl(\ket{0}^{\otimes r} \otimes \ket{\psi_\nu}\bigr)
\;=\; \sum_{\bar\omega \in S_{\omega_0}} \hat f(\bar\omega - \nu)\,\ket{\bar\omega} \otimes \ket{\psi_\nu},
\end{equation}
where $\hat f(\bar\omega) := \frac{1}{\sqrt{N}} \sum_{\bar t \in S_{t_0}} \ee^{-\ii \bar\omega \bar t} f(\bar t)$ is the discrete Fourier transform of $f(\bar{t})$. Clearly, the choice of the filter function $f$ can significantly influence the result of the phase estimation. Let us consider
the following two cases.

\medskip 
\noindent\textbf{(i) Uniform $f$ -- standard QPE.} To weight each label uniformly by $f(\bar t) = 1/\sqrt{N}$, the $\textsc{Prep}_f$ becomes a simultaneous application of \emph{Hadamard} gates on each of the time register ancillas. A Hadamard gate is defined as, 
$$
\mathsf{H} \;:=\; \tfrac{1}{\sqrt{2}} \begin{pmatrix} 1 & 1 \\ 1 & -1 \end{pmatrix},
$$
and it maps each ancilla $\ket 0 \mapsto (\ket 0 + \ket 1)/\sqrt 2$. The collective result is the uniform time register state $\frac{1}{\sqrt N}\sum_{\bar t}\ket{\bar t}$. The output kernel then becomes
\begin{equation}
    \hat f(\bar\omega - \nu) \;=\; \frac{1}{N}\,\frac{1 - \ee^{-\ii(\bar\omega-\nu)Nt_0}}{1 - \ee^{-\ii(\bar\omega-\nu)t_0}},
\end{equation}
or the squared form that we mostly see in QPE literature, 
\begin{equation} \label{eq:QPE-prob-kernel}
    \bigl|\hat f(\bar\omega-\nu)\bigr|^2 \;=\; \frac{\sin^2\!\bigl((\bar\omega-\nu)Nt_0/2\bigr)}{N^2\sin^2\!\bigl((\bar\omega-\nu)t_0/2\bigr)}
\end{equation}
i.e. the Dirichlet kernel, peaked at $\bar{\omega} = \nu$ with width $\sim 1 / (Nt_0)$ and polynomially decaying \cite[Fig. 8]{chen2023quantum}. Sampling the ancilla register returns an $r$-bit approximation of the scaled eigenphase $\varphi_\nu = \nu t_0/(2\pi)$. In order to resolve it up to $n$ accurate bits with failure probability $\delta$ requires, 
\begin{equation} \label{eq:qpe-ancilla-count}
    r \;=\; n + \left\lceil \log\!\left(2 + \frac{1}{2\delta}\right) \right\rceil
\end{equation}
ancilla qubits \cite[Sec.~5.2.1]{bible}, for a total controlled Hamiltonian evolution time $\sum_{k=0}^{r-1} 2^k t_0 = O(1/\varepsilon)$ at resolution $\varepsilon = 2^{-n}$.

\medskip
\noindent\textbf{(ii) Smoothly decaying $f$.} Replacing the uniform filter by a rapidly decaying one, e.g. a discretized Gaussian \eqref{eq:gaussian-filter-t}, or a Gevrey function \cite[text]{ding2025efficient}, turns $\hat{f}$ into a smooth kernel of tuneable width, smearing each eigenvalue into a packet of the window's width ($2\sigma$ for the Gaussian \eqref{eq:smoothed-jump}).

Resolution can be similarly computed: for a Gaussian width $2\sigma$, truncating the discretized window to relative $\ell_2$-tail $\varepsilon$ requires time range $Nt_0 \gtrsim \sigma^{-1}\sqrt{2\log(1/\varepsilon)}$ and frequency spacing $\omega_0 = 2\pi/(Nt_0) \lesssim \sigma$, giving
\begin{equation} \label{eq:gaussian-ancilla-count}
    r \;\geq\; \left\lceil \log\!\frac{2\pi}{\sigma t_0} \;+\; \tfrac{1}{2}\log\ln\frac{1}{\varepsilon} \right\rceil,
\end{equation}
and total evolution time $Nt_0 = \mathcal{O}\!\bigl(\sigma^{-1}\sqrt{\log(1/\varepsilon)}\bigr)$, which is linear in the inverse resolution up to $\sqrt{\log(1/\varepsilon)}$ (Gevrey windows replace the square root by $\mathrm{polylog}(1/\varepsilon)$ \cite[text]{ding2025efficient}).

Both case (i) and (ii) obey $T\cdot\Delta\omega \gtrsim 1$ with $T = Nt_0$ the total controlled Hamiltonian evolution time and $\Delta \omega$ the width of the frequency kernel $\hat{f}$. However, they have different qualities that are worth highlighting to understand what we gain (or lose) by choosing the latter one as our subroutine.

(i) Standard QPE is optimal for the classical-output task: input the eigenstate $\ket{\psi_\nu}$, return an $r$-bit estimate $\hat{\nu}$ of the eigenphase. For this task a matching quantum query lower bound holds: any algorithm that accesses $H$ only through controlled $U = \ee^{\pm \ii H t_0}$ and returns $\varepsilon$-accurate estimate with constant success probability must use total controlled-evolution time $T = \Omega(1/\varepsilon)$ \cite{giovannetti2006quantum}. This is the \emph{Heisenberg limit} $\Delta E \cdot T \sim 1$ of quantum metrology and is a genuine quantum speed-up over the ``standard quantum limit" $T\sim1/\varepsilon^2$ obtainable from $M$ independent depth $t_0$ measurements followed by classical averaging. QPE's $U^1, U^2, \ldots, U^{2^{r-1}}$ cascade attains $T = \mathcal{O}(1 / \varepsilon)$ \cite{bible} and thus saturates this lower bound. Although its Dirichlet kernel has polynomially-decaying tails \eqref{eq:QPE-prob-kernel}, and is \emph{not} a minimizer of frequency spread in the sense of case (ii), the median estimation error still decays as $1 / T$, which is all the estimation task demands.

(ii) The KMS detailed balanced Lindbladians of \hyperref[sec:kms-lindbladian]{Sec.~\ref*{sec:kms-lindbladian}} use Gaussian filtered phase estimates: by sandwiching a jump operator $A$ between filtered phases estimations, we can decompose jump operators into their Fourier components $\hat{A}(\bar{\omega})$ (\eqref{eq:oft}) and allow us to weigh them with the KMS transition rates $\gamma(\bar{\omega})$. The key metric is no longer estimation error but the $L^2$ width of $\bar{f}$. The governing bound then, is the classical Heisenberg inequality: for any $f\in L^2(\mathbb{R})$ of unit norm \cite{folland1997uncertainty}, 
$$
\sigma_t\,\sigma_\omega \;\geq\; \tfrac{1}{2},
$$
Under the circuit identification $\sigma_t \leftrightarrow T$, $\sigma_\omega \leftrightarrow \Delta\omega$ this becomes a lower bound on controlled-evolution time that no window can evade. The Gaussian \eqref{eq:gaussian-filter-t} uniquely saturates the above equality, while the Dirichlet does not. But this is a result coming from harmonic analysis, independent of any quantum metrological argument, and is not to be confused with the Heisenberg limit of (i). Although, recent results found that the Heisenberg-limit with shallower circuits are achievable \cite{ding2023even} (with just an ancilla combined with statistical phase-estimation protocols), these cannot be used in our setup since they collapse the ancilla to a classical sample, destroying the superposition of $\ket{\bar{\omega}}$.

Since we will work with OFT-based Lindbladian generators, the inherent obstacle that arises is the classical Fourier bound, not Heisenberg limit. As the Bohr spectrum $B_H$ of a non-trivial many-body Hamiltonian is exponentially dense in $n$, an attempt to exactly resolve Bohr frequencies would require $T$ exponentially growing with $n$, and thus no choice of filter function can escape this. The smoothened generators of \hyperref[sec:kms-lindbladian]{Sec.~\ref*{sec:kms-lindbladian}} take this impossibility as the starting point: they fix a controlled resolution width $\sigma$ and arrange the dissipator and Lamb-Shift such that KMS detailed balance survives even this finite-resolution smearing.

\smallskip
\label{sec:prelim-trotter}
\section{Trotterization}
\begin{figure}[ht]
  \centering
  \resizebox{\textwidth}{!}{%
  \begin{quantikz}[
    transparent, row sep=0.6cm, column sep=0.4cm,
    /tikz/operator/.append style={inner sep=3pt, minimum height=6mm, font=\small}
  ]
    \lstick{$q_1$}
  % A-block L (active)
      & \gate[2]{R_{X_1X_2}}
      & \gate[2]{R_{Y_1Y_2}}
      & \gate[2]{R_{Z_1Z_2}}
  % B-block L (q_1 idle)
      & \qw & \qw & \qw
  % C-block (boundary, q_1 leads the [3]-wide gate; R^2 = R applied twice = rotate by 2 dt)
      & \gate[3, label style={yshift=0.4cm}]{R_{X_3X_1}^{2}}
      & \gate[3, label style={yshift=0.4cm}]{R_{Y_3Y_1}^{2}}
      & \gate[3, label style={yshift=0.4cm}]{R_{Z_3Z_1}^{2}}
  % B-block R (q_1 idle)
      & \qw & \qw & \qw
  % A-block R (active)
      & \gate[2]{R_{X_1X_2}}
      & \gate[2]{R_{Y_1Y_2}}
      & \gate[2]{R_{Z_1Z_2}}
      & \qw \\
  %
    \lstick{$q_2$}
  % A-block L (empty: consumed by q_1's gate[2])
      & & &
  % B-block L (active, q_2 leads)
      & \gate[2]{R_{X_2X_3}}
      & \gate[2]{R_{Y_2Y_3}}
      & \gate[2]{R_{Z_2Z_3}}
  % C-block: q_2 passes through unaffected
      & \dashedthrough & \dashedthrough & \dashedthrough
  % B-block R (active, q_2 leads)
      & \gate[2]{R_{X_2X_3}}
      & \gate[2]{R_{Y_2Y_3}}
      & \gate[2]{R_{Z_2Z_3}}
  % A-block R (empty: consumed by q_1's gate[2])
      & & &
      & \qw \\
  %
    \lstick{$q_3$}
  % A-block L (q_3 idle)
      & \qw & \qw & \qw
  % B-block L (empty: consumed by q_2's gate[2])
      & & &
  % C-block (empty: consumed by q_1's gate[3])
      & & &
  % B-block R (empty: consumed by q_2's gate[2])
      & & &
  % A-block R (q_3 idle)
      & \qw & \qw & \qw
      & \qw
  \end{quantikz}%
  }
  \caption{One second-order Strang Trotter step for the $n{=}3$ periodic
  Heisenberg chain. The sum of Pauli strings in $H_\mathrm{Heis}$ turn into a product of 2-qubit unitary Pauli rotations,
  $R_{\sigma_j\sigma_k}(\theta) \equiv \ee^{-\ii\theta\,\sigma_j\sigma_k/2}$, where $\theta$ encodes coupling strengths of the Hamiltonian and the required time step.
  The boundary bond $(3,1)$ closes the ring and $q_2$ passes through unaffected
  (dashed wire).}
  \label{circ:trotter-strang}
\end{figure}
Up to this point we have concealed a fact that Hamiltonian time evolutions $U(t) = \ee^{\pm \ii H t}$ often cannot be exactly implemented by a quantum computer, and they have to be approximated by a product formula $S_p(t)$ instead. This is also true for the local (i.e. each term acts on a bounded number of qubits) and non-commuting Hamiltonians we care about: there is no known shallow circuit that would avoid the use of Trotterization. Here, we will give only the necessary details that are relevant for the thesis, based on the work of Childs et al. \cite{childs2021theory}.

Let us write the Hamiltonian as a sum of $\Gamma$ non-commuting summands
$$
H \;=\; \sum_{\ell=1}^{\Gamma} H_\ell,
$$
where each $H_\ell$ is either a single Pauli string or a commuting family of such strings whose exponential $\ee^{-\ii t H_\ell}$ is easy to implement (typically because all terms in $H_\ell$ share support disjointly or pairwise commute). A $p$-th order product formula $S_p(t)$ is an ordered product of single-summand evolutions $\ee^{-\ii t H_\ell}$ that together approximates the full Hamiltonian evolution $\ee^{-\ii t H}$ up to error $\mathcal{O}(t^{p + 1})$ for small $t$ \cite{childs2021theory}. The simplest one for reference, with $p=1$ is called the Lie-Trotter formula, 
\begin{equation}
    S_1(t) \;:=\; \ee^{-\ii t H_\Gamma}\cdots \ee^{-\ii t H_1},
\end{equation}
but the one we mainly use for the numerical analysis is the Strang splitting with $p=2$:
\begin{equation}\label{eq:strang}
    S_2(t) \;:=\; \Bigl(\prod_{\ell=1}^{\Gamma} \ee^{-\ii (t/2) H_\ell}\Bigr)\Bigl(\prod_{k=\Gamma}^{1} \ee^{-\ii (t/2) H_\ell}\Bigr).
\end{equation}
The former leads to errors $\|S_1(t) - \ee^{-\ii H t}\| = \mathcal{O}(t^2)$ while the latter $\|S_2(t) - \ee^{-\ii H t}\| = \mathcal{O}(t^3)$. There is also a structural difference, namely that Strang splitting is palindromic, i.e. $S_2(t)^\dagger = S_2(-t)$ which is worth preserving to avoid additional deviations from detailed balance \hyperref[rem:palindromic]{Remark~\ref*{rem:palindromic}}. The circuit realization of $S_2$ for $n=3$ qubit, periodic Heisenberg chain $H_\mathrm{Heis}$ \eqref{eq:H_heis} is given in \hyperref[circ:trotter-strang]{Fig.~\ref*{circ:trotter-strang}}. Since each $H_k$ is Hermitian by construction, each factor $\ee^{-\ii t H_\ell}$ is unitary, so $S_p(t)$ and its iterate $S_p^{(M)}(t) := S_p(t/M)^M$ are unitary for every choice of $p$ and $M$. Notably, we see that an $n$ qubit unitary $U$ is broken down into 2-qubit gates. This is important for implementation since they can be easily transpiled to native 1-qubit and 2-qubit gates that a quantum hardware uses. To motivate it even more: transpiling a general $n$ qubit unitary to native gates in a brute-force manner is currently only possible up to $n = 9$ qubits \cite{rakyta2024highly}.

Higher-order Suzuki formulas $S_{2r}$ can be built recursively from $S_2$ by the fractal construction $S_{2r}(t) = S_{2r-2}(u_r t)^2\,S_{2r-2}((1-4u_r)t)\,S_{2r-2}(u_r t)^2$ with $u_r = 1/(4 - 4^{1/(2r-1)})$, leading to errors of $\mathcal{O}(t^{2r+1})$ \cite{childs2021theory}.

A naive Taylor truncation only controls the error through $\|H_k\|$, ignoring the commutativity of the summands. The key result of \cite{childs2021theory} is that the error of any $p$-th order formula is in fact controlled by nested commutators of the summands. For some real $t$ the error is given by \cite[Theorem 6]{childs2021theory}, 
\begin{equation} \label{eq:childs-comm-bound}
    \begin{split}
        &\bigl\|S_p(t) - \ee^{-\ii H t}\bigr\|
        \;=\; \mathcal{O}\!\bigl(\tilde\alpha_{\mathrm{comm}}\,|t|^{p+1}\bigr) \\ 
        \text{with}\qquad &\tilde\alpha_{\mathrm{comm}}
        \;:=\; \sum_{\ell_1,\dots,\ell_{p+1}=1}^{\Gamma}
        \bigl\|\,[H_{\ell_{p+1}},\,\cdots\,[H_{\ell_2},H_{\ell_1}]\cdots]\,\bigr\|,
    \end{split}
\end{equation}
where the constant $\tilde{\alpha}_\mathrm{comm}$ is universal (independent of $\Gamma$, $\|H_\ell\|$, and the ordering fixed by the particular formula). More concretely, for the second-order Strang splitting,
\begin{equation} \label{eq:strang-explicit-error}
        \bigl\|S_2(t) - \ee^{-\ii H t}\bigr\|
    \;\leq\; \sum_{\ell_1 < \ell_2}\!\left(
    \frac{t^3}{12}\bigl\|[H_{\ell_2},[H_{\ell_2},H_{\ell_1}]]\bigr\|
    \;+\; \frac{t^3}{24}\bigl\|[H_{\ell_1},[H_{\ell_1},H_{\ell_2}]]\bigr\|
    \right)
    \;+\; \mathcal{O}(t^4).
\end{equation}
For longer time evolutions the target evolution $\ee^{-\ii H t}$ is split into $M$ Trotter steps of size $\delta t = t / M $. Then, a total error of $\varepsilon$ can be achieved if $M$ satisfies, 
\begin{equation} \label{eq:trott-num-bound}
    M \;\geq\; \Omega\!\left(\frac{\tilde\alpha_{\mathrm{comm}}^{1/p}\;t^{\,1+1/p}}{\varepsilon^{\,1/p}}\right).
\end{equation}
For $p=1$ this is the $\mathcal{O}(\tilde\alpha_{\mathrm{comm}}\, t^2/\varepsilon)$ worst-case count, and for Strang ($p = 2$) the $\mathcal{O}(\sqrt{\tilde\alpha_{\mathrm{comm}}}\, t^{3/2}/\sqrt{\varepsilon})$ bound is what is used for our specialized results in the Trotter error analysis of [REF] \todo{ref to trotter}

One of the Hamiltonians we use for the numerical analysis (the same we have a sketched Trotter step for in \hyperref[circ:trotter-strang]{Fig.~\ref*{circ:trotter-strang}}) is the periodic Heisenberg chain on $n$ qubits, 
\begin{equation} \label{eq:H_heis}
    H_{\mathrm{Heis}}
\;=\; J \sum_{i=1}^{n} \bigl( X_i X_{i+1} + Y_i Y_{i+1} + Z_i Z_{i+1} \bigr)
\;+\; \sum_{i=1}^{n} \zeta_i Z_i,
\end{equation}
with coupling $J$ (if positive - antiferromagnetic, if negative - ferromagnetic) and disordering external fields $\zeta_i \in [-\zeta,\zeta]$. 

In the following, we give more details about the commutant constant $\tilde{\alpha}_\mathrm{comm}$ for the family of \emph{geometrically local} Hamiltonians -- the family that the Heisenberg model \eqref{eq:H_heis} is also part of. These Hamiltonians can be written as
$$
 H \;=\; \sum_c h_c, \qquad \|h_c\| \le h_*, 
$$
where each local term $h_c$ is supported on at most $k$ sites, with maximum overlapping supports of degree $d$, making the Hamiltonian $H$ $(k, d)-$\emph{local} with $k, d = \mathcal{O}(1)$. Such an $H$ admits a partition into $\Gamma = \mathcal{O}(1)$ commuting classes, i.e. a sum of pair-wise commuting local terms. Expanding the $(p+1)$ nested commutator $[H_{\ell_{p+1}}, \cdots [H_{\ell_2}, H_{\ell_1}]]$ term-by-term over local terms, we obtain a sum over ordered $(p+1)$-tuples of bonds $(c_1, \dot, c_{p+1})$. The only surviving terms in this sum are the ones that have \emph{connected} bonds, i.e. where each $c_{i+1}$ overlaps with at least one of $c_1, \dots, c_i$, since the disjoint terms commute with each other. Therefore, an estimating $\tilde{\alpha}_\mathrm{comm}$ for $(k,d)$-local Hamiltonians reduces to the problem of counting these connected bonds. The number of connected $(p+1)$-tuples anchored at a pivot bond $c_1$ is at most $(2d)^p$, and there are $\Theta(n)$ choices of pivot bonds, giving
\begin{equation} \label{eq:alpha-comm}
\tilde\alpha^{(p)}_{\mathrm{comm}} \;\leq\; C_{d,p}\,n\,h_*^{p+1} \;=\; \mathcal{O}(n).
\end{equation}
Here, $C_{d,p} \le \Gamma^{p+1}\,(2d)^p$ is a combinatorial constant depending only on the bond-graph degree $d$ and the order $p$. The counting argument for this is present in \cite[Sec. V B]{childs2021theory} and goes as follows. The outer commuting-class grouping fixes $\Gamma = \mathcal{O}(1)$ and selects which classes contribute (e.g. even or odd sites), while inner bounded-degree bond counting yields the $\Theta(n)\cdot(2d)^p$ bound. It is the inner bond layer that carries the geometric locality assumption. For instance, an all-to-all interaction that has $\Gamma = 2$ summands but unbounded bond degree $d$, yields $\mathcal{O}(n^3)$ commutator constant rather than $\mathcal{O}(n)$ \cite[Sec. IV C]{childs2021theory}.

The 1D (periodic) Heisenberg chain \eqref{eq:H_heis} is a canonical case for $(k,d)$-local Hamiltonian. Its terms can be grouped into a couple of commuting classes: odd and even sites in case $n$ is even (or into three classes if $n$ is odd due to the boundary.) As the interaction degree on the chain is $d=2$, and the number of commuting classes $\Gamma = 2$ (or $3$), it is possible to find some estimate for the combinatorial factor in \eqref{eq:alpha-comm}. Although, with numerical heuristics \cite{childs2021theory}, we can get that constant value, for us the scaling is already sufficient, i.e. for $p=2$: $\tilde\alpha^{(2)}_{\mathrm{comm}}(H_{\mathrm{Heis}}) = \mathcal{O}(|J|^3 n)$. The same principle applies to the periodic 1D transverse-field Ising model and the 2D Heisenberg model on a square lattice, again with $\tilde\alpha^{(2)}_{\mathrm{comm}} = \mathcal{O}(n)$ scaling and constant prefactors that change only with the bond-interaction degree $d$.

\smallskip
\section{Linear Combination of Unitaries}\label{sec:prelim-LCU}

\begin{figure}[ht]
  \centering
  \begin{quantikz}[row sep=0.75cm, column sep=0.55cm]
    \lstick{$\ket{0^b}$}
      & \qwbundle{}
      & \gate{\text{PREP}}
      & \ctrl{1}
      & \gate{\text{PREP}^\dagger}
      & \rstick{$\bra{0^b}$} \qw
      \\
    \lstick{$\ket{\psi}$}
      & \qwbundle{}
      & \qw
      & \gate{U(t)}
      & \qw
      & \rstick{$\displaystyle\frac{1}{\|f\|_1}\sum_t f(t)\,U(t)\,\ket{\psi}$} \qw
  \end{quantikz}
  \caption{LCU block encoding $U_A$
  of $A = \sum_t f(t)\,U(t)$: \textsc{Prep} writes the filter state
  $\ket{f} = \frac{1}{\sqrt{\|f\|_1}}\sum_t \sqrt{f(t)}\ket{t}$, \textsc{Select}
  applies $U(t)$ controlled on the ancilla, and \textsc{Prep}$^\dagger$
  uncomputes. Post-selecting the ancilla on $\ket{0}$ leaves
  $A\ket{\psi}/\|f\|_1$ on the system, with success probability
  $\|A\ket{\psi}\|^2/\|f\|_1^2$.}
  \label{circ:lcu}
\end{figure}

Another key subroutine we will make use of is the linear combination of unitaries (LCU), along with the transformation of these block encoded operators introduced in the subsequent section. For a comprehensive introduction, please see \cite{childs2012hamiltonian,berry2015simulating}.

Fix a unitary family $\{U(t)\}_{t\in S_t}$ indexed by a discrete grid $S_t$ and a scalar weight $f : S \to \mathbb{R}_{\geq 0}$ with finite $\ell_1$ norm $\|f\|_1 = \sum_t f(t)$. The target operator
$$
A \;=\; \sum_{t \in S} f(t)\,U(t)
$$
is in general Hermitian but not unitary if $U(-t) = U(t)^\dagger$ and $f$ is real and symmetric. Its operator norm $\|A\|$ is at most $\|f\|_1$ by the triangle inequality. The LCU circuit in \hyperref[circ:lcu]{Fig.~\ref*{circ:lcu}} realises $A$ up to a known rescaling using an ancilla register of $b = \lceil \log|S_t|\rceil$ qubits and the following three ingredients: 
\begin{enumerate}
    \item \textsc{Prep} on the ancilla, prepares the filter state $\ket{f} := \frac{1}{\sqrt{\|f\|_1}}\sum_{t \in S} \sqrt{f(t)}\,\ket{t}$
    \item \textsc{Select} picks out each $U(t)$ for a given $t$, i.e. $\sum_{t}\ket{t}\!\bra{t}\otimes U(t)$
    \item $\textsc{Prep}^\dagger$ uncomputes the ancilla register.
\end{enumerate}
Writing the composite circuit as $U_A := (\textsc{Prep}^\dagger \otimes \idm)\, \textsc{Select}\, (\textsc{Prep} \otimes \idm)$, a direct calculation shows that its $\ket{0^b}$-$\ket{0^b}$ matrix element (the top left corner) is exactly, 
\begin{equation} \label{eq:lcu}
    \bra{0^b}\,U_A\,\ket{0^b} \;=\; \frac{A}{\|f\|_1},
\end{equation}
i.e. $U_A$ is a block encoding of $A$ with subnormalisation $Z_A := \|f\|_1$ \cite{low2019hamiltonian}. The subnormalisation is the price for packing a non-unitary $A$ inside a unitary: $U_A$ acts as $A / Z_A$ on the block projector $\ket{0^b}\!\bra{0^b}$ and as an unspecified completion on its orthogonal subspace, and all complexity statements about simulation times and success probabilities scale with $Z_A$.

This approach can be extended to cases where $f$ assumes signed or complex values. There are different solutions, but the one we end up using later is the following: split $\textsc{Prep}$ into an asymmetric left/right pair $\textsc{Prep}'$ / $\textsc{Prep}$. The prior writes the signed / phased amplitude $\sum_t \mathrm{sgn}(f(t))\sqrt{|f(t)|/\|f\|_1}\,\ket{t}$ on the left, and the latter writes the magnitude-only amplitude $\sum_t \sqrt{|f(t)|/\|f\|_1}\,\ket{t}$ on the right, so that the two amplitudes multiply to $f(t)/\|f\|_1$. \hyperref[circ:U_B_a]{Fig.~\ref*{circ:U_B_a}} adopts this with $\mathrm{sgn}(\cdot)$ is a unit-modulus phase in the complex case.

Applying $U_A$ on $\ket{0^b}\otimes\ket{\psi}$ and \emph{post-selecting} the ancilla on $\ket{0^b}$ leaves, up to normalisation, the state $A \ket{\psi} / Z_A$ on the system, with success probability $\|A\ket{\psi}\|^2/Z_A^2$. When $Z_A \leq 1$, the block encoding can be taken as a deterministic implementation of $A$. In the case when $Z_A \geq 1$, however, the post-selection succeeds only with $\mathcal{O}(1 / Z_A^2)$, in which case LCU is \emph{not} deterministic. For bare applications of LCU, there are already remedies for this, e.g. a repeat-until-success procedure with $\mathcal{O}(Z_A^2)$ of repetitions, or the so-called oblivious amplitude amplifications which involves additional reflections. But in our thesis, an indeterministic LCU can only happen when we create a block encoding for a Hermitian operator $B$ and then directly transform it to a coherent evolution $\ee^{-\ii B t}$. This transformation goes hand in hand with block-encodings, which is the next subroutine we will introduce. There, it becomes clear what role the subnormalization factor $Z_A$ plays in the algorithm's complexity.

\smallskip
\section{Quantum Signal Processing} \label{sec:prelim-QSP}

\begin{figure}[ht]
  \centering
  \resizebox{1.1\textwidth}{!}{%
  \begin{quantikz}[row sep=0.7cm, column sep=0.5cm]
  %
  % ==================================================================
  % Row 1: QSP ancilla.
  % Time order in this figure (left to right = forward in time):
  %   col 3:   opening rotation L_0 = R_MW(theta_0, phi_0, lambda_0)
  %   col 4-5: representative open-controlled slot, repeated d times
  %            (j = 1, ..., d, applied first in time after L_0)
  %   col 6-7: representative closed-controlled slot, repeated d times
  %            (j = d+1, ..., 2d, applied last in time)
  % j INCREASES left-to-right in this figure (matching MW2024 Eq. 42 with
  % the index convention "j = 1 is the first layer applied after L_0").
  % ==================================================================
    \lstick{$\ket{0}_\mathrm{QSP}$}
      & \qw
      & \gate{R(\theta_0,\phi_0,\lambda_0)}
      & \octrl{1}
        \gategroup[3,steps=2,style={dashed,rounded corners,fill=blue!6,inner sep=5pt},
                   background,label style={label position=above,yshift=0.25cm}]
          {\scriptsize Repeat $d$ times}
      & \gate{R(\theta_j,\phi_j,0)}
      & \qw
      & \ctrl{1}
        \gategroup[3,steps=2,style={dashed,rounded corners,fill=black!6,inner sep=5pt},
                   background,label style={label position=above,yshift=0.25cm}]
          {\scriptsize Repeat $d$ times}
      & \gate{R(\theta_{d+j},\phi_{d+j},0)}
      & \rstick{$\bra{0}_\mathrm{QSP}$}\qw
      \\
  %
  % ==================================================================
  % Row 2: Block-encoding ancilla (bundled).
  % Multi-row gates: \gate[2]{W} for the open-ctrl block (col 4),
  %                  \gate[2]{W^\dagger} for the closed-ctrl block (col 6).
  % NO uncontrolled tail in Form B.
  % PREP/PREP^dagger live INSIDE U_A inside each W and are not drawn here.
  % ==================================================================
    \lstick{$\ket{0^b}$}
      & \qwbundle{}
      & \qw
      & \gate[2]{W}
      & \qw
      & \qw
      & \gate[2]{W^\dagger}
      & \qw
      & \rstick{$\bra{0^b}$}\qw
      \\
  %
  % ==================================================================
  % Row 3: System register (bundled).
  % Cells under the multi-row gates (col 4: W, col 6: W^dagger) are LEFT
  % EMPTY so quantikz draws those gates spanning rows 2-3.
  % ==================================================================
    \lstick{$\ket{\psi}$}
      & \qwbundle{}
      & \qw
      &
      & \qw
      & \qw
      &
      & \qw
      & \rstick{$\approx \ee^{-\ii\delta A}\ket{\psi}$}\qw
  \end{quantikz}
  }
  \caption{GQSP circuit approximating
    $\ee^{-\ii\delta A}$ as a Laurent polynomial on the post-selected $\ket{0}_\mathrm{QSP}\otimes\ket{0^b}$
    block from the walk operator $W = R_b\,U_A$. The dashed boxes are repeated $d$ times: one with $W$ the other with $W^\dagger$ (each followed by a two-parameter rotation $R(\theta_\cdot,\phi_\cdot,0)$), with indices $j = 1,\dots,d$ on the first block and $d+j$ on the second. In a gate-level implementation the
    $\textsc{Prep}/\textsc{Prep}^\dagger$ inside $U_A$ are pulled out of
    the controlled walks (adjacent $\textsc{Prep}^\dagger\textsc{Prep}=I$) to avoid unnecessary controlled applications.}
  \label{circ:gqsp}
\end{figure}

Without spoiling too much of the coming theoretical discussions, let us simply claim, that for a detailed balanced Lindbladian we need to apply a coherent evolution of the form $\ee^{-\ii\delta A}$ with respect to some Hermitian operator $A = A^\dagger$. There are certainly many ways to apply such a coherent part in a Lindbladian evolution, the one we have decided upon is conceptually and in circuit construction terms the simplest while operating at the same approximating error order as the other parts of \hyperref[alg:main]{Algorithm~\ref*{alg:main}}, namely as the weak measurement scheme, see the next section.

The LCU circuit of the previous section \hyperref[circ:lcu]{Fig.~\ref*{circ:lcu}} delivers a block encoding $U_A$ on a $b$-qubit register of some Hermitian $A$ on an $n$-qubit register, 
$$
\langle 0^b |\,U_A\,|0^b \rangle \;=\; A/\alpha, \qquad \|A/\alpha\| \leq 1,
$$
with $A = \sum_j \alpha_j U_j$, $\alpha = \sum_j \alpha_j$, preparation unitary $\textsc{Prep}\ket{0} = \sum_j \sqrt{\alpha_j/\alpha}\,\ket{j}$, selection unitary $\textsc{Select} = \sum_j \ket{j}\!\bra{j}\otimes U_j$, and $U_A = (\textsc{Prep}^\dagger\otimes I)\,\textsc{Select}\,(\textsc{Prep}\otimes I)$. 

It is often the case however, that we do not purely want to apply $A$ to the system but rather some function of it, e.g. the coherent evolution $\ee^{-\ii\delta A}$, with $A = A^\dagger$. \emph{ Quantum signal processing} (QSP) is the technique that bridges the gap: given a block encoding like the LCU, it post-selectively realises an arbitrary polynomial transformation of the encoded operator's $A$ spectrum. There are a miriad of ways to implement QSP for a thorough review see \cite{skelton2025hitchhiker}. We have decided upon the \emph{Generalized QSP} (GQSP) variant of \cite{motlagh2024generalized}, which lifts the parity restriction of the original QSP construction and works on the unit circle by using more general $\mathrm{SU}(2)$ rotations. This makes the natural class for Hamiltonian simulation targets $\ee^{-\ii\delta\alpha\cos\theta}$, the class of all complex Laurent polynomials $P$ directly accessible, as we explain below.

Define the walk operator $W$ as the unitary block encoding $U_A$ followed by a reflection $R_b$ around the all-zero state of the block encoding ancilla \cite[Cor. 8]{motlagh2024generalized}
\begin{equation} \label{qsp-walk}
    W \;:=\; R_b\,U_A, \qquad R_b \;:=\; 2\,\ket{0^b}\!\bra{0^b} - I.
\end{equation}
The primitives $\textsc{Prep}/\textsc{Select}/\textsc{Prep}^\dagger$ will always be defined as part of the unitary block encoding $U_A$ -- both \hyperref[circ:lcu]{Fig.~\ref*{circ:lcu}} and \hyperref[circ:gqsp]{Fig.~\ref*{circ:gqsp}} adopt this convention. For each eigenvector 
$\ket\lambda$ of $A$ with $A\ket\lambda = \lambda\ket\lambda$, the two-dimensional subspace containing $\ket{0^b}\ket\lambda$ is invariant under $W$. Hence, in a suitable orthonormal basis of that subspace, $W$ acts as the $\mathrm{SU}(2)$ matrix, 
$$
W_\lambda \;=\; \begin{pmatrix} \lambda/\alpha & \sqrt{1-(\lambda/\alpha)^2} \\[2pt] -\sqrt{1-(\lambda/\alpha)^2} & \lambda/\alpha \end{pmatrix}.
$$
The eigenvalues are given by $\ee^{\pm\ii\theta_\lambda}$, where $\cos\theta \;=\; \lambda/\alpha$ with $\lambda \in [-\alpha,\alpha]$ and $\theta \in [0,\pi]$. Any spectral transformation on the $\lambda$-eigenspace of $A$ thus reduces to approximating the target function at $z = \ee^{\ii\theta_\lambda}$ on the unit circle $\mathbb{T} = \{z\in\mathbb{C} : |z|=1\}$.

For an ordinary polynomial $P \in \mathbb{C}[z]$ with $\|P\|_{\infty,\mathbb{T}} := \max_{z\in\mathbb{T}}|P(z)| \leq 1$, and an integer $0 \leq k \leq \deg P$, the GQSP circuit \hyperref[circ:gqsp]{Fig.~\ref*{circ:gqsp}} applies a polynomial transformation onto the qubitized subspace by interleaving walk operators and single-qubit rotations defined by 
$$                      
  R(\theta, \phi, \lambda) = \begin{pmatrix} \ee^{\ii(\phi+\lambda)}\cos\theta &
  \ee^{\ii\phi}\sin\theta \\ \ee^{\ii\lambda}\sin\theta & -\cos\theta \end{pmatrix}                  
$$
with angles found by \cite[Alg. 1]{motlagh2024generalized}. Concretely:
\begin{enumerate}
    \item Start with an initial rotation $L_0 = R(\theta_0,\phi_0)\cdot\mathrm{diag}(\ee^{\ii\lambda_0}, 1)$ on the QSP ancilla, with $\lambda_0\in \mathbb{R}$ an absorbed phase parameter distinct from the spectral values $\lambda$.
    \item This is followed by a $\deg P$ controlled walk interleaved with the QSP ancilla rotations $R(\theta_j, \phi_j, 0)$, where $j = 1, \ldots, \deg P$.
\end{enumerate}
Importantly, the $\deg P$ controlled walks can be split into two groups of polarity in order to realise negative powers of polynomials. The first $\deg P - k$ slots ($j = 1, \ldots, \deg P - k$) as open-controlled forward walks, 
$$
\Lambda(W) \;:=\; \ket{0}\!\bra{0}_{\mathrm{QSP}}\otimes W \;+\; \ket{1}\!\bra{1}_{\mathrm{QSP}}\otimes I,
$$
and the last $k$ slots ($j = \deg P - k + 1, \ldots, \deg P$) are closed-controlled inverse walks, 
$$
\Lambda'(W) \;:=\; \ket{0}\!\bra{0}_{\mathrm{QSP}}\otimes I \;+\; \ket{1}\!\bra{1}_{\mathrm{QSP}}\otimes W^\dagger.
$$
Post-selecting both ancillas on $\ket{0}_{\mathrm{QSP}}\ket{0^b}$  leads to the operator \cite[Thm. 6]{motlagh2024generalized}, 
\begin{equation} \label{eq:qsp-poly}
    P'(W_\lambda) \;=\; W_\lambda^{-k}\,P(W_\lambda),
\end{equation}
which on the unit circle $z = \ee^{\ii\theta_\lambda} \in \mathbb{T}$  realises, for each qubitised two-dimensional subspace, the Laurent polynomial:
$$
P'(z) \;=\; z^{-k}\,P(z) \;=\; \sum_{m=0}^{\deg P} c_m\,z^{m - k} \;\in\; \mathbb{C}[z, z^{-1}].
$$
The monomials here, span the bilateral range $\{z^{-k}, z^{1-k}, \ldots, z^{\deg P - k}\}$. The two slot polarities are precisely what lifts the accessible class from ordinary into Laurent polynomials. The $k$ closed-controlled $\Lambda'(W)$ slots together contribute the negative-exponent prefactor $z^{-k}$, while the $\deg P - k$ open-controlled $\Lambda(W)$ slots, paired with the QSP-ancilla rotations $\{R(\theta_j,\phi_j,0)\}$ and the opening rotation $L_0$, build up the ordinary factor $P(z)$. The angle vector $(\theta_0, \phi_0, \lambda_0; \{(\theta_j, \phi_j)\}_{j=1}^{\deg P})$ depends only on $P$, while $k$ enters purely through the slot pattern. The two are chosen independently subject to $0 \leq k \leq \deg P$.

Letting $(P, k)$ range over the constraint set $\|P\|_{\infty,\mathbb{T}} \leq 1$, $0 \leq k \leq \deg P$, GQSP allows access to the full class
$$
\mathcal{P}_{\text{GQSP}} \;=\; \bigl\{P' \in \mathbb{C}[z, z^{-1}] \;:\; P'(z) = z^{-k}P(z),\; P \in \mathbb{C}[z],\; \|P\|_{\infty,\mathbb{T}} \leq 1,\; 0 \leq k \leq \deg P \bigr\}
$$
of Laurent polynomials whose ordinary partner is sup-norm-bounded on $\mathbb{T}$ — strictly richer than the ordinary polynomials accessible from an all-$\Lambda(W)$ circuit ($k = 0$). This Laurent class is exactly what the coherent target $\ee^{-\ii\delta A}$ requires.

The Hamiltonian simulation target can be cast in the form $\ee^{-\ii\delta\alpha\cos\theta}$, which is a Laurent function of $\ee^{\ii\theta}$, \cite[Thm. 6]{motlagh2024generalized}. Using the Jacobi-Anger expansion on this form yields, 
\begin{equation} \label{eq:jacobi-anger}
    \ee^{-\ii\delta\alpha\cos\theta} \;=\; \sum_{k=-\infty}^{\infty} (-\ii)^k\,J_k(\delta\alpha)\,\ee^{\ii k\theta}, \qquad \theta \in \mathbb{R},
\end{equation}
with the $k$-th Bessel functions $J_k$. As it turns out, truncating this series at $|k|\leq d < \infty$ can still give a good approximation to the target exponential function. We find the Laurent polynomials $L_d$ of degree $d$ in the expression:
$$
L_d(\ee^{\ii\theta}) \;:=\; \sum_{k=-d}^{d} (-\ii)^k\,J_k(\delta\alpha)\,\ee^{\ii k\theta} \;\in\; \mathbb{C}[z,z^{-1}],
$$
that have a tail bound 
$$
\bigl\|\ee^{-\ii\delta\alpha\cos\theta} - L_d\bigr\|_{\infty,\mathbb{T}} \;\leq\; 2\!\!\sum_{|k|>d}\!|J_k(\delta\alpha)| \;\leq\; \mathcal{O}\!\left(\!\left(\frac{\ee\delta\alpha}{2(d+1)}\right)^{d+1}\right),
$$
justifying the use of a truncation. The trick to efficiently realise $L_d$, is to first introduce an underlying polynomial $P$ as, 
$$
P(z) \;:=\; z^{d}\,L_{d}(z) \;=\; \sum_{m=0}^{2d} c_m\,z^m, \qquad c_m \;=\; (-\ii)^{m-d}\,J_{m-d}(\delta\alpha),
$$
of degree $2d$. Since $|z^d|=1$ on $\mathbb{T}$, $\|P\|_{\infty,\mathbb{T}} = \|L_d\|_{\infty,\mathbb{T}} \leq 1 + \mathcal{O}((\ee\delta\alpha/(2(d+1)))^{d+1})$, and a rescaling $(1-\varepsilon_\text{QSP})P$ enforces \eqref{eq:qsp-poly}. Applying \cite[Thm. 6]{motlagh2024generalized} with $\deg P = 2d$ and $k = d$ produces
$$
P'(z) \;=\; z^{-d}\,P(z) \;=\; L_d(z),
$$
so the post-selected top-left block directly realises the Jacobi–Anger Laurent truncation, giving $L_d(W_\lambda) = \ee^{-\ii\delta\alpha\cos\theta_\lambda}\,I_{2\times 2} + \mathcal{O}(\varepsilon_\text{QSP}) = \ee^{-\ii\delta\lambda}\,I_{2\times 2} + \mathcal{O}(\varepsilon_\text{QSP})$ on each qubitization subspace, and hence $\ee^{-\ii\delta A} + \mathcal{O}(\varepsilon_\text{QSP})$ on the system register.

The truncation index $d$ that suffices for $\varepsilon_\mathrm{QSP}$-accuracy on $\mathbb{T}$ is \cite[Thm. 7]{motlagh2024generalized}, 
\begin{equation} \label{eq:qsp-degree}
    d \;=\; \mathcal{O}\!\left(\delta\alpha \;+\; \frac{\log(1/\varepsilon_\text{QSP})}{\log\log(1/\varepsilon_\text{QSP})}\right).
\end{equation}
Consequently, one GQSP layer applies $d$ controlled-$W$ calls, $d$ controlled-$W^\dagger$ calls, together with $2d + 1$ single-qubit rotations on the QSP ancilla. In the regime $\delta\alpha \lesssim 1$, relevant for our coherent step, already $d=1$ suffices:
$$
L_1(\ee^{\ii\theta}) \;=\; J_0(\delta\alpha) \,-\, 2\ii\,J_1(\delta\alpha)\,\cos\theta.
$$
This can be realised by one controlled-$W$, one controlled-$W^\dagger$, and three rotations $(R_0, R_1, R_2)$. Inserting $J_0(z) = 1 + \mathcal{O}(z^2)$ and $2J_1(z) = z + \mathcal{O}(z^3)$ yields
$$
L_1(\ee^{\ii\theta}) \;=\; 1 \,-\, \ii\,\delta\alpha\,\cos\theta \,+\, \mathcal{O}((\delta\alpha)^2) \;=\; \ee^{-\ii\delta\alpha\cos\theta} \,+\, \mathcal{O}((\delta\alpha)^2).
$$
Larger subnormalization $\alpha$ or a larger time step $\delta$ at fixed accuracy could push $d \geq 2$ \eqref{eq:qsp-degree}, and the cost grows linearly with $\delta \alpha$. 

Given the polynomial coefficients $\{c_m\}_{m=0}^{2d}$, the angles $\{(\theta_0,\phi_0,\lambda_0)\}\cup\{(\theta_j,\phi_j)\}_{j=1}^{2d}$ are classically precomputed for the quantum circuit. First step is to construct a complementary polynomial $Q\in\mathbb{C}[z]$ with $\deg Q = 2d$ and on $\mathbb{T}$ satisfying
$$
|P(z)|^2 + |Q(z)|^2 \;=\; 1.
$$
The existence of $Q$ for a given $P$ is proven by \cite[Thm. 4]{motlagh2024generalized}. But we deviate from their proposed way, and use the contour-integral / Fast Fourier Transform construction of \cite{berntson2025complementary}. They do not use optimization to solve the problem, moreover, they can evaluate $Q$ with explicit error bounds up front, and work in classical time $\tilde{O}(d)$. Once we have $Q$, using the recursive extraction of \cite[Alg. 1]{motlagh2024generalized} results in the $2d + 1$ rotation triplets from $(P, Q)$ in $\mathcal{O}(d^2)$ arithmetic operations. Both stages are performed once, before any quantum computation, for each value $\delta \alpha$. The same angles are then reused at every quantum call of the layer. 

\smallskip
\section{Weak measurement} \label{sec:prelim-weak-meas}

\begin{figure}[H]
  \centering
  \begin{quantikz}[row sep=0.85cm, column sep=0.55cm]
    \lstick{$\ket{0}_\delta$}
      & \qw
      & \qw
      & \qw
      & \gate{Y_\delta}
      & \octrl{1}
      & \meter{}
      \\
    \lstick{$\ket{0^b}$}
      & \qwbundle{}
      & \gate[2]{U_L}
      & \qw
      & \octrl{-1}
      & \gate[2]{U_L^{\dagger}}
      & \rstick{$\bra{0^b}$} \qw
      \\
    \lstick{$\rho$}
      & \qwbundle{}
      &
      & \qw
      & \qw
      &
      & \rstick{$\approx\ \rho + \tfrac{\delta}{Z_L^{2}}\,\mathcal{D}_L(\rho) + \mathcal{O}(\delta^{2})$} \qw
  \end{quantikz}
  \caption{Weak measurement as a Lindblad dissipator, with
  $\bra{0^b} U_L \ket{0^b} = L/Z_L$ a block encoding of the jump
  $L$. The $Y_\delta$ rotation is controlled on the block register
  $\ket{0^b}$, so that it fires only on the post-selectable branch of
  $U_L$. Post-selecting $b$ on $\ket{0^b}$ and measuring $q_{\delta}$
  realises one Euler step of $\mathcal{D}_L(\rho) = L\rho L^{\dagger}
  - \tfrac{1}{2}\{L^{\dagger}L,\rho\}$ at rate
  $\gamma = \delta/(Z_L^{2}\,\dd t)$.}
  \label{circ:weak-measurement}
\end{figure}

There is no one-size-fits-all quantum simulation of Lindbladian evolutions. Based on the Lindbladian's specifics, one method could work while another one not. The \emph{weak measurement} scheme we introduce here, established by \cite[Thm. III.1]{chen2023quantum}, fits indeed best to the quantum algorithm in the centre of this thesis, \hyperref[alg:main]{Alg.\ref*{alg:main}}.
A Lindbladian dissipator for a jump operator $L$ is given by $\mathcal{D}_L(\rho) := L\rho L^\dagger - \tfrac{1}{2}\{L^\dagger L,\rho\}$. It is a combination of a transition term and a reflection term, that together leads to a trace-preserving map $\rho \mapsto \ee^{t\mathcal{D}_L}\rho$. A unitary block-encoding $U_L$ like the LCU \hyperref[circ:lcu]{Fig.~\ref*{circ:lcu}} could already act as the transition term $L\rho L^\dagger$: post-select on the block-encoding ancillas as $|0^b\rangle$ to have $U_L$ imprint $L\ket{\psi} / Z_L$ on the system (with $\|L/Z_L\|\leq 1$). Clearly, the same technique could not deliver the remaining reflection part of the Lindbladian. Chen et al. solves this in \cite{chen2023quantum} by weaving a weak measurement ancilla $\ket{0}_\delta$ between $U_L$ and a conditional $U_L^\dagger$ in such a way that post-selecting on all ancillas to be zero implements an Euler step of $\ee^{\delta\mathcal{D}_L/Z_L^2}$ on $\rho$ up to $\mathcal{O}(\delta^2)$ (\hyperref[circ:weak-measurement]{Fig.~\ref*{circ:weak-measurement}}).

The circuit walkthrough after initialising the three registers in $\ket{0}_\delta \otimes \ket{0^b} \otimes \rho$ is the following:
\begin{enumerate}[(i)]
    \item \emph{Block encoding.} $U_L$ acts on $(b, S)$ s.t. on the $\ket{0^b}$ block of the state imprints the amplitude $L/Z_L$ on $S$, leaving an orthogonal artefact on $(b, S)$:
    $$
    |\phi^\perp\rangle := \left(\idm - |0^b\rangle\!\langle0^b|\!\otimes\idm\right)\,U_L\,|0^b\rangle\!\ket\psi.
    $$
    \item \emph{Ancilla rotation.} $Y_\delta = R_Y(2\arcsin\sqrt\delta)$ is controlled on the block encoding being successful, i.e. $\ket{0^b}$, and leads to
    \begin{equation}
        \begin{split}
            Y_\delta\,\left(\ket{0}_\delta\otimes|0^b\rangle\right) &\;=\; \left(\sqrt{1-\delta}\,\ket{0}_\delta + \sqrt{\delta}\,\ket{1}_\delta\right)\otimes |0^b\rangle, \\
            Y_\delta\,\left(\ket{0}_\delta\otimes |\phi^\perp\rangle\right) &\;=\; \ket{0}_\delta\otimes|\phi^\perp\rangle.
        \end{split}
    \end{equation}
    In other words, the weak measurement ancilla acquires amplitude $\sqrt{\delta}$ on $\ket{1}_\delta$ precisely where $U_L$ succeeded, and carries $\ket{0}_\delta$ unchanged on the leaked $\ket{\phi^\perp}$ branch.
    \item \emph{Controlled undo.} The open-controlled $U_L^\dagger$ fires only on the $\ket{0}_\delta$ branch. On the $\ket{0^b}$-block of $U_L\ket{0^b}\ket\psi$ this uncomputes to $\ket{0^b}\ket\psi$. The artefact part $\ket{\phi^\perp}$ survives without any change to it. On the $\ket{1}_\delta$ branch $U_L^\dagger$ is absent, thus $\ket{0^b}\otimes(L/Z_L)\ket\psi$ stands.
    \item \emph{Measure and post-select.} Finally, measure all ancillas, and post-select only on the block encoding ones as $\ket{0^b}$.
\end{enumerate}

We can understand what is happening in this construction by Kraus operators \eqref{eq:kraus}: in the \emph{no-jump} branch  ($\ket{0}_\delta$ outcome, $\ket{0^b}$ post-selected), the uncomputation $U_L^\dagger U_L = \idm$ acts on the block, while the $\ket{\phi^\perp}$ artefact has $\ket{0^b}$ post-selection failure probability $\|L\ket\psi\|^2/Z_L^2$. In the \emph{jump} branch ($\ket{1}_\delta$ outcome, $\ket{0^b}$ post-selected) the full amplitude $L\ket\psi/Z_L$ survives, weighted by the $\sqrt{\delta}$ of the weak rotation. The results for a given Lindblad operator $L$ derived by \cite{chen2023quantum} can be summarized as the Kraus jumps
\begin{equation}
    M_0 \;=\; \sqrt{1-\delta}\,\idm \;-\; \tfrac{1}{2}\,\tfrac{\delta}{Z_L^2}\,L^\dagger L \;+\; \mathcal{O}(\delta^2), \qquad M_1 \;=\; \sqrt{\delta}\;\frac{L}{Z_L}.
\end{equation}
By tracing out the ancillas (summing over the branches) and expanding in $\delta$,
\begin{equation} \label{eq:weak-meas-euler}
    \rho \;\mapsto\; M_0\,\rho\,M_0^\dagger + M_1\,\rho\,M_1^\dagger \;=\; \rho \;+\; \frac{\delta}{Z_L^{2}}\Bigl(L\rho L^\dagger - \tfrac{1}{2}\{L^\dagger L,\rho\}\Bigr) \;+\; \mathcal{O}(\delta^{2}).
\end{equation}
which is precisely one Euler step of $\dot\rho = (1/Z_L^2)\,\mathcal{D}_L(\rho)$ up to $\mathcal{O}(\delta^2)$ errors. The post-selection on $b$ succeeds with probability $1 - \delta\,\|L\ket\psi\|^2/Z_L^2 + \mathcal{O}(\delta^2) = 1 - \mathcal{O}(\delta)$, i.e. unlike the post-selection loss ($1/Z_L^2$) of LCU \eqref{eq:lcu}, here there is no need for oblivious amplitude amplificiation or any other trick to boost the success probability. The smallness of $\delta$ is enough to make this into a deterministic mechanism. Or put differently, even the artefact blocks contribute to the Lindbladian correctly -- as the reflection part -- up to $\mathcal{O}(\delta^2)$ errors. The extension to multiple jumps is straight-forward, which is precisely what we can see in \cite{chen2023quantum,chen2023efficient}, and also in our slightly modified variant \hyperref[alg:diss]{Alg.~\ref*{alg:diss}}.

The Euler-step accuracy of \eqref{eq:weak-meas-euler} is only first order in $\delta$ so reaching a diamond-norm error $\varepsilon$ over the total time $T$ costs $T/\delta = \mathcal{O}(T^2 / \varepsilon)$ calls to $U_L$, i.e. polynomial in $1 / \varepsilon$. One might therefore ask why we do not instead pick a higher-order Lindblad simulator, such as the one proposed by \cite{li2022simulating} that uses Duhamel's principle. This could simulate a generator $\mathcal{L}(\rho) = -\ii[H,\rho] + \sum_{j=1}^m \mathcal{D}_{L_j}(\rho)$ with $\mathcal{O}(\tau\,\mathrm{polylog}(t/\varepsilon))$ block-encoding queries to $U_H$ and to $U_{L_j}$ and $\mathcal{O}(m\tau \mathrm{polylog}(t / \varepsilon))$ additional 1- and 2-qubit gates, where $\tau := t (\alpha_0 + \tfrac{1}{2}\sum_{j=1}^m \alpha_j^2)$. Accordingly, that higher-order simulators are meant to take in a \emph{finite} list $\{L_j\}_{j=1}^m$ of jump operators, each separately block encoded by $U_{L_j}$, leading to a cost linear in $m$. This breaks down in case the generator is of the form \eqref{eq:ckg-L}
$$
\mathcal{L}(\rho) \;=\;\int \dd\omega\;\gamma(\omega)\,\mathcal{D}_{L_j(\omega)}(\rho),
$$
with the continuously parametrized jump operators $L_j(\omega)$, that is used later to estimate the Bohr-frequencies jump components for the detailed balanced Davies generator \eqref{eq:davies-generator}. Even after a quadrature discretization of $\int \dd \omega$ to $\omega_0\sum_{\bar{\omega} \in \Omega}$, $m$ is at least the size of the frequency grid $\Omega$, and each grid point would need its own block encoding rather than being addressed coherently through the $\Omega$ register of the phase estimation \hyperref[circ:qpe]{Fig.~\ref*{circ:qpe}}. The weak measurement primitive is exactly what one uses when the integral cannot be separated from the block encodings. A single $U_L$ call already block-encodes the full sum of jumps $\{L_j(\omega)\}_{\omega\in\Omega}$. Then, one Euler step of \eqref{eq:weak-meas-euler} realises one step of the \emph{whole} Lindbladian per call, with cost that scales with the resolution of the phase estimations rather than with the number of effective jumps.

Undeniably, there is a structural difference between the two Lindbladian simulations. The higher-order series scheme buys $\mathrm{polylog}(1/\varepsilon)$ accurate channel implementation via LCUs, but only when the jump set is finite, and preferably not a large set. In this sense, the continuously parametrized jumps $\{L_j(\omega)\}_{\omega\in\Omega}$ simply does not compress into this scheme. An alternative KMS detialed balanced Lindbladian construction \cite{ding2025efficient} that came after the Chen et al. results, manages to confine the jumps to a finite set. While that allows them to utilize the higher-order simulators of \cite{li2022simulating}, their construction introduces more cost at other parts of the algorithm, making it less clear which solution is better or worse in general. We discuss this in more details in [REF] \todo{IF you write the comparison section, ref to it.}


% \part{Gibbs Sampling on a Quantum Computer}
%\chapter{Quantum Metropolis Sampling}
%We have now come to the point where we can turn our attention to recent developments in quantum Gibbs sampling. The moment we mentioned Gibbs sampling in the classical setting, we also learnt the recipe that researchers used for the past decades in all kinds of scenarios. So much so, that the \hyperref[parag:MH]{Metropolis-Hastings algorithm} is now a ubiquitous word of the sampling jargon, and it became one of the most important algorithms of the 20th century \cite{dongarra2000guest}. 
%While we have no doubt the classical Metropolis algorithm will continue to advance and have even more successes, it naturally has scaling limitations, just as most algorithms do. The dimensional explosion of the state space in quantum systems and similarly in large classical systems can pose a problem for even the state-of-the-art implementations of the algorithm. \textcolor{purple}{[cite]} Now that quantum computers have been on the horizon for some years, it could be beneficial to explore its capabilities for (Gibbs) sampling. Then one of the natural questions to ask is,
%$$\textit{Can we generalize the Metropolis algorithm to the quantum setting?}$$
%The answer has only been recently given in a foundational work \cite{temme2011quantum}, that established the notion of the \textit{quantum Metropolis algorithm.} 

%If we recall the steps of the classical algorithm, it becomes quickly evident why the quantum generalization kept the community waiting. We start in some initial state and then propose a new state: if the energy of the new configuration is lower accept it, otherwise reject it with probability given by the (local) Metropolis rule, $\exp(\beta(E_\text{old} - E_\text{new}))$. Classically, we can compute the energy of the new state, and if it has to be rejected, then we throw it away and use a copy of the previous state. In quantum mechanics however, we cannot copy, or \textit{clone} states \cite{wootters1982single}. This could pose a problem, because some operations in a quantum computer, like the energy measurement of states, are irreversible. One way to get around this, is to disturb the system as weakly as possible when we measure, such that the rejected state is still in some relation with the initial state. If some overlap between them survives, then using a Marriott-Watrous amplification scheme \cite{marriott2005quantum} could rewind back iteratively to the initial state. Crucially, this procedure together with the rest of the quantum subroutines, cannot break the detailed balanced condition, or else there is no guarantee for sampling from the Gibbs distribution. 

\chapter{Dissipative Quantum Gibbs Sampling}

In the previous chapter we have seen and analysed the very first breakthrough algorithm in quantum Gibbs sampling. The quantum Metropolis algorithm has kicked off a myriad of research, trying to better understand it, use it as a subroutine for other problems \todo{cites}, and also to possibly circumvent its main drawbacks, namely, the costly rewinds and the lack of exact detailed balance. There is no clear-cut, perfect solution to either of these problems. 

One way to get rid of the rewinds, is to work with a continuous-time Markov sampling instead of a discrete one. In this case, we never have to worry about a subpar jump and its rewind, but rather worry about the rate of the jumps. As we will see in this chapter, we are constrained to use low rates for the jumps, or else the approximations would get out of hand. Nevertheless, this might be an improvement to the Metropolis algorithm, since the rewind itself required a weak measurement, i.e. a weakly disturbed system, as well. 

Concerning the question of detailed balance, it must be stated that it is not a priori clear how and if at all $exact$ detailed balance is necessary for Gibbs sampling \textit{in practice}. Of course, if we have detailed balance we can work with well established techniques for reversible Markov semigroups; it gives guarantees that the sampling eventually reach the Gibbs state and thus helps with rigorous mixing time bound proofs. However, implementing an exact detailed balanced algorithm on a computer could easily break that nice, yet fragile symmetry. Even the algorithm we will introduce here, will break its proven detailed balance the moment we put it in quantum circuit form. The hope is that once there is an exact detailed balance framework established, it can be used as a blueprint for follow-up algorithms that break the symmetry in a controlled and rigorous way. 

The goal for this chapter is to (i) introduce a recently proposed approximate and exact detailed balanced Lindbladians and their implementations from Chen et al. \cite{chen2023quantum,chen2023efficient}; (ii) Analyse it numerically to understand their $real$ complexity, e.g. mixing times for given systems, gate complexities, temperature dependencies; (iii) Work out how best can we implement them on a quantum device.

\smallskip
\subsection{Approximate GNS Lindbladian.} After learning that the Davies generator has a proven GNS detailed balance in \hyperref[subsec:davies]{Section \ref*{subsec:davies}}, it is clear that the crux is not only about having detailed balance, but also about having an efficient implementation for it. The main obstacle with the Davies generator was that it required perfect energy estimates for the jumps. For small systems this might even be possible, and we could distinguish each and every Bohr frequency component of a given jump, then weigh each of them  with the prescribed transition rates (that satisfy KMS condition). But as we increase the system size and consider more interesting systems, the Bohr frequencies get exponentially close to each other. This poses an inherent metrology problem for implementations, because of the energy-time uncertainty principle. Which means that the more precise an energy value needs to be estimated, the longer Hamiltonian time evolution is necessary to do in the quantum computer (e.g. with QPE), thus leading to exponential costs.

A natural question to ask is, what happens with detailed balance if we allow a controlled error around the Bohr frequencies we need to estimate. As it turns out from \cite{chen2023quantum}, these smoothened Davies generators still \textit{approximately} satisfy detailed balance. To see why, let us start from the Davies generator, which fulfils exact GNS detailed balance \eqref{eq:davies-gns-db}. There, it was useful to work with the Heisenberg evolution of the jumps $A$:
$$
    A(t) = \ee^{\ii Ht} A \ee^{-\ii Ht} = \sum_{\nu\in B_H} A_\nu \ee^{\ii \nu t},
$$
where the Bohr frequencies $\nu\in B_H$ are all the energy differences in the Hamiltonian's spectrum, $\text{Spec}(H)$. The key step introduced in the Chen et al. papers is to use a filter function $f: \mathbb{R}\rightarrow \mathbb{C}$ and define the \textit{Operator Fourier Transform} (OFT) of a jump $A$ as
\begin{equation}
    \label{eq:oft}
    \hat{A}(\omega) = \frac{1}{\sqrt{2\pi}} \int_{-\infty}^{\infty} f(t) A(t) \ee^{-\ii \omega t} \dd t = \sum_{\nu \in B_H} A_\nu \hat{f}(\omega - \nu),
\end{equation}
with the Fourier transform of $f(t)$,
$$
    \hat{f}(\omega) = \frac{1}{\sqrt{2\pi}} \int_{-\infty}^{\infty} f(t) \ee^{-\ii \omega t} \dd t.
$$
As a reminder: in a quantum computer we always have to work in the time domain, and the Fourier transform is to the tool that brings us to the energy domain. In order to implement the Davies generator, one would need a Fourier transform that maps exactly to $A_\nu$. That can be also understood as an OFT, but with $f(t) = 1$ and thus $\hat{f}(\omega) = \sqrt{2\pi} \delta(\omega)$:
\begin{equation}
    \hat{A}_{f=1}(\omega) = \sqrt{2\pi}\sum_{\nu \in B_H} A_\nu \delta(\omega - \nu).
\end{equation}
This can then only be exactly achieved if the Fourier transform is computed over the full, infinitely large time domain. Hence, the idea here is to choose an $f(t)$ such that it has some decaying tail that effectively truncates the Fourier transforms to some more manageable subdomains. From \eqref{eq:oft} we can see the effect of the non-trivial filter function: the OFT leads to $\hat{A}(\omega)$ instead of $A_\nu$, which includes neighbouring energy contributions centred around $A_\nu$. What happens with the Davies generator and its GNS detailed balance if we swap out the jump operators, to their smoothened out form? In \cite{chen2023quantum} it was shown, that it leads to an approximately GNS detailed balance Lindbladian. 

There is some freedom in how we choose our filter function, but it necessarily has to be well-behaving $L^1$ integrable distribution. Clearly, for a not too costly implementation this function $f(t)$ has to decay rapidly in the time domain as $|t|\rightarrow \infty$, and its Fourier integral needs to be efficiently amenable to quadrature techniques. Moreover, its Fourier transform $\hat{f}(\omega)$ has to be smooth too.

One possible choice that fulfils all the above requirements is, to use a Gaussian filter function,
\begin{equation}
    \label{eq:gaussian-filter-t}
    f(t) := \left(\frac{2 \sigma^2}{\pi}\right)^{\frac{1}{4}}  \exp(-\sigma^2 t^2) \quad\text{with}\quad \int_{-\infty}^{\infty} |f(t)|^2 \dd t = 1.
\end{equation}
The width of the Gaussian $1 / (2 \sigma)$ is a tuneable parameter, that determines how strongly we smoothen the jumps and thus the generator, and it is directly affecting the cost of the quantum algorithm. Gaussians have the nice property, that their Fourier transforms are also Gaussians, but with inverse width $2\sigma$,
\begin{equation} \label{eq:gaussian-filter-w}
    \hat{f}(\omega) = \left(2\pi \sigma^2 \right)^{-\frac{1}{4}} \exp\left(-\frac{\omega^2}{4\sigma^2}\right).
\end{equation}
Then in this case, the smooth jump operators that the Fourier transforms are mapping to, can be understood as Gaussian distributions centred around each Bohr frequency $\nu$ with an uncertainty of $2\sigma$,
\begin{equation}
    \label{eq:smoothed-jump}
    \hat{A}(\omega) = \left(2\pi \sigma^2 \right)^{-\frac{1}{4}} \sum_{\nu \in B_H} \exp\left(- \frac{(\omega - \nu)^2}{4\sigma^2}\right)A_\nu.
\end{equation}
\todo{figure for Gaussian smoothing around $\nu$}
Note, that a Gaussian filter is not the only function that could be used here. Gevrey functions also fulfil the necessary requirements and have been successfully used for KMS detailed balanced Lindbladians in \cite{ding2025efficient} that we will look at later. Until then let us keep ourselves to Gaussians. 

The Davies generator with the smoothened jump operators \eqref{eq:smoothed-jump} we will call after the authors, \textit{Chen-Kastoryano-Brandão-Gilyén} or \textit{CKBG Lindbladian}, and it takes the following form in the Schrödinger picture,
\begin{equation}
    \label{eq:L-ckbg}
    \mathcal{L}_\text{CKBG}(\cdot) = \sum_{a} \int_{-\infty}^{\infty} \gamma(\omega)\left(\hat{A}^a(\omega)(\cdot)\hat{A}(\omega)^\dagger -\frac{1}{2}\left\{\hat{A}^a(\omega)^\dagger \hat{A}^a(\omega), \cdot \right\}\right)\dd\omega .
\end{equation}

The sum runs over the set of jumps, $\{A^a\}_{a\in\mathcal{A}}$, that can be chosen in any way that ensures ergodicity. Usually, with a varied enough set (breaking any symmetries that the Hilbert space might have) where each jump also has its adjoint included works well,
\begin{equation}
    \label{eq:ergodic-set-jumps}
    \{A^a\}_{a\in\mathcal{A}} = \{(A^a)^\dagger\}_{a\in\mathcal{A}} \quad\text{and}\quad \left\|\sum_a A^{a\dagger} A^a\right\| \leq 1.
\end{equation}
Here, the norm will be used to normalize the jumps and its purpose is (i) to make block encodings (i.e. implementation) possible \todo{cite to block encoding} and (ii) it also helps in analysis since it fixes a variable, namely the "strength" of the Lindbladian for different number of jump operators.
The naive set of single-site Pauli jumps work already quite well, and since they are unitary we can avoid using some block encodings and excess ancillas in the circuits. (In this case, the normalization factor in \eqref{eq:ergodic-set-jumps} becomes the total number of jumps $|\mathcal{A}|$.) Even if we choose a single-site jump as the core, the Kraus operators of the Lindbladian, $\{L_j\}_{j\in\mathcal{J}} = \{\sqrt{\gamma(\omega)}\hat{A}(\omega)\}$ will not be constraint to just that one site. We can see this from the OFT construction \eqref{eq:oft}, where it becomes clear, that as the Pauli operators evolve over time in the Heisenberg picture, its support will also spread out spatially. We can estimate this size by the Lieb-Robinson bound to see that for the small system sizes that we can numerically simulate, these jumps spread out over the whole lattice. See the Appendix for more details [APPENDIX] \todo{appendix LR bound}.

The transition rates $\gamma(\omega)$ can be any bounded function, $\gamma(\omega) \in [0, 1]$, that satisfy the KMS condition \eqref{eq:KMS-condition}, which now has the form
\begin{equation}
    \label{eq:KMS-condition-2}
    \frac{\gamma(\omega)}{\gamma(-\omega)} = \ee^{-\beta \omega}.
\end{equation}
Let us first choose a tuneable Gaussian for these transition rates,
\begin{equation} \label{eq:gaussian-transition}
    \gamma(\omega) = \exp\left( -\frac{(\omega + \omega_\gamma)^2}{2\sigma_\gamma^2} \right),
\end{equation}
where the centre of the Gaussian is at $\omega_\gamma$ and its width is given by $\sqrt{2}\sigma_\gamma.$ In other words, most of the jumps are suppressed in this case and only the ones that lead to an energy difference with $\omega \approx \omega_\gamma$ contribute meaningfully. The KMS condition \eqref{eq:KMS-condition-2} for this Gaussian $\gamma(\omega)$ actually fixes the following relation between the parameters $\beta, \omega_\gamma, \sigma_\gamma$
\begin{equation} \label{eq:beta-approx-gns-gaussian}
    \beta := \frac{2\omega_\gamma}{\sigma_\gamma^2}
\end{equation}

In [FIGURE] \todo{add figure of gammas} you can see how constraining it is compared to transitions that we will consider later on. But before we could do that, we need to see what kind of detailed balance, if any, we have with these Gaussian smoothed jump operators \eqref{eq:oft}. 

In its current form \eqref{eq:L-ckbg} has been shown to satisfy \textit{approximate GNS detailed balance} in \cite{chen2023quantum}. Quantitatively, we have understood the severity of breaking detailed balance as the norm of the anti-Hermitian part $\|\mathcal{A}(\rho, \mathcal{L})\|_{2\rightarrow 2}$ of the generator's discriminant $\mathcal{D}(\rho, \mathcal{L})$, see \hyperref[def:approx-db]{Definition \ref*{def:approx-db}}. Thus, if we can choose algorithm parameters such that the norm of the anti-Hermitian part is below some $\varepsilon$, we can also bound how well does the fixed point $\rho_\text{fix}(\mathcal{L})$ approximate the target state $\rho$.

\begin{corollary}
    \textup{(Fixed point accuracy \cite{chen2023quantum})} If a Lindbladian $\mathcal{L}$ satisfies $\varepsilon$-approximate $\rho$-detailed balance condition, then its fixed point $\rho_\text{fix}(\mathcal{L})$ deviates from $\rho$ by at most $$\|\rho_\text{fix}(\mathcal{L}) - \rho\|_1 \leq 20 t_\text{mix}(\mathcal{L})\,\varepsilon .$$
\end{corollary}
Chen et al. provides the proof of this, which is indeed less than obvious. The core thought is that perturbations to the parent Hamiltonian $\mathcal{H}$ with $\mathcal{A}$, leads to only small changes for the eigenvectors with separated eigenvalues. In fact, one can show that the top eigenvalue of $\mathcal{H}$ will not deviate further from $0$ (its detailed balanced value) than $$|\lambda_1(\mathcal{H})| \leq \|\mathcal{A}\|_{2\rightarrow 2}.$$ We can combine all of this with the CKBG Lindbladian and find that it is approximately GNS detailed balanced.
\begin{theorem}
    \textup{(Approximate detailed balance for $\mathcal{L}_\text{CKBG}$  \cite{chen2023quantum})}. Any Lindbladian $\mathcal{L}_\textup{CKBG}$ of the form \eqref{eq:L-ckbg} with jumps \eqref{eq:ergodic-set-jumps}, satisfying the KMS condition \eqref{eq:KMS-condition-2}, and having Gaussian filter functions \eqref{eq:gaussian-filter-t} has an approximate detailed balance. Hence, its fixed point $\rho_\textup{fix}(\mathcal{L_\textup{CKBG}})$ is approximately the Gibbs state $\rho_\beta$
    $$
    \left\|\rho_\textup{fix}(\mathcal{L}_\textup{CKBG}) - \rho_\beta\right\|_1 = \mathcal{O}\left(\sigma \beta \sqrt{\log\frac{1}{\sigma \beta}}\,t_\textup{mix}(\mathcal{L}_\textup{CKBG})\right).
    $$
\end{theorem}
The main takeaway of these results is that if some error appears on the generator level, it will get amplified with time, and can have significant impact on Gibbs sampling results.
This is even worse for cases (usually the interesting ones) where fast or rapid mixing has not been proven, and thus can very well make $t_\text{mix}(\mathcal{L})$ exponentially scaling with system sizes and with inverse temperature $\beta$.

One of the main culprits leading to deviations from the Gibbs state is the presence of non-secular terms in the CKBG Lindbladian \eqref{eq:L-ckbg}. For the \hyperref[subsec:davies]{Davies generator} we had to apply the secular approximation in order to get a correct, physical generator. But here, with the smoothened jump operators \eqref{eq:oft} we have reintroduced some off-diagonal entries in the Kossakowski matrix. To better see this, let us derive its analytical form by computing the integral over the energies in \eqref{eq:L-ckbg} for the transition part of the Lindbladian.
\begin{equation} \label{eq:gauss-transition-part-of-lindbladian}
    \begin{split}
        \mathcal{T} = &\int_{-\infty}^\infty \gamma(\omega) \hat{A}(\omega) (\cdot) \hat{A}(\omega)^\dagger \dd \omega \\
        &= \frac{1}{\sigma \sqrt{2\pi}} \sum_{\nu_1, \nu_2 \in B_H} \int_{-\infty}^{\infty} \ee^{-\frac{(\omega + \omega_\gamma)^2}{2\sigma_\gamma^2}}\ee^{-\frac{(\omega - \nu_1)^2}{4\sigma^2}}\ee^{-\frac{(\omega - \nu_2)^2}{4\sigma^2}}A_{\nu_1}(\cdot)A_{\nu_2}^\dagger \dd \omega \\
        &=: \sum_{\nu_1, \nu_2 \in B_H} \alpha_{\nu_1, \nu_2} A_{\nu_1}(\cdot)A_{\nu_2}^\dagger \dd \omega
    \end{split}
\end{equation}
with the Kossakowski matrix,
\begin{equation} \label{eq:kossakowski-gaussian}
    \alpha_{\nu_1,\nu_2} = \frac{\sigma_\gamma}{\sqrt{\sigma^2 + \sigma_\gamma^2}} \exp\left(-\frac{(\nu_1 + \nu_2 + 2\omega_\gamma)^2}{8(\sigma^2 + \sigma_\gamma^2)}\right)\exp\left(-\frac{(\nu_1 - \nu_2)^2}{8\sigma^2}\right).
\end{equation}
There are a few things we can read off from this expression: (i) There are non-secular terms present $\nu_1 \neq \nu_2$ due to the Gaussian smearing around the $\nu$'s; (ii) It is a positive semi-definite matrix (and Hermitian), which makes the generator $\mathcal{L}$ to be a Lindbladian generator; (iii) 
The Gaussian filters are not only breaking the secular structure but also raising the temperature. To see this latter point, let us reconsider what we are aiming to approximate: the GNS detailed balance that was physically derived by retaining diagonal terms, and drops all off-diagonal ones via the secular approximation. Setting $\nu_1 = \nu_2 = \nu$ in \eqref{eq:kossakowski-gaussian} yields, 
\begin{equation}
    \alpha_{\nu, \nu} = \frac{\sigma_\gamma}{\sqrt{\sigma^2 + \sigma_\gamma^2}} \exp\left(-\frac{(\nu + \omega_\gamma)^2}{2(\sigma^2 + \sigma_\gamma^2)}\right).
\end{equation}
These diagonal entries have to satisfy the KMS condition,
\begin{equation} \label{eq:kossakowski-diag-condition}
    \frac{\alpha_{\nu, \nu}}{\alpha_{-\nu, -\nu}} = \exp\left(-\frac{2\nu\omega_\gamma}{\sigma^2 + \sigma_\gamma^2}\right) =: e^{-\beta_\text{eff} \nu}
\end{equation}
with 
\begin{equation} \label{eq:beta-eff}
    \beta_\text{eff} = \frac{2 \omega_\gamma}{\sigma^2 + \sigma_\gamma^2} < \beta = \frac{2\omega_\gamma}{\sigma_\gamma^2}.
\end{equation}
In other words, Gibbs sampling with the $\mathcal{L}_\text{CKBG}$ from \eqref{eq:L-ckbg} and with $\sigma > 0$, effectively leads to some higher temperature Gibbs state. Though because of the additional non-secular errors, the error structure is not as clear-cut as just "raised temperatures". It can now also be clearly seen, that by taking the $\sigma \rightarrow 0$ limit, all errors to the GNS detailed balance vanish as well: $\beta_\text{eff} \rightarrow \beta$ and as the Gaussian filters become delta peaks in this limit, they also eliminate any off-diagonal elements and only $\nu_1 = \nu_2$ remains. 

\label{sec:kms-lindbladian}
\smallskip
\subsection{KMS Lindbladian} While the previous approximate GNS detailed balanced Lindbladian was a result on its own, its improvement is what got researchers (with myself included) thrilled about quantum Gibbs sampling. In \cite{chen2023efficient} Chen et al. have managed to come up with a structure for the Lindbladian where they do not fight the above-mentioned errors, but rather lean into them and fix them all at once with an additional coherent term. The result is an exact KMS detailed balanced Lindbladian instead of an approximation of the GNS detailed balance and the Davies generator with high costs. In order to better see how terms in the Lindbladian interact with each other to get KMS detailed balance, it will be useful to write the Lindbladian in this general form:
\begin{equation}
    \mathcal{L} = \mathcal{T} - \frac{1}{2}\mathcal{R} - \ii \mathcal{B}
\end{equation}
with the transition $\mathcal{T}$, decay $\mathcal{R}$ and coherent $\mathcal{B}$ parts, 
\begin{alignat}{2} \label{eq:ckg-L-parts}
\mathcal{T}(\cdot) 
  &= \sum_a\sum_{\nu_1,\nu_2} \alpha_{\nu_1,\nu_2} A^a_{\nu_1}(\cdot)(A^a_{\nu_2})^\dagger
  &\quad&\text{with}\quad
  \mathcal{T}^\dagger(\cdot) = \sum_a\sum_{\nu_1,\nu_2} \bar\alpha_{\nu_1,\nu_2}(A^a_{\nu_1})^\dagger(\cdot)A^a_{\nu_2},
  \\
\mathcal{R}(\cdot) 
  &= \sum_a\sum_{\nu_1,\nu_2} \alpha_{\nu_1,\nu_2}\{(A^a_{\nu_2})^\dagger A^a_{\nu_1},\,\cdot\,\}
  &\quad&\text{with}\quad
  \mathcal{R}^\dagger = \mathcal{R},
  \\
\mathcal{B}(\cdot) 
  &= \sum_a[B_a,\,\cdot\,]
  &\quad&\text{with}\quad
  (-i\mathcal{B})^\dagger = i\mathcal{B}.
\end{alignat}
For $\mathcal{R}^\dagger$ we implicitly took the Kossakowski matrix to be positive semidefinite and thus Hermitian, $\alpha_{\nu_1, \nu_2} = \bar\alpha_{\nu_2, \nu_1}$. Though, this needs to be checked for new constructions, without it, we would not have a GKSL Lindbladian.

Other than the additional coherent term $\mathcal{B}$, the rest of the setup is unchanged: we take the same Gaussian filter functions for the jump operators \eqref{eq:oft}; we keep the Gaussian transition weights $\gamma(\omega)$ \eqref{eq:gaussian-transition}; Consequently, we have the exact same Kossakowski matrix $\alpha_{\nu_1, \nu_2}$ too \eqref{eq:kossakowski-gaussian} with one small but pivotal difference, the temperature. We set the temperature such that the diagonal entries of $\alpha_{\nu_1,\nu_2}$ satisfy the KMS condition, i.e. as $\beta := \beta_\text{eff}$. This instantly makes $\gamma(\omega)$ itself break the original KMS condition, however, the goal here is to satisfy the condition \textit{with} the Gaussian filters together. From the relation \eqref{eq:beta-eff}, it is clear that the Gaussian parameters and the temperature are intertwined. In the KMS case, there is a convenient way to choose them to keep the Kossakowski matrix necessarily skew-symmetric, by 
\begin{equation}
    \sigma = \omega_\gamma = \sigma_\gamma = \frac{1}{\beta}.
\end{equation}
Nevertheless, we will keep the following expressions general, because it is not a priori clear how best to choose these parameters for a faster Gibbs state preparation. It is best to keep in mind, that since the mixing time and these Gaussian parameters are all temperature dependent, they can affect one another. This is one of the questions we have tackled in our numerical analysis, see \todo{ref to results: mixing time dependence on more or less off-diagonal terms.}

Previously the off-diagonal terms $\nu_1 \neq \nu_2$ were either eliminated via secular approximations (Davies) or were the reason for deviation from GNS detailed balance (CKBG Lindbladian). Now these terms are part of the design. But that also means we must have a condition on how these off-diagonal terms relate to each other just like what we had for the diagonal part in \eqref{eq:kossakowski-diag-condition}, in order to attain KMS detailed balanced.
In the following we will see parts of the Lindbladian \eqref{eq:ckg-L-parts} behave differently in this KMS space. In fact, while the transition part $\mathcal{T}$ keeps its structure under Gibbs conjugations, the decay $\mathcal{R}$ and coherent $\mathcal{B}$ parts mix together.

Let us start with the transition part: for it to satisfy the KMS detailed balance condition, it needs to be self-adjoint with respect to the KMS inner product, or equivalently: 
\begin{equation} \label{eq:transition-self-adjoint}
    \mathcal{T}^\dagger = \Gamma_\rho^{-1}\, \circ \, \mathcal{T}\, \circ\, \Gamma_\rho^{1},
\end{equation}
with $\Gamma_\rho(\cdot) := \rho_\beta^{\frac{1}{2}} (\cdot) \rho_\beta^{\frac{1}{2}}$. The equality indeed holds if each jump $A^a$ has its Hermitian conjugate in the jumps set $\mathcal{A}$, and most importantly if the Kossakowski matrix has a specific skew-symmetry:
\begin{equation} \label{eq:skew-symmetry}
    \alpha_{\nu_1, \nu_2} = \alpha_{-\nu_2, -\nu_1} e^{-\beta(\nu_1 + \nu_2) / 2},
\end{equation}
which can even be thought of as the off-diagonal extension of the diagonal condition we found previously \eqref{eq:kossakowski-diag-condition}.
This skew-symmetry has to be satisfied for any given combination of filter function $f(t)$, acceptance rates $\gamma(\omega)$ and $\beta$. For our current Gaussian case \eqref{eq:kossakowski-gaussian}, it was shown that the inverse temperature has to take its already introduced form, $\beta = \beta_\text{eff}$ from \eqref{eq:beta-eff}. 

Now let us look at the remaining parts of the Lindbladian, $\mathcal{R}$ and $\mathcal{B}$. As a quick check, we can conjugate $\mathcal{R}$ with the Gibbs state to see what kinds of terms we get: 
\begin{equation} \label{eq:decay-gibbs-conjugation}
    \Gamma_\rho^{-1}\, \circ \, \mathcal{R}\, \circ\, \Gamma_\rho^{1}(\cdot) =\sum_{\nu_1, \nu_2} \left( e^{\beta(\nu_1 - \nu_2)/2}R_{\nu_1, \nu_2}(\cdot) + (\cdot) R_{\nu_1, \nu_2} e^{-\beta(\nu_1 - \nu_2)/2}\right),
\end{equation}
with $R_{\nu_1, \nu_2} := \sum_a \alpha_{\nu_1, \nu_2} (A^a_{\nu_2})^\dagger A^a_{\nu_1}$. This clearly shows, that the two terms pick up a different factor in the anti-commutator of $\mathcal{R}$ and thus no simple condition on $\alpha_{\nu_1, \nu_2}$ like the skew-symmetry, can lead to KMS detailed balance alone. However, the coherent term $\mathcal{B}$ leads to similar terms under the KMS conjugation as the decay part, 
\begin{equation} \label{eq:coherent-gibbs-conjugation}
    \Gamma_\rho^{-1}\circ\mathcal{B}\circ\Gamma_\rho(\cdot) = \sum_\nu \left(e^{+\beta\nu/2}\,B_\nu\, (\cdot) \;-\; (\cdot)\,B_\nu e^{-\beta\nu/2}\right),
\end{equation}
with the usual Bohr frequency decomposition $B = \sum_\nu B_\nu$. Then, one can dream up a form for $B$ such that all the detailed balance violating terms in the decay part in \eqref{eq:decay-gibbs-conjugation} cancel (while the skew-symmetry \eqref{eq:skew-symmetry} for the Kossakowski matrix holds). Which is exactly what Chen et al. did in \cite{chen2023efficient}.

Then, the key equation that has to hold is
\begin{equation}
    \left(-\frac{1}{2}\mathcal{R} + \ii \mathcal{B}\right) = \Gamma_\rho^{-1}\circ\left(-\frac{1}{2}\mathcal{R} - \ii \mathcal{B}\right)\circ\Gamma_\rho.
\end{equation}
With the Hilbert-Schmidt adjoint on the LHS and the Gibbs conjugation on the RHS. In other words, we are looking for an expression of $\mathcal{B}$ such that $\mathcal{R}$ and $\mathcal{B}$ together are self-adjoint with respect to the KMS inner product. If we write all these terms out using \eqref{eq:decay-gibbs-conjugation} and \eqref{eq:coherent-gibbs-conjugation}, and identify $\nu = \nu_1 - \nu_2$, then we can match the prefactors of the same type of terms on both hands. For instance, the left-multiplication prefactors give: 
\begin{equation}
    -\frac{1}{2}R_{\nu} + \ii B_\nu = -\frac{1}{2} \ee^{\beta\nu / 2}R_{\nu} - \ii \ee^{\beta\nu / 2} B_\nu
\end{equation}
which rearranges to
\begin{equation}
    \ii B_\nu\left(1 + \ee^{\beta\nu / 2}\right) = -\frac{1}{2}\left(\ee^{\beta\nu / 2} - 1\right) R_{\nu}.
\end{equation}
This determines $B_\nu$ uniquely at each Bohr frequency,
\begin{equation}
    B_\nu = \frac{1}{2\ii}\,\frac{1 - \ee^{\beta\nu /2}}{1 + \ee^{\beta\nu /2}}\;R_\nu = \frac{1}{2\ii}\tanh\!\left(-\frac{\beta\nu}{4}\right)\,R_\nu,
\end{equation}
where we used the identity: $\frac{1 - \ee^x}{1 + \ee^x} = \tanh(-x / 2).$ The prefactors for the right-multiplications lead to the same result. Restoring the indices and summing up over them leads to the full coherent term
\begin{equation} \label{eq:B-bohr-domain}
    B = \frac{1}{2\ii}\sum_a\sum_{\nu_1,\nu_2}\tanh\!\left(-\frac{\beta(\nu_1-\nu_2)}{4}\right)\;\alpha_{\nu_1\nu_2}\;(A^a_{\nu_2})^\dagger A^a_{\nu_1}.
\end{equation}
Concluding this section, we have successfully found a unique and Hermitian $B$ which makes the decay and coherent parts of the Lindbladian, $\mathcal{R}$ and $\mathcal{B}$ KMS detailed balanced. Since $\alpha_{\nu_1, \nu_2}$ has been kept skew-symmetric, the translation part $\mathcal{T}$ is also correct. Combining everything together, we get a KMS detailed balanced Lindbladian, that we will call from now on as the CKG Lindbladian,
\begin{equation}\label{eq:ckg-L}
    \mathcal{L}_{\text{CKG}}(\cdot) = -\ii[B, \cdot\,] + \sum_a \int_{-\infty}^\infty \gamma(\omega) \left(\hat{A}^a(\omega)(\cdot)\hat{A}(\omega)^\dagger -\frac{1}{2}\left\{\hat{A}^a(\omega)^\dagger \hat{A}^a(\omega), \cdot \right\}\right)\dd\omega .
\end{equation}

\smallskip
\subsection{Designing transition weights}
\todo{plot for gammas}
Only allowing transitions with energies concentrated as a Gaussian around some $\omega_\gamma$ can be very restricting, which inevitably leads to slow thermalization. The question is, \textit{can we come up with more accepting rates $\gamma(\omega)$ while keeping the KMS detailed balance of the CKG Lindbladians?} The answer is yes, by finding a clear connection between the condition on the transition weights $\gamma(\omega)$ and the necessary skew-symmetry of the Kossakowski matrix $\alpha_{\nu_1, \nu_2}$.
\begin{lemma} \label{lem:shifted-kms-condition}
    \textup{(Shifted KMS condition for detailed balance, \cite{ramkumar2024mixing})} If $\tilde{\gamma}(\omega)$ satisfies the KMS condition $\tilde{\gamma}(\omega) = \tilde{\gamma}(-\omega)e^{-\beta \omega},$ then the Kossakowski matrix 
    $$
    \alpha_{\nu_1, \nu_2} = \int_{-\infty}^\infty \gamma(\omega) \hat{f}(\omega - \nu_1) \hat{f}(\omega - \nu_2) \dd \omega \quad\textup{with}\quad \gamma(\omega) = \tilde{\gamma}\left(\omega + \frac{\beta \sigma^2}{2}\right),
    $$
    will be skew-symmetric by \eqref{eq:skew-symmetry}. Here, $\hat{f}$ is any filter with $|\hat{f}(\omega)|^2 = |\hat{f}(-\omega)|^2$, like the usual Gaussian filter from \eqref{eq:gaussian-filter-w}.
\end{lemma}
\noindent The proof from the authors follows by direct verification of the self-adjointness condition \eqref{eq:transition-self-adjoint} for $\mathcal{T}$. 

Hence, we can take any transition weight function $\tilde{\gamma}(\omega)$ that satisfies the KMS condition, shift it to $\gamma(\omega)$ and find a KMS detailed balanced CKG Lindbladian. This is quite interesting, since we can import canonically well-working transition weights that has been used for Gibbs sampling in the past. One such function, is the \textit{Metropolis function}, which takes the following form after the prescribed shift from \hyperref[lem:shifted-kms-condition]{Lemma~\ref*{lem:shifted-kms-condition}}: 
\begin{equation} \label{eq:shifted-metro-func}
    \gamma_\text{M}(\omega)= \exp\left(-\beta \max\left(\omega + \frac{\beta\sigma^2}{2}, 0\right)\right).
\end{equation}

From Fig. \todo{fig} it becomes apparent why the Metropolis function would lead to better mixing compared to the previous Gaussian case. It accepts almost all transitions that lower the energy with probability 1, and only suppresses energy increasing jumps. Since the Metropolis function is the centrepiece of the Metropolis-Hastings algorithm \cite{metropolis1953equation, hastings1970monte} which is already decades old, the community had time to rigorously analyse why it works so well. 
Peskun showed in \cite{peskun1973optimum} that among all reversible chains with the same proposal and target, the one with the pointwise maximal off-diagonal transition rates (between states $i \neq j$), has the smallest asymptotic variance for any estimator. The (unshifted) Metropolis acceptance is the unique pointwise maximum satisfying detailed balance with $\|\gamma\|_\infty\leq 1$. Then, Tierney extends these results in \cite{tierney1998note} by finding that the spectral gap among these reversible chains, the spectral gap is also the largest for the Metropolis transition, hence it is the fastest mixing of them all.

These are strong results, but they apply only for classical Markov chains and as it turns out, they do not generally hold in the quantum case. It is an open question if there exists a quantum variant of the above results. But at the moment it is not clear if and how the off-diagonal elements of a Kossakowski matrix $\alpha_{\nu_1, \nu_2}$ improve upon the mixing times of quantum Markov semigroups. There are recent numerical evidences in the general KMS detailed balanced framework \cite{chen2024boosting,hahn2025efficient} suggesting that there is not one choice that fits all problems. But also shows a gap in analytical results.

We can however look at just the CKG Lindbladian case and ask, for given Gaussian filter functions \eqref{eq:gaussian-filter-w}, would the Metropolis transition rate $\gamma_M$ lead to the fastest mixing?  

\begin{proposition} \label{prop:optim-metro}
    \textup{(Optimality of Metropolis weights for CKG Lindbladians)} Assume that the CKG Lindbladian $\mathcal{L}_\textup{CKG}^{\gamma_M}$ with transition weight $\gamma_M$ \eqref{eq:shifted-metro-func} is primitive. For any $\gamma: \mathbb{R} \rightarrow [0, \infty)$ that satisfies the shifted KMS condition \hyperref[lem:shifted-kms-condition]{Lemma~\ref*{lem:shifted-kms-condition}} and $\|\gamma\|_\infty \leq 1$: 
    $$\textup{Gap}\left(\mathcal{L}_\textup{CKG}^{\gamma_M}\right) \geq \textup{Gap}\left(\mathcal{L}_\textup{CKG}^{\gamma}\right) \geq 0.$$
\end{proposition}

\begin{proof}
    Define $\Delta := \gamma_M - \gamma$. We claim $\Delta \geq 0$. For the piece of the function where $\omega \geq -\beta\sigma^2 / 2$, the shifted KMS condition gives
    $$
    \gamma(\omega) = \gamma(-\omega - \beta\sigma^2)e^{-\beta(\omega + \beta\sigma^2/2)} \leq e^{-\beta(\omega + \beta\sigma^2 / 2)} \equiv \gamma_M(\omega),
    $$
    using $\|\gamma\|_\infty \leq 1$. For the other piece, $\omega < -\beta\sigma^2/2$, we directly have $\gamma(\omega) \leq 1 = \gamma_M(\omega)$. This argument is analogous to Peskun \cite{peskun1973optimum}. The CKG shifted KMS condition is a linear identity in $\gamma$: since both $\gamma_M, \gamma$ satisfy it, so does $\Delta$. Thus, $\mathcal{L}_\textup{CKG}^\Delta$ is a valid GKLS generator which satisfies KMS detailed balance.

    Since the map $\gamma\rightarrow\mathcal{L}^\gamma$ is linear: for any non-negative $\gamma_1, \gamma_2$, both satisfying the CKG shifted KMS condition, and any $c_1, c_2 \geq 0$, we have $\mathcal{L}_\textup{CKG}^{c_1\gamma_1 + c_2\gamma_2} = c_1 \mathcal{L}_\textup{CKG}^{\gamma_1} + c_2 \mathcal{L}_\textup{CKG}^{\gamma_2}$. The reason behind this is that every construction in the CKG Lindbladian involves an integral linear in $\alpha_{\nu_1, \nu_2}$ and correspondingly linear in $\gamma$. Then we can write: $$\mathcal{L}_\textup{CKG}^{\gamma_M} = \mathcal{L}_\textup{CKG}^\gamma + \mathcal{L}_\textup{CKG}^\Delta, 
    \quad\textup{or}\quad
     \left(\mathcal{L}_\textup{CKG}^{\gamma_M}\right)^\dagger = \left(\mathcal{L}_\textup{CKG}^\gamma\right)^\dagger + \left(\mathcal{L}_\textup{CKG}^\Delta\right)^\dagger.$$
    Recall that the spectral gap is defined as
    $\textup{Gap}\left(\mathcal{L}^\dagger\right) = \inf_{X} \mathcal{E}(X, X) / \|X\|^2_\textup{KMS}$ with $X\neq 0$ and $\tr(\rho_\beta X) = 0$, and with the Dirichlet form $\mathcal{E}^\gamma(X, X) := -\langle X, (\mathcal{L}^\gamma)^\dagger(X) \rangle_\textup{KMS}$. See also [REF] \todo{refer to preliminaries} and note that the Dirichlet form is defined in the Heisenberg picture, hence the $\mathcal{L}^\dagger$. Consequently, the proof boils down to showing that $$\mathcal{E}^\Delta(X, X) \geq 0 \quad\forall X \in \mathcal{B(H)},$$
    which is true since: (i) the Schrödinger and Heisenberg operators share the same spectra; (ii) $\mathcal{L}_\textup{CKG}^\Delta$ is a GKLS generator and thus has $\textup{Re}(\lambda) \leq 0$ for all eigenvalues $\lambda$; (iii) $\mathcal{L}_\textup{CKG}^\Delta$ has KMS detailed balance, i.e. it is self-adjoint with respect to the KMS inner-product which makes all eigenvalues on this space real. Combining all of this leads to $\langle X, (\mathcal{L}^\Delta)^\dagger(X) \rangle_\textup{KMS} \leq 0$ and finally to the above statement.
    
    The last step is to relate the Dirichlet forms and the spectral gaps. For all $X$ with $\tr(\rho_\beta X) = 0$, $X \neq 0$:
    $$  
    \mathcal{E}^{\gamma_M}(X, X) = \mathcal{E}^{\gamma}(X, X) + \mathcal{E}^{\Delta}(X, X) \geq \mathcal{E}^{\gamma}(X, X).
    $$
    Since $\|X\|^2_\textup{KMS}$ is independent of $\gamma$, dividing with it and taking the infimum yields:
    $$
    \textup{Gap}\left(\mathcal{L}_\textup{CKG}^{\gamma_M}\right) \geq \textup{Gap}\left(\mathcal{L}_\textup{CKG}^{\gamma}\right)
    $$
    The primitivity assumption ensures $\textup{Gap}\left(\mathcal{L}_\textup{CKG}^{\gamma_M}\right) > 0$.
\end{proof}

Although, we cannot say anything definite about optimal Kossakowski matrices, with the above proposition, we have at least the motivation to numerically analyse CKG Lindbladians with the Metropolis function $\gamma_M$ or (risking foreshadowing here) with its smoothened variants.

There is however an interesting technique that is worth mentioning, as it can help us design transition functions and Kossakowski matrices.
In \cite{chen2023efficient} the authors found this shifted Metropolis function $\gamma_M$ by different means. The idea is to take a linear combination of detailed balanced Lindbladians in order to speed up the mixing. Due to linearity, a linear combination of skew-symmetric Kossakowski matrices also respects the symmetry,
\begin{equation}\label{eq:lin-comb-alpha}
    \alpha_{\nu_1,\nu_2}^{(s, g)} := \int_{\frac{\beta\sigma^2}{2}(1 + s)}^\infty g(x) \alpha_{\nu_1,\nu_2}^{(\omega_\gamma, \sigma_\gamma)} \dd x \quad\text{with}\quad (\omega_\gamma(x), \sigma_\gamma(x)) = \left(x, \sqrt{\frac{2x}{\beta} - \sigma^2}\right).
\end{equation}
Where $g \in \ell_1(\mathbb{R})$ is necessary to keep the resulting Kossakowski matrix Hermitian. This is a linear combination of the skew-symmetric Gaussian-case Kossakowski matrices \eqref{eq:kossakowski-gaussian} by sweeping with the centre of the transition Gaussians $\omega_\gamma$. Hence, we find a transition that admits a larger set of energy differences.

By choosing the $g$ as the inverse $\ell_1$ weight of the Gaussians, $g(x) = \frac{1}{\sqrt{2\pi}\sigma_\gamma}$ and $s = 0$, Chen et al. managed to find the shifted Metropolis function \eqref{eq:shifted-metro-func}, after applying the linear combination to the Gaussian transition weight $\gamma_G(\omega)$ with the same $\omega_\gamma(x), \sigma_\gamma(x)$ parameters as in \eqref{eq:lin-comb-alpha},
\begin{equation} \label{eq:gamma_M-from-lin-comb-of-G}
    \gamma_M(\omega) = \int_{\beta\sigma^2 / 2}^\infty g(x)\gamma_G(\omega) \dd x.
\end{equation}
But while this leads to better mixing than the Gaussian case, it has some problems the moment we want to implement it. 

The source of all these problems is the fact that though $\gamma_M$ is Lipschitz continuous, it has a kink at $\omega = -\beta\sigma^2 / 2$ (and the unshifted variant at $\omega = 0$), see \todo{plot}. To be clear up front, the problems that this fact leads to are not dominating the asymptotic error / cost scalings of the algorithm, but they seem to be unnecessary costs, that for implementations prior fault-tolerance can be important. 
\begin{itemize}
    \item On one side this would lead to extra quantum circuit synthesis costs, since in the quantum algorithm at one point we have to apply a rotation onto an ancilla qubits based on $\gamma(\omega)$ \todo{refer to q. alg and circuit}. There are multiple ways to do this, but one is to use QSP, which has to approximate $\gamma_M$ via some polynomials. Due to the kink however we would need to have higher degree polynomial for approximations, and thus a higher gate cost \todo{refer to something possibly}. Chen et al. notes \cite[footnote 33.]{chen2023quantum} that we could do this also piecewise, if we check the energy register via a comparator circuit block (to decide between left and right of the kink) and then apply QSP on smoother pieces of functions. But the costs are higher if the kink is not exactly at $\omega = 0$, since the comparator circuit then has to check more bits in the register, not just the sign bit. 
    \item The other problem is more important however: since there is an integral over $\omega$ in the CKBG and CKG Lindbladians, we would have to approximate this with a quadrature sum in the quantum computer. A classical result of Fourier analysis \cite[Chapter 2.]{stein2011fourier} states that if a function $f$ is $k$ times continuously differentiable with $f^{(k)}$, its Fourier transform decays as $\hat{f}(t) = \mathcal{O}(|t|^{-(k+1)})$. The Metropolis $\gamma_M$ has a discontinuous first derivative at the kink, and thus its Fourier transform decays only as $\mathcal{O}(1 / t)$. Of course the convolution with the Gaussian filters with width $ \sigma^{-1}$ smoothen this kink and improves the scaling, but the algebraic prefactor persists $\mathcal{O}(\frac{e^{-\sigma^2 t^2}}{\sigma^2 t^2})$. Hence, this algebraic tail requires a larger truncation range and a finer grid (more estimating ancilla qubits) to achieve the same precision as an analytic transition weight.
    \item Finally, the existence of the kink is tightly connected to having a constant piece of the Metropolis function that equates to 1 for $\omega \leq \beta\sigma^2 / 2$. When passing to the time-domain representation of the Lindbladian (the one that runs in the quantum computer) via a Fourier transform, this constant piece produces a $\delta(t)$ singularity in the integral kernel. In practice this singularity has to be tamed by some regularization technique which incurs more costs as well. \todo{refer to Metropolis tame and maybe our taming}
\end{itemize}

But we can have our cake and eat it too, almost that is. We can have close to optimal transitions while not stumbling into all the additional implementation costs. The solution is to smoothen the kink of the Metropolis function $\gamma_M$, which eliminates the first two previously mentioned problems. After a long battle of trying to improve on \cite{chen2023quantum}'s way of regularizing the singularity in the time domain, it seems that the incurring additional costs cannot be made better. The remnants of the battle can be found in [REF to APPENDIX] \todo{appendix where parameter a is back in the lin comb.}.

To designing a better transition function $\gamma$, we will make use of the linear combination that was introduced in \eqref{eq:lin-comb-alpha}. The key is to choose $s \neq 0$ to shift the lower limit of the integration away from the kink, in a tuneable manner, i.e. the larger $s$ is the stronger we smoothen. Explicitly writing out $\gamma_G$ and $g(x)$ yields,
\begin{equation} \label{eq:smooth-metro-int-defi}
    \gamma_M^{(s)}(\omega) = \int_{\frac{\beta\sigma^2}{2}(1 + s)}^\infty \frac{1}{\sqrt{2\pi \sigma_\gamma^2(x)}} \exp\left(\frac{-(\omega + x)^2}{2 \sigma_\gamma^2(x)}\right) \dd x.
\end{equation}
The resulting the transition function cleanly separates into the Metropolis function $\gamma_M$ with an additional factor dependent on $s, \beta, \sigma$:
\begin{equation} \label{eq:smooth-metro}
    \gamma_M^{(s)}(\omega) = \gamma_M^{(s=0)}(\omega) \times \frac{1}{2}\left[ \erfc\left(z_-\right) + e^{\beta |\tilde{\omega}|} \erfc\left(z_+\right) \right]
\end{equation}
with $\tilde{\omega} = \omega + \beta\sigma^2 / 2$ and 
$$
z_\pm := \frac{\beta \sigma \sqrt{s}}{2\sqrt{2}} \pm \frac{|\tilde{\omega}|}{\sigma\sqrt{2s}} 
$$
\todo{check with numerics, if these were the correct factors here.}
Of course, there is a drawback here, since \hyperref[prop:optim-metro]{Proposition~\ref*{prop:optim-metro}} we know that any other choice of transition function $\gamma$ could lead to a smaller spectral gap. More concretely in our case, all $\gamma^{(s)}$ are pointwise smaller or equal to $\gamma_M$ since the right multiplying factor in \eqref{eq:smooth-metro} is $\leq 1$ for $s > 0$. But in practice we only need a small $s$, which is enough to tame quadrature errors (and circuit costs), while keeping almost the same spectral gap of the CKG Lindbladian, see the numerical analysis [REF] \todo{ref}. The parameter $s$ creates a family of Metropolis-like transitions, with $s\rightarrow 0^+$ recovering the Metropolis function $\gamma_M$ (the square bracket with the error functions equals 2 in this limit). Also, we get the Glauber-like transition weights if we choose $s = 2$, which has been found by Chen et al. in \cite{chen2023quantum}. But that is a stark smoothening, which leads to noticeably worse mixing times in the simulations.

Let us state rigorously in what sense the smooth Metropolis $\gamma^{(s > 0)}_M$ improves upon its kinky version $\gamma^{(s = 0)}_M$. First, we have to recall the definition of Gevrey functions \cite{adwan2017global}.
\begin{definition}
    \textup{(Gevrey functions).} Let $\Omega \subseteq \mathbb{R}$. A $C^\infty$ function $h: \Omega \to \mathbb{C}$ is a Gevrey function of order $\kappa \geq 0$ if there exist constants $C_1, C_2 > 0$ such that
    \[
    \|h^{(n)}\|_{L^\infty(\Omega)} \leq C_1\, C_2^n\, n^{n \kappa} \quad \text{for all } n \geq 0,
    \]
    For the given constants $C_1, C_2, \kappa$, the set of Gevrey functions is denoted by $\mathcal{G}_{C_1, C_2}^\kappa(\Omega)$.
\end{definition}
Being in a Gevrey set tells us something about the regularity of a function. For instance, with order $\kappa = 0$ the set of real analytic functions is recovered. Gevrey order with finite $\kappa$ provides a quantitative interpolation: integrals over $\mathbb{R}$ admit quadrature errors decaying as $\exp(-c N^{1 / \kappa})$ under optimal truncation of the Euler-Maclaurin expansion, where $N$ is the number of quadrature nodes \cite{trefethen2014exponentially}.
Conversely, if the integrand has only finite regularity, for instance due to a kink, the quadrature errors converge only at a polynomial rate $\mathcal{O}(N^{-k})$ determined by the order of smoothness $k$ \cite{trefethen2014exponentially}.

In the following proposition we will denote the integrand in \eqref{eq:smooth-metro-int-defi} as $F(\omega, x)$, i.e. $\gamma_M^{(s)}(\omega) = \int_{x_0}^\infty F(\omega, x)\,dx$, with the lower limit $x_0 := \beta \sigma^2 (1 + s) / 2$.
\begin{proposition} \label{prop:smooth-metro-gevrey}
    \textup{(Smooth Metropolis is Gevrey-1/2).} For $s > 0, \beta > 0$, the smooth Metropolis weight $\gamma^{(s)}_M$ defined by \eqref{eq:smooth-metro-int-defi} satisfies:
    \begin{enumerate}[(i)]
        \item For all $n\geq 1$ and $\omega \in \mathbb{R}$:
        \[
        \left|\left(\gamma_M^{(s)}\right)^{(n)}(\omega)\right| \leq \frac{1}{2\pi}\left(\frac{2\beta^2}{s}\right)^{n/2} \Gamma\!\left(\frac{n}{2}\right),
        \]
        with $f^{(n)}$ denoting the $n$-th derivative, and $\Gamma(z) = \int_{0}^{\infty} t^{z - 1} \ee^{-t} \dd t$ is the Euler gamma function.
        \item Consequently, $\gamma_M^{(s)} \in \mathcal{G}^{1/2}_{C_1, C_2}(\mathbb{R})$ with 
        \[
        C_1 = 1, \qquad C_2 = \frac{1}{\sigma\sqrt{\ee\,s}}.
        \]
        %\item For $s = 0$, the Metropolis weight $\gamma_M^{(0)}(\omega) = e^{-\beta\max(\omega + 1/(2\beta),\, 0)}$ is Lipschitz but not $C^1$, hence not in any Gevrey class.
    \end{enumerate}
\end{proposition}
\begin{proof}
    For a fixed $x \geq x_0$, the function $\omega \mapsto F(\omega, x)$ is a Gaussian with mean $-x$ and variance $\sigma_\gamma^2(x) \geq s \sigma^2 / 2 > 0$. Its Fourier transform in $\omega$ is
    \[
    \mathcal{F}[F(\cdot, x)](\xi) = \frac{1}{\sqrt{2\pi}}\,\ee^{\ii x\xi}\, \ee^{-\sigma_\gamma^2(x)\, \xi^2/2}.
    \]
    Substituting $\sigma_\gamma^2(x) = 2x / \beta  - \sigma^2$:
    \[
    \mathcal{F}[F(\cdot, x)](\xi) = \frac{1}{\sqrt{2\pi}}\,\ee^{\ii x\xi - (x/\beta - \sigma^2/2)\,\xi^2} = \frac{1}{\sqrt{2\pi}}\,\ee^{\sigma^2\xi^2/2}\, \ee^{x(\ii\xi - \xi^2/\beta)}.
    \]
    Then for, $n \geq 1$, differentiation under the integral is justified by the locally uniform integrability of $|\partial_\omega^n F(\omega, x)|$ over $x \in [x_0, \infty)$: the Gaussian decay $F(\omega, x) \lesssim e^{-\beta x/4}$ provides the domination, and $\sigma_\gamma > 0$ prevents singularities. Therefore:
    \[
    \left(\gamma_M^{(s)}\right)^{(n)}(\omega) = \int_{x_0}^\infty \frac{\partial^n F}{\partial \omega^n}(\omega, x)\, \dd x.
    \]
    Taking the Fourier transform of both sides in $\omega$ and applying Fubini's theorem:
    We compute the integral for the linear combination. For real $\xi \neq 0$
    \begin{equation*}
        \begin{split}
            \mathcal{F}\left[\left(\gamma_M^{(s)}\right)^{(n)}\right](\xi) &= \frac{(\ii\xi)^n}{\sqrt{2\pi}} \int_{x_0}^\infty \ee^{\sigma^2\xi^2/2}\, \ee^{x(\ii\xi - \xi^2/\beta)}\, \dd x. \\
            &= \frac{(i\xi)^n \cdot \beta}{\sqrt{2\pi}} \cdot \frac{e^{\sigma^2\xi^2/2 + x_0(i\xi - \xi^2/\beta)}}{\xi(\xi - i\beta)}.
        \end{split}
    \end{equation*}
    Where for computing the integral, $\mathrm{Re}(\xi^2/\beta - i\xi) = \xi^2/\beta > 0$, thus we do not run into singularities.
    Substituting $x_0$ and doing some arithmetics we arrive at
    \[
    \mathcal{F}[(\gamma_M^{(s)})^{(n)}](\xi) = \frac{(i\xi)^{n-1} \cdot i\beta}{\sqrt{2\pi}} \cdot \frac{e^{-s\,\sigma^2\xi^2/2\; +\; i\beta\sigma^2(1+s)\xi/2}}{\xi - i\beta},
    \]
    with modulus:
    \[
    \left|\mathcal{F}[(\gamma_M^{(s)})^{(n)}](\xi)\right| = \frac{\beta\, |\xi|^{n-1}}{\sqrt{2\pi}\,\sqrt{\xi^2 + \beta^2}}\; e^{-s\,\sigma^2\xi^2/2}.
    \]
    Since $\mathcal{F}[(\gamma_M^{(s)})^{(n)}] \in L^1(\mathbb{R})$ for $n \geq 1$, we can now apply the Fourier inversion formula and bound the derivatives of the smooth Metropolis weights,
    \[
    \left|(\gamma_M^{(s)})^{(n)}(\omega)\right| \leq \frac{1}{\sqrt{2\pi}} \int_{\mathbb{R}} \left|\mathcal{F}[(\gamma_M^{(s)})^{(n)}](\xi)\right| d\xi = \frac{\beta}{2\pi} \int_{\mathbb{R}} \frac{|\xi|^{n-1}\, e^{-s\,\sigma^2\xi^2/2}}{\sqrt{\xi^2 + \beta^2}}\, d\xi.
    \]
    Using $\sqrt{\xi^2 + \beta^2} \geq \beta$ for all $\xi\in\mathbb{R}$ (which is a loose bound, but a tighter solution still would lead to the same Gevrey order):
    \[
    \left|(\gamma_M^{(s)})^{(n)}(\omega)\right| \leq \frac{1}{2\pi} \int_{\mathbb{R}} |\xi|^{n-1}\, e^{-s\,\sigma^2\xi^2/2}\, d\xi = \frac{1}{\pi} \int_0^\infty \xi^{n-1}\, e^{-s\,\sigma^2\xi^2/2}\, d\xi.
    \]
    This Gaussian integral can be evaluated with $\alpha := s \sigma^2 / 2$ and the identity $\int_0^\infty \xi^{n-1} e^{-\alpha \xi^2}\, d\xi = \frac{1}{2}\,\alpha^{-n/2}\,\Gamma(n/2)$:
    \[
    \left|\left(\gamma_M^{(s)}\right)^{(n)}(\omega)\right| \leq \frac{1}{2\pi}\left(\frac{2}{s\,\sigma^2}\right)^{n/2} \Gamma\!\left(\frac{n}{2}\right),
    \] 
    proving the part $(i)$ of the proposition.
    The last step is to show that 
    \[
    \frac{1}{2\pi}(2/(s\sigma^2))^{n/2}\,\Gamma(n/2) \leq C_1\, C_2^n\, n^{n/2}
    \]
    for suitable constants $C_1, C_2$. Let us handle the different cases of $n$ separately: 
    \begin{itemize}
        \item $n = 0$: $\|\gamma_M^{(s)}\|_\infty \leq 1$ since $\gamma_M^{(s)} \leq \gamma_M^{(s=0)} \leq 1$.
        \item $n = 1$: Substituting $\Gamma(1 / 2) = \sqrt{\pi}$ into (i) gives $\| (\gamma_M^{(s)})'\|_\infty \leq \frac{1}{\sqrt{2\pi} \sigma s}$ 
        \item $n \geq 2$: in these cases the Stirling upper bound gives $\Gamma(z) \leq \sqrt{2\pi} (z / \ee)^z$. Applying this with $z = n /2 \geq 1$ to (i) gives, 
        \[
        \| (\gamma_M^{(s)})^{(n)}\|_\infty \leq \frac{1}{\sqrt{2\pi}}\left(\frac{1}{\ee s \sigma^2}\right)^{n / 2} n^{n / 2}
        \]
    \end{itemize}
    From which we find that $C_1 = 1$ and $C_2 = \frac{1}{\sigma \sqrt{\ee s}}$ (where we could absorb the prefactor of $C_2$, since $1 / \sqrt{2\pi} < C_1$).
\end{proof}
\noindent An important remark here is that the Metropolis function $\gamma_M^{(s = 0)}$, has a left derivative at the kink $\omega + \beta\sigma^2 / 2$ equal to $0$, and a right derivative equal to $-\beta$, which makes it only $C^0$ but not $C^1$, and in particular it is not a Gevrey function of any order. This qualitative difference between transition rates will become important for our quadrature error discussions [REF] \todo{ref}.

\smallskip
\subsection{Improved Kossakowski matrices.}
For bookkeeping and later analysis, it is worth to write down what form the Kossakowski matrix takes if instead of the Gaussian transition $\gamma_G$ we use the smoothened Metropolis functions $\gamma_M^{(s)}$. Similarly to \eqref{eq:kossakowski-gaussian}, integrating over $\omega$ with the kernel consisting of $\gamma_M^{(s)}$ and the Gaussian filters $\hat{f}(\omega - \nu)$, we find, 
\begin{equation} \label{eq:kossakowski-smooth-metro}
    \alpha_{\nu_1, \nu_2}^{(s)} = \frac{1}{2}e^{-\frac{\nu_-^2}{8\sigma^2}}e^{-\frac{\beta}{4}(\nu_+ +|\nu_+|)}\left[\text{erfc}(\zeta_-) + e^{\frac{\beta}{2}|\nu_+|}\text{erfc}(\zeta_+)\right]
\end{equation}
with $\nu_\pm := \nu_1 \pm  \nu_2$ and 
$$
\zeta_\pm := \frac{\beta \sigma \sqrt{1 + s}}{\sqrt{8}} \pm \frac{|\nu_+|}{\sigma\sqrt{8(1 + s)}}.
$$
A few things to note here:
\begin{enumerate}[(i)]
    \item $\alpha^{(s)}_{\nu_1,\nu_2}$ is positive semidefinite, since $\gamma_M^{(s)}\geq 0$ and $\hat{f}\geq 0$, guaranteeing GKSL form for the Lindbladian. It also satisfies the skew-symmetry \eqref{eq:skew-symmetry}, i.e. leads to a KMS detailed balanced Lindbladian as it was promised via linear combinations of detailed balanced Lindbladians.
    \item In the limit of $s\rightarrow 1$ we again recover the Metropolis Kossakowski matrix that Chen et al. found \cite[Appendix B]{chen2023quantum}, exactly how we previously recovered $\gamma_M$ from $\gamma_M^{(s)}$.
    \item In the $\sigma\rightarrow 0$ limit consistently finds the diagonal Kossakowski matrix, i.e. the Davies generator, since the off-diagonal suppression $e^{-\nu_-^2 / (8\sigma^2)}$ forces all terms with $\nu_1 \neq \nu_2$ to vanish.
    \item Due to the factorizing Gaussian filter functions, $\alpha_{\nu_1,\nu_2}^{(s)}$ can be separated into to two parts, that each depend on either $\nu_+$ or $\nu_-$. This will be important for later, when we derive time-domain functions, since it allows for a double Fourier transformation. Note, that this is also true for the case with the Gaussian transition weight \eqref{eq:kossakowski-gaussian}.
    \item Another implication of the factorizing Gaussian filters is a simple shift relation between the Kossakowski matrices for the GNS and KMS detailed balanced Lindbladians (i.e. between CKBG and CKG Lindbladians), analogous to \hyperref[lem:shifted-kms-condition]{Lemma~\ref*{lem:shifted-kms-condition}}:
\end{enumerate}

\begin{corollary} \label{cor:kossakowski-shift}
    \textup{(Kossakowski shift between CKG and CKBG).}
    Let $\tilde{\gamma}(\omega) = \tilde{\gamma}(-\omega)\,e^{-\beta\omega}$, and let
    $\hat{f}$ be the Gaussian filter with width~$\sigma$ from
    \eqref{eq:gaussian-filter-w}.
    Denote by $\alpha^{\textup{KMS}}_{\nu_1,\nu_2}$ the Kossakowski matrix
    of the CKG Lindbladian, which uses the shifted transition weight
    $\gamma(\omega) = \tilde{\gamma}(\omega + \beta\sigma^2/2)$ from
    \hyperref[lem:shifted-kms-condition]{Lemma~\ref*{lem:shifted-kms-condition}}:
    %
    $$
      \alpha^{\textup{KMS}}_{\nu_1,\nu_2}
      \;=\; \int_{-\infty}^{\infty}
        \tilde{\gamma}\!\bigl(\omega + \tfrac{\beta\sigma^2}{2}\bigr)\,
        \hat{f}(\omega - \nu_1)\,\hat{f}(\omega - \nu_2)\,\dd\omega.
    $$
    %
    Then the Kossakowski matrix of the CKBG Lindbladian, which uses
    $\tilde{\gamma}$ directly,
    %
    $$
      \alpha^{\textup{GNS}}_{\nu_1,\nu_2}
      \;=\; \int_{-\infty}^{\infty}
        \tilde{\gamma}(\omega)\,
        \hat{f}(\omega - \nu_1)\,\hat{f}(\omega - \nu_2)\,\dd\omega,
    $$
    %
    satisfies
    %
    \begin{equation} \label{eq:kossakowski-shift}
      \alpha^{\textup{GNS}}_{\nu_1,\nu_2}
      \;=\;
      \alpha^{\textup{KMS}}_{\,\nu_1 - \beta\sigma^2/2,\;\,\nu_2 - \beta\sigma^2/2}.
    \end{equation}
    %
    Equivalently, with $\nu_{\pm} := \nu_1 \pm \nu_2$, the off-diagonal
    suppression factor $e^{-\nu_-^2/(8\sigma^2)}$ is invariant under the
    shift, while every occurrence of $\nu_+$ in
    $\alpha^{\textup{KMS}}$ is replaced by $\nu_+ - \beta\sigma^2$.
\end{corollary}

\begin{proof}
    Substitute $u = \omega + \beta\sigma^2/2$ in the defining integral of
    $\alpha^{\textup{KMS}}_{\nu_1 - \beta\sigma^2/2,\;\nu_2 - \beta\sigma^2/2}$:
    %
    \begin{align*}
      \alpha^{\textup{KMS}}_{\nu_1 - \beta\sigma^2/2,\;\nu_2 - \beta\sigma^2/2}
      &= \int_{-\infty}^{\infty}
        \tilde{\gamma}\!\bigl(\omega + \tfrac{\beta\sigma^2}{2}\bigr)\,
        \hat{f}\!\bigl(\omega - \nu_1 + \tfrac{\beta\sigma^2}{2}\bigr)\,
        \hat{f}\!\bigl(\omega - \nu_2 + \tfrac{\beta\sigma^2}{2}\bigr)\,
        \dd\omega \\[4pt]
      &= \int_{-\infty}^{\infty}
        \tilde{\gamma}(u)\,
        \hat{f}(u - \nu_1)\,\hat{f}(u - \nu_2)\,\dd u
      \;=\; \alpha^{\textup{GNS}}_{\nu_1,\nu_2}.
    \end{align*}
    %
    For the $\nu_\pm$ decomposition, the Gaussian filter gives
    %
    $$
      \hat{f}(\omega - \nu_1)\,\hat{f}(\omega - \nu_2)
      \;\propto\;
      \exp\!\Bigl(-\frac{(\omega - \nu_+/2)^2}{2\sigma^2}\Bigr)\,
      \exp\!\Bigl(-\frac{\nu_-^2}{8\sigma^2}\Bigr).
    $$
    %
    The shift $(\nu_1, \nu_2) \mapsto
    (\nu_1 - \beta\sigma^2/2,\;\nu_2 - \beta\sigma^2/2)$ sends
    $\nu_- \mapsto \nu_-$ and $\nu_+ \mapsto \nu_+ - \beta\sigma^2$,
    leaving the $\nu_-$-dependent factor unchanged, and shifting $\nu_+$ as prescribed.
\end{proof}
\noindent Though it is a simple corollary to \hyperref[lem:shifted-kms-condition]{Lemma~\ref*{lem:shifted-kms-condition}}, we will make use of it in our numerical analysis, where we look at the mixing performances of the approximately GNS detailed balanced CKBG Lindbladian [REF to num analysis of GNS] \todo{ref}.

\smallskip
\subsection{Time, time, time.} \todo{add norms and LCU subnormalizations here.}
As we have seen, the energy-domain representation is natural for analysing detailed balance, but if we want to run any algorithm in a quantum computer, we have to understand and analyse the given algorithm also in the time-domain. We have already seen one of the core subroutines in the time-domain that \cite{chen2023efficient,chen2023quantum} used, the Operator Fourier Transform, \eqref{eq:oft} and its circuit form [REF to Quantum Circuits chapter] \todo{ref}. With this the dissipative part of the Lindbladians is more or less covered, though how one implements the corresponding open quantum system evolution, $\ee^{t\mathcal{L}}(\rho)$, will be left for the next paragraph.

Then, we are left to work out what form the coherent part $G$ \eqref{eq:B-bohr-domain} takes in the time-domain. The factorization property that we have noted in the previous chapter is the key to derive the time-domain representation. For any transition weight $\gamma$ (e.g. $\gamma_G, \gamma_M^{(s)}$) and with the Gaussian filters \eqref{eq:gaussian-filter-w} the Kossakowski matrix splits as
\begin{equation}
    \frac{1}{2\pi}\alpha_{\nu_1,\nu_2} = \hat{f}_+(\nu_+)\,\hat{f}_-(\nu_-), \qquad \nu_{\pm} := \nu_1 \pm \nu_2,
\end{equation}
where the $\hat{f}_-$ factor encodes the off-diagonal suppression and the coherent-term structure ($\tanh$ piece), while $\hat{f}_+$ carries all dependence on the transition weight $\gamma$. By \cite[Corollary~A.1]{chen2023efficient}, the two-dimensional Fourier transform over the Bohr frequencies $(\nu_1,\nu_2)$ then decouples into two independent one-dimensional transforms, yielding a nested double integral in the time domain. 

After executing the two independent Fourier transforms on $\hat{f}_\pm(\nu_\pm)$, we find that the coherent term $G$ takes the following general form in the time-domain \cite[Corollary III.1]{chen2023efficient}:
\begin{equation} \label{eq:coherent-time-domain}
    B = \sum_{a \in \mathcal{A}} \int_{-\infty}^{\infty} b_-(t)\, e^{-iHt / \sigma} \left[\int_{-\infty}^{\infty} b_+(t')\, A^{a\dagger}(\beta t')\, A^a(-\beta t')\,\dd t' \right] e^{iHt/\sigma}\,\dd t
\end{equation}
Here $b_-(t)$ comes from $\hat{f}_-(\nu_-)$, and it is universal across all $\gamma$ choices. It takes the following form:
\begin{equation} \label{eq:b_minus}
    b_-(t) = \frac{2\sqrt{\pi}}{\beta \sigma}e^{\beta^2\sigma^2/8}\left[\frac{1}{\cosh(\frac{2\pi t}{\beta \sigma})}\ast_t \sin(-\beta \sigma t) e^{-2t^2}\right] \quad\text{with}\quad \|b_-\|_1 \leq \frac{\pi}{\sqrt{2}}e^{\beta^2 \sigma^2 / 8},
\end{equation}
where $\ast_t$ denotes the convolution over the variable $t$, 
$$
    (f \ast_t g)(t) := \int_{-\infty}^{\infty} f(s) g(t - s)\dd s.
$$
Importantly, $b_-(t)$ is real function that has a support over $t = \mathcal{O}(\beta \sigma)$ due to the width of the $1 / \cosh$ kernel. The time evolution of the outer integral is $e^{\pm \beta H t / \sigma}$, which implies Hamiltonian simulation time of $\mathcal{O}(\beta)$.

Since $b_+$ depends on $\gamma$, we will have to look at the Gaussian $\gamma_G$ and the smooth Metropolis $\gamma_M^{(s)}$ separately. The Gaussian case is definitely the simpler one out of the two. It takes the following form:
\begin{equation} \label{eq:b_plus-gauss}
    b_+(t) = \frac{2 \sigma_\gamma}{\pi\sqrt{\pi}}e^{-4t^2\beta^2 \omega_\gamma - 2it\beta\omega_\gamma} \quad\text{with}\quad \|b_+\|_1 = \frac{\beta}{\omega_\gamma} \sqrt{\frac{\sigma_\gamma}{2\pi}}.
\end{equation}
This is a rapidly decaying function with width $\mathcal{O}(1 / \sqrt{\beta \omega_\gamma}) = \mathcal{O}(1)$, if we use that $\omega_\gamma$ has a natural unit of $1 / \beta$. Hence, the inner integral also needs an $\mathcal{O}(\beta)$ amount of Hamiltonian simulation time in this case.

The Metropolis case is slightly more complicated. But this should not surprise us after the foreshadowing in the discussions of \eqref{eq:gamma_M-from-lin-comb-of-G}. Since in the energy-domain the Metropolis function has a constant piece for most of the negative energies $\omega < -\beta\sigma^2 / 2$, the Fourier transform produces a $\delta(t)$ contribution in the time domain. In the smooth case $s \neq 0$, $\gamma_M^{(s)}$ is only asymptotically reaching that flat $\gamma = 1$ portion. Nevertheless, the near singularity behaviour also persists in these smoother cases. Despite the divergent $\ell_1$ norm of $b_+^{(s)}$ the operator $B_M$ is actually bounded. The key idea is that the $1/t$ singularity, when sandwiched between operator valued functions $A^\dagger(t)A(t)$, can be absorbed into $H \textup{sinc}(2Ht)$ terms that are bounded at $t=0$, \cite[Proposition B.2]{chen2023efficient}. Though mathematically this is sound, and implementing it would also be asymptotically efficient with the Quantum Singular Value Transformation (QSVT), but would also be cumbersome. Moreover, it would be a significant change to the LCU quantum circuit [ref]\todo{ref to quantum circuits section} that we use for the Gaussian case coherent term $B_G$. 
For LCU circuits we care about the subnormalization of the operator we want to implement, as it determines the quantum circuit cost. For the exact Metropolis term $B_M$ it bounded by \cite[Proposition B.2]{chen2023efficient}
\begin{equation}
    \|B^{(s=0)}\| \equiv\|B_M\| \leq \left\|\sum_a A^{a\dagger}A^a\right\| \mathcal{O}\left(\log\left(\beta \| H \|\right)\right)
\end{equation}
Instead of implementing the Metropolis coherent term exactly, we can regularize $b_+^{(s)}$ around $t=0$ by replacing the singularity near $|t|\leq\eta$ with a smooth cutoff \cite[Proposition B.1]{chen2023efficient}, which leads to the following expression:
\begin{equation} \label{eq:b_plus-s-eta}
    b_+^{(s, \eta)}(t) = \frac{1}{2\sqrt{2}\pi^2}\frac{\ee^{-\sigma^2\beta^2 t(2t + \ii)(1 + s)} + \mathbb{I}(|t|\leq\eta)\ii(2t + \ii)}{t(2t + \ii)}, 
\end{equation}
with the $\ell_1$ norm, 
$$
\|b_+^{(s, \eta)}\|_1 < \frac{2}{5 \pi^{3/2}} - \frac{\ln(1 + s)}{2\sqrt{2}\pi^2}+ \frac{1}{\sqrt{2}\pi^2}\ln(1 / \eta).
$$

\todo{write somewhere: B M as A'A additional term in it}
Consistency check: Chen's norm bound is recovered for $s=0$, albeit with the corrected prefactors. There was a slight prefactor discrepancy in their paper, the corrected norm is smaller by a factor of $1 / \pi^{3 / 2}$, which improves their implementation by a constant factor. For the $s$-family the log-term is negative for $s > 0$, the norm is strictly smaller than the Metropolis case with $s = 0$, which also helps by a $\ln(1 + s)$ decrease for the LCU implementations. 

In \eqref{eq:b_plus-s-eta} it also becomes clear that smoothing with $s$ indeed does not shift the poles of the expression above $(t = 0, -\ii / 2)$, and thus the regularization is still necessary. 
In \cite[Corollary III.2]{chen2023efficient} they have shown that the error between the regularized coherent term with $B^{(s = 0)}_\eta$ and the non-regularized one $B^{(s = 0)}$ is bounded. Their result also adapts to the smooth Metropolis variants since the error is purely due to the regularized region $|t| \leq \eta$ where the smoothening plays no role, 
\begin{equation}
    \|B^{(s)} - B_\eta^{(s)} \| \leq \left\|\sum_a A^{a\dagger}A^a\right\| \min\left(\frac{\eta\beta\|H\|}{\sqrt{2}\pi}, \mathcal{O}\left(\left(\eta \beta \|H\|\right)^3\right)\right).
\end{equation}
To achieve an $\varepsilon$ accuracy in this approximation, we would have to set
\begin{equation}
    \eta = \frac{\varepsilon}{\beta \|H\| \|\sum_a A^{a\dagger}A^a\|},
\end{equation}
which in turn leads to the following subnormalization for the LCU implementation
\begin{equation} \label{eq:subnorm-eta-B}
    \|b_-\|_1 \|b_+^{(s,\eta)}\|_1 = \mathcal{O}\left(\log\frac{\beta \|H\| \|\sum_a A^{a\dagger}A^a\| }{\varepsilon}\right).
\end{equation}
Combining it all, we end up with a Hamiltonian simulation time required for the coherent term $B^{(s, \eta)}$ of $\mathcal{O}(\beta \log(1 / \varepsilon))$. This is a $\log$ factor more than what we have seen for the Gaussian case

We could try to use different regularization techniques, like the $\ii \varepsilon$ prescription from quantum field theory, i.e. to shift the pole away from the real time values to the imaginary ones, such that the integral of the time-domain would not run into it. In our case this would materialize by an exponentially decaying factor multiplying the weights of the linear combination, i.e. $\ee^{- \varepsilon \beta x} g(x)$. But while it does make both poles imaginary, avoiding Chen's regularization and its additional $\log$ costs, it does it at the price of an $e^{-\varepsilon \beta / 2}$ reduction in the transition weight near the kink, which degrades the spectral gap. As it turns out, this degradation would lead to higher costs asymptotically, than the $\eta$ regularization \eqref{eq:subnorm-eta-B}. Hence, our regularization has been tucked away in the appendix as a failed but hopefully constructive attempt \todo{ref}. For the main text, we will stick to the above form for $b_+^{(s)}$ \eqref{eq:b_plus-s-eta}.

\label{sec:quad-errors-diss}
\smallskip
\subsection{Quadrature errors for $\mathcal{L}_\text{diss}$.}
The next step towards implementation is the realization that every integral expression in the CKG Lindbladian \eqref{eq:ckg-L}, both the coherent term $B$ and the dissipative terms, have to be approximated via quadrature sums in the quantum computer. This introduces two types of errors that are interconnected with each other:
\begin{enumerate}[(i)]
    \item \emph{Truncation errors} occur due to the restriction of infinite integrals to a finite window $t\in[-T/2, T/2)$ and $\omega \in [-W/2, W/2)$. These errors are controlled by the decaying tail of the respective integrands.
    \item \emph{Discretization / Aliasing errors} occur since the integrals are replaced by a sum on the finite grids. These are controlled by the smoothness of the integrand in the given truncated domain.
\end{enumerate}
We will tackle these errors for the dissipative and coherent parts separately. First let us cast the CKG Lindbladian \eqref{eq:ckg-L} into its discretized form. The dissipative part becomes,
\begin{equation} \label{eq:discr-dissiative-L}
    \mathcal{\bar{L}}_\text{diss}(\cdot) = \sum_a \sum_{\bar{\omega} \in S_{\omega_0}} \gamma(\bar{\omega}) \hat{A}(\bar{\omega})(\cdot)\hat{A}(\bar{\omega})^\dagger - \frac{1}{2}\left\{\hat{A}(\bar{\omega})^\dagger \hat{A}(\bar{\omega}), \cdot\right\},
\end{equation}
with the discretized OFT, 
\begin{equation} \label{eq:oft-discr}
    \hat{A}(\bar{\omega}) = \frac{1}{\sqrt{N}}\sum_{\bar{t}\in S_{t_0}} \bar{f}(\bar{t})\ee^{-\ii \bar{\omega} \bar{t}} A(\bar{t}) =  \frac{1}{\sqrt{N}}\sum_{\bar{t}\in S_{t_0}} \bar{f}(\bar{t})\ee^{-\ii \bar{\omega} \bar{t}} \ee^{\ii H \bar{t}} A \ee^{-\ii H \bar{t}}.
\end{equation}
Here, we merged the Riemann-sum prefactor $t_0$ into the \emph{discretized} filter function $\bar{f}(t) := \sqrt{t_0} f(t)$. This lets us work with the natural Discrete Fourier Transform prefactor $1 / \sqrt{N}$, with $N$ the number of grid points. (It is also the reason why there is no explicit $\omega_0$ in \eqref{eq:discr-dissiative-L}.)
The infinite integrals are now replaced by discretized sums, where we will denote the discrete values by $\bar{t}, \bar{\omega}$ whose sums run over
\begin{equation*}
    \begin{split}
    S^{\lceil N \rfloor} &:= \left\{-\lceil(N-1)/2\rceil, \dots, -1, 0, 1, \dots, \lfloor (N - 1) / 2 \rfloor\right\} \\
    S^{\lceil N \rfloor}_{\omega_0} &:= \omega_0 S^{\lceil N \rfloor}, \quad S^{\lceil N \rfloor}_{t_0} := t_0 S^{\lceil N \rfloor},
\end{split}
\end{equation*}
with $\omega_0t_0 = \frac{2\pi}{N}$, where $N = 2^r$ is determined by the number of estimating qubits $r$ in the quantum circuit. Hence, the resolution can be tuned by choosing a larger $r$. But there is an interplay between time- and energy-domains due to the Fourier duality. Since we have integrals in both domains, the ideal scenario is when the integrands and their Fourier transforms are both well-behaving, i.e. smooth and rapidly decaying, functions.

In order to quantify the quadrature errors for the dissipative part of the Lindbladian (but also for the coherent part), we follow the strategy of \cite[Appendix C]{chen2023quantum}. All parts of the algorithm rely on the OFT, thus it is essential to understand how it behaves under discretization / truncation \eqref{eq:oft-discr}. The central difficulty here can be immediately revealed if we look at its decomposition $\hat{A}(\omega) = \sum_{\nu_\in B(H)} \hat{f}(\omega - \nu) A_\nu$. A sum over the Bohr frequencies $B(H)$ means, that a quadrature error in $\hat{f}$ would aggregate over all pairs of Bohr frequencies $|B(H)|^2$ -- a potentially exponentially large number. 

Chen et al. circumvent this problem by introducing a \emph{rounding Hamiltonian} $\bar{H}$ whose eigenvalues are the result of rounding each eigenvalue of $H$ to the nearest integer multiple of some $\eta >0$. Then $\|H - \bar{H}\| \leq \eta$ and $|B(\bar{H})| \leq 4 \|H\| / \eta + 1$. To be clear, this is only a proof technique, neither the quantum circuit nor our numerical simulations ever use $\bar{H}$; both work with the true (but rescaled to $\|H\| = 0.5$ for the QFT) Hamiltonian $H$.

Then, the total error between the continuous $\mathcal{L}_\text{diss}$ and the discretized $\mathcal{\bar{L}}_\text{diss}$ is bounded by the following triangle inequality,
\begin{equation} \label{eq:main-lindbladian-tri-ineq-discr}
    \begin{split}
    \|\mathcal{L}_{(f,H)} &- \bar{\mathcal{L}}_{(f,H)}\|_{1 \to 1}
\;\leq\; \\
&\underbrace{\|\mathcal{L}_{(f,H)} - \mathcal{L}_{(f,\bar{H})}\|_{1 \to 1}}_{\text{(i) continuous: }H \to \bar{H}}
+\;
\underbrace{\|\mathcal{L}_{(f,\bar{H})} - \bar{\mathcal{L}}_{(f,\bar{H})}\|_{1 \to 1}}_{\text{(ii) continuous} \leftrightarrow \text{discrete for }\bar{H}}
+\;
\underbrace{\|\bar{\mathcal{L}}_{(f,\bar{H})} - \bar{\mathcal{L}}_{(f,H)}\|_{1 \to 1}}_{\text{(iii) discrete: }\bar{H} \to H}.
\end{split}
\end{equation}
Terms (i) and (iii) are perturbation bounds comparing two Lindbladians that use the same grid but different Hamiltonians. They are controlled by \cite[Corollary A.1, A.2]{chen2023quantum},
\begin{equation}
    \text{(i), (iii)} \;\leq\; 
8\big(\|f - f_T\|_2 + T\|H - \bar{H}\|\big) 
\;\leq\; 8\big(\|f - f_T\|_2 + T\eta\big),
\end{equation}
where $f_T(t) := f(t) \mathbf{1}(t \in [-T/2, T/2))$ captures the truncation error. These terms are independent of the choice of $\gamma$, and depend only on the Gaussian filter tail and the rounding resolution.

The term (ii) is where the smoothness of $\gamma$ enters. By \cite[Lemma C.1]{chen2023quantum}, the operator error decomposes into $|B(\bar{H})|^2$ - many scalar integrals of the form
\begin{equation}
    \begin{split}
        &\left|\alpha_{\nu_1, \nu_2} - \bar{\alpha}_{\nu_1, \nu_2}\right| = \\
        &\left|\int_{-\infty}^{\infty} \gamma(\omega)\,\hat{f}^*(\omega - \nu_1)\,\hat{f}(\omega - \nu_2)\,\mathrm{d}\omega
\;-\;
\sum_{\bar{\omega} \in S_{\omega_0}} 
\tilde{\gamma}(\bar{\omega})\,
\bar{\mathcal{F}}\!\big[\bar{f}_T \cdot e^{i\nu_1\bar{t}}\big]^*\!(\bar{\omega})\;
\bar{\mathcal{F}}\!\big[\bar{f}_T \cdot e^{i\nu_2\bar{t}}\big](\bar{\omega})
        \right|,
    \end{split}
\end{equation}
with the discrete Fourier transform $\mathcal{\bar{F}}$.
In other words, the problem reduces to how well we can estimate the exact Kossakowski matrix $\alpha_{\nu_1, \nu_2}$ by quadrature sums. The norm of this error can be split further into three parts via another triangle inequality: frequency-tail truncation due to the $W$ window, i.e. $\gamma_W(\omega) := \gamma(\omega)\mathbf{1}(\omega \in [-W/2, W/2))$ in both the continuous and discretized settings; and the Riemann-sum error for the truncated integral. The convergence of the Riemann sum is governed by the regularity of the integrand 
\begin{equation}\label{eq:g-integrand}
    g_{\nu_1, \nu_2}(\omega) = \gamma(\omega)\hat{f}^*(\omega - \nu_2) \hat{f}(\omega - \nu_1).
\end{equation}
Since $\hat{f}$ is our smooth Gaussian filter, all comes down to $\gamma$ and its smoothness. 

We now go through each choice for the transition weight $\gamma$, clarifying what regularity each possesses and what quadrature error scaling it implies for their respective dissipative Lindbladians. In each case, we require the total discretization error to satisfy the following bound,
\begin{equation}
    \|\mathcal{L}^{(f,H)}_\text{diss} - \bar{\mathcal{L}}^{(f,H)}_\text{diss}\|_{1\to 1} \leq \varepsilon.
\end{equation}

\begin{enumerate}[(a)]
    \item \textbf{Gaussian weight $\gamma_G$}
    
    For the Gaussian weight, the integrand \eqref{eq:g-integrand} in (ii) is a product of Gaussians. Since the Gaussian is Gevrey-1/2 \cite[Proposition B.1]{hoepfner2019global} and self-dual under the Fourier transform, the Riemann sum converges exponentially in the number of grid points $N$. Likewise, the Gaussian filter tail error $\|f - f_T\|_2$ decays as $\mathcal{O}(e^{-T^2/2\sigma^2})$, and so does the frequency-domain tail \cite[Proposition A.9]{chen2023quantum}. Then, for the dissipative part of the Lindbladian to achieve an $\varepsilon$ precision, it requires 
    $$
    r = \mathcal{O}\!\left(\log\!\left(|A|(\|H\|{+}1)(\sigma{+}1/\sigma)(\beta{+}1)\right) + \log^2(1/\varepsilon)\right)
    $$
    estimating qubits, i.e. a polylogarithmic in $1/\varepsilon$ \cite[Corollary C.2]{chen2023quantum}. The dependence on $\log(1/\varepsilon)$ arises from the interplay between the truncation window $T = \Theta(\sigma \sqrt{\log(1/\varepsilon)})$, the frequency bandwidth $W = \Theta(\sigma^{-1}\sqrt{\log(1/\varepsilon)} + 2\|H\|)$, and the Fourier label condition on $\omega_0, t_0$. 
    
   That $\log(1/\varepsilon)$ is optimal follows from the classical uncertainty principles. Any filter achieving $\varepsilon$ truncation error on both time and frequency domains simultaneously, the product $T\cdot W$ scales at least as $\log(1 \varepsilon)$. That the Gaussian uniquely saturates this uncertainty bound is a classical result, see \cite{folland1997uncertainty}.
    
    \item \textbf{Metropolis weight $\gamma_M^{(s=0)}$}
    
    For the Metropolis weight $\gamma_M$ (with $s = 0$), the Lipschitz kink at $\omega = -\beta\sigma^2 / 2$, with Lipschitz constant $\beta$, limits $\gamma_M$ to $C^{0, 1}$ (in Hölder's notation), i.e. Lipschitz but not $C^1$. 
    
    By \cite[Proposition C.1]{chen2023quantum}, the Lipschitz constant $\beta$ enters the grid spacing constraint as
    $$
        \omega_0 = \mathcal{O}\!\left(\frac{\varepsilon}{\beta\, T\, |B(\bar{H})|^2}\right),
    $$
    which forces the number of estimating qubits to scale polynomially in $1 / varepsilon$ \cite[Corollary C.2]{chen2023quantum}: 
    $$
        r = \mathcal{O}\left(\log\tfrac{|A|(\|H\|+1)(\sigma + 1/\sigma)(\beta+1)}{\varepsilon}\right).
    $$
    Thus, the kink causes a qualitative degradation: polynomial many estimating qubits (rather than polylogarithmic) are needed to achieve the same precision.
    
    To see the mechanism concretely: after splitting the integral at $\omega_\text{kink} = -\beta\sigma^2/2$ and applying the Euler-Maclaurin formula on each half, the trapezoidal error for the $(\nu_1, \nu_2)$ Kossakowski entry is 
    $$
       \frac{\omega_0^2}{12}\,|\Delta g'| = \frac{\beta\,\omega_0^2}{12}\,|\hat{f}(\omega_{\mathrm{kink}} - \nu_1)|\hat{f}(\omega_{\mathrm{kink}} - \nu_2)|, 
    $$
    Summing over rounded Bohr-frequency pairs via \cite[Lemma C.1]{chen2023quantum} and using that $\sum_\nu \|A_\nu\|^2 \leq 1$ together with $|\hat{f}| \leq \|\hat{f}\|_\infty$ this aggregates to a total quadrature error of $\mathcal{O}(\beta \omega_0^2)$. Every other error contributions in the triangle inequalities (e.g. Gaussian aliasing, time truncation, spectral leakage from off-grid Bohr frequencies) otherwise converges exponentially in $1 / \omega_0$, which makes this algebraic quadrature error scaling stick out.

    \item \textbf{Metropolis weight $\gamma_M^{(s > 0)}$}
    
    In \hyperref[prop:smooth-metro-gevrey]{Proposition~\ref*{prop:smooth-metro-gevrey}}, we have shown that the smooth Metropolis function with $s > 0$ belongs to the Gevrey-1/2 class, the same regularity class as the Gaussian transition weights $\gamma_G$, albeit with different Gevrey constants $C_1, C_2$.
    
    An important property of Gevrey functions, is that their product is also Gevrey \cite[Lemma 26]{ding2025efficient}. Consequently, the Kossakowski integrand $g_{\nu_1, \nu_2}$ as a product of Gevrey-1/2 functions -- smooth Metropolis with $C_{1,2}^{(\gamma)}$ and Gaussian filters $\hat{f}$ with $C_{1,2}^{(\hat{f})}$ -- satisfies
    \begin{equation} \label{eq:gevrey-product-integrand}
        \begin{split}
            \hat{g}_{\nu_1,\nu_2} \in G^{1/2}_{C_1^{(\hat{g})},\, C_2^{(\hat{g})}} \qquad\text{with}& \\
            C_2^{(\hat{g})} = C_2^{(\gamma)} + 2\,C_2^{(\hat{f})} = \frac{1}{\sigma\sqrt{\ee\,s}} + \frac{\sqrt{2}}{\sigma \sqrt{\ee}} \quad\text{and}&\quad C_1^{(\hat{g})} = C_1^{(\gamma)}\left(C_1^{(\hat{f})}\right)^2.
        \end{split}
    \end{equation}
    Where $C_1^{(\hat{g})}$ is just the squared normalization of the Gaussian filters $\hat{f}$, since $C_1^{(\gamma)} = 1$.

    The trapezoidal quadrature error for this integrand is governed by the Poisson summation formula \cite[Chapter 5]{trefethen2014exponentially}, specifically for $g \in C^\infty(\mathbb{R}) \cap L^1(\mathbb{R})$ with all derivatives in $L^1(\mathbb{R})$ is
    \begin{equation}
        E_{\nu_1,\nu_2} := \omega_0 \sum_{n \in \mathbb{Z}} \hat{g}_{\nu_1,\nu_2}(n\omega_0) - \int_{\mathbb{R}} \hat{g}_{\nu_1,\nu_2}(\omega)\,\dd\omega = \sum_{k \neq 0} g_{\nu_1,\nu_2}\!\left(\frac{2\pi k}{\omega_0}\right).
    \end{equation}
    Bounding the Fourier transform $|g(\xi)|$ via $2K$-fold integration by parts gives $|g(\xi)| \leq |\xi|^{-2K}\,\|\hat{g}^{(2K)}\|_{L^1}$, hence
    \begin{equation} \label{eq:remainder-g-discr-1}
        |E_{\nu_1,\nu_2}| \leq 2\sum_{k=1}^{\infty}\left(\frac{\omega_0}{2\pi k}\right)^{2K}\|\hat{g}^{(2K)}\|_{L^1} \leq 4\left(\frac{\omega_0}{2\pi}\right)^{2K}\|\hat{g}^{(2K)}\|_{L^1},
    \end{equation}
    where we used $\sum_{k=1}^\infty k^{-2K} \leq \zeta(2) \leq 2$ for $K \geq 1$. 

    But in its current form with $L^1$ we cannot use Gevrey properties yet which are formulated in $L^\infty$. Since the Gaussian filter factor confine the integrand $\hat{g}_{\nu_1, \nu_2}$ to an effective support of width $W_\text{eff} = \mathcal{O}(\sigma)$ (which is $K$ independent), we have $\|\hat{g}^{(2K)}\|_{L^1} \leq W_\mathrm{eff}\,\|\hat{g}^{(2K)}\|_{L^\infty}$. The Gevrey-1/2 bound from $g \in \mathcal{G}^{1/2}_{C_1^{(\hat{g})}, C_2^{(\hat{g})}}$ gives 
    $$
       \|\hat{g}^{(2K)}\|_{L^\infty} \leq C_1^{(\hat{g})}\,(C_2^{(\hat{g})})^{2K}\,(2K)^{K}. 
    $$
    Combining it with \eqref{eq:remainder-g-discr-1}, 
    \begin{equation}
        |E_{\nu_1,\nu_2}| \leq 4\,W_\mathrm{eff}\,C_1^{(\hat{g})}\left(\frac{C_2^{(\hat{g})}\,\omega_0}{2\pi}\right)^{2K}(2K)^K.
    \end{equation}
    Then, we can optimize in $K$ by the standard saddle-point method by first setting 
    $$
        \Phi(K) := 2K\ln\!\left(\frac{C_2^{(g)}\omega_0}{2\pi}\right) + K\ln(2K),
    $$
    and then using the condition $\Phi'(K^*) = 0$ leads to 
    $$
        2\ln\!\left(\frac{C_2^{(g)}\omega_0}{2\pi}\right) + \ln(2K^*) + 1 = 0 \qquad\Longrightarrow\qquad 2K^* = \frac{1}{e}\!\left(\frac{2\pi}{C_2^{(g)}\omega_0}\right)^{\!2}.
    $$
    Inserting $K^*$ back we find that the optimum evaluates to
    $$
        \Phi(K^*) = -K^* = -\frac{2\pi^2}{e\,(C_2^{(g)}\omega_0)^2}.
    $$
    Hence, the per $(\nu_1, \nu_2)$ entry Riemann-sum error satisfies:
    \begin{equation}\label{eq:per-entry-error-smooth-metro}
    |E_{\nu_1,\nu_2}| = \mathcal{O}\!\left(\exp\!\left(-\frac{2\pi^2}{e\,(C_2^{(g)}\,\omega_0)^2}\right)\right).
    \end{equation}

    In the regime $s \ll 1$ (close to the Metropolis limit), the transition weight contribution $C_2^{(\gamma)} = 1/(\sigma\sqrt{es})$ dominates the filter contribution $2C_2^{(\hat{f})} = \sqrt{2}/(\sigma\sqrt{e})$, since $C_2^{(\gamma)}/2C_2^{(\hat{f})} = 1/(\sqrt{2s}) \gg 1$. In this regime $C_2^{(g)} \approx C_2^{(\gamma)}$ and the bound simplifies to
    \begin{equation}\label{eq:smooth-metro-dominant-regime}
        |E_{\nu_1,\nu_2}| = \mathcal{O}\!\left(\exp\!\left(-\frac{2\pi^2\,s\,\sigma^2}{\omega_0^2}\right)\right),
    \end{equation}
    where the factor $e$ from $C_2^{(\gamma)} = 1/(\sigma\sqrt{es})$ cancels the $1/e$ from the saddle-point optimization. 

    For the general regime where both $C_2^{(\gamma)}$ and $C_2^{(\hat{f})}$ contribute comparably, the full additive constant $C_2^{(g)} = C_2^{(\gamma)} + 2C_2^{(\hat{f})}$ must be retained. However, in general we have found that the smooth Metropolis requires the same polylogarithmic estimating qubit scaling as the Gaussian case (a), 
    \begin{equation}
        r = \mathcal{O}\!\left(\log\!\left(|A|(\|H\|{+}1)(\sigma{+}1/\sigma)(\beta{+}1)\right) + \frac{1}{s}\log^2(1/\varepsilon)\right).
    \end{equation}

    \begin{table}[ht]
    \centering
    \begin{tabular}{lll}
    \toprule
    Transition weight & Regularity & Quadrature error in $\omega_0$ \\
    \midrule
    Gaussian $\gamma_G$ & Gevrey-$1/2$ & $\mathcal{O}(e^{-c_G/\omega_0^2})$ — Gaussian \\
    Smooth Metropolis $\gamma_M^{(s)}$, $s > 0$ & Gevrey-$1/2$ & $\mathcal{O}(e^{-c_s/\omega_0^2})$ — Gaussian ($c_s \leq c_G$) \\
    Metropolis $\gamma_M$ ($s = 0$) & Lipschitz, not $C^1$ & $\mathcal{O}(\beta\,\omega_0^2)$ — algebraic \\
    \bottomrule
    \end{tabular}
    \caption{Convergence hierarchy of the trapezoidal quadrature error for the three transition weights.}
    \label{tab:quadrature-hierarchy}
    \end{table}
\end{enumerate}

The above discussions can be summarized in a table of the convergence hierarchy \hyperref[tab:quadrature-hierarchy]{Table~\ref*{tab:quadrature-hierarchy}}. The Gaussian and smooth Metropolis fall in the same qualitative convergence class. However, the Gaussian case with $\gamma_G$ has an optimal constant $c_G$ because the Kossakowski integrand $\hat{g}_{\nu_1, \nu_2} = \gamma_G \cdot |\hat{f}|^2$ is itself a Gaussian, whose Poisson summation aliases are exactly Gaussian and therefore optimally localized. For the smooth Metropolis, the integrand is note a pure Gaussian, so $c_s$ is generally smaller. A qualitative difference can be found between the kinky Metropolis $\gamma_M$ with $s=0$ and its smooth variant: the latter converges super-algebraically (restoring polylogarithmic qubit scaling), while the former is limited to an algebraic convergence due to its derivative discontinuity. Thus, for the dissipative Lindbladian, the smooth Metropolis $\gamma_M^{(s)}$ with $s > 0$ needs less amount of estimating qubits in general, while maintaining the broad Metropolis-like transition window, which keeps the spectral gap (and thus mixing time) favourable.

% Remark (comparison with DLL24). Ding, Li, and Lin [DLL24, Proposition 16] establish an analogous exponential quadrature rate $\exp(-c'((M{-}1)\tau)^{1/s}/(C\,e))$ for their Lindbladian, using compactly supported Gevrey filters and the Poisson summation formula for exact alias cancellation [DLL24, Lemma 31]. Their proof does not use Euler–Maclaurin optimization; instead, compact support gives exact cancellation of all aliases within the support bandwidth, so only the Paley–Wiener tail of the Fourier-domain kernel contributes [DLL24, Lemma 29]. Our setting differs: the Gaussian filter $\hat{f}$ is not compactly supported, so the DLL24 alias-cancellation technique yields no asymptotic gain over the CKBG23 rounding approach for our case ($a = 0$). However, both routes yield the same qualitative conclusion — Gevrey-$1/2$ regularity implies super-algebraic (Gaussian-type) convergence in the grid spacing.

The following remark has to be made on the frequency register size $N$. In the algorithm from Chen et al. the DFT gird parameters $N, \omega_0, t_0$ are set by: (i) the frequency range $N\omega_0/2 \geq 2\|H\| + \mathcal{O}(\sigma^{-1}\sqrt{\log(1/\varepsilon)})$, (ii) the Gaussian truncation in time $Nt_0/2 \geq \mathcal{O}(\sigma\sqrt{\log(1/\varepsilon)})$, and (iii) the DFT aliasing condition \cite[Proposition A.7]{chen2023quantum}. None of these conditions involve $\gamma$. This is to say, that for analytically prescribed (fine) grids where all error sources are already exponentially small, the number of frequency-register qubits $\lceil\log_2 N\rceil$ is the same for all three transition weights. However, in practice quadrature errors can be slightly below the algorithmic errors e.g. prescribed by $\delta$ in the weak measurement scheme, and yet achieve good results. If we operate on such a coarser grid (larger $\omega_0$, fewer qubits), then the exponential terms remain small while $\mathcal{O}(\beta\omega_0^2)$ dominates for the Metropolis case, making the kink the actual bottleneck. Smoothing with $s > 0$ removes this bottleneck entirely.

\todo{we might needs different structure than this paragraph / subsection..}
\label{sec:quad-errors-B}
\smallskip
\subsection{Quadrature errors for $B$.}
The remaining part in the CKG Lindbladian \eqref{eq:ckg-L} that has to be discretized is the coherent term $B$, which takes the following form
\begin{equation} \label{eq:B-discr}
    \bar{B} = \sum_a t_0\sum_{\bar{t}\in S_{t_0}} b_-(\bar{t}) \ee^{-\ii H \bar{t} / \sigma} \left[t_0'\sum_{\bar{t}' \in S_{t_0'}} b_+(\bar{t}') A^{a\dagger}(\beta \bar{t}') A^a(-\beta \bar{t}')\right]\ee^{\ii H \bar{t} / \sigma}.
\end{equation}
Note, that for the coherent term we keep the Riemann-sum factor explicit (since here there is no Fourier factors present). However, in quantum computer, the time registers have to then encode e.g. $\sqrt{t_0 b_-(\bar{t})} / \sqrt{\alpha_-}$, where $\alpha_- = t_0 \sum_{\bar{t}} |b_-(\bar{t})| \approx \|b_-\|_1$ is the subnormlaization of $b_-$.
Due to the nested integral form (which arose from the factorizing Gaussian filters), the quadrature error $\|B - \bar{B}\|$ decomposes into independent contributions from the outer ($b_-$) and inner ($b_+$) integrals via a triangle inequality. Concretely, define the continuous inner operator,
\begin{equation}
    I_+^a := \int_{-\infty}^{\infty} b_+(t')\, A^{a\dagger}(\beta t')\, A^a(-\beta t')\,\dd t',
\end{equation}
whose discretized counterpart is $\bar{I}_+^a$, and its outer conjugated operator $O_a(t) := \ee^{-\ii H t / \sigma} I_+^a \ee^{\ii H t}$. Then, the quadrature error for the full nested integral can be written as
\begin{equation} \label{eq:B-error-factorize}
    \begin{split}
        &\|B - \bar{B}\| \;\leq\; \\
    &\underbrace{\sum_a\left\|\int_\mathbb{R} b_-(t)\,O_a(t)\,dt - t_0\sum_{\bar{t}\in S_{t_0}} b_-(\bar{t})\,O_a(\bar{t})\right\|}_{\text{outer quadrature error (with exact } I_+^a\text{)}}
    \;+\;
    \underbrace{\|b_-\|_1 \cdot \sum_a \|I_+^a - \bar{I}_+^a\|.}_{\text{inner quadrature error, weighted by } \|b_-\|_1}
    \end{split}
\end{equation}
Although the errors factorize, a coupling between them arises nevertheless: 
the outer Riemann-sum error involves $\|I^a_+\| \leq \|b_+\|_1 \|A^a\|^2 \leq \|b_+\|_1$ through the norm of the operator being conjugated via $O_a$. Thus, the outer grid requirement inherits the subnormalization from the inner integral. 

We proceed by analysing the integrals one by one and then combine the results.

\medskip
\noindent\textbf{Outer integral ($b_-$).}
The kernel $b_-(t)$ is universal across all transition weight choices $\gamma$ and takes the form \eqref{eq:b_minus}. It is a convolution of two entire functions: $1/\cosh(2\pi t / (\beta\sigma))$, which decays exponentially as $\sim 2\ee^{-2\pi |t| / (\beta\sigma)}$, and $\sin(-\beta\sigma t)\,\ee^{-2t^2}$, which decays as a Gaussian. The convolution inherits the slower (exponential) decay  rate, giving and effective width $\mathcal{O}(\beta\sigma)$ as noted in \cite[footnote 13]{chen2023efficient}.

The outer integral involves Hamiltonian evolution $\ee^{\pm \ii H \bar{t} / \sigma}$, requiring simulation time of $\mathcal{O}(|t| / \sigma)$. Truncating to $|t| \leq T_-  / 2$ introduces a tail error that decays exponentially:
\begin{equation}
    \text{(outer truncation error)} = \mathcal{O}\!\left(\ee^{-c_-\, T_- / (\beta\sigma)}\right),
\end{equation}
for a constant $c_- > 0$ that depends on the $1/\cosh$ decay rate. Setting $T_- = \Theta(\beta\sigma\log(1/\varepsilon))$ controls this error to $\varepsilon$, corresponding to a Hamiltonian simulation time of $\mathcal{O}(\beta\log(1/\varepsilon))$ in the outer integral.

For the Riemann-sum error, we apply the operator-valued discretization bound from \cite[Lemma III.1]{chen2023efficient}:
\begin{equation}\label{eq:riemann-sum-operator}
    \left\|\int_{-T/2}^{T/2} O(t)\,f(t)\,\dd t - \sum_{\bar{t}\in S_{t_0}} O(\bar{t})\,f(\bar{t})\,t_0\right\| \;\leq\; \frac{T^2}{2|S_{t_0}|}\!\left(\|[O, H_\text{eff}]\|\sup_t|f(t)| + \|O\|\sup_t|f'(t)|\right),
\end{equation}
with $O(t) = \ee^{-\ii H_\text{eff} t} \,O\, \ee^{\ii H_\text{eff} t}$. Specifically, for the outer integral we have $f = b_-$, $H_\text{eff} = H / \sigma$ and $O = \bar{I}^a_+$, i.e. the discretized inner block. So the outer operator-norm error $\delta_-^a$ for each $a$ is:
\begin{equation}
    \delta^a_- \;\leq\; \frac{T_-^2}{2N_-}\left( \left\|[\bar{I}^a_+, H / \sigma]\right\|\|b_-\|_\infty + \|\bar{I}^a_+ \|\; \|b_-'\|_\infty\right).
\end{equation}
Using the commutator bound $\left\|[\bar{I}^a_+, H / \sigma]\right\|\leq 2 \|H\|\; \|\bar{I}^a_+\| / \sigma$, and the norm bound of $\|\bar{I}^a_+\|$ we find the full discretization error for the outer integral after summing over $a$: 
\begin{equation}
    \delta_- \leq \frac{T_-^2 |\mathcal{A}| \|b_+\|_1}{2 N_-} \left(\frac{2\|H\|}{\sigma}\|b_-\|_\infty + \|b_-'\|_\infty\right).
\end{equation}
Setting $\delta \leq \varepsilon$ with $T_- = \Theta(\beta\sigma\log(1/\varepsilon))$:
\begin{equation} \label{eq:r_-}
    r_- = \mathcal{O}\!\left(\log\!\left(\frac{\beta\sigma\|H\|\,|\mathcal{A}|\,\|b_+\|_1 \log(1/\varepsilon)}{\varepsilon}\right)\right).
\end{equation}
For the Gaussian transition weight $\|b_+\|_1 = \mathcal{O}(1)$, and thus the inner integral does not lead to additional costs. However, for the Metropolis-like transitions (both smooth and non-smooth) we have $\|b_+\|_1 = \mathcal{O}(\log(\beta\|H\|/\varepsilon))$, contributing as an additive $\log\log(\beta\|H\|/\varepsilon)$ term to $r_-$ (which is indeed not the end of the world).

\medskip
\noindent\textbf{Inner integral ($b_+$).}
Before discussing the different cases, we first have to establish the Riemann-sum bound for the inner integral \eqref{eq:B-discr}. Since 
$$
A^{a\dagger}(\beta t')\,A^a(-\beta t') = e^{i\beta H t'}\,A^{a\dagger}\,e^{-2i\beta H t'}\,A^a\,e^{i\beta H t'}
$$
is not of the Heisenberg conjugation form we cannot directly apply \cite[Lemma III.1]{chen2023efficient} like before. Instead, we use the general scalar Riemann-sum bound
\begin{equation}
    \begin{split}
        &\left\|\int_{-T_+/2}^{T_+/2} b_+(t')A^{a\dagger}(\beta t')\,A^a(-\beta t')\,dt' - \sum_{S_{t_0'}} b_+(\bar{t}')A^{a\dagger}(\beta \bar{t}')\,A^a(-\beta \bar{t}')\,t_0'\right\| \\
        &\leq \frac{T_+^2}{2N_+}\sup_{t'}\left\|\left(b_+(t')A^{a\dagger}(\beta t')\,A^a(-\beta t')\right)'(t')\right\|.
    \end{split}
\end{equation}
Where the norm of the derivative can be upper-bounded by the derivative of the parts. Specifically if we use that $\|A^{a\dagger}(\beta t')\,A^a(-\beta t')\| \leq \|A^a\|^2 \leq 1$ and that via the Leibniz rule we get
\begin{equation}
    \left\|\left(A^{a\dagger}(\beta t')\,A^a(-\beta t')\right)'(t')\right\| \leq 4\beta\|H\|\,\|A^a\|^2,
\end{equation}
then combining everything leads to the bound, 
\begin{equation}
    \left\|\left(b_+(t')A^{a\dagger}(\beta t')\,A^a(-\beta t')\right)'(t')\right\| \leq \|A^a\|^2\left(|b_+'(t')| + 4\beta\|H\|\,|b_+(t')|\right).
\end{equation}
This yields a per-$a$ discretization error $\delta_+^a$:
\begin{equation}
    \delta_+^a \leq \frac{T_+^2\,\|A^a\|^2}{2N_+}\sup_{t'}\!\left(|b_+'(t')| + 4\beta\|H\|\,|b_+(t')|\right).
\end{equation}
Summing over $a$:
\begin{equation} \label{eq:inner-discr-error}
    \delta_+ \leq \frac{T_+^2\,|A|}{2N_+}\sup_{t'}\!\left(|b_+'(t')| + 4\beta\|H\|\,|b_+(t')|\right).
\end{equation}
\begin{enumerate}[(a)]
    \item \textbf{Gaussian case}
    
    For the Gaussian weight $\gamma_G$, the time-domain kernel is given by a Gaussian \eqref{eq:b_plus-gauss} (surprise!) with width $\mathcal{O}(1)$. Hence, $\|b_+^{(G)}\|_\infty, \|(b_+^{(G)})'\|_\infty = \mathcal{O}(1)$.

    Correspondingly, the Gaussian decay gives 
    $$
    \text{(inner truncation)} = \mathcal{O}\!\left(e^{-c_G\, T_+^2}\right),
    $$
    so $T_+ = \Theta(\sqrt{\log(1/\varepsilon)})$ suffices. The Hamiltonian simulation time for the inner integral is $\beta T_+ = \Theta(\beta\sqrt{\log(1/\varepsilon)})$.

    Inserting $\|b_+^{(G)}\|_\infty, \|(b_+^{(G)})'\|_\infty$ into the general discretization error bound above \eqref{eq:inner-discr-error}:
    \begin{equation}
        \delta_+^{(G)} \leq \frac{T_+^2\,|\mathcal{A}|}{2N_+}\,\mathcal{O}(\beta\|H\|).
    \end{equation}
    Setting $\delta^{(G)}_+ \leq \varepsilon$, 
    \begin{equation} \label{eq:r_+_gaussian}
        r_+^{(G)} = \mathcal{O}\!\left(\log\!\left(\frac{\beta\|H\|\,|\mathcal{A}|\,\log(1/\varepsilon)}{\varepsilon}\right)\right).
    \end{equation}
    \item \textbf{Metropolis cases $(s \geq 0)$}
    
    Unlike for the dissipative case $\mathcal{\bar{L}}_\text{diss}$, here we find no qualitative difference between the smooth or kinky Metropolis variants, since both of them require the same $\eta$-regularization in the time-domain, which becomes the dominant factor. 

    However, for completeness and also to ground all of our later numerical results via analytic forms, let us derive the necessary quantities here.

    For $|t'|$ large enough that the indicator $\mathbb{I}(|t'|\leq\eta)$ is inactive, the kernel $b_+^{(s,\eta)}(t')$ is dominated by:
    $$
    |b_+^{(s,\eta)}(t')| \leq \frac{1}{2\sqrt{2}\pi^2}\,\frac{e^{-2\sigma^2\beta^2(1+s)\,t'^2}}{|t'|\,|2t'+i|} \quad\text{for } |t'| > \eta.
    $$
    The Gaussian decay rate $2\sigma^2\beta^2(1+s)$ increases with $s$, narrowing the effective support:
    \begin{equation} \label{eq:T_+_metro}
        T_+ = \Theta\!\left(\frac{\sqrt{\log(1/\varepsilon)}}{\sigma\beta\sqrt{1+s}}\right).
    \end{equation}

    Then, the required Hamiltonian simulation time for the inner integral in \eqref{eq:B-discr} is $\beta T_+$.

    The indicator function $\mathbb{I}(|t'|\leq\eta)$ in $b_+^{(s,\eta)}$ introduces jump discontinuities at $|t'| = \eta$. Since the Riemann-sum bound requires smoothness of the integrand, we split the inner integral at $t' = \pm\eta$ into three pieces:

    \textit{Central piece $[-\eta, \eta]$}: On the interval regularization removes the $1 / t'$ singularity by design, making $b_+^{(s,\eta)}$ bounded near the origin. Both $b_+^{(s,\eta)}$ and its derivatives are $\mathcal{O}(1)$ on $[-\eta,\eta]$. With a uniform grid of step size $t_0' = T_+/N_+$, the number of grid points in $[-\eta,\eta]$ is $\sim 2\eta N_+/T_+$, and the Riemann-sum error from this piece is
    $$
    \mathcal{O}\!\left(\frac{(2\eta)^2}{2\cdot (2\eta N_+/T_+)}\cdot \mathcal{O}(\beta\|H\|)\right) = \mathcal{O}\!\left(\frac{\eta T_+\beta\|H\|}{N_+}\right),
    $$
    which is a factor $\sim \eta/T_+$ of the outer-piece error and therefore negligible.

    \textit{Outer pieces $[-T_+/2, -\eta] \cup [\eta, T_+/2]$}: On these intervals $b_+^{(s,\eta)}(t')$ is smooth, with the Gaussian envelope providing rapid decay. The kernel has the scaling $|b_+^{(s,\eta)}(\eta^+)| = \mathcal{O}(1/\eta)$, while its derivative satisfies $\|(b_+^{(s,\eta)})'\|_{L^\infty(|t'|>\eta)} = \mathcal{O}(1/\eta^2)$ near $t' = \eta$ with Gaussian decay for larger $|t'|$. Applying the general discretization bound \eqref{eq:inner-discr-error} to the outer pieces the dominant contribution (the derivative) gives:
    \begin{equation}
        \delta_+^a \leq \frac{T_+^2\,\|A^a\|^2}{2N_+}\,\mathcal{O}\!\left(\frac{1}{\eta^2}\right).
    \end{equation}
    Summing over $a$, substituting $T_+$ and setting $\eta = \varepsilon/(\beta\|H\|\|\sum_a A^{a\dagger}A^a\|)$ leads to the scaling of the necessary estimating qubits for the inner integral of $B$:
    \begin{equation} \label{eq:r_+_metro}
        r_+^{(s)} = \mathcal{O}\!\left(\log\!\left(\frac{\|H\|\left\|\sum_a A^{a\dagger}A^a\right\| |\mathcal{A}|\,\log(1/\varepsilon)}{\sigma(1+s)\,\varepsilon}\right)\right).
    \end{equation}
    This is a polylogarithmic in $1 / \varepsilon$ for all $s \geq 0$. The $s$-dependence appears only as a mild $\log(1 + s)$ reduction inside the algorithm.
\end{enumerate}
As we will see later in the algorithm section \todo{ref}, to implement the block encoding of the coherent term $B$, we need two time registers to compute the nested integrals in \eqref{eq:B-discr}. Accordingly, the total number of estimating qubits is the sum of the above, separately derived expressions, $r = r_- + r_+$. At this discretization / estimating-qubit level, the outer $r_-$ and inner $r_+$ integrals contribute essentially the same way. Where it will matter more that the Metropolis cases had to be regularized with $\eta$ is in their implementation via the LCU method. There the complexity depends on the subnormalizations for which $\|b^{(s)}_+\|_1$ indeed leads to additional costs. \todo{ref}

% Comment on Ding's Poission technique that works for smooth variants but not for time singular Metropolis functions

\smallskip
\subsection{Hamiltonian Trotterization.} 

\todo{could add a comment about childs constants for Heis or Isin}
\todo{Add that discussion into the Trotter prelims}
The CKG Lindbladian \eqref{eq:ckg-L} is abundant with exact Hamiltonian time evolutions $\ee^{\pm \ii H t}$. Any implementation on a quantum computer must approximate these using some product formula when $H = \sum_{k = 1}^K H_k$ has non-commuting terms. With this subsection we have managed to reach the deepest point of the dark, error-analysis-cave, where we quantify the incurring errors due to the replacement of $e^{\pm iH\bar{t}}$ in the discretized OFT and coherent term with a product formula $S_p^{(M)}(\bar{t}) := S_p(\bar{t}/M)^M$ of order $p$ with $M$ Trotter steps. 

In \hyperref[sec:prelim-trotter]{Section~\ref*{sec:prelim-trotter}} the general Trotterization strategy and the approximation errors they lead to have been introduced using \cite{childs2021theory} as the main reference. Here, we make use of those results, namely that by \cite[Theorem 6]{childs2021theory} a $p$-th order product formula $S_p^{(M)}(t)$ for $\ee^{\ii Ht}$ satisfies the commutator-scaling bound:
\begin{equation} \label{eq:childs-trotter-bound}
    \bigl\|S_p^{(M)}(t) - e^{iHt}\bigr\|
    = \mathcal{O}\!\left(\frac{\tilde{\alpha}_{\mathrm{comm}}\, |t|^{p+1}}{M^p}\right),
\end{equation}
where
$$
\tilde{\alpha}_{\mathrm{comm}}
:= \sum_{k_1, k_2, \ldots, k_{p+1} = 1}^{K}
\bigl\|\bigl[H_{k_{p+1}}, \cdots [H_{k_2}, H_{k_1}]\cdots\bigr]\bigr\|
$$
is the commutator-scaling constant. With $M$ Trotter steps of size $\delta t = t / M$, the total error is $\|S_p^{(M)}(t) - e^{iHt}\| \leq \varepsilon_p(t, M)$ with $\varepsilon_p(t, M) = O(\tilde{\alpha}_{\mathrm{comm}}\, |t|^{p+1}/M^p)$. 

We use the following notation throughout: $\hat{A}(\bar{\omega})$ is the discretized OFT operator with \emph{exact} Hamiltonian evolution; $\tilde{A}(\bar{\omega})$ is the same object with every $\ee^{\pm \ii H\bar{t}}$ replaced by $S_p(\pm \bar{t})$. Similarly, $\mathcal{\bar{L}}_\text{diss}$ is built from $\hat{A}(\bar{\omega})$, while $\mathcal{\tilde{L}}_\text{diss}$ is its Trotterized counterpart built from $\tilde{A}(\bar{\omega})$.

First, we tackle the dissipative part of the Lindbladian, i.e. the Trotter errors propagating from the OFTs \eqref{eq:oft-discr} within. The Trotterized version replaces $\ee^{\ii H\bar{t}} A_a \ee^{-\ii H\bar{t}}$ with $S_p^{(M)}(\bar{t})\, A_a\, S_p^{(M)}(\bar{t})^\dagger$:
\begin{equation} \label{eq:trotter-oft}
    \tilde{A}_a(\bar{\omega})
    = \frac{1}{\sqrt{N}} \sum_{\bar{t} \in S_{t_0}} \ee^{-\ii\bar{\omega}\bar{t}}\, \bar{f}(\bar{t})\, S_p^{(M)}(\bar{t})\, A_a\, S_p^{(M)}(\bar{t})^\dagger.
\end{equation}

\begin{proposition} \label{prop:trotter-diss}
    \textup{(Trotter error for $\mathcal{L}_\textup{diss}$).} Let $H = \sum_{k=1}^K H_k$ with non-commuting terms, and let $\{A_a\}_{a \in \mathcal{A}}$ satisfy the normalization $\|\sum_a A_a^\dagger A_a\| \leq 1$. Let $\bar{f} : S_{t_0} \to \mathbb{R}$ be some discretized filter function with $\|\bar{f}\|_2^2 := \sum_{\bar{t} \in S_{t_0}} |\bar{f}(\bar{t})|^2 = 1$, and let the transition weight $\gamma$ satisfy $\|\gamma\|_\infty \leq 1$. With the Trotter product formula defined by \eqref{eq:childs-trotter-bound} and the corresponding Trotterized OFT \eqref{eq:trotter-oft} the total incurring Trotter error for the $\mathcal{\bar{L}}_\textup{diss}$ can be bounded as:
    $$
    \bigl\|\bar{\mathcal{L}}_{\mathrm{diss}} - \tilde{\mathcal{L}}_{\mathrm{diss}}\bigr\|_{1 \to 1}
    \leq 8\,\sqrt{\sum_{\bar{t} \in S_{t_0}} |\bar{f}(\bar{t})|^2\, \varepsilon_p^2(\bar{t}, M)},
    $$
    where $\varepsilon_p(\bar{t}, M) := \|S_p^{(M)}(\bar{t}) - \ee^{\ii H\bar{t}}\|$ is the product-formula error at time $\bar{t}$.

    For the second-order Strang splitting ($p = 2$), define the explicit commutator-scaling constant
    $$
    \tilde{\alpha}_{\mathrm{comm}}^{(2)}
    := \sum_{k_1 < k_2}
    \left(
      \frac{1}{12}\bigl\|[H_{k_2},[H_{k_2},H_{k_1}]]\bigr\|
    + \frac{1}{24}\bigl\|[H_{k_1},[H_{k_1},H_{k_2}]]\bigr\|
    \right),
    $$
    which is the leading-order prefactor from the tight second-order analysis of \cite[Proposition 10]{childs2021theory}. Then at leading order in $1 / M$:
    $$
    \bigl\|\bar{\mathcal{L}}_{\mathrm{diss}}
          - \tilde{\mathcal{L}}_{\mathrm{diss}}\bigr\|_{1 \to 1}
    \;\leq\;
    \frac{\sqrt{15}\;\tilde{\alpha}_{\mathrm{comm}}^{(2)}}
         {M^2\,\sigma^3}
    \;+\; \mathcal{O}\!\left(\frac{1}{M^3\,\sigma^4}\right).
    $$
\end{proposition}
\begin{proof}
    \noindent\textbf{(i) Bound $\bigl\|\bar{\mathcal{L}}_{\mathrm{diss}} - \tilde{\mathcal{L}}_{\mathrm{diss}}\bigr\|_{1 \to 1}$.}
    Define the Stinespring isometries with the discretized OFT \eqref{eq:oft-discr} and Trotterized OFT \eqref{eq:trotter-oft}, 
    $$
    \bar{V} := \sum_{a,\bar{\omega}} \sqrt{\gamma(\bar{\omega})}\, |\bar{\omega}, a\rangle\langle \bar{0}|\, \hat{A}_a(\bar{\omega}), \qquad
    \tilde{V} := \sum_{a,\bar{\omega}} \sqrt{\gamma(\bar{\omega})}\, |\bar{\omega}, a\rangle\langle \bar{0}|\, \tilde{A}_a(\bar{\omega}).
    $$
    Both $\mathcal{\bar{L}}_\text{diss}$ and $\mathcal{\tilde{L}}_\text{diss}$ admit the Stinespring factorization 
    $$
    \mathcal{L}[{\rho}] = \tr_{\mathrm{anc}}(V \rho\, V^\dagger) - \frac{1}{2}\tr_{\bar{0}}\{V^\dagger V, \rho\}
    $$
    Writing $\delta V := \bar{V} - \tilde{V}$ and expanding both the transition and decay terms, a triangle inequality on the $1 \to 1$ norm gives (following the proof technique of \cite[Lemma A.1]{chen2023quantum}):
    $$
    \bigl\|\bar{\mathcal{L}}_{\mathrm{diss}} - \tilde{\mathcal{L}}_{\mathrm{diss}}\bigr\|_{1 \to 1}
    \leq \|\delta V\|\bigl(\|\bar{V}\| + \|\tilde{V}\|\bigr) + \|\delta V\|\bigl(\|\bar{V}\| + \|\tilde{V}\|\bigr) = 2\,\|\delta V\|\bigl(\|\bar{V}\| + \|\tilde{V}\|\bigr).
    $$
    Both $\bar{V}$ and $\tilde{V}$ are sub-isometries with $\|V\| \leq 1$: this follows from the operator Parseval identity (\cite[Proposition A.1]{chen2023quantum}), $\|\gamma\|_\infty \leq 1$, $\|\bar{f}\|_2 = 1$, and the CKG normalization $\|\sum_a A_a^\dagger A_a\| \leq 1$. (For $\tilde{V}$, the same bound holds since $S_p^{(M)}(\bar{t})$ is unitary, so the Parseval identity applies identically with $S_p^{(M)}$ in place of $U$.) Therefore,
    $$
    \bigl\|\bar{\mathcal{L}}_{\mathrm{diss}} - \tilde{\mathcal{L}}_{\mathrm{diss}}\bigr\|_{1 \to 1}
    \leq 4\,\|\bar{V} - \tilde{V}\|.
    $$
    Note, that the key thought here is to use the perturbation techniques introduced in Chen et al. for the error analysis of discretization, with the understanding that Trotterization is just another kind of perturbation in this picture.

    \medskip
    \noindent\textbf{(ii) Bound $\|\bar{V} - \tilde{V}\|$}. We claim that,
    $$
    \|\bar{V} - \tilde{V}\|^2 \leq \sum_{\bar{t} \in S_{t_0}} |\bar{f}(\bar{t})|^2 \, \left\|\sum_a \delta C_a(\bar{t})^\dagger\, \delta C_a(\bar{t})\right\|,
    $$
    where
    $$
    \delta C_a(\bar{t}) := U(\bar{t})^\dagger A_a\, U(\bar{t}) - S_p^{(M)}(\bar{t})^\dagger A_a\, S_p^{(M)}(\bar{t})
    $$
    is there per-time-point Heisenberg-picture error due to Trotterization.

    By linearity, we can write,
    $$
    \bar{V} - \tilde{V} = \sum_{a,\bar{\omega}} \sqrt{\gamma(\bar{\omega})}\, |\bar{\omega}, a\rangle\langle \bar{0}|\, \delta A_a(\bar{\omega}),
    $$
    with $\delta A_a(\bar{\omega}) := \hat{A}_a(\bar{\omega}) - \tilde{A}_a(\bar{\omega}) = \frac{1}{\sqrt{N}} \sum_{\bar{t}} e^{-i\bar{\omega}\bar{t}}\, \bar{f}(\bar{t})\, \delta C_a(\bar{t})$
    Using the orthonormality of the ancillas $\langle \bar{\omega}', a' | \bar{\omega}, a\rangle = \delta_{aa'}\delta_{\bar{\omega}\bar{\omega}'}$, we find
    $$
    (\bar{V} - \tilde{V})^\dagger(\bar{V} - \tilde{V}) = \sum_{a,\bar{\omega}} \gamma(\bar{\omega})\, \delta A_a(\bar{\omega})^\dagger\, \delta A_a(\bar{\omega}) \preceq \sum_{a,\bar{\omega}} \delta A_a(\bar{\omega})^\dagger\, \delta A_a(\bar{\omega}),
    $$
    In the last step we used $\gamma(\omega) \leq 1$. Then we can apply Parseval's identity from \cite[Proposition A.1]{chen2023quantum},
    $$
    \sum_{\bar{\omega}} \delta A_a(\bar{\omega})^\dagger\, \delta A_a(\bar{\omega}) = \sum_{\bar{t}} |\bar{f}(\bar{t})|^2\, \delta C_a(\bar{t})^\dagger\, \delta C_a(\bar{t}).
    $$
    Inserting this back we get, 
    $$
    (\bar{V} - \tilde{V})^\dagger(\bar{V} - \tilde{V}) \preceq \sum_{\bar{t}} |\bar{f}(\bar{t})|^2 \sum_a \delta C_a(\bar{t})^\dagger\, \delta C_a(\bar{t}).
    $$
    Each summand $|\bar{f}(\bar{t})|^2 \sum_a \delta C_a(\bar{t})^\dagger\, \delta C_a(\bar{t})$ is positive semidefinite, and after using triangle inequality, we find the claimed bound on $\|\bar{V} - \tilde{V}\|^2$. 

    \medskip
    \noindent\textbf{(iii) Bound $\left\|\sum_a \delta C_a(\bar{t})^\dagger\, \delta C_a(\bar{t})\right\|$}. We claim that,
    $$
    \Bigl\|\sum_a \delta C_a(\bar{t})^\dagger\, \delta C_a(\bar{t})\Bigr\| \leq 4\,\varepsilon_p^2(\bar{t}, M).
    $$
    Define $\Delta(\bar{t}) := U(\bar{t}) - S_p^{(M)}(\bar{t})$, and its norm (the Trotter error) satisfies $\|\Delta(\bar{t})\| \leq \varepsilon_p(\bar{t}, M)$. Decomposing $\delta C_a(\bar{t})$:
    $$
    \begin{aligned}
        \delta C_a(\bar{t}) &= U^\dagger A_a\, U - (S_p^{(M)})^\dagger A_a\, S_p^{(M)} \\
    &= (S_p^{(M)} + \Delta)^\dagger A_a\, U - (S_p^{(M)})^\dagger A_a (U - \Delta) \\
    &= \Delta^\dagger A_a\, U + (S_p^{(M)})^\dagger A_a\, \Delta \\
    &=: X_a + Y_a,
    \end{aligned}
    $$
    Since $(X_a - Y_a)^\dagger(X_a - Y_a) \succeq 0$ implies $X_a^\dagger Y_a + Y_a^\dagger X_a \preceq X_a^\dagger X_a + Y_a^\dagger Y_a$, we get
    $$
    \delta C_a^\dagger\, \delta C_a = (X_a + Y_a)^\dagger(X_a + Y_a) \preceq 2\bigl(X_a^\dagger X_a + Y_a^\dagger Y_a\bigr).
    $$
    We can bound $\sum_a X_a^\dagger X_a$ by using that $\Delta\Delta^\dagger \preceq \|\Delta\|^2\, \mathbf{1} \preceq \varepsilon_p^2\, \mathbf{1}$, that $U$ is unitary and $\|\sum_a A_a^\dagger A_a\| \leq 1$,
    $$
    \sum_a X_a^\dagger X_a = \sum_a U^\dagger A_a^\dagger\, \Delta\Delta^\dagger\, A_a\, U
    \preceq \varepsilon_p^2 \sum_a U^\dagger A_a^\dagger A_a\, U
    \preceq \varepsilon_p^2\, \mathbf{1}.
    $$
    The result is the same for $\sum_a Y_a^\dagger Y_a$ but by using that $S_p^{(M)}$ is unitary.
    Combining them gives, 
    $$
    \sum_a \delta C_a^\dagger\, \delta C_a \preceq 2\bigl(\varepsilon_p^2\, \mathbf{1} + \varepsilon_p^2\, \mathbf{1}\bigr) = 4\,\varepsilon_p^2\, \mathbf{1}.
    $$
    Taking the operator norm finally proves the upper-bound on $\left\|\sum_a \delta C_a(\bar{t})^\dagger\, \delta C_a(\bar{t})\right\|$.

    \medskip
    \noindent\textbf{(iv) Assemble the cake.} Putting steps (i)-(iii) together, we find that the Trotter error on the full superoperator $\mathcal{\bar{L}}_\text{diss}$ is bounded by
    \begin{align*}
        \bigl\|\bar{\mathcal{L}}_{\mathrm{diss}} - \tilde{\mathcal{L}}_{\mathrm{diss}}\bigr\|_{1 \to 1}
    &\leq 4\,\|\bar{V} - \tilde{V}\| \tag{i}\\
    &\leq 8 \,\sqrt{\sum_{\bar{t}} |\bar{f}(\bar{t})|^2\, \left\|\sum_a \delta C_a(\bar{t})^\dagger\, \delta C_a(\bar{t})\right\|} \tag{ii}\\
    &\leq 8 \,\sqrt{\sum_{\bar{t}} |\bar{f}(\bar{t})|^2\, \varepsilon_p^2(\bar{t}, M)}. \tag{iii}
    \end{align*}

    \medskip
    \noindent\textbf{(v) Strang splitting and Gaussian filters.} In our numerical analysis we will use Strang splitting, i.e. Trotterization of order $p = 2$. In this case, the error bound specializes as follows.
    \begin{align*}
    \bigl\|S_2(\delta t) - \ee^{\ii H \delta t}\bigr\|
    &\leq \sum_{k_1 < k_2}
    \left(
      \frac{1}{12}\bigl\|[H_{k_2},[H_{k_2},H_{k_1}]]\bigr\|
    + \frac{1}{24}\bigl\|[H_{k_1},[H_{k_1},H_{k_2}]]\bigr\|
    \right)
    + \mathcal{O}\bigl(|\delta t|^4\bigr) \\
    &= \tilde{\alpha}_{\mathrm{comm}}^{(2)}\,|\delta t|^3
    + \mathcal{O}\bigl(|\delta t|^4\bigr),
    \end{align*}
    where the leading-order coefficients $1/12$ and $1/24$ are tight, and the $\mathcal{O}(|\delta t|^4)$ correction arises from the multi-term bootstrapping of \cite[Proposition~9]{childs2021theory}.
    With $M$ Trotter steps of size $\delta t = \bar{t}/M$ and a triangle inequality
    over the $M$ steps:
    $$
    \varepsilon_2(\bar{t}, M)
    \leq \tilde{\alpha}_{\mathrm{comm}}^{(2)}\,
        \frac{|\bar{t}|^3}{M^2}
    + \mathcal{O}\!\left(\frac{|\bar{t}|^4}{M^3}\right).
    $$
    Substituting the leading term into the assembled bound and evaluating the sixth moment of $|\bar{f}|^2$, where in the CKG case the filter is the Gaussian $f(t) = \ee^{-\sigma^2 t^2}\big/\bigl(\pi/(2\sigma^2)\bigr)^{1/4}$ with squared density $|f(t)|^2 \propto \ee^{-2\sigma^2 t^2}$, whose even moments evaluate to
    $$
    \int_{-\infty}^{\infty} |f(t)|^2\,t^{2k}\,\dd t
    = \frac{(2k-1)!!}{(2\sigma)^{2k}}.
    $$
    The discrete sum $\sum_{\bar{t}} |\bar{f}(\bar{t})|^2\,|\bar{t}|^6$
    agrees with its continuous counterpart up to an exponentially small
    correction (by the quadrature analysis of the preceding subsection).
    Setting $k = 3$ yields:
    $$
    \bigl\|\bar{\mathcal{L}}_{\mathrm{diss}}
        - \tilde{\mathcal{L}}_{\mathrm{diss}}\bigr\|_{1 \to 1}
    \leq \frac{8\,\tilde{\alpha}_{\mathrm{comm}}^{(2)}}{M^2}\;
        \sqrt{\frac{15}{64\,\sigma^6}}
    + \mathcal{O}\!\left(\frac{1}{M^3\,\sigma^4}\right)
    = \frac{\sqrt{15}\;\tilde{\alpha}_{\mathrm{comm}}^{(2)}}
        {M^2\,\sigma^3}
    + \mathcal{O}\!\left(\frac{1}{M^3\,\sigma^4}\right).
    $$
    The $\mathcal{O}(1/(M^3\sigma^4))$ correction arises from the $\mathcal{O}(|\bar{t}|^4/M^3)$ sub-leading Trotter error propagated through the eighth moment of $|\bar{f}|^2$.
\end{proof}
\noindent The main takeaway here is that the Trotter errors from the OFT propagate to the superoperator level in an expected manner, but with a crucial realization: Since the natural unit of the Gaussian filter width is the temperature $\sigma \sim 1 / \beta$, that means the lower the temperature is the more the  number of Trotter steps $M$ have to compensate to control the errors. Not even higher-orders $ p>2$ can help here, as this relation between $M$ and $\sigma$ is independent of the order. This was to be expected since the lower the temperature, the better the resolution in the energy-domain has to be, and thus the longer the respective time-evolutions are in the time-domain. Nevertheless, this has not been written out explicitly for a Trotter-domain before. 

\medskip
With the Trotter errors for the dissipative part of the Lindbladian $\mathcal{L}_\text{diss}$ fully understood, it remains to see how Trotterization affects the coherent term $B$. Replacing all exact Hamiltonian time evolutions in \eqref{eq:B-discr} with the Trotterized evolution $S_p^{(M)}(-\bar{t})$ leads to
\begin{equation} \label{eq:B-trotter}
    \begin{split}
         \tilde{B} &= \sum_a t_0\sum_{\bar{t}} b_-(\bar{t}) S_p^{(M)}(-\bar{t} / \sigma) \left[t_0'\sum_{\bar{t}'} b_+(\bar{t}') \tilde{A}_{a\dagger}(\beta \bar{t}') \tilde{A}_a(-\beta \bar{t}')\right]S_p^{(M)}(\bar{t} / \sigma) \\
         &=: \sum_{a} t_0 \sum_{\bar{t} \in S_{t_0}} b_-(\bar{t})\, S_p^{(M)}(-\bar{t}/\sigma)\, \tilde{O}_a\, S_p^{(M)}(\bar{t}/\sigma),
    \end{split}
\end{equation}
with $\tilde{O}_a$ denoting the inner Trotterized integral, and the Trotterized Heisenberg evolution of order $p$ and steps $M$, 
$$
\tilde{A}_a(t) := S_p^{(M)}(t) A_a S_p^{(M)}(-t).
$$
In the following we will use $\|b_-\|_{1} := t_0 \sum_{\bar{t}} |b_-(\bar{t})|$ and $\|b_+\|_{1} := t_0' \sum_{\bar{t}'} |b_+(\bar{t}')|$ for the discrete $\ell_1$ norms (with quadrature weights).

\todo{ref to impl}
There are a few ways to implement the evolution via the CKG Lindbladian, especially in terms of how we add the coherent evolution via $B$, but in all cases the core error boils down to controlling $\|\bar{B} - \tilde{B}\|$.
\begin{proposition} \label{prop:trotter-B}
    \textup{(Trotter error for the coherent term $B$).} Let $H = \sum_{k=1}^K H_k$ with non-commuting terms, let the jump operators satisfy $\bigl\|\sum_a A_a^\dagger A_a\bigr\| \leq 1$, and let $b_\pm$ be the time-domain weight functions of the coherent term with finite discrete norms $\|b_-\|_{1}, \|b_+\|_{1} < \infty$. Let $S_p^{(M)}(t)$ be a $p$-th order product formula with $M$ Trotter steps satisfying the commutator-scaling bound $\varepsilon_p(t, M) = \mathcal{O}(\tilde{\alpha}_{\mathrm{comm}}\,|t|^{p+1}/M^p)$.

    Then the operator-norm Trotter error of the coherent term satisfies, 
    $$
    \|\bar{B} - \tilde{B}\|
    \;\leq\;
    2\,\|b_+\|_{1}\; t_0 \sum_{\bar{t}} |b_-(\bar{t})|\, \varepsilon_p\!\bigl(|\bar{t}|/\sigma,\, M\bigr)
    \;+\;
    4\,\|b_-\|_{1}\; t_0' \sum_{\bar{t}'} |b_+(\bar{t}')|\, \varepsilon_p\!\bigl(\beta|\bar{t}'|,\, M\bigr).
    $$
    Here, the first term captures the outer Trotter error (replacing the outer exact time evolution $\ee^{\pm \ii H \bar{t} / \sigma}$), while the second captures the inner Trotterization error (replacing $\ee^{\pm \ii H\beta\bar{t}'}$ inside $A_a(\beta \bar{t}')$).

    For the special case of Strang splitting ($p = 2$), substituting $\varepsilon_2(t, M) \leq \tilde{\alpha}_{\mathrm{comm}}^{(2)}\,|t|^3/M^2 + \mathcal{O}(|t|^4/M^3)$ gives the following result at leading order:
    $$
    \|\bar{B} - \tilde{B}\|
    \;\leq\;
    \frac{\tilde{\alpha}_{\mathrm{comm}}^{(2)}}{M^2}
    \left(
    \frac{2\,\|b_+\|_1}{\sigma^3}\, \mu_3(|b_-|)
    \;+\;
    4\,\beta^3\,\|b_-\|_1\, \mu_3(|b_+|)
    \right)
    + \mathcal{O}\!\left(\frac{\beta^4}{M^3}\right),
    $$
    where $\mu_3(|b_\pm|) := t_0^{(\prime)} \sum |b_\pm(\bar{t})|\,|\bar{t}|^3$ denotes the third absolute moment of the respective weight function (with $t_0$ or $t_0'$).
\end{proposition}
\begin{proof}
    \hfill\break
    \noindent\textbf{(i) Decompose the Trotter error.}
    For each outer grid point $\bar{t}$ and jump operator $a$, define:
    $$
    \bar{Q}_a(\bar{t}) := U(-\bar{t}/\sigma)\, \bar{O}_a\, U(\bar{t}/\sigma), \qquad
    \tilde{Q}_a(\bar{t}) := S(-\bar{t}/\sigma)\, \tilde{O}_a\, S(\bar{t}/\sigma).
    $$
    Then, to bound an error between these two expressions, follow these steps in order: use the three-term telescoping identity for products of the form $PQR - P'Q'R'$, simplify, sum over $a$, use that $\|U\| = 1 = \|S\|$. At this point we arrive at
    $$
    \Bigl\|\sum_a \bigl[\bar{Q}_a(\bar{t}) - \tilde{Q}_a(\bar{t})\bigr]\Bigr\|
    \leq
    \varepsilon_p(|\bar{t}|/\sigma, M)\, \Bigl\|\sum_a \bar{O}_a\Bigr\|
    + \Bigl\|\sum_a \bigl[\bar{O}_a - \tilde{O}_a\bigr]\Bigr\|
    + \Bigl\|\sum_a \tilde{O}_a\Bigr\|\, \varepsilon_p(|\bar{t}|/\sigma, M),
    $$

    \noindent\textbf{(ii) Bound $\Bigl\|\sum_a \bar{O}_a\Bigr\|$ and $\Bigl\|\sum_a \tilde{O}_a\Bigr\|$.} The inner operator $O_a$ has the following general form: $A_a^\dagger W A_a$ for some unitary $W$ in both the Trotterized case \eqref{eq:B-trotter} and discretized causes \eqref{eq:B-discr}. Since, $-\mathbf{1} \preceq W \preceq \mathbf{1}$, we also have $-A_a^\dagger A_a \preceq A^\dagger_a W A_a \preceq A_a^\dagger A_a$. Summing over $a$ and taking norm then leads to:
    $$
    \Bigl\|\sum_a A_a^\dagger\, W\, A_a\Bigr\| \leq \Bigl\|\sum_a A_a^\dagger A_a\Bigr\| \leq 1,
    $$
    Therefore:
    $$
    \Bigl\|\sum_a \bar{O}_a\Bigr\| \leq t_0' \sum_{\bar{t}'} |b_+(\bar{t}')| = \|b_+\|_1.
    $$
    The identical bound holds for $\bigl\|\sum_a \tilde{O}_a\bigr\| \leq \|b_+\|_1$ since $S$ is also unitary.

    \noindent\textbf{(ii) Bound $\Bigl\|\sum_a \bigl[\bar{O}_a - \tilde{O}_a\bigr]\Bigr\|$.} We have
    $$
    \sum_a [\bar{O}_a - \tilde{O}_a] = t_0' \sum_{\bar{t}'} b_+(\bar{t}') \sum_a \Bigl[\hat{A}_a^\dagger(\beta\bar{t}')\, \hat{A}_a(-\beta\bar{t}') - \tilde{A}_a^\dagger(\beta\bar{t}')\, \tilde{A}_a(-\beta\bar{t}')\Bigr],
    $$
    so it suffices to bound $\bigl\|\sum_a \bigl[\hat{A}_a^\dagger(\beta\bar{t}')\hat{A}_a(-\beta\bar{t}') - \tilde{A}_a^\dagger(\beta\bar{t}')\tilde{A}_a(-\beta\bar{t}')\bigr]\bigr\|$ for each $\bar{t}'$. Define $\delta_a(t) := \hat{A}_a(t) - \tilde{A}_a(t)$ for the per-operator Heisenberg-picture Trotter error. Then,
    $$
    \sum_a \Bigl[\hat{A}_a^\dagger(\beta\bar{t}')\hat{A}_a(-\beta\bar{t}') - \tilde{A}_a^\dagger(\beta\bar{t}')\tilde{A}_a(-\beta\bar{t}')\Bigr]
    = \sum_a \delta_a(\beta\bar{t}')^\dagger\, \hat{A}_a(-\beta\bar{t}') + \sum_a \tilde{A}_a^\dagger(\beta\bar{t}')\, \delta_a(-\beta\bar{t}').
    $$
    We bound each sum using the operator Cauchy-Schwarz inequality, which in general can be written for operators $\{X_a\}, \{Y_a\}$ as
    $$
    \Bigl\|\sum_a X_a Y_a\Bigr\|^2 \leq \Bigl\|\sum_a X_a X_a^\dagger\Bigr\| \cdot \Bigl\|\sum_a Y_a^\dagger Y_a\Bigr\|.
    $$
    The first sum with $X_a = \delta_a(\beta\bar{t}')^\dagger$ and $Y_a = \hat{A}_a(-\beta\bar{t}')$ is bounded by two corresponding factors. The left one is
    $$
    \bigl\|\sum_a \delta_a(\beta\bar{t}')^\dagger\, \delta_a(\beta\bar{t}')\bigr\| \leq 4\,\varepsilon_p^2(\beta|\bar{t}'|, M),
    $$
    which was shown in step (iii) of \hyperref[prop:trotter-diss]{Proposition~\ref*{prop:trotter-diss}}. The right factor for the bound is 
    $$
    \bigl\|\sum_a \hat{A}_a(-\beta\bar{t}')^\dagger\, \hat{A}_a(-\beta\bar{t}')\bigr\| = \bigl\|\sum_a A_a^\dagger A_a\bigr\| \leq 1.
    $$
    Therefore, $\bigl\|\sum_a \delta_a(\beta\bar{t}')^\dagger\, \hat{A}_a(-\beta\bar{t}')\bigr\| \leq 2\,\varepsilon_p(\beta|\bar{t}'|, M).$ An identical bound holds for the second sum as well. Combining them via triangle inequality leads to,
    $$
    \Bigl\|\sum_a [\bar{O}_a - \tilde{O}_a]\Bigr\|
    \leq 4\, t_0' \sum_{\bar{t}'} |b_+(\bar{t}')|\, \varepsilon_p(\beta|\bar{t}'|, M).
    $$
    \noindent\textbf{(iv) Assemble another cake.} From steps (i)-(iii):
    $$
    \Bigl\|\sum_a [\bar{Q}_a(\bar{t}) - \tilde{Q}_a(\bar{t})]\Bigr\|
    \leq 2\,\|b_+\|_1\, \varepsilon_p(|\bar{t}|/\sigma, M)
    + 4\, t_0' \sum_{\bar{t}'} |b_+(\bar{t}')|\, \varepsilon_p(\beta|\bar{t}'|, M).
    $$
    Taking the outer $b_-$-weighted sum of $\|\|\bar{B} - \tilde{B}\|\|$, then inserting the above bound in yields our result
    \begin{align*}
        \|\bar{B} - \tilde{B}\|
    &= \Bigl\|t_0 \sum_{\bar{t}} b_-(\bar{t})\, \sum_a [\bar{Q}_a(\bar{t}) - \tilde{Q}_a(\bar{t})]\Bigr\|
    \leq t_0 \sum_{\bar{t}} |b_-(\bar{t})| \cdot \Bigl\|\sum_a [\bar{Q}_a(\bar{t}) - \tilde{Q}_a(\bar{t})]\Bigr\| \\
    & \leq 2\,\|b_+\|_1\, t_0 \sum_{\bar{t}} |b_-(\bar{t})|\, \varepsilon_p(|\bar{t}|/\sigma, M)
    + 4\,\|b_-\|_1\, t_0' \sum_{\bar{t}'} |b_+(\bar{t}')|\, \varepsilon_p(\beta|\bar{t}'|, M).
    \end{align*}    
    \noindent\textbf{(v) Strang splitting specialization.} For $p = 2$, the leading-order Trotter error is $\varepsilon_2(t, M) \leq \tilde{\alpha}_{\mathrm{comm}}^{(2)}|t|^3/M^2$. Define the third absolute moment of a weight function $w$ on grid $S_{\tau}$:
    $$
    \mu_3(|w|) := \tau \sum_{\bar{t} \in S_\tau} |w(\bar{t})|\, |\bar{t}|^3.
    $$
    Then the outer term becomes, $t_0 \sum_{\bar{t}} |b_-(\bar{t})|\, \varepsilon_2(|\bar{t}|/\sigma, M) \leq \frac{\tilde{\alpha}_{\mathrm{comm}}^{(2)}}{M^2\, \sigma^3}\, \mu_3(|b_-|)$, while the inner term, $t_0' \sum_{\bar{t}'} |b_+(\bar{t}')|\, \varepsilon_2(\beta|\bar{t}'|, M) \leq \frac{\tilde{\alpha}_{\mathrm{comm}}^{(2)}\, \beta^3}{M^2}\, \mu_3(|b_+|)$. Hence, up to leading order we find
    $$
    \|\bar{B} - \tilde{B}\|
    \leq \frac{\tilde{\alpha}_{\mathrm{comm}}^{(2)}}{M^2}
    \left(
    \frac{2\,\|b_+\|_1}{\sigma^3}\, \mu_3(|b_-|)
    + 4\,\beta^3\, \|b_-\|_1\, \mu_3(|b_+|)
    \right)
    + \mathcal{O}\!\left(\frac{\beta^4}{M^3}\right).
    $$
\end{proof}
The two terms present in the Trotter error bound of $\|\bar{B} - \tilde{B}\|$ arise from the nested outer and inner time-sums in $B$ \eqref{eq:B-discr}. Qualitatively they contribute similarly. The $\log$ scaling of $\|b^{(s, \eta)}_+\|_1$ \eqref{eq:b_plus-s-eta} that appears for Metropolis-like transitions (due to the $\eta$-regularization) is the only difference between the errors. Note that at leading-order, the third absolute moments $\mu_3(|b_\pm|)$ are $\mathcal{O}(1)$ for both Gaussian and (smooth) Metropolis-like weights: the reason is that the $|t|^3$ factor within renders the $1 / t$ singularity integrable. Thus, the moment decays dominantly as a Gaussian.

In principle, we could compose the two Trotter error results of \hyperref[prop:trotter-diss]{Proposition~\ref*{prop:trotter-diss}} and \hyperref[prop:trotter-B]{Proposition~\ref*{prop:trotter-B}} into one complete generator level error bound by a simple triangle inequality. But this gives no additional insight: in our implementation we Trotter split the generator and apply the coherent evolution by $B$ separate from the dissipative part $\mathcal{L}_\text{diss}$. Moreover, in implementation we have the flexibility of choosing different number of Trotter steps for these two parts since to control Trotter errors on OFT's might be simpler than controlling it on some less regular functions $b_\pm$ in $B$.

It is worth clarifying which structural properties of the CKG Lindbladian survive discretization and Trotterization. The three levels of approximation are: 
\begin{itemize}
    \item \textbf{Continuous} $\mathcal{L}$: exact KMS DB by construction  \hyperref[sec:kms-lindbladian]{KMS Lindbladian} based on \cite[Theorem I.1]{chen2023efficient}.
    \item \textbf{Discretized} $\mathcal{\bar{L}}$: approximate KMS DB. Quadrature errors break detailed balance. The size of the violation is controlled by the analysis in \hyperref[sec:quad-errors-diss]{Quadrature errors for $\mathcal{L}_\text{diss}$} and \hyperref[sec:quad-errors-B]{Quadrature errors for $B$}.
    \item \textbf{Trotterized} $\mathcal{\tilde{L}}$: further approximates KMS DB, with additional Trotter errors from \hyperref[prop:trotter-diss]{Proposition~\ref*{prop:trotter-diss}} and \hyperref[prop:trotter-B]{Proposition~\ref*{prop:trotter-B}} on top of the quadrature errors.
\end{itemize}
The exact KMS DB of the continuous Lindbladian relies on several ingredients acting in unison: the shifted KMS condition on the transition weight $\gamma$, the unique construction of the coherent term $B$ cancelling the anti-Hermitian residual of the quantum discriminant, the properties of the OFT construction. Not all of these are relevant at the discretized level, but one structural identity does remain exact here -- and the question is whether the imminent Trotterization preserves it. We cover this in the following remark.

\begin{remark} \label{rem:palindromic}
    \textup{(Palindromic product formulas).} The identity that survives discretization and can be made to survive Trotterization as well, is the \emph{OFT adjoint identity}:
    $$
    \hat{A}_{a'}(-\bar{\omega})^\dagger = \hat{A}_a(\bar{\omega}), \qquad \text{with } A_{a'} = A_a^\dagger \in \mathcal{A}.
    $$
    At the discretized level with exact Hamiltonian evolution, this identity holds exactly because the proof only requires that, (i) $\bar{f}(\bar{t}) \in \mathbb{R}$, (ii) $\bar{f}(-\bar{t}) = \bar{f}(\bar{t})$ (the used Gaussian filter is even), and (iii) $U(\bar{t})^\dagger = U(-\bar{t})$. All three hold for the discretized OFT \eqref{eq:oft-discr}. This identity ensures that the discriminant $\mathcal{D}_\beta$ from \eqref{eq:discriminant} satisfies $\mathcal{D}_\beta = \mathcal{D}_\beta^\dagger$, i.e. its anti-Hermitian part $\mathcal{A}(\rho_\beta, \mathcal{\bar{L}})$
    -- the measure of DB violation -- from \hyperref[def:approx-db]{Definition~\ref*{def:approx-db}} vanishes.

    At the Trotterized level, input (iii) must be replaced by $S_p^{(M)}(\bar{t})^\dagger = S_p^{(M)}(-\bar{t})$, which holds if and only if the product formula is \emph{palindromic} (or adjoint-symmetric).
\end{remark}
\begin{proof}
    A product formula $S_p(\delta t) = \prod_{j=1}^{J} e^{-ic_j H_{\pi(j)} \delta t}$ is palindromic if $c_{J+1-j} = c_j$ and $\pi(J+1-j) = \pi(j)$ for all $j$. For such a formula:
    $$
    S_p(\delta t)^\dagger = \prod_{j=J}^{1} e^{ic_j H_{\pi(j)} \delta t} = \prod_{j=1}^{J} e^{ic_{J+1-j} H_{\pi(J+1-j)} \delta t} = \prod_{j=1}^{J} e^{ic_j H_{\pi(j)}\delta t} = S_p(-\delta t),
    $$
    where the second step is the palindromic reindexing. This naturally extends to $M$ steps $S_p^{(M)}(\bar{t})^\dagger = S_p^{(M)}(-\bar{t})$. Strang splitting (and all other even Trotter formulas) satisfies this by construction.
    Then, we can explicitly check for $\tilde{A}_a(\bar{\omega})$ from \eqref{eq:trotter-oft}:
    \begin{align*}
        \tilde{A}_a(\bar{\omega})^\dagger &= \frac{1}{\sqrt{N}} \sum_{\bar{t} \in S_{t_0}} e^{i\bar{\omega}\bar{t}}\, \bar{f}(\bar{t})\, S_p^{(M)}(\bar{t})\, A_a^\dagger\, S_p^{(M)}(\bar{t})^\dagger \\
        &= \frac{1}{\sqrt{N}} \sum_{\bar{t}} e^{-i\bar{\omega}\bar{t}}\, \bar{f}(\bar{t})\, S_p^{(M)}(\bar{t})^\dagger\, A_{a'}\, S_p^{(M)}(\bar{t}),
    \end{align*}
    where we sent $\bar{t} \to -\bar{t}$ on the symmetric grid $S_{t_0}$, applied $\bar{f}(-\bar{t}) = \bar{f}(\bar{t})$ and the palindromic identity $S_p^{(M)}(-\bar{t}) = S_p^{(M)}(\bar{t})^\dagger$. We can also compute $\tilde{A}_{a'}(-\bar{\omega})$ separately, 
    \begin{align*}
        \tilde{A}_{a'}(-\bar{\omega}) &= \frac{1}{\sqrt{N}} \sum_{\bar{t}} e^{i\bar{\omega}\bar{t}}\, \bar{f}(\bar{t})\, S_p^{(M)}(\bar{t})\, A_{a'}\, S_p^{(M)}(\bar{t})^\dagger \\
        &= \frac{1}{\sqrt{N}} \sum_{\bar{t}} e^{-i\bar{\omega}\bar{t}}\, \bar{f}(\bar{t})\, S_p^{(M)}(\bar{t})^\dagger\, A_{a'}\, S_p^{(M)}(\bar{t}),
    \end{align*}
    and see that the two expressions coincide, $\tilde{A}_a(\bar{\omega})^\dagger = \tilde{A}_{a'}(-\bar{\omega})$.
\end{proof}
The anti-Hermitian part $\mathcal{A}(\rho_\beta, \mathcal{\bar{L}})$ of the quantum discriminant receives contributions from all sources of error: quadrature, Trotter, and possibly from non-palindromic product formulas (e.g. Lie-Trotter) violating the OFT adjoint identity. A palindromic formula eliminates the last channel entirely, leaving the Trotter contribution to the DB violation coming solely from the failure of $S_p^{(M)}(\bar{t})$ to commute with $\rho_\beta^{1/2}$ (defined with the exact Hamiltonian $H$).
For a non-palindromic formula (necessarily odd-order) formula of order $p$, the violation can be quantified with the commutator constant $\tilde{\alpha}_\text{comm}$ from \cite{childs2021theory} as 
$$
\|S_p^{(M)}(\bar{t})^\dagger - S_p^{(M)}(-\bar{t})\| = \mathcal{O}\left(\frac{\tilde{\alpha}_\text{comm} |\bar{t}|^{p + 1}}{M^p}\right), 
$$
so the violation enters at the same order as the Trotter error itself. In other words, the asymptotic scaling of the total incurring Trotter errors $\mathcal{O}(1 / M^p)$ is the same for both palindromic and non-palindromic formulas, but the constant prefactor is strictly smaller in the former case.

For commuting Hamiltonians, the product formula is exact for any $M$ and there is no Trotterization error. But the CKG framework was designed to handle non-commuting Hamiltonians. The commuting case is more naturally handles by the Davies generator directly \cite[text]{kastoryano2016quantum,bardet2023rapid}, for which the CKG Gaussian-smearing and coherent term $B$ are unnecessary.

\smallskip
\subsection{Generator splitting.}
In the implementation we use two distinct decomposition of the CKG generator \eqref{eq:ckg-L}. 

(i) First, the full generator is decomposed into per-jump contributions $\mathcal{L}_a$,
$$
\mathcal{L} = \sum_a \mathcal{L}_a.
$$
There are two ways to do this, each of them preserves the Gibbs state as the fixed point. We can either sweep over $A_a \in \mathcal{A}$ or randomly sample the jumps $A_a$ with probability $p_a$, where each contribution is rescaled by $1 / p_a$ such that after many samples we find the full generator $\mathcal{L}$ on average. Asymptotically both strategies give the same complexity, but there is some qualitative differences between the two. While the sequential Lie-Trotter splitting breaks KMS DB in general, at leading order it does so by an antisymmetric term which does not change the spectral gap. On the other hand, the random sampling is essentially a convex combination of KMS DB Lindbladians which leads to an KMS DB evolution on average [REF] \todo{ref}. Even though this sounds advantageous, it also means that the leading order errors do perturb the spectral gap. Since these errors are of the same order as errors from other sources in the algorithm, it made no significant difference which strategy we chose for our numerical simulations. But it is hard to argue against the better error scaling for an implementation, so we will mainly discuss the sequential Lie-Trotter splitting. (For other theoretical works, where implementation via discretization, Trotterization, etc. is not relevant, then keeping KMS DB might be more appealing since it keeps the quantum discriminant Hermitian.)

(ii) Second, for each fixed $a$, the generator $\mathcal{L}_a$ is further split into a coherent and a dissipative part, each implemented as their own, separate primitive:
$$
 \mathcal L_a = \mathcal L_{a,\mathrm{coh}} + \mathcal L_{a,\mathrm{diss}},
    \qquad
    \mathcal L_{a,\mathrm{coh}}(\cdot) := - i [B_a,\cdot].
$$
These two splittings are qualitatively different from each other: while the jump-wise splitting preserves the Gibbs state as the fixed point, the coherent-dissipative splitting generally does not. Consequently, the former introduces errors in the spectrum that can slow the mixing time, while the latter produces a genuine fixed-point bias. In either case, we analyse the incurring errors here and find a suiting way to control them for our purposes.

It is important to keep in mind that for the implementation cost, the relevant baseline error comes from the weak measurement simulation of the dissipative part $\mathcal{L}_\mathrm{diss}$. This primitive is only first-order accurate in the step size, i.e. for a step $\delta$ it has a local error $\mathcal{O}(\delta^2)$, and over the total simulated time $t_\text{mix}$ the accumulated error is $\mathcal{O}(t_\text{mix} \delta)$. Accordingly, all the generator decomposition we introduce in (i) and (ii) should respect this scaling as well.

\medskip
\noindent\textbf{(i) Jump-wise splitting.} Each of the Lindbladians $\mathcal{L}_a$ individually satisfies exact KMS DB for $\rho_\beta$, i.e. they are self-adjoint with respect to the KMS inner product ($\mathcal{L}_a^* = \mathcal{L}_a$). Consequently, the Gibbs state is the fixed point of each generator $\mathcal{L}_a(\rho_\beta) = 0$, and hence also of each semigroup channel $e^{\delta\mathcal{L}_a}(\rho_\beta) = \rho_\beta$. Note that, the kernel of an individual $\mathcal{L}_a$ is generically larger than $\mathrm{Span}\{\rho_\beta\}$. Primitivity is a property of the full generator $\mathcal{L} = \sum_a \mathcal{L}_a$ only, not of each summand.

The simplest (yet sufficient) way to split the generator sequentially, is to use the first-order Lie-Trotter splitting:
\begin{equation}
    \Phi_{\mathcal{A}} := \ee^{\delta \mathcal{L}_{M_\mathcal{A}}} \circ \cdots \circ \ee^{\delta \mathcal{L}_2} \circ \ee^{\delta \mathcal{L}_1},
\end{equation}
with $M_\mathcal{A} := |\mathcal{A}|$ denoting the number of jump operators. By the standard first-order product formula bound, applied to the superoperators $\mathcal{L}_a$, we can write  
\begin{equation} \label{eq:gen-split-approx-error}
    \left\|\Phi_{\mathcal{A}} - e^{\delta \mathcal{L}}\right\|_{1\to 1} \leq \frac{\delta^2}{2} \left\|\sum_{a < b} [\mathcal{L}_a, \mathcal{L}_b]\right\|_{1\to 1} + O(\delta^3).
\end{equation}
Under the normalization $\|\sum_a A_a^\dagger A_a\| \leq 1$, each per-jump Lindbladian has norm
\begin{equation}
    \|\mathcal{L}_a\|_{1 \to 1} = O\!\left(\frac{\log(\beta\|H\|)}{M_\mathcal{A}}\right),
\end{equation}
where the $\log(\beta\|H\|)$ factor originates from the Metropolis-like coherent term $B_a$ [REF within thesis]\todo{ref}. The dissipative part alone satisfies $\|\mathcal{L}_{a,\mathrm{diss}}\|_{1\to 1} \leq 4\|A_a^\dagger A_a\|$, so the coherent contribution dominates. Applying the triangle inequality to the commutator sum in \eqref{eq:gen-split-approx-error}:
\begin{equation}
    \left\|\sum_{a<b} [\mathcal{L}_a, \mathcal{L}_b]\right\|_{1\to 1} 
    \leq \binom{M_\mathcal{A}}{2} \cdot 2 \max_{a,b} \|\mathcal{L}_a\|_{1\to 1}\|\mathcal{L}_b\|_{1\to 1} 
    = O\!\left(\log^2(\beta\|H\|)\right),
\end{equation}
for which the $M_\mathcal{A}$ dependence cancelled. The per-step approximation error is therefore
\begin{equation} \label{eq:gen-split-explicit-bound}
    \left\|\Phi_{\mathcal{A}} - e^{\delta \mathcal{L}}\right\|_{1\to 1} = O\!\left(\delta^2 \log^2(\beta\|H\|)\right),
\end{equation}
independently of $M_\mathcal{A}$. For Gaussian transition weights, where the $\eta$-regularization of the coherent term is unnecessary, the $\log^2$ factor is absent and the bound simplifies to $O(\delta^2)$.

There is some structure to this error however, that is worth exploring.
Since each $\ee^{\delta \mathcal{L}_a}$ is KMS DB, and the KMS adjoint reverses the composition, $({\Phi_1 \circ \Phi_2})^* = \Phi_2^* \circ \Phi_1^*$, we obtain: 
\begin{equation}
    \Phi_{\mathcal{A}}^* = e^{\delta \mathcal{L}_1} \circ e^{\delta \mathcal{L}_2} \circ \cdots \circ e^{\delta \mathcal{L}_{M_\mathcal{A}}} \neq \Phi_{\mathcal{A}}.
\end{equation}
This is the same kind of DB violation as what we have seen for non-palindromic Trotter products for the Hamiltonian time evolutions in the previous section: the choice of ordering in the product formula breaks time-reversal symmetry. The DB violation is quantified by the difference between the forward and reverse orderings. By Baker-Campbell-Hausdorff (BCH), the effective generators of the two orderings differ in the sign of the leading commutator term, 
$$
\delta \sum_a \mathcal{L}_a \pm \frac{\delta^2}{2}\sum_{a < b}[\mathcal{L}_a, \mathcal{L}_b] + O(\delta^3),
$$
yielding the DB violation at leading order, 
\begin{equation}
    \left\|\Phi_{\mathcal{A}} - \Phi_{\mathcal{A}}^*\right\|_{1\to 1} = O\!\left(\delta^2 \left\|\sum_{a<b}[\mathcal{L}_a, \mathcal{L}_b]\right\|_{1\to 1}\right),
\end{equation}
which is of the same order as the approximation error \eqref{eq:gen-split-approx-error}. Crucially, even though the sequential product breaks detailed balance, it still has the Gibbs state as its fixed point exactly: since each factor fixes $\rho_\beta$, so does $\Phi_{\mathcal{A}}(\rho_\beta) = \rho_\beta$. Moreover, the DB violation has a definite algebraic structure: the commutator of any two KMS self-adjoint generators is KMS anti-self-adjoint: 
$$
[\mathcal{L}_a, \mathcal{L}_b]^* = [\mathcal{L}_b^*, \mathcal{L}_a^*] = -[\mathcal{L}_a, \mathcal{L}_b].
$$
Hence, the leading BCH correction involving the commutator is a purely antisymmetric perturbation of the generator. This has a direct spectral consequence. Since $\mathcal{L}$ is KMS self-adjoint, its eigenvalues in the KMS inner product are real. By standard perturbation theory, an anti-self-adjoint term can only lead to a purely imaginary perturbation to the spectrum $\{\lambda_i\}$ of $\mathcal{L}$ or written out:
\begin{equation} \label{eq:gen-split-imaginary-delta2-errors}
    \Delta\lambda_k = \frac{\delta}{2}\langle \psi_k, \sum_{a<b}[\mathcal{L}_a, \mathcal{L}_b](\psi_k) \rangle_{\mathrm{KMS}} \in \ii\mathbb{R},
\end{equation}
where $\{\psi_k\}$ are the KMS-orthonormal eigenvectors of $\mathcal{L}$. Therefore, the jump-wise generator splitting is less harmful than a generic first-order product formula: for the $\delta$ step channel it lets the spectral gap of $\Phi_\mathcal{A}$ agree with that of $\ee^{\delta \mathcal{L}}$ up to $\mathcal{O}(\delta^3)$.

The sequential product $\Phi_\text{seq}$ leading to this specific structure for the error terms, raises a question, namely: could this anti-KMS-adjoint error term accelerate the convergence to the Gibbs state? The only reason we would even ask such a question is because this is precisely the structure that has been shown to accelerate mixing in the continuous generator setting. For classical diffusions, \cite{hwang2005accelerating} showed that such an antisymmetric addition can only improve the spectral gap. A similar fact for the quantum case was established by \cite{li2025speeding}, showing that for a DB QMS this could mean an up to quadratic speed-up in mixing. But none of these papers are one-to-one applicable to our case, nor did the simulations show any noticeable speed-up. It seems that turning these pesky errors into an advantage keeps eluding us.
\todo{ref, we could check this numerically.}

As a final note to the jump-wise splitting: higher order product formulas would make no meaningful difference here, especially if the weak measurement scheme within $\mathcal{L}_a$ already introduces errors of $\mathcal{O}(\delta^2)$ per step.

\medskip
\noindent\textbf{(ii) Coherent--Dissipative}
Splitting the coherent term from the dissipative one and implementing their respective evolutions separately from each other, is a qualitatively different process compared to the jump-wise splitting. The coherent part is chosen precisely such that the sum of the coherent and dissipative parts are exactly KMS DB. Unlike in the previous case (i), now the separate terms do not satisfy KMS DB, thus the emerging errors already shift the fixed point at leading orders. To see this let us write out the first-order splitting:
\begin{equation}
    \ee^{\delta \mathcal L_{a,\mathrm{coh}}}\ee^{\delta \mathcal L_{a,\mathrm{diss}}}
    = \exp\!\left(
        \delta \mathcal L_a
        + \frac{\delta^2}{2}[\mathcal L_{a,\mathrm{coh}},\mathcal L_{a,\mathrm{diss}}]
        + O(\delta^3)
    \right),
\end{equation}
so that the effective generator is 
\begin{equation} \label{eq:gen-split-coh-diss}
    \mathcal L_{a,\mathrm{eff}}
    = \mathcal L_a
    + \frac{\delta}{2}[\mathcal L_{a,\mathrm{coh}},\mathcal L_{a,\mathrm{diss}}]
    + O(\delta^2).
\end{equation}
Since neither $L_{a,\mathrm{coh}}$ nor $L_{a,\mathrm{diss}}$ satisfy KMS DB, their commutator does not become KMS anti-self-adjoint. Therefore, the fixed point is now shifted, and the leading perturbation now has an impact even on the real part of the spectrum. 

The already used and convenient way to express this is via the fixed-point perturbation bound of \cite[Lemma II.1]{chen2023quantum} \todo{ref to where we used this} 
\begin{equation}
        \|\rho_{\mathrm{fix}}(\mathcal L_1)-\rho_{\mathrm{fix}}(\mathcal L_2)\|_1
    \le 4\, \|\mathcal L_1-\mathcal L_2\|_{1\to1}\, t_{\mathrm{mix}}(\mathcal L_1).
\end{equation}
Applying this with $\mathcal L_1=\mathcal L_a$ and $\mathcal L_2=\mathcal L_{a,\mathrm{eff}}$ yields a fixed-point shift of order $\mathcal O\left(t_\text{mix}\, \delta\right)$. If the substeps (coherent evolution via $B$ and weak measurement scheme for $\mathcal{L}_\text{diss}$) were implemented exactly, then a Strang splitting would improve this BCH error to second order in the effective generator:
\begin{equation}
    e^{\frac{\delta}{2}\mathcal L_{a,\mathrm{coh}}}
    e^{\delta \mathcal L_{a,\mathrm{diss}}}
    e^{\frac{\delta}{2}\mathcal L_{a,\mathrm{coh}}}
    = \exp\!\left(\delta \mathcal L_a + O(\delta^3)\right),
\end{equation}
which would reduce the induced fixed-point bias to $\mathcal O(t_{\mathrm{mix}}\delta^2)$ over the full evolution. However, this improvement is not useful in our implementation, since any more accuracy beyond $\mathcal O\left(t_\text{mix}\, \delta\right)$ is washed away by the errors due to the weak measurement subroutine. 

For this reason, higher-order coherent-dissipative product formulas do not improve the asymptotic end-to-end scaling in our setting. If the weak measurement scheme were to be replaced by some more accurate variant, then we could always adjust all the above product formulas to a higher order variant to comply with the baseline error. \eqref{eq:relative-entropy}

% Use Trotter splitting also for coherent and dissipative parts. t_mix scaling becomes poly with this compared to linear for Chen's, but doesn't need QSVT, doesn't need a ton of repetitions due to subnorm, coherent ancillas are independent and incoherent from dissipative ancillas. It is a valid realization with NISQ in mind, even if asymptotically it's worse.

\smallskip
\label{sec:algorithm}
\subsection{Algorithm.}

\begin{algorithm}[H]
\caption{\textsc{CKG-GibbsSampler}}
\label{alg:main}
\begin{algorithmic}[1]
\small
\Require
  System Hamiltonian $H = \sum_{k=1}^{K} H_k$ on $n$ qubits;
  jump set $\{A_a\}_{a\in\mathcal{A}}$ with $\bigl\lVert\sum_a A_a^\dagger A_a\bigr\rVert \le 1$;
  inverse temperature $\beta$;
  Gaussian filter width $\sigma$ for filter function $f$ \eqref{eq:gaussian-filter-t};
  transition weight $\gamma$ (e.g.\ $\gamma_M^{(s)}$ from~\eqref{eq:smooth-metro});
  functions $b_\pm$ for the coherent term \eqref{eq:b_plus-s-eta} and \eqref{eq:b_minus};
  step size $\delta$;
  target trace-distance accuracy $\varepsilon$
\Ensure System register $S$ in state
  $\rho \approx \rho_\beta = \ee^{-\beta H}/\mathrm{tr}(\ee^{-\beta H})$
\Statex
\State $L \gets \bigl\lceil t_{\mathrm{mix}}(\mathcal{L})\,
       \log(2/\varepsilon)\,/\,\delta \bigr\rceil$
       \Comment{total Lindbladian steps}
\State Initialise $S$ in an arbitrary state $\rho_0$
\For{$\ell = 1, \ldots, L$}
  \For{$a = 1, \ldots, M_{\mathcal{A}}$}
       \Comment{sequential sweep over jumps}
    \State \Call{CoherentStep}{$S, a, \delta, b_\pm$}
           \Comment{$\ee^{-\mathrm{i}\delta B_a}$, Alg.~\ref{alg:coh}}
    \State \Call{DissipativeStep}{$S, a, \delta, f, \gamma$}
           \Comment{$\ee^{\delta\mathcal{L}_{a,\mathrm{diss}}}$, Alg.~\ref{alg:diss}}
  \EndFor
\EndFor
\State \Return $S$
\end{algorithmic}
\end{algorithm}


\begin{figure}[ht]
  \centering
  \resizebox{\textwidth}{!}{%
  \begin{quantikz}[row sep=0.6cm, column sep=0.5cm]
  %
  % Cell-by-cell layout (15 cells per row):
  %   1: \lstick          (left label)
  %   2: wire decl        (\qw or \qwbundle{r})
  %   3: R_0 / qw         (coherent: first QSP rotation)
  %   4: octrl / U_{B_a}  (Lambda(U_{B_a}))
  %   5: octrl / R_T      (Lambda(R_T))
  %   6: R_k / qw         (coherent: inner-loop QSP rotation)
  %   7: meter / qw       (terminate coherent ancilla)
  %   8: setwiretype{n}   (transition cell A: hide outgoing wire)
  %   9: new lstick       (transition cell B: re-init wire as dissipative ancilla)
  %  10: U_{diss} / qw    (dissipative forward)
  %  11: Y_delta / octrl  (weak measurement rotation)
  %  12: U_{diss}^dag     (dissipative reverse, 0-ctrl by q_delta)
  %  13: meter / rstick   (terminate dissipative ancilla / system output)
  %
  % Meter and setwiretype{n} are placed in SEPARATE cells so that the
  % wire leading into the meter stays visible (otherwise quantikz2 applies
  % {n} retroactively to the incoming wire and breaks the meter).
  %
  % ==================================================================
  % Row 1:  q_QSP  (coherent, cells 3-7)  -->  q_delta  (dissipative, cells 9-13)
  % ==================================================================
    \lstick{$\ket{0}_{\mathrm{QSP}}$}
      & \qw
      & \gate{R(\theta_0,\phi_0,\lambda_0)}
      & \octrl{1}
      \gategroup[4,steps=3,style={dashed,rounded corners,fill=blue!6,inner sep=5pt},
                 background,label style={label position=above,yshift=0.30cm,anchor=south}]
        {\scriptsize \textsc{CoherentStep}: repeat $d$ times with angles $(\theta_k,\phi_k)\Rightarrow\ee^{-\ii\delta B_a}$}
      & \octrl{1}
      & \gate{R(\theta_k,\phi_k,0)}
      & \meter{}
      & \setwiretype{n}
      & \lstick{$\ket{0}_\delta$} \setwiretype{q} \wireoverride{n}
      & \qw
      & \qw
      \gategroup[4,steps=4,style={dashed,rounded corners,fill=red!6,inner sep=5pt},
                 background,label style={label position=below,yshift=-0.40cm,anchor=north}]
        {\scriptsize \textsc{DissipativeStep}: weak measurement simulation of $\ee^{\delta\mathcal{L}_{a,\mathrm{diss}}}$}
      & \gate{Y_\delta}
      & \octrl{1}
      & \meter{}
      \\
  %
  % ==================================================================
  % Row 2:  T_+  -->  q_gamma  (holds the multi-row gate labels in both phases:
  %         U_{B_a}/R_T on coherent, U_{diss}/U_{diss}^dag on dissipative.
  %         q_gamma sits directly under q_delta so the weak measurement
  %         control is the minimal \octrl{-1}.)
  % ==================================================================
    \lstick{$\ket{0}_{T_+}$}
      & \qwbundle{r_+}
      & \qw
      & \gate[3]{U_{B_a}}
      & \gate[2]{R_T}
      & \qw
      & \meter{}
      & \setwiretype{n}
      & \lstick{$\ket{0}_\gamma$} \setwiretype{q} \wireoverride{n}
      & \qw
      & \gate[3]{U_{\mathrm{diss}}}
      & \octrl{-1}
      & \gate[3]{U_{\mathrm{diss}}^\dagger}
      & \meter{}
      \\
  %
  % ==================================================================
  % Row 3:  T_-  -->  Omega  (bundled throughout; cells 4, 5, 10, 12 empty as
  %         they are consumed by the multi-row gates sitting on row 2 above.)
  % ==================================================================
    \lstick{$\ket{0}_{T_-}$}
      & \qwbundle{r_-}
      & \qw
      &
      &
      & \qw
      & \meter{}
      & \setwiretype{n}
      & \lstick{$\ket{0}_\Omega$} \setwiretype{q} \wireoverride{n}
      & \qwbundle{r}
      &
      & \qw
      &
      & \meter{}
      \\
  %
  % ==================================================================
  % Row 4:  S  (system register, active throughout with no wire-type change)
  % ==================================================================
    \lstick{$\rho$}
      & \qwbundle{n}
      & \qw
      &
      & \qw
      & \qw
      & \qw
      & \qw
      & \qw
      & \qw
      &
      & \qw
      &
      & \qw
      & \rstick{$\ee^{\delta\mathcal{L}_a}(\rho)+\mathcal{O}(\delta^2)$}\qw
  \end{quantikz}%
  }
  \caption{Algorithm~\ref*{alg:main}: one $\delta$-step of the CKG Lindbladian
  $\ee^{\delta\mathcal{L}_a} = \ee^{-\ii\delta B_a}\cdot
  \ee^{\delta\mathcal{L}_{a,\mathrm{diss}}}+\mathcal{O}(\delta^2)$ for a fixed
  jump $A_a$, split into a GQSP \emph{coherent block} (blue, $d$ repetitions
  of $\Lambda(U_{B_a}),\Lambda(R_T),R(\theta_k,\phi_k,0)$) and a
  weak measurement \emph{dissipative block} (red).  Mid-circuit measurements
  mark ancilla reuse: $q_{\mathrm{QSP}}, T_\pm \!\to q_\delta, q_\gamma, \Omega$.}
  \label{circ:algorithm-delta-step}
\end{figure}

It is time to assemble the subroutines discussed in the preceding subsections into a complete quantum algorithm for Gibbs state preparation via the CKG Lindbladian \eqref{eq:ckg-L}. The resulting procedure, Algorithm~\eqref{alg:main}, simulates the continuous-time Markov semigroup $\ee^{t \mathcal{L}}$ by repeated application of small $\delta$-step channels that approximate $\ee^{\delta \mathcal{L}_a}$ for each jump $a\in\mathcal{A}_a$. Its structure draws on three key ingredients: the \emph{weak measurement simulation scheme} of \cite{chen2023quantum}, which implements the dissipative channels $\ee^{\delta \mathcal{L}_{a, \text{diss}}} + \mathcal{O}(\delta^2)$; The \emph{coherent correction term} $B_a$ \eqref{eq:coherent-time-domain} introduced in \cite{chen2023efficient} (or its Trotterized version \eqref{eq:B-trotter}), that makes the CKG generator $\mathcal{L}$ exactly KMS detailed balanced; and \emph{Generalized Quantum Signal Processing} (GQSP) from \cite{motlagh2024generalized}, which provides a near-optimal method to realize the unitary $\ee^{-\ii \delta B_a}$ up to $\mathcal{O}(\delta^2)$ errors from the block encoding of $B_a$.
The end result is a quantum channel that, after $L$ iterations of the outer loop, produces a state $\rho$ that is at most $\varepsilon$ trace distance from the Gibbs state $\rho_\beta$.

We now give a few comments on Algorithm~\eqref{alg:main}. 
\begin{itemize}
    \item The system register $S$ can be initalised in any arbitrary state $\rho_0$, due to the primitivity of $\mathcal{L}$: $\rho_\beta$ is the unique fixed point and the semigroup is ergodic, thus any initial state converges to $\rho_\beta$ at the same asymptotic rate. In practice however, it is ideal to choose a state that has a reasonable overlap with $\rho_\beta$ (and is not too costly to prepare), e.g. the maximally mixed state $\mathbf{1} / d$, to eliminate part of the initial sampling steps, related to the factor $1/\sqrt{\lambda_{\min}(\rho_\beta)}$.
    \item From the mixing time definition \eqref{eq:mixing-time-defi} and the spectral bound \eqref{eq:quantum-spectral-bound}, after evolving some initial state $\rho_0$ for a total Lindbladian time $t = t_{\mathrm{mix}}(\mathcal{L})\,\log(2/\varepsilon)$ we are guaranteed to have $\|e^{t\mathcal{L}}(\rho_0) - \rho_\beta\|_1 \leq \varepsilon$.
    \item For each $\delta$ step of the outer-loop we perform a sequential sweep over all $M_\mathcal{A} = |\mathcal{A}|$ jumps. This is a Lie-Trotter splitting of $e^{\delta\mathcal{L}}$ to $\Phi_\mathcal{A} := e^{\delta\mathcal{L}_{M_\mathcal{A}}} \circ \cdots \circ e^{\delta\mathcal{L}_1}$, introducing approximation errors of $\mathcal{O}(\delta^2)$ \eqref{eq:gen-split-approx-error} but perturbing the real part of the spectrum only at $\mathcal{O}(\delta^3)$ \eqref{eq:gen-split-imaginary-delta2-errors}.
    \item In the inner-loop for a fixed $a$ we split the generator further as $\mathcal{L}_a = \mathcal{L}_{a,\mathrm{coh}} + \mathcal{L}_{a,\mathrm{diss}}$, introducing an error to the channel $\ee^{\delta \mathcal{L}_a}$ of order $\mathcal{O}(\delta^2)$ order \eqref{eq:gen-split-coh-diss}.
    \item \textsc{CoherentStep} is the subalgorithm that implements the unitary $\ee^{-\ii \delta B_a} + \mathcal{O}(\delta^2)$ via GQSP, see Alg.~\eqref{alg:coh}.
    \item \textsc{DissipativeStep} is the subalgorithm that implements the dissipation $e^{\delta\mathcal{L}_{a,\mathrm{diss}}}(\rho) + \mathcal{O}(\delta^2)$ via the weak measurement scheme, see Alg.~\eqref{alg:diss}.
\end{itemize}
\medskip 

%% ---------------------------------------------------------------
\setcounter{algorithm}{0}
\renewcommand{\thealgorithm}{\arabic{algorithm}a}
\begin{algorithm}[!htbp]
\caption{\textsc{CoherentStep}: Hamiltonian simulation of
         $e^{-\mathrm{i}\delta B_a}$ via block encoding and QSP}
\label{alg:coh}
\begin{algorithmic}[1]
\small
\Require System register $S$, jump index $a$, step size $\delta$
\Ensure  $S$ evolved by $e^{-\mathrm{i}\delta B_a}$
  up to QSP and Trotter precision of $\mathcal{O}(\delta^2)$
\Statex
\State \textbf{Ancillas:}
  outer time register $T_-$ ($r_-$ qubits),
  inner time register $T_+$ ($r_+$ qubits),
  QSP signal qubit $q_{\mathrm{QSP}}$
\Statex
\Statex\textit{--- Block encoding $U_{B_a}$ of
  $B_a/\alpha$~\eqref{eq:B-discr}, with
  $\alpha := \lVert b_-\rVert_1\,\lVert b_+\rVert_1$ ---}
\Statex\textit{Forward pass:}
\State \textsc{StatePrep}$(T_-)$:
  $|0\rangle \to |b_-\rangle :=
   \displaystyle\sum_{\bar{t}}
   \sqrt{\frac{t_0\,|b_-(\bar{t})|}{\alpha_-}}\;
   \mathrm{sgn}\bigl(b_-(\bar{t})\bigr)\,|\bar{t}\rangle$
  \Comment{outer kernel~\eqref{eq:b_minus}; $\alpha_- = \lVert b_-\rVert_1$}
\State \textsc{Ctrl-Trott}$(T_-, S)$: apply
  $S_2^{(M_-)}(-\bar{t}/\sigma)$ on $S$, controlled on $T_-$
  \Comment{outer: $\approx e^{-\mathrm{i}H\bar{t}/\sigma}$}
\State \textsc{StatePrep}$(T_+)$:
  $|0\rangle \to |b_+\rangle :=
   \displaystyle\sum_{\bar{t}'}
   \sqrt{\frac{t_0'\,|b_+(\bar{t}')|}{\alpha_+}}\;
   \mathrm{sgn}\bigl(b_+(\bar{t}')\bigr)\,|\bar{t}'\rangle$
  \Comment{inner kernel~\eqref{eq:b_plus-s-eta}; $\alpha_+ = \lVert b_+\rVert_1$}
\Statex
\Statex \textit{Inner Heisenberg evolution}
  $$\tilde{A}_a^\dagger(\beta\bar{t}')\,\tilde{A}_a(-\beta\bar{t}')
  = S_2^{(M_+)}(\beta\bar{t}')\,A_a^\dagger\,
    S_2^{(M_+)}(-2\beta\bar{t}')\,A_a\,
    S_2^{(M_+)}(\beta\bar{t}'),$$
  \textit{controlled on $T_+$}:
\State \textsc{Ctrl-Trott}$(T_+, S)$: apply
  $S_2^{(M_+)}(+\beta\bar{t}')$ on $S$
  \Comment{$\approx e^{+\mathrm{i}\beta H\bar{t}'}$}
\State Apply $A_a^\dagger$ on $S$
\State \textsc{Ctrl-Trott}$(T_+, S)$: apply
  $S_2^{(M_+)}(-2\beta\bar{t}')$ on $S$
  \Comment{$\approx e^{-2\mathrm{i}\beta H\bar{t}'}$}
\State Apply $A_a$ on $S$
\State \textsc{Ctrl-Trott}$(T_+, S)$: apply
  $S_2^{(M_+)}(+\beta\bar{t}')$ on $S$
  \Comment{$\approx e^{+\mathrm{i}\beta H\bar{t}'}$}
\Statex
\State \textsc{StatePrep}$^\dagger(T_+)$:
  $|b_+\rangle \to |0\rangle$
  \Comment{reflect / uncompute inner register}
\State \textsc{Ctrl-Trott}$(T_-, S)$: apply
  $S_2^{(M_-)}(+\bar{t}/\sigma)$ on $S$
  \Comment{undo outer evolution}
\State \textsc{StatePrep}$^\dagger(T_-)$:
  $|b_-\rangle \to |0\rangle$
  \Comment{reflect / uncompute outer register}
\Statex \Comment{Block encoding:
  $\langle 0|_{T_-}\!\langle 0|_{T_+}\;
   U_{B_a}\;
   |0\rangle_{T_-}\!|0\rangle_{T_+}
   \;=\; B_a\,/\,\alpha$
  acting on $S$}
\Statex
\Statex \textit{--- QSP: implement $e^{-\mathrm{i}\delta B_a}$
  from the block encoding ---}
\State Apply QSP~\cite{motlagh2024generalized} with precomputed rotation angles     
    $\{(\theta_k,\phi_k)\}_{k=0}^{d}$ for the Laurent polynomial $P(e^{i\theta}) \approx e^{-\mathrm{i}\delta\alpha\,\sin\theta}$
\Statex \Comment{Degree                                        
    $d = 1$ calls to $U_{B_a}$}
\State Measure and reset ancillas $q_{\mathrm{QSP}}$, $T_\pm$
\end{algorithmic}
\end{algorithm}
\renewcommand{\thealgorithm}{\arabic{algorithm}}
%% ---------------------------------------------------------------
\begin{figure}[ht]
  \centering
  \resizebox{1.05\textwidth}{!}{%
  \begin{quantikz}[row sep=0.9cm, column sep=0.42cm]
  % Row 1:  T_-  (outer time register; controls cells 4 and 12)
  % ==================================================================
    \lstick{$\ket{0}_{T_-}$}
      & \qwbundle{r_-}
      & \gate{\text{PREP}'_-}
      & \ctrl[style=halfctrl]{2}
      & \qw
      & \qw
      & \qw
      & \qw
      & \qw
      & \qw
      & \qw
      & \ctrl[style=halfctrl]{2}
      & \gate{\text{PREP}_-^\dagger}
      & \rstick{$\bra{0}_{T_-}$}\qw
      \\
  %
  % ==================================================================
  % Row 2:  T_+  (inner time register; controls cells 6, 8, 10;
  %              prepared in cell 5, uncomputed in cell 11)
  % ==================================================================
    \lstick{$\ket{0}_{T_+}$}
      & \qwbundle{r_+}
      & \qw
      & \qw
      & \gate{\text{PREP}'_+}
      \gategroup[2,steps=7,style={dashed,rounded corners,fill=blue!6,inner sep=5pt},
                 background,label style={label position=below,yshift=-0.35cm,anchor=north}]
        {\scriptsize \textsc{Inner sum}:
         $\sum_{\bar{t}'} b_+(\bar{t}')\tilde{A}_a^\dagger(\beta\bar t')\,\tilde{A}_a(-\beta\bar t')$
         controlled on $T_+$}
      & \ctrl[style=halfctrl]{1}
      & \qw
      & \ctrl[style=halfctrl]{1}
      & \qw
      & \ctrl[style=halfctrl]{1}
      & \gate{\text{PREP}_+^\dagger}
      & \qw
      & \qw
      & \rstick{$\bra{0}_{T_+}$}\qw
      \\
  %
  % ==================================================================
  % Row 3:  S  (system register; target of every controlled evolution
  %             and of the two jump applications A_a^dag, A_a)
  % ==================================================================
    \lstick{$\ket{\psi}$}
      & \qwbundle{n}
      & \qw
      & \gate{S_2\left(\frac{\bar t}{\sigma}\right)}
      & \qw
      & \gate{S_2(\beta\bar t')}
      & \gate{A_a}
      & \gate{S_2(-2\beta\bar t')}
      & \gate{A_a^\dagger}
      & \gate{S_2(\beta\bar t')}
      & \qw
      & \gate{S_2\left(-\frac{\bar t}{\sigma}\right)}
      & \qw
      & \rstick{$\tilde{B}_a\ket{\psi} / \alpha$}\qw
  \end{quantikz}%
  }
  \caption{Block encoding $U_{B_a}$ of the Trotterized coherent
  Metropolis term $\tilde{B}_a/\alpha$ with $\alpha=\lVert b_-\rVert_1\lVert b_+\rVert_1$
  (Alg.~\ref*{alg:coh}, \eqref{eq:B-block-encoding}), on outer/inner time
  registers $T_\mp$. Half-filled controls denote uniform (SELECT) control
  over the bundled time registers.
  $\text{PREP}'_\pm$ prepares the signed state
  $\sum_{\bar t}[b_\pm(\bar t)/\sqrt{|b_\pm(\bar t)|}]\ket{\bar t}/\sqrt{\lVert b_\pm\rVert_1}$
  and $\text{PREP}_\pm$ the magnitude state
  $\sum_{\bar t}\sqrt{|b_\pm(\bar t)|/\lVert b_\pm\rVert_1}\ket{\bar t}$,
  so that amplitudes multiply to the signed kernel $b_\pm(\bar t)$. Controlled
  evolutions are Trotterized with the second-order Strang splitting $S_2$ with $M_\pm$ Trotter steps respectively.}
  \label{circ:U_B_a}
\end{figure}

\medskip
\noindent \textsc{CoherentStep:} \textbf{Coherent evolution of $B$ with QSP.} The \textsc{CoherentStep}, \hyperref[alg:coh]{Algorithm~\ref*{alg:coh}}, implements the unitary evolution $\ee^{-\ii \delta B_a}$ generated by the per-jump Hermitian term $B_a$. This is the component that distinguishes the CKG algorithm from the CKBG one: without it (and without the shifted transitions weights) the algorithm reduces to the CKBG Gibbs sampler [ref to text \todo{ref}], which satisfies only approximate GNS detailed balance. Recall from the \hyperref[sec:kms-lindbladian]{KMS Lindbladian~\ref*{sec:kms-lindbladian}} subsection that the coherent term $B$ was derived \eqref{eq:B-bohr-domain} to cancel the emergent anti-Hermitian part of the quantum discriminant, upgrading the Lindbladian to exact KMS detailed balance. Its per-jump contribution $B_a$ enters the coherent-dissipative splitting \eqref{eq:gen-split-coh-diss} as the generator $\mathcal{L}_{a,\mathrm{coh}}(\cdot) = -i[B_a, \cdot\,]$, with the first-order splitting error of $\mathcal{O}(\delta^2)$ already matched by the weak measurement error of the dissipative step. Of course, discretization and Trotterization errors analysed in the preceding subsections will re-introduce controlled violations of the detailed balance, but these are precisely the errors we have learned to bound.

The subalgorithm breaks down to two steps: first we construct a block encoding $U_B$, satisfying
\begin{equation} \label{eq:B-block-encoding}
    \langle 0|_{T_-}\langle 0|_{T_+}\,U_{B_a}\,|0\rangle_{T_-}|0\rangle_{T_+} = B_a/\alpha,
\end{equation}
where $\alpha = \|b_-\|_1\,\|b_+\|_1$ is the subnormalisation factor of $B_a$. Then, we use GQSP \cite{motlagh2024generalized} (in most cases with degree $d = 1$) to implement $\ee^{-\ii \delta B_a}$ from $U_{B_a}$. In the following, we will give a few more details about each of the steps.

The block encoding of the coherent term \eqref{eq:B-block-encoding} containing the nested time integrals of \eqref{eq:coherent-time-domain} is done by a nested LCU. [Ref to prelims \todo{ref}]. This has been established in the CKG paper \cite[Proposition III.1]{chen2023efficient}, where they also prove its asymptotic scaling. The full nested double-integral structure with two time registers has been left implicit by them. Here we unpack it completely.

The starting point is the nested integral form of $B_a$ \eqref{eq:coherent-time-domain} which is discretized by nested Riemann-sums \eqref{eq:B-discr} with prefactors $t_0, t_0'$. In the LCU framework \cite{childs2012hamiltonian,berry2015simulating} this manifests as a state preparation / controlled-evolution / unpreparation on two separate ancilla registers: an outer time register $T_-$ of $r_-$ qubits, and an inner time register $T_+$ of $r_+$ qubits. 

Steps 2 and 4 of \hyperref[alg:coh]{Algorithm~\ref*{alg:coh}} prepare these registers by
\begin{equation}
    \begin{split}
        |0\rangle_{T_-} &\;\longrightarrow\; |b_-\rangle := \sum_{\bar t}\, \sqrt{\frac{t_0\,|b_-(\bar t)|}{\alpha_-}}\;\mathrm{sgn}\!\big(b_-(\bar t)\big)\,|\bar t\rangle, \qquad \alpha_- = \|b_-\|_1, \\
        |0\rangle_{T_+} &\;\longrightarrow\; |b_+\rangle := \sum_{\bar t'}\, \sqrt{\frac{t_0'\,|b_+(\bar t')|}{\alpha_+}}\;\mathrm{sgn}\!\big(b_+(\bar t')\big)\,|\bar t'\rangle, \qquad \alpha_+ = \|b_+\|_1.
    \end{split}
\end{equation}
Here, the absolute values go into the amplitudes (ensuring normalization $\langle b_\pm | b_\pm \rangle = 1$) and the signs go into the phases. For the outer kernel $b_-$ \eqref{eq:b_minus}, which is real-valued, $\mathrm{sgn}(b_-(\bar t)) \in \{-1,+1\}$ is the usual sign function. For the inner kernel $\mathrm{sgn}(b_+(\bar t'))$ \eqref{eq:b_plus-s-eta}, which is complex-valued in general, the notation $\mathrm{sgn}(b_+(\bar t'))$ should be understood as the unit-modulus phase $b_+(\bar t')/|b_+(\bar t')|$. The same convention applies to the Gaussian case \eqref{eq:b_plus-gauss}. 

The state preparation costs for $\ket{b_\pm}$ depend on the register sizes and the regularity of the kernels. For actual costs it is best to look at our numerical analysis [REF, \todo{ref}], but we can already say something qualitative here based on our quadrature error results for $r_-$ \eqref{eq:r_-}, and $r_+$ in the Gaussian case \eqref{eq:r_+_gaussian}, in the Metropolis case \eqref{eq:r_+_metro}. For our parameter regimes (moderate temperatures $\beta \leq 20$ and system sizes $n \leq 12$) the register sizes $r_- = 5-8$ and $r_+ = 10-15$ are already sufficient for our purposes to block encode $B_a$. There are several ways to prepare a quantum state, and as usual, the question boils down to space-time trade-offs: either with ancillas and fewer gates, or no ancillas and more gates. Since we are aiming to specify details for an implementation on a partially fault-tolerant quantum hardware which would likely not have an abundance of qubits, the ideal state preparation is one that uses no ancillas and its gate complexity is subdominant to the other parts of the Gibbs preparation. One such algorithm is to use uniformly controlled rotations \cite{mottonen2004transformation,shende2008cnot,plesch2011quantum} at cost $\mathcal{O}(2^r)$ gates. At $r \leq 12$, this is at most a few thousand gates, definitely below the cost of implementing controlled Hamiltonian evolutions. But for lower temperatures, or higher precisions for Gibbs sampling (smaller $\delta$) we would likely find ourselves in the regimes where the exponential scaling of the state preparation becomes dominant. In this case, we would suggest using the QSP-based method from \cite{mcardle2022quantum} that achieves $\mathcal{O}(r \cdot \log(1 / \varepsilon))$ gate complexity at the cost of 3-4 ancilla qubits by exploiting the (possibly piecewise) smoothness of $b_\pm$ via polynomial approximation. [REF \todo{ref to prelims}]

After preparing the time registers, we proceed to apply the controlled Hamiltonian evolutions. As written in \hyperref[alg:coh]{Algorithm~\ref*{alg:coh}}, all of these time evolutions are implemented by second order Trotterization [Ref to prelims \todo{ref}]: $S_2^{(M_-)}(\pm \bar{t}/\sigma)$ for the outer-, and $S_2^{(M_+)}(\pm \beta\bar{t})$ for the inner time sum (with $M_\pm$ many Trotter steps). For the desired quadrature errors, we need to retain contributions to the time sums up to $T_- = \Theta(\beta\sigma\log(1/\varepsilon))$ and $T_+ = \Theta(\sqrt{\log(1/\varepsilon)}/(\sigma\beta\sqrt{1+s}))$ (in the Metropolis case \eqref{eq:T_+_metro}). Hence, the maximum Hamiltonian simulation times are $\mathcal{O}(\beta\log(1/\varepsilon))$ and $\mathcal{O}(\sqrt{\log(1/\varepsilon)}/(\sigma\sqrt{1+s}))$ respectively. Note, that we decided for the Strang splitting opposed to the Lie-Trotter splitting because: (i) for these time scales its second order nature actually leads to less gate complexity, and (ii) it also avoids some additional errors to the detailed balance since it is palindromic $S_p(t)^\dagger = S_p(-t)$, see \hyperref[rem:palindromic]{Remark~\ref*{rem:palindromic}}.

The subnormalisation $\alpha = \|b_-\|_1\,\|b_+\|_1$ is a critical parameter: it enters both the QSP degree (see below) and the success probability of the block encoding. Its value depends on the choice of transition weight $\gamma$. More concretely, $b_-$ is universal for all cases and its norm is upper bounded, $\|b_-\|_1 \leq (\pi/\sqrt{2})\,e^{\beta^2\sigma^2/8}$ \eqref{eq:b_minus}. While this is ideal, and so is $b_+$ in the Gaussian case with $\|b_+\|_1 = (\beta/\omega_\gamma)\sqrt{\sigma_\gamma/(2\pi)}$ \eqref{eq:b_plus-gauss}, the same cannot be said about $b_+$ in the Metropolis cases \eqref{eq:b_plus-s-eta}. There, due to the singularity in the time-domain, we have an additional factor, and have
$$
\alpha = \|b_-\|_1\,\|b_+^{(s,\eta)}\|_1 = \mathcal{O}\!\left(\log\frac{\beta\|H\|\,\|\textstyle\sum_a A^{a\dagger}A^a\|}{\varepsilon}\right).
$$
How this effectively enters the total simulation time for the coherent evolution, we will see in the following QSP section.

\medskip
Given the block encoding $U_{B_a}$, with $\langle 0|U_{B_a}|0\rangle = B_a/\alpha$, we wish to implement the unitary $\ee^{-\ii\delta B_a}$. Since $U_{B_a}$ is a unitary that block-encodes a Hermitian operator $B_a/\alpha$ with $\|B_a/\alpha\| \leq 1$, its eigenvalues come in conjugate pairs $\ee^{\pm \ii\theta}$ where $\sin\theta = \lambda$ for each eigenvalue $\lambda$ of $B_a/\alpha$.  Thus, implementing $\ee^{-\ii\delta B_a}$ on the eigenspaces of $U_{B_a}$ reduces to finding a Laurent polynomial $P(\ee^{\ii\theta}) \approx \ee^{-\ii\delta\alpha\sin\theta}$ on the unit circle. We use GQSP of \cite{motlagh2024generalized} for this task. The basics were already covered in [REF to prelims \todo{ref}]. Here we only clarify the specific features to our application.

The degree of the Laurent polynomial required for $\varepsilon_\text{QSP}$-approximation of $\ee^{-\ii\delta\alpha\sin\theta}$ is given by \cite[Theorem 7.]{motlagh2024generalized} as
$$
d = \mathcal{O}\!\left(\delta\,\alpha + \frac{\log(1/\varepsilon_\mathrm{QSP})}{\log\log(1/\varepsilon_\mathrm{QSP})}\right),
$$
which in turn means this amount of applications of $U_{B_a}$. Indeed, $\delta$ in our cases is quite small, so much so, that even for lower temperatures $\beta = 20$ (for which $\alpha$ becomes larger) we can effectively use $d = 1$ and still have better precisions than what is set as a baseline by the weak measurement scheme. Of course, for larger systems moderate temperatures would manifest for larger $\beta$ than for smaller systems. Meaning, that though for the humble system sizes that we can classically simulate it seems that degree $d = 1$ is sufficient for the QSP, the same would not hold for larger systems. There, the irregularity of the Metropolis case function $b_+^{(s, \eta)}$ would eventually require higher QSP degrees and with it longer simulation times.

\setcounter{algorithm}{0}
\renewcommand{\thealgorithm}{\arabic{algorithm}b}
\begin{algorithm}[!htbp]
\caption{\textsc{DissipativeStep}: weak measurement simulation of
         $e^{\delta\mathcal{L}_{a,\mathrm{diss}}}$}
\label{alg:diss}
\begin{algorithmic}[1]
\small
\Require System register $S$, jump index $a$, step size $\delta$, filter function $f$, transition weight $\gamma$
\Ensure  $S$ evolved by
  $e^{\delta\mathcal{L}_{a,\mathrm{diss}}}(\rho) + \mathcal{O}(\delta^2)$
\Statex
\State \textbf{Ancillas:}
  frequency register $\Omega$ ($r$ qubits),
  Boltzmann qubit $q_\gamma$,
  weak measurement qubit $q_\delta$
  \Comment{all initialised to $|0\rangle$}
\Statex
\Statex \hfill \textit{--- Forward: Trotterized OFT~\eqref{eq:trotter-oft} ---}
\State \textsc{StatePrep}$(\Omega)$:
  $|0\rangle^{\otimes r} \;\longrightarrow\;
   |f\rangle := \sum_{\bar{t}\in S_{t_0}^{[N]}}
   \bar{f}(\bar{t})\,|\bar{t}\rangle$
   \Comment{Gaussian filter amplitudes~\eqref{eq:gaussian-filter-t}}
\State \textsc{Ctrl-Trott}$^{-}(\Omega, S)$:
  $|\bar{t}\rangle\,|\psi\rangle \;\longrightarrow\;
   |\bar{t}\rangle\; S_p^{(M)}(\bar{t})^\dagger \,|\psi\rangle$
   \Comment{$\approx e^{-\mathrm{i}H\bar{t}}$}
\State Apply $A_a$ on $S$
   \Comment{jump operator (e.g.\ single-site Pauli)}
\State \textsc{Ctrl-Trott}$^{+}(\Omega, S)$:
  $|\bar{t}\rangle\,|\psi'\rangle \;\longrightarrow\;
   |\bar{t}\rangle\; S_p^{(M)}(\bar{t})\,|\psi'\rangle$
   \Comment{$\approx e^{+\mathrm{i}H\bar{t}}$}
\State \textsc{QFT}$(\Omega)$
  \Comment{time $\to$ energy: register now encodes $|\bar{\omega}\rangle$}
\Statex \Comment{Joint state:
  $\displaystyle\sum_{\bar\omega}
   |\bar\omega\rangle \otimes
   \tilde{A}_a(\bar\omega)\,|\psi\rangle$}
\Statex
\Statex \textit{--- Boltzmann filter ---}
\State Controlled on $|\bar\omega\rangle$, rotate $q_\gamma$:\;
  $R_Y\!\Bigl(2\arcsin\!\sqrt{1-\gamma(\bar\omega)}\Bigr)$
  \Comment{see [REF to text]}
\Statex \Comment{$|0\rangle_{q_\gamma} \to
   \sqrt{\gamma(\bar\omega)}\,|0\rangle
   + \sqrt{1 - \gamma(\bar\omega)}\,|1\rangle$}
\Statex
\Statex \textit{--- Weak measurement ---}
\State Controlled on $q_\gamma = |0\rangle$, rotate $q_\delta$:\;
  $R_Y\!\bigl(2\arcsin\!\sqrt{\delta}\bigr)$
  \Comment{step-size acceptance}
\Statex \Comment{In the $q_\gamma\!=\!|0\rangle$ branch:
  $|0\rangle_{q_\delta} \to
   \sqrt{1-\delta}\,|0\rangle + \sqrt{\delta}\,|1\rangle$;
  \quad net jump probability $= \gamma(\bar\omega)\,\delta$}
\Statex
\Statex \textit{--- Reverse: uncompute $U$ (controlled on $q_\delta = |0\rangle$) ---}
\Statex \Comment{No-jump branch ($q_\delta\!=\!|0\rangle$): full uncomputation;
  jump branch ($q_\delta\!=\!|1\rangle$): $\Omega, q_\gamma, S$ remain entangled}
\State Undo Boltzmann rotation on $q_\gamma$
  \Comment{undo step~7}
\State \textsc{QFT}$^\dagger(\Omega)$
  \Comment{undo step~6}
\State \textsc{Ctrl-Trott}$^{-}(\Omega, S)$
  \Comment{undo step~5}
\State Apply $A_a^\dagger$ on $S$
  \Comment{undo step~4; for Pauli jumps $A_a^\dagger = A_a$}
\State \textsc{Ctrl-Trott}$^{+}(\Omega, S)$
  \Comment{undo step~3}
\State \textsc{StatePrep}$^\dagger(\Omega)$:
  $|f\rangle \to |0\rangle^{\otimes r}$
  \Comment{undo step~2}
\Statex
\Statex \textit{--- Discard ancillas ---}
\State Measure and reset $q_\delta$, $q_\gamma$, and $\Omega$
  \Comment{discard outcomes}
\Statex \Comment{Net channel on $S$:
  $\;\rho \;\mapsto\; e^{\delta\mathcal{L}_{a,\mathrm{diss}}}(\rho)
   + \mathcal{O}(\delta^2)$}
\end{algorithmic}
\end{algorithm}
\renewcommand{\thealgorithm}{\arabic{algorithm}}
\setcounter{algorithm}{1}
%% ---------------------------------------------------------------
\begin{figure}[ht]
  \centering
  \resizebox{1.05\textwidth}{!}{%
  \begin{quantikz}[row sep=0.9cm, column sep=0.42cm]
  % ==================================================================
    \lstick{$\ket{0}_\gamma$}
      & \qw
      & \qw
      & \qw
      & \qw
      & \qw
      & \qw
      & \gate{Y_{1-\gamma}}
      & \rstick{$\bra{0}_\gamma$}\qw
      \\
  %
  % ==================================================================
  % Row 2:  Omega  (time/frequency register, middle wire;
  %                  controls cells 4, 6 (downward to S) and cell 8 (upward to q_gamma),
  %                  prepared in cell 3, QFT'd in cell 7)
  % ==================================================================
    \lstick{$\ket{0}_\Omega$}
      & \qwbundle{r}
      & \gate{\text{PREP}}
      \gategroup[2,steps=5,style={dashed,rounded corners,fill=red!6,inner sep=5pt},
                 background,label style={label position=below,yshift=-0.40cm,anchor=north}]
        {\scriptsize $U_\mathrm{OFT}$}
      & \ctrl[style=halfctrl]{1}
      & \qw
      & \ctrl[style=halfctrl]{1}
      & \gate{\text{QFT}}
      & \ctrl[style=halfctrl]{-1}
      & \rstick{$\ket{\bar\omega}_\Omega$}\qw
      \\
  %
  % ==================================================================
  % Row 3:  S  (system register, bottom wire; target of the two controlled
  %             Trotters and of the bare jump A_a in between)
  % ==================================================================
    \lstick{$\ket{\psi}_S$}
      & \qwbundle{n}
      & \qw
      & \gate{S_2(-\bar t)}
      & \gate{A_a}
      & \gate{S_2(\bar t)}
      & \qw
      & \qw
      & \rstick{$\sqrt{\gamma(\bar{\omega})}\tilde A_a(\bar\omega)\ket{\psi}$}\qw
  \end{quantikz}%
  }
  \caption{Dissipative unitary $U_{\mathrm{diss}}$ (Alg.~\ref*{alg:diss}) on
  $q_\gamma\otimes\Omega\otimes S$ with PREP that prepares the time register as the Gaussian filter $\ket{\bar{f}(\bar{t})}_\Omega$, and with Boltzmann rotation
  $Y_{1-\gamma(\bar{\omega})}$.
  The red dashed block is the OFT,
  $\bra{\bar\omega}_\Omega U_{\mathrm{OFT}}\ket{0}_\Omega = \tilde A_a(\bar\omega)$.}
  \label{circ:u-diss}
\end{figure}

\medskip
\noindent\textsc{DissipativeStep:} \textbf{Weak measurement simulation.}
The \textsc{DissipativeStep}, described in \hyperref[alg:diss]{Algorithm~\ref*{alg:diss}}, implements the dissipative channel $\ee^{\delta\mathcal{L}_{a,\mathrm{diss}}}(\rho) + \mathcal{O}(\delta^2)$ for a single jump $A_a$. This part is inherited from \cite[Theorem III.1]{chen2023quantum} up to some small changes that we will highlight here. 

The forward pass (steps 2--6) that encodes the Trotterized OFT $\tilde{A}_a(\bar{\omega})$ \eqref{eq:trotter-oft} is denoted by the unitary $U_\mathrm{diss}$ that acts on the joint system $\Omega \otimes S$. Concretely, 
\begin{equation}
    U_{\mathrm{diss}} := \mathrm{QFT}(\Omega) \cdot \mathrm{Ctrl\text{-}Trott}^+(\Omega, S) \cdot A_a \cdot \mathrm{Ctrl\text{-}Trott}^-(\Omega, S) \cdot \mathrm{StatePrep}(\Omega),
\end{equation}
where $\mathrm{StatePrep}(\Omega)$ prepares the Gaussian filter state $\ket{f}$ on the time/frequency register, and the Heisenberg evolutions are implemented by Trotterized evolutions $S_2^{(M)}(\bar t)$. The net effect is then, 
\begin{equation}
    |0\rangle_\Omega |\psi\rangle_S \mapsto U_{\mathrm{diss}}\,|0\rangle_\Omega\,|\psi\rangle_S = \sum_{\bar\omega} |\bar\omega\rangle_\Omega \otimes \tilde{A}_a(\bar\omega)\,|\psi\rangle_S.
\end{equation}
The unitary $U_\mathrm{diss}$ is the dissipative analogue of the block encoding of the coherent part $U_{B_a}$, with a small difference: it is a unitary on the full $\Omega \otimes S$ space, and projecting onto a given $\ket{\bar{\omega}}$ selects the corresponding filtered jump component $\tilde{A}_a(\bar{\omega})$, i.e. $\langle \bar\omega|U_{\mathrm{diss}}|0\rangle = \tilde{A}_a(\bar\omega)$. 

\todo{this maybe at the end?}
Our construction is a simplified version of the one presented in \cite{chen2023quantum}. We split the generator jump-wise and implement dissipation of just $A_a$ in a given $\delta$ step, and use only single-site Pauli operators that are unitary ($A^2_a = \mathbf{1}$) and Hermitian ($A_a^\dagger = A_a$) jumps. Hence, we  do not need ancillas to block encode the sum of jumps, or to block encode non-unitary jumps either. A caveat here is that we would not apply the sum of jumps $\sum_a A_a$ at once, which has an effect on the mixing of the effective generator, but these errors are controlled to relevant orders \eqref{eq:gen-split-approx-error}, \eqref{eq:gen-split-imaginary-delta2-errors}. The other caveat, is that we certainly do not know what is a \emph{good} set of jump operators, ones that would mix the system quicker than other sets. However, after trying out different jump sets and seeing no real difference for our setup, we ended up sticking with the naive choice of single Pauli jumps. Then, the only ancillas we need are the ones for the time/frequency register $\Omega$, the Boltzmann qubit $q_\gamma$, and the weak measurement qubit $q_\delta$, a total of $r+2$ ancilla qubits.

Given the energy-resolved state after $U_\mathrm{diss}$, the remaining steps are the following:
\begin{enumerate}
    \item \emph{Boltzmann filter} (step 7): controlled on the energy register $\ket{\bar{\omega}}$, rotate the Boltzmann ancilla qubit $q_\gamma$ by $R_Y(2\arcsin\sqrt{1 - \gamma(\bar\omega)})$. This maps 
    $$|0\rangle_{\gamma} \to \sqrt{\gamma(\bar\omega)}\,|0\rangle_\gamma + \sqrt{1-\gamma(\bar\omega)}\,|1\rangle_\gamma,
    $$
    encoding the transition weight $\gamma$ into the amplited of the $\ket{0}_\gamma$ branch.
    \item \emph{Weak measurement} (step 8): controlled on $q_\gamma = \ket{0}_\gamma$ rotate the weak measurement ancilla qubit $q_\delta$ by $R_Y(2\arcsin\sqrt\delta)$. In the acceptance branch,
    $$
    |0\rangle_{\delta} \to \sqrt{1-\delta}\,|0\rangle_\delta + \sqrt\delta\,|1\rangle_\delta. 
    $$
    The net probability of the "jump" outcome ($q_\gamma = \ket{1}_\gamma$) is $\gamma(\bar{\omega}) \cdot \delta$, i.e. first-order in $\delta$ which is required by the weak measurement approximation $e^{\delta\mathcal{L}_{a,\mathrm{diss}}} = \mathrm{id} + \delta\mathcal{L}_{a,\mathrm{diss}} + \mathcal{O}(\delta^2)$. 
\end{enumerate}
The small jump probability is precisely what makes this a \emph{weak} measurement: the system state is only weakly disturbed per step, such that the accumulated effect of $L \cdot M_\mathcal{A}$ steps reproduces the dissipative semigroup to the desired accuracy.  The $\mathcal{O}(\delta^2)$ local error per step is the baseline that all other error sources (Trotterisation, quadrature, generator splitting) are tuned to match.

A comment has to be made regarding the Boltzmann rotation, as it could also lead to additional circuit costs. It requires a circuit that, controlled on the $r$-qubit energy register $\ket{\bar{\omega}}$ rotates $q_\gamma$ by the angle $\theta(\bar\omega) = 2\arcsin\sqrt{1 - \gamma(\bar\omega)}$. Crucially, this is a \emph{function-controlled} rotation: the classical function $\gamma$ must be compiled into the quantum circuit and its cost depends on the regularity of $\gamma$ and on the register size $r$.

For large $r$ (say $r \geq 15$) the standard approach is to use QSP/QSVT \cite{gilyen2019quantum} to implement a polynomial approximation of $\sqrt{\gamma(\bar{\omega})}$ as a function of the block-encoded normalised energy $\bar\omega / (\omega_0 N/2)$. The required polynomial degree $d_\gamma$ depends on which transition weight $\gamma$ we pick:
\begin{itemize}
    \item For the Gaussian $\gamma_\mathrm{G}$ (analytic) or the smooth Metropolis $\gamma_\mathrm{M}^{(s)}$ with $s>0$ (Gevrey-1/2, \hyperref[prop:smooth-metro-gevrey]{Propositon~\ref*{prop:smooth-metro-gevrey}}): $d_\gamma = \mathcal{O}(\mathrm{polylog}(1/\varepsilon_\gamma))$, since Gevrey-1/2 functions admit super-algebraically converging polynomial approximations.
    \item For the kinky Metropolis $\gamma_\mathrm{M}$ ($s = 0$, only Lipschitz at $\omega = -\beta\sigma^2 / 2$): $d_\gamma = \mathcal{O}(\mathrm{poly}(1/\varepsilon_\gamma))$, since Lipschitz functions require polynomially many terms. This can indeed be circumvented by splitting the domain piecewise (via a comparator circuit block on $\ket{\bar{\omega}}$) as \cite{chen2023quantum} notes, but that is still some additional circuit cost compared to the smooth Metropolis family $s>0$.
\end{itemize}

For moderate $r$ (around $r \leq 12)$, one can bypass QSP and just use uniformly controlled $R_Y$ rotations \cite{mottonen2004transformation,shende2008cnot}, the same we proposed for $\ket{b_\pm}$ in the \textsc{CoherentStep}. Although the number of CNOT gates in this case is exponentially growing $\mathcal{O}(2^r)$, at these moderate regimes, the cost is still subdominant to the other parts of the circuit. This is the approach that we use in practice. 

In either case the Boltzmann rotation is not a bottleneck: for smooth $\gamma$ and large $r$, QSP gives polylogarithmic overhead. For moderate $r$, the brute-force decomposition is cheap. 

\todo{maybe need better structuring with subsections...}
After the weak measurement, we apply the reverse pass $U_{\mathrm{diss}}^\dagger$ controlled on $q_\delta = |0\rangle$. Then after measuring and resetting all ancillas, it has been shown in \cite[Theorem III.1]{chen2023quantum} that the resulting state is 
$$
\rho \mapsto \rho + \delta\,\mathcal{L}_{a,\mathrm{diss}}(\rho) + \mathcal{O}(\delta^2),
$$
which is the first-order expansion of the dissipative semigroup channel we want to implement. 

Compared to the CoherentStep, the DissipativeStep is structurally simpler: it uses a single time/frequency register $\Omega$ (versus the nested $T_-, T_+$ of the coherent term), involves no QSP, and its Hamiltonian simulation time per step is $\mathcal{O}(\sigma\sqrt{\log(1/\varepsilon)})$ (set by the Gaussian filter truncation $T = \Theta(\sigma\sqrt{\log(1/\varepsilon)})$ and the controlled Trotter evolutions). 

\todo{add circuit figures}

\smallskip
\subsection{Parent Hamiltonian.} Beyond its role as a design principle for the algorithm, KMS detailed balance carries another analytical payoff that we have already touched upon: the similarity transformation leads to a \emph{Hermitian} quantum discriminant $\mathcal{D}(\rho_\beta, \mathcal{L})$ \eqref{eq:discriminant}. This transformation preserves the spectrum of the Lindbladian $\mathcal{L}$, which also keeps it negative semi-definite. Hence, to define a PSD parent Hamiltonian, we flip the signs, 
\begin{equation} \label{eq:parent-ham-defi}
    \mathcal{H}_\beta \;:=\; -\mathcal{D}(\rho_\beta, \mathcal{L}) \;\geq\; 0.
\end{equation}
Its ground state is the purified Gibbs state (since $\mathcal{H}_\beta$ acts on the double Hilbert space $\mathcal{H}\otimes\mathcal{H}$),
\begin{equation}
    |\sqrt{\rho_\beta}\rangle\rangle := \frac{1}{\sqrt{Z}}\sum_i e^{-\beta E_i/2}|\psi_i\rangle
   \otimes |\psi_i^*\rangle \in \mathcal{H}\otimes\mathcal{H},
\end{equation}
where $\{|\psi_i\rangle\}$ are the energy eigenstates and $|\psi_i^*\rangle$ denotes the entry-wise complex conjugate. As the spectrum is preserved, the spectral gap of the parent Hamiltonian equals that of the Lindbladian, $\lambda_\text{gap}(\mathcal{H}_\beta) = \lambda_\text{gap}(\mathcal{L})$. Consequently, the question of mixing time, reduces to bounding the spectral gap of a Hermitian PSD operator, based on the quantum spectral bound \eqref{eq:quantum-spectral-bound},
\begin{equation}\label{eq:parent-ham-gap-mixing}
    t_\text{mix}(\varepsilon) \;\leq\; \frac{1}{\lambda_\text{gap}(\mathcal{H}_\beta)} \ln \left( \frac{1}{\varepsilon\sqrt{\lambda_{\min}(\rho_\beta)}} \right).
\end{equation}
By reformulating the problem as a Hermitian spectral gap, we can tap into a well-developed research area of Hamiltonian spectral theory, e.g. perturbation theory for gapped self-adjoint operators \cite{kato1966perturbation}, detectability lemma \cite{aharonov2009detectability,anshu2016simple}, or specialized variational (min-max) principles \cite{bhatia2013matrix}.

But having a Hermitian parent Hamiltonian for detailed balanced generators is nothing new [REF]. But the parent Hamiltonians of CKG Lindbladians were also shown to have an additional structural property \cite[Proposition I.1]{chen2023efficient}: they are \emph{frustration-free}. Intuitively, this can be understood by the fact that the CKG Lindbladian decomposes as $\mathcal{L} = \sum_{a\in\mathcal{A}} \mathcal{L}_a$, where each per-jump term $\mathcal{L}_a$ individually satisfies the KMS DB condition $\mathcal{L}_a(\rho_\beta) = 0$. The discriminant transformation carries this over to the parent Hamiltonian as well,
\begin{equation}
    \mathcal{H}_\beta = \sum_{a\in\mathcal{A}} \mathcal{H}_\beta^a, \qquad \mathcal{H}_\beta^a \geq 0, \qquad \mathcal{H}_\beta^a\,|\sqrt{\rho_\beta}\rangle\!\rangle = 0 \quad\text{for each}\;\; a\in\mathcal{A}.
\end{equation}
In other words, every local term $\mathcal{H}_\beta^a$ is PSD and has (only) the purified Gibbs state in its kernel. For lattice Hamiltonians with local jumps $A_a$, the parent Hamiltonian even inherits the quasi-locality of the Lindbladian, making each $\mathcal{H}_\beta^a$ to act on a patch of radius $\tilde{\mathcal{O}}(\beta)$ of radius around the support of $A_a$. Frustration-freeness has the following three concrete analytical uses:
\begin{enumerate}
    \item \emph{Spectral gap stability.} In \cite{rouze2025efficient} Rouzé et al. proved that for 1D lattice Hamiltonians, the CKG parent Hamiltonian satisfies local topological quantum order at $\beta = 0$ where the gap is explicitly computable. Using the Michalakis-Zwolak stability theorem for frustration-free gapped Hamiltonians \cite{michalakis2013stability}, they showed that the gap persists at all sufficiently small inverse temperatures $\beta$, establishing the first unconditional rapid mixing results for a KMS DB Lindbladian with non-commuting $H$. This was recently strengthened by \cite{bergamaschi2024quantum}, who showed that 1D spin chains mix rapidly at \emph{all} finite temperatures ($\beta < \infty$), not just in the high-temperature regime. Without frustration-freeness, the stability theorem simply does not apply.
    \item \emph{Detectability lemma and gap amplification.} Since the ground state projector factorizes into local projectors (one per kernel of $\mathcal{H}_\beta^a$), the detectability lemma of \cite{anshu2016simple} provides a way to amplify the spectral gap via coarse-graining: grouping the local terms into $\mathcal{O}(1)$ layers of commuting projectors yields a product of projections whose spectral gap lower-bounds that of $\mathcal{H}_\beta$ up to a constant factor. This is also a key ingredient in the proof of \cite{rouze2025efficient}.

    In an even more recent paper \cite{fang2026quantum} the authors took this one step further, using the detectability lemma \cite{aharonov2009detectability,anshu2016simple}. They could circumvent the Lindbladian simulation and prepare the Gibbs state from the parent Hamiltonian. Their construction yields an $\mathcal{O}(|\mathcal{A}|)$-factor gate saving over simulation-based methods in general, and when the parent Hamiltonian is strictly local (e.g. for commuting $H$), a quadratic improvement in the spectral-gap dependence via QSVT. 

    \item \emph{Gap-to-MLSI upgrade and rapid mixing.} As we have discussed in the preliminaries [REF] the spectral bound \eqref{eq:parent-ham-gap-mixing} could in principle be upgraded to a tighter MLSI bound (albeit the MLSI constant has only been found to a few special cases so far). The frustration-free decomposition of $\mathcal{H}_\beta$ can help with this conversion, which is exactly what the authors used in \cite{kochanowski2025rapid}: they proved such a gap-to-MLSI equivalence for 1D Davies dynamics (GNS DB) of commuting Hamiltonians at all positive temperatures, strengthening earlier 1D rapid-mixing results. Indeed, the analogous statement in the non-commuting KMS DB setting remains open.
\end{enumerate}

Beyond proving mixing time bounds, the Hermitian and frustration-free structure of $\mathcal{H}_\beta$ also allows for some numerical techniques that could help with classical simulations as well. 

Without any special structure the Lindblad dynamics can already be simulated without having to retain the full eigenspace: in \cite{minganti2022arnoldi} the authors apply Arnoldi iteration directly to the non-Hermitian Liouvillian $\mathcal{L}$ of dimension $d^2$, maintaining a full upper-Hessenberg matrix and reorthogonalising against all $K$ Krylov basis vectors accumulated basis vectors at every step. This procedure has a cost of $\mathcal{O}(K d^2)$. Non-normality of $\mathcal{L}$ however further complicates the relationship between its spectral gap and the iteration count. 

\emph{Hermeticity}, obtained by the discriminant transformation $\mathcal{L} \mapsto \mathcal{H}_\beta$ reduces Arnoldi to the Lánczos iteration. The Hessenberg matrix becomes tridiagonal, and the iterations involve only the last two basis vectors, reducing the per-step cost to $\mathcal{O}(d^2)$. Moreover, the convergence of the Ritz values is now controlled directly by the spectral gap $\lambda_\text{gap}(\mathcal{H}_\beta)$ and the overlap of the starting vector with the kernel. 

\emph{Frustration-freeness} adds a structural layer on top due to the detectability lemma \cite{aharonov2009detectability,anshu2016simple}. Because the purified Gibbs state lies in the common kernel of the PSD local terms $\{\mathcal{H}_\beta^a\}$, the lemma turns this collection into an explicit operator (as the product of the corresponding local kernel projectors) whose repeated action drives any state towards $|\rho_\beta \rangle\rangle$ at a rate set by $\lambda_\text{gap}(\mathcal{H}_\beta)$. As we have already noted, on the quantum side this has been recently used by \cite{fang2026quantum}. On the classical side, the same lemma can be used in tensor network methods, where it is the analytical backbone of the polynomial-time ground state algorithms for 1D gapped local Hamiltonians \cite{landau2015polynomial, arad2017rigorous}.

\newpage


% --- DOCUMENT INFO ---
\title{Quantum Gibbs Sampling: Theory, Algorithms, and Numerical Analysis}
\author{Bence Temesi}
\date{\today}

% --- DOCUMENT BODY ---
\begin{document}

\maketitle
\clearpage
\thispagestyle{empty}

% Abstract ----
\clearpage
\chapter*{Abstract}
Your abstract text here...

% Borges epigraph ----
\clearpage
\vspace*{\fill}
\vspace{3cm} % push slightly below center

\hfill
\begin{minipage}{0.65\textwidth}
\raggedleft
\textit{%
``I imagined a labyrinth of labyrinths, amaze of mazes, a twisting, turning, ever-widening labyrinth that contained both past and future and somehow implied the stars.''}

\vspace{0.8em}
\normalfont
-- Jorge Luis Borges, \textit{The Garden of Forking Paths}
\end{minipage}

\vspace*{\fill}

% Acknowledgment ----
\clearpage
\chapter*{Acknowledgment}

\tableofcontents
% --- CHAPTERS ---
\part{Preliminaries}

\todo{this can be skipped pretty sure:}
As with any topic, we need to first understand the basic notions and tools that build the foundation. After the foundation is set, we can see what is possible to our current knowledge. Both the title and the introduction suggest that we would like to understand how a quantum system evolves to its thermal equilibrium state, the Gibbs state. For this, it is crucial to define the quantum dynamics that are at play here. But also to introduce the properties of the Gibbs state we would like to prepare. As we will see, these properties will set the difficulty level of the equilibration process. The underlying physics is however only one side of the problem, the other is the method with which we would like to prepare this target state. The tool we would \textit{eventually} like to use is a quantum computer, and accordingly we will need a quantum algorithm that could run on such a machinery. In classical computation Markov chain sampling was the quintessential technique to prepare and sample from the Gibbs distributions, and thus it comes as no surprise that we will work with quantum Markov chains in this thesis. Understanding how the quantum side works is definitely one aim of this chapter, but we also want to make the differences between the classical and quantum pictures as clear as possible. This should complement later analysis of the parameter regimes in which we can hope for a quantum advantage.

\chapter{Introduction}
\newpage
\chapter{Quantum Dynamics}

\label{sec:quantum-formalism}
\section{Quantum Formalism and Notation}
In order to have a self-contained thesis, we ought to start with the necessary concepts in quantum information. The aim here is not to give a complete introduction, but rather to provide all the definitions that are used later on in the text. In case the reader needs more details and pedagogic examples, we would advise them to open up already well-established textbooks and lecture notes like \cite{bible, wolf2012lecture, breuer2002theory}.

\smallskip
\label{sec:quantum-states}
\subsection{Quantum states} The state of a finite-dimensional quantum system is described within a complex vector space known as the Hilbert space, $\mathcal{H}$. While pure states are represented by vectors in $\mathcal{H}$, the most general description of a quantum state is given by a density operator, $\rho$.

These bounded operators belong to the space of all linear operators on $\mathcal{H}$, which we denote by $\mathcal{B}(\mathcal{H})$. For a physical state, an operator $\rho \in \mathcal{B}(\mathcal{H})$ must satisfy the following properties:
\begin{itemize}
    \item \textit{Hermiticity.} $\rho^\dagger = \rho$
    \item \textit{Positive semi-definiteness.} $\rho \geq 0$
    \item \textit{Normalization.} $\tr[\rho] = 1$
\end{itemize}
The set of all operators that fulfil these conditions forms the space of physical states, which we will denote by $\mathcal{D}(\mathcal{H})$. This space is a convex subset of $\mathcal{B}(\mathcal{H})$, i.e. any convex combination (or mixture) of density matrices is also a density matrix, $\sum_i \lambda_i \rho_i = \rho \in \mathcal{D}(\mathcal{H}).$

The states in $\mathcal{D}$ can be characterized by their purity, a measure of mixedness defined as $\tr [\rho^2]$. At one end there are the \textit{pure states}, for which the purity takes its maximal value of $1$. In this case the state is given by a rank-one projector onto a state vector $\ket{\psi} \in \mathcal{H}$
$$
\rho = \ket{\psi}\bra{\psi}
$$
States that are not pure, are called \textit{mixed states}. At this end we find the maximally mixed state, $\rho = 1 / \dim(\mathcal{H})$, which has the minimum possible purity of $\idm / \dim(\mathcal{H})$.
We can decompose any mixed state into a convex sum of orthogonal pure states as well:
\begin{equation}
    \rho = \sum_i \lambda_i \ket{\psi_i} \bra{\psi_i},
\end{equation}
where $\{\ket{\psi_i}\}$ is an orthonormal basis on $\mathcal{H}$ and are the eigenvectors of $\rho$. This is a powerful decomposition, as it tells us that mixed states can be represented by an ensemble of pure states with some probabilities $\lambda_i$. This will be quite useful later on for our numerical computations, as it is much easier to simulate the evolution of pure states than mixed states.

The task of preparing a specific state, such as the thermal Gibbs state, on a quantum computer involves navigating this state space $\mathcal{D}(\mathcal{H})$ through a sequence of physical operations. Once a state is prepared, we can then measure an observable quantity of interest, such as the system's energy.

\textbf{TODO: maybe correlation and entanglement and distances}

\smallskip
\subsection{Quantum measurements} In order to extract valuable information about the physical system, we still need to apply measurements to the system. In classical theory, predicting the outcome of a measurement should be possible, granted that we know the initial condition of the system and have the right theory. But due to the inherently probabilistic nature of quantum mechanics, such a prediction is just not possible. What is possible, is to predict probabilities and thus also the resulting measurement statistics that we expect for our statistical experiments. It is useful to first separate two notions, one is \textit{what} we can measure, and the other is \textit{how} we can measure it. A physical observable $A \in \mathcal{B}(\mathcal{H})$ that we can measure is described by a Hermitian matrix. The expectation value of the measurement can be then only understood with respect to a state $\rho$, and we can write it as $\langle A \rangle = \tr [\rho A]$. 
We can give this expression slightly more depth if we label the possible outcomes of the measurement by $\{a_\alpha\}$ and the probability with which we measure them by $\{p_\alpha\}$. Consequently, the expectation value is $\langle A \rangle = \sum_\alpha p_\alpha a_\alpha$. 

Now it is just a question of what do we really understand by the probabilities $p_\alpha$, which is also the same question of \textit{how} do we do measurements. In principle, we can always associate an outcome to a positive-valued operator $0 < E_\alpha \leq \idm$ such that,
$$p_\alpha = \tr [\rho E_\alpha].$$ 
If these operators (called \textit{effects}) have the normalization property $\sum_\alpha E_\alpha = \idm$, then they define a \textit{positive operator-valued measure} or POVM. This is still quite a general way of understanding quantum measurements. The way we connect back to the physical observable $A$ and its measurement statistics, is to find the spectral decomposition, $A = \sum_\alpha a_\alpha P_\alpha$, in which case $\{P_\alpha\}$ forms a special kind of POVM, a \textit{projection valued measure} (PVM). The naming is due to the fact that now each $E_\alpha = P_\alpha$ satisfy $P_\alpha^\dagger P_\alpha =  P_\alpha$, and thus an orthogonal projection to the eigenspace of $A$ with eigenvalue $a_\alpha$. While the POVM formalism is perfectly suited for describing measurement statistics, it does not care about the quantum state that we would find post measurement. However, in the PVM case, we definitely know in which state the system is after the measurement, as the operator projects to one of the eigenspaces of $A$:
\begin{equation}
    \label{eq:post_meas_state}
    \rho_\alpha = \frac{P_\alpha\rho P_\alpha}{\sqrt{\tr{\left[P_\alpha \rho\right]}}} = \frac{P_\alpha\rho P_\alpha}{\sqrt{p_\alpha}}.
\end{equation}
Due to the projection operator's property, if we were to repeat this measurement we would get the same $a_\alpha$ result and the same state $\rho_\alpha$. Famously, this is called as the collapse of the wave function, as the projection singles out a component of the superposition. Clearly not all measurements are repeatable in quantum mechanics, so the projective measurement must really be a special case. For instance when we measure a photon, we destroy it in the process and there is no way to repeat that measurement again. The post measurement state in the most general case can still be written up, and its form is similar to \eqref{eq:post_meas_state}, but without the projective property of the operator,
\begin{equation}
    \label{eq:post_meas_state_general}
    \rho_\alpha' = \frac{M_\alpha \rho M_\alpha^\dagger}{\sqrt{\tr\left[M_\alpha \rho M_\alpha^\dagger\right]}} = \frac{M_\alpha \rho M_\alpha^\dagger}{\sqrt{p_\alpha}}\quad\text{with}\quad M_\alpha M_\alpha^\dagger = E_\alpha.
\end{equation}
The difference here is that $\rho_\alpha'$ is in general not an eigenstate of $A$ and if we were to measure the same operator $M_\alpha$ we would get a different post measurement state.

\todo{THIS NEEDS REWRITING FOR THE NEW STRUCTURE}
\section{Closed systems}
With the introduction of quantum measurements in the previous chapter, we have already stumbled upon the notion of how a quantum system $\rho$ can change from one moment in time to another. This process can be described by an \textit{open quantum system evolution}, where we, the observers or (if we take away the human part) the environment, act on a quantum system. From the point of view of the observed system, this is an irreversible process, i.e. there is no physical map that would reverse the dynamics for all possible initial states. Irreversibility arises because we only observe a subsystem and lose the information of the rest. But the dynamics of the whole \textit{closed} system we can describe by a reversible process. 

Before we describe closed and open system dynamics in more detail, this is also a good place to introduce two frameworks in which one can understand evolutions. Since in quantum mechanics the only physically measurable quantities are the expectation values of some observables, $\langle A \rangle (t)$, and a physical process as we have seen is divided into a state preparation and a measurement, we have the freedom to choose which side we would like to evolve. In the \textit{Schrödinger picture} we evolve the system $\rho(t)$ and the operators $A$ are considered to be static, while in the \textit{Heisenberg picture} we keep the state $\rho$ static and evolve the operators $A(t)$. Naturally these pictures are equivalent to each other and give the same physical expectation values,
\begin{equation}
    \label{eq:schrödinger-heis}
    \langle A \rangle (t) = \tr \left[A \rho(t) \right] = \tr \left[ A(t) \rho \right].
\end{equation}

A reversible evolution of a closed system is described by a unitary map $u(\rho)$
\begin{equation}
    \label{eq:unitary-map}
    \rho \mapsto \rho(t) = \mathcal{U}(\rho) = U \rho U^\dagger,
\end{equation}
in the Schrödinger picture, or as the dual map $\mathcal{U}^\dagger(A)$ in the Heisenberg picture. These maps are naturally related via \eqref{eq:schrödinger-heis}, and by using the cyclic property of the trace we can see that $A(t) = \mathcal{U}^\dagger(A) = U^\dagger A U$. From Stone's theorem, we know that if these unitaries build a homogeneous group and have a continuous time parameter, i.e. multiple time evolutions can always be written as one: $U(t) U(t') = U(t + t'),\,\forall t, t'\in\mathbb{R}$, then there exists a Hermitian $H$, such that
$$
U(t) = \exp (-\ii H t).
$$
If we let a closed system evolve naturally, it will do so with respect to the Hermitian $H$ we call the Hamiltonian of the system. In principle, we can also evolve a system coherently via a unitary given by some other Hermitians, which we will see when we come to the quantum circuits. It is important to note that if we act with a unitary on a state, it will preserve its norm, which means that in a closed system evolution we can be certain that our evolving quantum state remains a density matrix. The same for open system evolutions cannot be said in general, however as we will see, we will always use map that will preserve all density matrix properties, for which reason we will call the map \textit{physical.}
\section{Open systems} \label{sec:OQS}
While closed system evolutions can describe a lot of scenarios, it requires for us to know the \textit{whole} system. In many cases this is impossible to do, either because the number of degrees of freedom is intractable, or because the interaction with the environment is not even known. 
For this exact reason, there has been a lot of research on how to describe the evolution of just a subsystem, while discarding the information of the environment. Naturally, this will not be a reversible process any more, because the discarded information is generally irrecoverable by any physical operation on the subsystem alone. 

A physical meaningful evolution in the Schrödinger picture is given by a linear map $\mathcal{E}: \mathcal{D}(\mathcal{H}) \rightarrow \mathcal{D}(\mathcal{H}')$ that preserves all the density matrix properties. 
This means that $\mathcal{E}$ has to be \textit{trace preserving} (TP), $\tr[\mathcal{E}(A)] = \tr[A]$ and also \textit{completely positive} (CP). Positivity for a map means, that it would map positive operators to positive operators, i.e. $\mathcal{E}(A^\dagger A) \geq 0$. Imagine that we only act with $T$ on just part of the system and rest we leave alone, $\mathcal{E}\otimes \id_n$. We will also require this map to be positive, for all $n \in \mathbb{N}$, which is exactly what we call complete positivity. If a linear map fulfils both of these conditions we call it a CPTP map or a \textit{quantum channel}. 

We can also define here the dual map $\mathcal{E}^\dagger : \mathcal{B}(\mathcal{H'})\rightarrow \mathcal{B}(\mathcal{H})$ that describes the same evolution in the Heisenberg picture, which we can always find from $T$ by $\tr[\mathcal{E}(\rho) A] = \tr[\rho \mathcal{E}^\dagger(A)]$. The only difference here is that the trace preserving property manifests itself as the \textit{unitality} of the map, $\mathcal{E}^\dagger(\idm) = \idm$.

While the unitary map, discussed for the closed system earlier, is a CPTP map, we can now introduce more general maps too. A quantum channel can be brought into different forms and representations that will be useful in different settings. 

\smallskip
\subsection{Kraus representation.} A linear map $\mathcal{E}: \mathcal{B}(\mathcal{H}) \rightarrow \mathcal{B}(\mathcal{H}')$ is completely positive iff there exists $r$ operators $K_j: \mathcal{H}\rightarrow\mathcal{H}'$ such that
\begin{equation}
    \label{eq:kraus}
    \mathcal{E}(A) = \sum_{j=1}^r K_j A K_j^\dagger,
\end{equation}
and $\mathcal{E}$ is trace preserving iff $\sum_j K_j^\dagger K_j = \idm$. The minimum number of Kraus operators needed to describe a map is called the Kraus rank. 
Now we can also better understand quantum measurements, as they are a very important example of irreversible evolution. A quantum channel can be understood as a general measurement process where the classical outcome is discarded. There is a clear similarity, between the post-measurement state we found in \eqref{eq:post_meas_state_general} and the Kraus representation: we find that the Kraus operators $\{K_j\}$ are precisely the measurement operators $\{M_\alpha\}$. From \eqref{eq:kraus} we also see, that the resulting state after applying all the measurements, or $\mathcal{E}(\rho)$, is the ensemble average over all possible post-measurement states.

This operator-sum representation is useful when our aim is to find out what state we get after applying a quantum channel to an initial state. However, sometimes we are interested in the map itself and want to analyse its static properties. The next representation could provide some help with that. 

\smallskip
\todo{rewrite this a bit more rigorously, mention isometry}
\subsection{Choi-Jamiolkowski representation.} A channel that acts on $d = \dim(\mathcal{H})$ dimensional operators, can be expressed in terms of a $d^2\times d^2$ matrix. But before we turn maps into matrices, let us first turn matrices into vectors. \textit{Vectorization} of a $d\times d$ matrix $A$ is the $d^2 \times 1$ column vector, denoted by $\mathrm{vec}(A)$, formed by stacking the columns of $A$ on top of one another:
$$
A = 
\begin{pmatrix}
    a &\! b \\
    c &\! d
\end{pmatrix}
\quad \Rightarrow \quad 
\mathrm{vec}(A) \equiv \ket{A}\rangle = 
\begin{pmatrix}
    a \\
    c \\
    b \\
    d
\end{pmatrix}.
$$
The way linear maps $\mathcal{E}$ are related to operators $\tau$ is slightly more complicated. Define a bipartite system $\mathrm{AB}$ with the joined Hilbert space $\mathcal{H}_A \otimes \mathcal{H}_B$. Then prepare this system in the maximally entangled state $\ket{\Omega}$
\begin{equation}
    \ket{\Omega} = \frac{1}{\sqrt{d}} \sum_{j=1}^d \ket{j}_A \otimes \ket{j}_B \equiv \frac{1}{\sqrt{d}} \sum_{j=1}^d \ket{jj}_{AB}.
\end{equation}
The Choi matrix $\tau$ is then found by applying the map $\mathcal{E}$ only on one part of the maximally entangled system $\tau = (\mathcal{E} \otimes \id_d)(\ket{\Omega}\bra{\Omega})$, and it is in one-to-one correspondence with the linear map $\mathcal{E}$ via
\begin{equation}
    \tr[A \mathcal{E}(B)] = d \tr[\tau A \otimes B^T].
\end{equation}
After some entry reshuffling (see \textbf{Appendix}) we can find a key vectorization identity:
\begin{equation}
    \mathrm{vec} (AXB) = (B^T \otimes A) \ket{X}\rangle.
\end{equation}
For instance if applied to the Kraus representation of the channel \eqref{eq:kraus}
we find that
\begin{equation}
    \label{eq:vectorized_kraus}
    \ket{\mathcal{E}(A)}\rangle = \left(\sum_j \bar K_j \otimes K_j\right)\ket{A}\rangle =: \hat{\mathcal{E}}\ket{A}\rangle.
\end{equation}

\smallskip
\subsection{Lindbladian generator.} So far we have described open quantum system evolutions with discrete maps that would transform some initial state to a final state over a fixed time. However, we are often interested in what hides in gaps between now and then, i.e. in continuous evolutions. Just as a single parameter, time, parametrized the unitary evolution of closed systems, it can also parametrize a family of quantum channels $(\mathcal{P}_t)_{t\geq 0}$: $\mathcal{D}(\mathcal{H}) \rightarrow \mathcal{D}(\mathcal{H})$. We call such a family a \textit{quantum dynamical semigroup} if it satisfies the following properties for all $t, s \in \mathbb{R}_+$
$$
\mathcal{P}_t\mathcal{P}_s = \mathcal{P}_{t + s} \quad\text{and}\quad \mathcal{P}_0 = \id.
$$
Mathematically, these are the properties that describe a continuous, homogeneous and \textit{Markovian} process. Physically, we can understand it as a memoryless process, i.e. that a future state depends only on the present state and not on any of the past states. In a moment we will see what physical assumptions are needed for the open quantum system to have its evolution described by a physical quantum dynamical semigroup. But first let us look into from the mathematical side.

An important result from the theory of one-parameter semigroups is that any such continuous family of maps is uniquely determined by a time-independent operator 
known as the generator, denoted by $\mathcal{L}$. The generator is a linear map on the space of operators, $\mathcal{L}: \mathcal{B}(\mathcal{H})\rightarrow \mathcal{B}(\mathcal{H})$, and is formally defined as the derivative of the dynamical map at $t=0$:, 
\begin{equation}
\mathcal{L} = \lim_{t \to 0^+} \frac{\mathcal{P}_t - \mathrm{id}}{t}.
\end{equation}
The generator thus describes the infinitesimal change that drives the dynamics. The naming also becomes clear since if we have $\mathcal{L}$, we can generate the semigroup $(\mathcal{P}_t)_{t\geq 0}$ through the exponentiation of the superoperator
\begin{equation}
    \mathcal{P}_t = \exp (t \mathcal{L}).
\end{equation}
For our purposes of course $\mathcal{L}$ cannot take up any arbitrary form. The evolution it generates $\mathcal{P}_t$ must be a valid physical process for all times $t \geq 0$, i.e. each $\mathcal{P}_t$ must be a CPTP map. In this case we call the semigroup a \textit{quantum Markov semigroup} (QMS) and we can describe the continuous evolution of a subsystem within an environment starting from some initial state $\rho(0)$ by
\begin{equation}
    \label{eq:oqs-evo-with-L}
    \rho(t) = \exp(t \mathcal{L})(\rho(0)).
\end{equation}
Or conversely,
\begin{equation}
    \frac{\dd}{\dd t}\rho(t) = \mathcal{L}(\rho(t)).
\end{equation}

What is the general form of $\mathcal{L}$ that generates physical evolutions? This was answered by Gorini, Kossakowski, Sudarshan, and Lindblad (GKSL) by a fundamental theorem they proved in \cite{gorini1976completely,lindblad1976generators}. The way to get to their result, starts by pointing out that any generator $\mathcal{L}$ of a QMS can be written in the form:
\begin{equation}
    \label{eq:qms_generator}
\mathcal{L}(\rho) = \Psi(\rho) + K\rho + \rho K^\dagger,
\end{equation}
where $\Psi: \mathcal{B}(\mathcal{H})\rightarrow \mathcal{B}(\mathcal{H})$ is a completely positive map and $K$ is an operator in $\mathcal{B}(\mathcal{H})$. The separation between the completely positive part and the remaining part allows us to reveal some structure, if we decompose the operator $K$ into its Hermitian and anti-Hermitian parts
\begin{equation}
    K = -V - \ii H,
\end{equation}
with the Hermitian operators: $H = \frac{\ii}{2}(K - K^\dagger)$ and $V = \frac{\ii}{2}(K + K^\dagger)$. Substituting this into the expression \eqref{eq:qms_generator} yields
\begin{equation}
\mathcal{L}(\rho) = -i[H, \rho] + \Psi(\rho) - \{V, \rho\}.
\end{equation}
Up until now the decomposition is completely general. But this is where the GKSL theorem comes in and finds a central constraint for $\mathcal{L}$ to be a generator of a CPTP semigroup: the map $\Psi$ must have a Kraus form $\Psi(\rho) = \sum_k L_k \rho L_k^\dagger$, and for the dynamics to be trace-preserving, the operator $V$ must be given by $\frac{1}{2}\sum_k L_k^\dagger L_k$. Putting these pieces together, we arrive at the celebrated GKSL form of the Lindbladian generator
\begin{equation}
    \label{eq:lindblad-generator}
    \mathcal{L}(\rho) = -\ii [H, \rho] + \sum_{k}\gamma_k \left(L_k^\dagger \rho L_k - \frac{1}{2}\left\{L_k^\dagger L_k, \rho\right\}\right).
\end{equation}
We can also clearly separate the different dynamics in this picture. The first term is the $\textit{coherent part}$, generated by the Hermitian $H$. Note, that $H$ does not necessarily have to be the Hamiltonian of the system, it could also be just some generator that coherently drives the system (foreshadowing). The second term is the \textit{dissipative part}, which describes the various ways of decoherence and dissipation via the Lindblad jump operators $\{L_k\}$ with a dissipation or jump rate given by $\gamma_k \geq 0$.

\smallskip
\subsection{Microscopic picture.} While it is nice to arrive at a mathematically rigorous result, it is also insightful to derive the same GKSL form of the generator from the physical principles and assumptions. For this we need to work with the microscopic model of a system $S$, interacting with a large environment or reservoir, $E$. The combined system is described by the following Hamiltonian
\begin{equation}
    H = H_S + H_E + H_I,
\end{equation}
where $H_S$ is the system Hamiltonian, $H_E$ the environment Hamiltonian, and $H_I$ defines the interaction between them. The total system is assumed to be closed, i.e. the evolution is purely unitary and governed by the von Neumann equation for the total density matrix $\rho$,
\begin{equation}
    \label{eq:von-neumann}
    \frac{\dd}{\dd t}\rho(t) = -\ii [H(t), \rho(t)].
\end{equation}
Note, that this is the same evolution as \eqref{eq:unitary-map} but written in the differential form. But of course, we are interested in the effective equation of motion of only the subsystem $\rho_S$. In order to derive it, we first have to translate the von Neumann equation to the interaction picture. This is another picture we can work in, and it is equivalent to the previously mentioned Schrödinger and Heisenberg pictures. The difference is that here we split the time evolutions between states and observables, and say that the states are evolved by the unitary generated by only the interaction Hamiltonian $H_I$ while the observables by $H_S + H_E$. In this picture the von Neumann equation in the Schrödinger picture \eqref{eq:von-neumann} will simply turn into
\begin{equation}
    \label{eq:von-neumann-int}
    \frac{\dd}{\dd t}\rho(t) = -\ii [H_I(t), \rho(t)],
\end{equation}
or in the integral form
\begin{equation}
    \rho(t) = \rho(0) - \ii \int_0^t \dd s [H_I(s), \rho(s)].
\end{equation}
Inserting this back into the von Neumann equation \eqref{eq:von-neumann-int} we find
\begin{equation}
    \label{eq:starting-integral-von-neumann}
    \frac{\dd}{\dd t}\rho_S(t) = -\int_0^t \dd s \tr_E \left[H_I(t), \left[H_I(s), \rho(s)\right]\right].
\end{equation}
Here, $\tr_E[\cdot]$ denotes the partial trace over the Hilbert space $\mathcal{H}_E$, which leads us to the system state, $\rho_S = \tr_E[\rho]$. Without loss of generality we assumed above that the interaction does not generate any dynamics in the environment from the initial state,
\begin{equation}
    \tr_E\left[H_I(t), \rho(0)\right] = 0.
\end{equation}
To proceed we need to use our first approximation, called the \textit{Born-Markov approximation}, which describes systems that interact weakly with their environment. Thus, the state of the environment is negligibly affected by the interaction and can be approximately separated from the system state, 
\begin{equation}
    \label{eq:sys-env-separable}
    \rho(t) \approx \rho_S(t) \otimes \rho_E.
\end{equation}
Though it might sound trivial, it is crucial to separate the system from the environment, because otherwise any statement about the system's dynamics alone is ill-defined. In \eqref{eq:sys-env-separable} we also assume that the environment is in a fixed stationary state $\rho_E$, which is essential if we want to derive a time-independent Lindblad generator \eqref{eq:lindblad-generator}. But it raises the question on which stationary state should we choose for $\rho_E$, and if this reflects at all what happens in Nature. In the next chapter we will see that choosing the environment to be the Gibbs state is necessary to derive thermalizing dynamics for the system $S$.
The Born-Markov approximation also assumes a timescale separation, in which the excitations in the environment decay much faster than what the system $\rho_S$ could ever resolve. This allows the replacement of $\rho(s)$ with $\rho(t)$, but also lets us replace $s$ with $t-s$ and send the upper limit to infinity (since the integrand decays quickly above the timescale of the environment). Applying all these steps to \eqref{eq:starting-integral-von-neumann} gets us
\begin{equation}
    \label{eq:born-markov-me}
    \frac{\dd}{\dd t}\rho_S(t) = -\int_0^\infty \dd s \tr_E \left[H_I(t), \left[H_I(t-s), \rho_S(t)\otimes\rho_E\right]\right].
\end{equation}
However, there is a problem here: these last steps indeed bring us closer to find a time-independent generator, it also fails to generate completely positive maps. Works from \cite{davies1974markovian,dumcke1979proper} illuminated why, and also provided the rigorous approximations we need to derive the Lindblad generator. Since the techniques derived in these works are necessary for us to understand the algorithms in this thesis, we give here a bit more thorough review. 

Let us write the interaction Hamiltonian in the Schrödinger picture as a tensor product of Hermitian system operators $A^a$ and Hermitian environment operators $B^a$
\begin{equation}
    H_I = \sum_a A^a \otimes B^a.
\end{equation}
The Bohr frequencies of the Hamiltonian $H_S$ are given by 
\begin{equation}
    B_{H_S} = \{\nu = \lambda_i - \lambda_j : \lambda_i,\lambda_j \in \mathrm{Spec}(H_S)\},
\end{equation}
where $\mathrm{Spec}(H_S)$ is the spectrum of $H_S$. Then we can write the Bohr decomposition of an operator $A$ as
\begin{equation}
    \label{eq:bohr-decomp}
    A = \sum_{\lambda_i, \lambda_j \in \mathrm{Spec}(H_S)} \Pi_i A \Pi_j = \sum_{\nu \in B_{H_S}} A_\nu,
\end{equation}
with
\begin{equation}
    A_\nu := \sum_{\lambda_i - \lambda_j = \nu} \Pi_i A \Pi_j.
\end{equation}
If we now do a unitary rotation, we find the operator's form in the interaction picture, 
\begin{equation}
    A(t) = e^{\ii H_S t} A e^{-\ii H_S t} = \sum_\nu e^{\ii \nu t} A_\nu \quad\text{or}\quad [H_S, A_\nu] = \nu A_\nu.
\end{equation}
Generally the Bohr spectrum $B_{H_S}$ is degenerate, meaning that there could be always different parts of an operator $A$ that would lead to the same energy difference $\nu$. The interaction Hamiltonian in the interaction picture follows
\begin{equation}
    H_I(t) = \sum_{a, \nu} e^{\ii \nu t} A^a_\nu \otimes B^a(t),
\end{equation}
with $B(t) = e^{\ii H_E t} B e^{-\ii H_E t}$. This finally lets us reinterpret the Born-Markov master equation \eqref{eq:born-markov-me} as
\begin{equation}
    \label{eq:bohr-decomposed-born-markov-me}
    \begin{split}
        \frac{\dd}{\dd t}\rho_S(t) =& \sum_{\nu, \nu'}\sum_{a, b}e^{\ii(\nu - \nu')t}\Gamma_{ab}(\nu)\left(A^b_\nu\, \rho_S(t)\, (A^{a}_{\nu'})^\dagger - (A^a_{\nu'})^\dagger A_{\nu'}^b\, \rho_S(t)\right) \\
        &+ \text{h.c.}
    \end{split}
\end{equation}
Here, h.c. stands for the Hermitian conjugate part. $\Gamma_{ab}(\nu)$ obtained by carrying out $\int \dd s$ in \eqref{eq:born-markov-me}, or as it turns out the half Fourier transforming the reservoir correlation functions:
\begin{equation}
    \label{eq:half-fourier-reservoir-corr-fn}
    \begin{split}
        \Gamma_{ab}(\nu) &:= \int_{0}^{\infty} \dd s\, e^{\ii\nu s} \tr_E\left[(B^a)^\dagger(t)B^b(t-s)\rho_E\right] \\
        &\equiv \int_{0}^{\infty} \dd s\,e^{\ii\nu s}\langle (B^a)^\dagger(t)B^b(t-s) \rangle \\
        &= \int_{0}^{\infty} \dd s\,e^{\ii\nu s}\langle (B^a)^\dagger(s)B^b(0) \rangle.
    \end{split}
\end{equation}
In the last equation we used the assumption that the bath is in a stationary state, $[H_E, \rho_E] = 0$, which leads to the correlation functions time translation invariant, with which we see that $\Gamma_{ab}(\nu)$, the function the determines the jump rates in the system-environment interactions is time-independent. 

For the Born-Markov form we also assumed a timescale separation between the system and the environment, which manifests here by assuming that the reservoir correlation functions decay rapidly $\lim_{s\rightarrow\infty}\langle (B^a)^\dagger(s)B^b(0) \rangle = 0$, before the system could ever be affected by it. But this also elucidate why a large or infinite environment is also required: if $\mathrm{Spec}(H)$ is not continuous then there will be quasi-periodic remnants left from the Fourier transform. From the physics point of view, in the case of a smaller, finite environment, the information that flowed from the system to the environment over the past interactions, would never die off, and after some time (Poincarré recurrence time) it would reappear and affect the system. Hence, ruining the Markovianity of the dynamics.

But there are still problems left after the Born-Markov approximation: $(i)$ there is still time-dependency left in the generator, $(ii)$ by assuming the coefficient matrix \eqref{eq:half-fourier-reservoir-corr-fn} to be constant over all times, it can in general break the positivity of the whole generator in some time regions as well. Thus, it cannot yet be a generator of a QMS. The solution to both of these problems is the \textit{secular approximation}. In \eqref{eq:bohr-decomposed-born-markov-me}, the secular terms are those that have $\nu = \nu'$, and they correspond to resonant processes where the total energy of the system and the environment is conserved. The non-secular terms, with $\nu \neq \nu'$ are off-resonant transitions to virtual states for which the total energy is temporarily not conserved. In the master equation they are also accompanied by the oscillating prefactor $e^{\ii(\nu - \nu')t}$. This oscillation is happening on the system's timescale, $\tau_S \sim \nu^{-1} \sim |\nu - \nu'|^{-1}$. Another timescale that is relevant here, is the relaxation time $\tau_R$ over which the system state changes through the interaction with the reservoir. It is inversely proportional to the coupling strength or jump rates, i.e. quite long in the weak-coupling limit we are working with. The secular approximation then says that we are only interested in cases where
\begin{equation}
    \label{eq:secular-approx}
    \tau_S \gg \tau_R \quad\text{or}\quad {|\nu - \nu'|} \gg 1.
\end{equation}
This effectively means that we can drop all the off-diagonal terms with $\nu \neq \nu'$ in \eqref{eq:bohr-decomposed-born-markov-me} and find
\begin{equation}
    \label{eq:after-secular-me}
    \frac{\dd}{\dd t}\rho_S(t) = \sum_{\nu}\sum_{a, b}\Gamma_{ab}(\nu)\left(A^b_\nu\, \rho_S(t)\, (A^{a}_{\nu})^\dagger - (A^a_{\nu})^\dagger A_{\nu}^b\, \rho_S(t)\right) + \text{h.c.},
\end{equation}
which is finally a QMS generator! To bring it into the Lindblad form we need to decompose the coefficient matrix into its Hermitian and anti-Hermitian parts:
\begin{equation}
    \Gamma_{ab}(\nu) = \frac{1}{2}\gamma_{ab}(\nu) + \ii S_{ab}(\nu)
\end{equation} 
with
\begin{equation}
    \gamma_{ab}(\nu) = \Gamma_{ab}(\nu) + \Gamma_{ab}^\dagger(\nu)
    = \int_{-\infty}^{\infty}\dd s\, e^{\ii\nu s}\langle (B^a(s))^\dagger B^b(0)\rangle
\end{equation}
and
\begin{equation}
       S_{ab}(\nu) = \frac{1}{2\ii}\left(\Gamma_{ab}(\nu) - \Gamma_{ab}^\dagger(\nu)\right).
\end{equation}
$S_{ab}(\nu)$ is a Hermitian and $\gamma_{ab}(\nu)$ is called the \textit{Kossakowski matrix} and it contains all jump rate coefficients; it is positive since now the Fourier transform over the reservoir correlations is a full Fourier (Bochner's theorem).
By inserting these into \eqref{eq:after-secular-me}, we have finally derived the (non-diagonal) Lindblad generator from microscopic principles and more-or-less physical assumptions:
\begin{equation}
    \label{eq:lindblad-microscopic}
    \mathcal{L}(\rho_S) = -\ii[H_{LS}, \rho_S] +  \sum_{ab}\sum_\nu \gamma_{ab}(\nu)\left(A^b_\nu\, \rho_S (A^a_\nu)^\dagger - \frac{1}{2}\left\{(A^a_\nu)^\dagger A^b_\nu, \rho_S\right\}\right).
\end{equation}
Comparing it to the general Lindblad generator \eqref{eq:lindblad-generator} we find that the Lindblad jumps are $L^a_\nu = \sqrt{\gamma_a(\nu)} A^a_\nu$, and that there is a contributing coherent term of the form:
\begin{equation}
    H_{LS} = \sum_{ab}\sum_\nu S_{ab}(\nu)(A_\nu^a)^\dagger A_\nu^b.
\end{equation}
This is the Lamb shift Hamiltonian, and it originates from the interactions with the environment which effectively shift the system's energy levels (since $[H_S, H_{LS}] = 0$). As the system $S$ evolves coherently with respect to its Hamiltonian $H_S$, the Lamb shift will add to it and lead to an evolution with $H_\text{eff} = H_S + H_{LS}$. 

There is one more comment to be made about \eqref{eq:lindblad-microscopic}: in its current form it is non-diagonal in the indices $a, b$. However, it is always possible to diagonalize the Kossakowski matrix $\gamma_{ab}(\nu)$, since it is Hermitian. Thus, we can always find some unitary $U$ such that $\gamma_a(\nu) = U \gamma_{ab}(\nu) U^\dagger$. We can accordingly find the rotated jump operators for which we can work in this diagonal picture by
\begin{equation}
    \label{eq:diagonalize-jumps-in-generator}
    L^a_\nu = \sum_b U_{ab} A^b_\nu.
\end{equation} 
Such a diagonalization `trick' can be quite useful for efficient numerical analysis as we will later see. 

Up to this point, we have introduced the main character of Markovian open quantum system dynamics: the Lindbladian generator. We have derived it from mathematical and physical principles, and have found that they generate CPTP maps, i.e. as the subsystem evolves within an environment, it preserves its density matrix properties. Of course, we are not only interested in how the system evolves but also where it is headed. Generally, these evolutions inch the system towards some stationary state(s). But a priori it is not clear if the system ever reaches that equilibrium (within a practical amount of time), since it can stumble upon a metastable state, where the evolution can get stuck. In order to understand this side of the dynamics, we will introduce the theory of Markov chains in the next chapter.
\chapter{Detailed Balanced Markov Dynamics}
A Markov chain is a stochastic process that describes the evolution of a probability distribution over a set of states $\mathcal{X}$. The dynamics are governed by a stochastic map that has the Markov property, which dictates that the probability of transitioning to any future state depends only upon the current state. Once we get to know the fundamentals of Markov chains, we can ask questions about to which state or states does the Markov chain evolve to, and how fast it reaches its goal. For more details on this topic, we refer to the indispensable source from Levin and Peres, \cite{levin2017markov}.

\section{Classical Markov Semigroups}
It is instructive to start the discussion with the classical theory of Markov chains. Even if the notions have to be generalized later to the quantum case, introducing them first classically, sheds a lot of light on the path ahead.

In classical theory, a probability distribution over the states $x\in\mathcal{X}$ at some time $t$ is given by a row vector $\mu_t$, where $\mu_t(x)$ gives the probability of finding the system in state $x$. The evolution of the probability distribution is described by a stochastic transition matrix $P$ as
$$
\mu_{t+1} = \mu_t P.
$$
The stochasticity of $P$, i.e. having non-negative entries and each of its row sums to $1$, ensures that $\mu_t$ stays a well-defined probability distribution at all times during the evolution.

In this work, we will be interested in Markov chains that for long enough times $t\rightarrow\infty$ let the distribution $\mu_t$ converge to an equilibrium distribution, known as the \textit{stationary distribution} $\pi$. This distribution is a fixed point of the dynamics, 
\begin{equation}
    \pi = \pi P.
\end{equation}
The existence of a unique stationary distribution to which an initial distribution converges, is guaranteed if the chain is \textit{irreducible} and \textit{aperiodic} due to the Perron-Frobenius theorem. Irreducible means that it is possible to reach any state from another state, while aperiodic means that the chain does not get stuck in a cycle of states. If both of these conditions hold, we will also call the chain \textit{ergodic}. As an example, if the system has some symmetry and the transitions prescribed by $P$ are all within that symmetric subspace, then the Markov chain would never be able to discover the states in the whole space, and would never reach its global stationary state $\pi$.

In Markov chain applications, the goal can be simply stated as: construct a transition matrix $P$ that prepares a desired probability distribution $\pi$ from some initial distribution(s). Naturally, we would first ask, why is there even a need for this, couldn't we directly sample from the target distribution $\pi$ if we know its form already?
In general the answer is no, and there are cases where it is indeed easier to walk through a Markov chain defined by $P$ that converges to $\pi$. Such a case is exactly what we have at our hands in this thesis, finding and sampling from the Gibbs (or Boltzmann) distribution
\begin{equation}
    \label{eq:cl-gibbs}
    \pi_\beta(x) = \frac{e^{-\beta E(x)}}{Z}, \quad  Z = \tr [e^{-\beta E(x)}]
\end{equation}
With $E(x)$ the energy of the state $x$, and $Z$ the partition function, which is a sum over all the possible energies in a system... which is a Herculean task to compute for larger system sizes.

Thankfully, there is a universal way to construct Markov chains with a given stationary distribution $\pi$. (Though universal, it does not always mean an efficient Markov chain.) The main concept for these processes, is \textit{detailed balance}. A chain is said to satisfy the detailed balance condition if, at equilibrium $\mu_t = \pi$, the probability flow from state $x$ to state $y$ is equal to the flow from $y$ to $x$:
\begin{equation}
    \label{eq:db-classical}
    \pi(x) P(x, y) = \pi(y) P(y, x), \quad\forall x, y\in\mathcal{X}.
\end{equation}
These chains are also called \textit{reversible}. This refers to a time-reversal symmetry of the dynamics, like a musical palindrome, for which we could not distinguish if it is played forwards or backwards. From a thermodynamic point of view these are processes that produce zero entropy at the equilibrium.

Important to note that while the detailed balance condition \eqref{eq:db-classical} is a sufficient condition for $\pi$ to be a stationary distribution, it is not strictly necessary. A weaker version of the condition, the \textit{global balance condition}, does not require for each pair of states to have reversible probability flow, but rather the net probability flow for each state $\pi(x)$ should remain zero under $P$:
\begin{equation}
    \sum_{x\in\mathcal{X}} \pi(x) P(x, y) = \sum_{x\in\mathcal{X}} \pi(y) P(y, x) = \pi(y),\quad\forall y\in\mathcal{X}.
\end{equation}
This latter condition however, is much less utilized in applications, since in terms of states it is not a local condition any more.

There is also another, operational, way of understanding detailed balance, which will help with the quantum detailed balance generalization later on. For this, we will work in the space of real-valued functions that act on the state space, $L^2(\pi)$, equipped with the $\pi$\textit{-weighted inner product}, $\langle\cdot,\cdot\rangle_\pi$. For two functions $f: \mathcal{X}\rightarrow \mathbb{R}$ and $g: \mathcal{X}\rightarrow \mathbb{R}$, it is defined as:
\begin{equation}
    \langle f, g\rangle_\pi := \sum_{x\in\mathcal{X}} f(x) g(x) \pi(x),
\end{equation}
where $\pi$ is still the stationary distribution of the chain. Since $\pi(x) > 0$ for all $x$ in an irreducible chain, the inner product is positive definite. Basically, it weighs each state's contribution by how likely that state is to occur at equilibrium. An inner product induces the natural norm in this space, $\|f\|_{2} = \sqrt{\langle f, f \rangle_\pi}$.

An operator (our transition matrix $P$) is then said to be self-adjoint with respect to the $\pi$-weighted inner product  if for any two functions $f$ and $g$, the following holds:
\begin{equation}
    \label{eq:P-selfadjoint-cl}
    \langle f, Pg\rangle_\pi = \langle Pf, g\rangle_\pi.
\end{equation}
Both sides written out gives,
\begin{equation}
    \sum_{x, y\in\mathcal{X}} f(x) g(y) \pi(x) P(x, y) = \sum_{x, y\in\mathcal{X}} f(x) g(y) \pi(y) P(y, x),
\end{equation}
which is exactly the detailed balance condition \eqref{eq:db-classical}, since it holds for all $f, g$. After showing the converse, we find a fundamental equivalence that, even if not in its exact form but in concept, survives in the quantum regime: a Markov chain satisfies detailed balance if and only if its transition matrix $P$ is a self-adjoint with respect to the $\pi$-weighted inner product.  
A consequence that comes from \eqref{eq:P-selfadjoint-cl} is that all the eigenvalues $\lambda_i(P)$ of $P$ are real, which allows us to order all eigenvalues on the real line, 
$$
1 = \lambda_1(P) > \lambda_2(P) \geq ... \geq -1,
$$
and have a clean definition for the spectral gap $\gamma(P) := 1 - \lambda_2(P)$. Intuitively, a large gap $\gamma(P)$ means that all terms in the distribution that are not the stationary state decay quickly. Conversely, a vanishing gap implies very slow convergence. We will soon see how to bound the mixing time of a chain with the spectral gap, but let us first explain how to construct a Markov chain that prepares a desired target distribution.

\label{parag:MH}
\smallskip
\subsection{Metropolis-Hastings algorithm.} The generic recipe for constructing a transition matrix $P$ that satisfies the detailed balance condition for any target distribution $\pi$. It was pioneered and was run on the MANIAC computer as early as 1953 \cite{metropolis1953equation}, and the algorithm has been generalized by Hastings in \cite{hastings1970monte}. The recipe is as follows:
\begin{enumerate}
    \item Decide upon an initial state $x\in\mathcal{X}$ where the Markov chain starts from. 
    \item Propose a new state $y\in\mathcal{X}$ based on a \textit{proposal matrix} $Q(x, y)$, which is some irreducible Markov chain on the state space $\mathcal{X}$. This can be as simple as randomly picking a state in the space. 
    \item Accept the state with probability given by the \textit{acceptance matrix} $A(x,y)$.
\end{enumerate}
Therefore, the transition matrix $P$ for the resulting Markov chain is a combination of a proposal $Q$ and acceptance $A$.
\begin{equation}
    P(x,y) =
\begin{cases}
    Q(x,y)A(x,y), \quad &\text{if}\; y\neq x, \\
    1 - \sum_{z\neq x} Q(x, z)A(x, z), \quad &\text{if}\;y = x.
\end{cases} 
\end{equation}
The key here is to define the acceptance probability $A(x, y)$ such that the final chain $P$ satisfied detailed balance condition \eqref{eq:db-classical} with respect to $\pi$, 
\begin{equation}
    \pi(x) Q(x, y) A(x, y) = \pi(y) Q(y, x) A(y, x),
\end{equation}
for which the choice that maximizes the acceptance rates $A$ (thus allowing for more efficient exploration in the state space) is given by the Metropolis-Hastings form: 
\begin{equation}
    \label{eq:MH-acceptance}
    A(x,y) = \min\left(1, \frac{\pi(y)Q(y, x)}{\pi(x)Q(x,y)}\right).
\end{equation}
It is straightforward to show that this choice of $A(x,y)$ forces the detailed balance condition to hold for $P$. The special case, when the proposal matrix is symmetric, i.e. $Q(x, y) = Q(y, x)$, is the Metropolis algorithm that Hastings later generalized to \eqref{eq:MH-acceptance}.

Finally, it becomes clear why this algorithm is so powerful. The acceptance $A$ and thus the full transition $P$ depends only on the \textit{ratio} of the target probabilities. For sampling from the Gibbs distribution $\pi(x) = e^{-\beta E(x)} / Z$, we would never have to compute the intractable partition function $Z$. The sole information the Markov chain needs in this case, is to the energy difference between $x\rightarrow y$: $\Delta E = E(y) - E(x)$. 

Note, that the Metropolis-Hastings algorithm makes no guarantee about the speed of convergence. It heavily depends on the state space, the proposal  matrix $Q$, and its interplay with the target distribution $\pi$. Quantifying this efficiency is the subject of mixing time analysis, which will be crucial for us for later findings as well. 

We have now come to the point where we can turn our attention to recent developments in quantum Gibbs sampling. The moment we mentioned Gibbs sampling in the classical setting, we also learnt the recipe that allowed researchers for decades to sample Gibbs states. So much so, that the \hyperref[parag:MH]{Metropolis-Hastings algorithm} is now an indispensable word of the sampling jargon. 

\smallskip
\subsection{Mixing time.} The go-to quantity that defines how fast the Markov chain converges to the target distribution. First, we need to define a distance between two probability distributions $\mu$ and $\nu$ using the \textit{total variation distance}:
\begin{equation}
    \label{eq:tv-distance}
    \|\mu - \nu\|_{TV} = \frac{1}{2} \sum_{x\in\mathcal{X}} |\mu(x) - \nu(x)|.
\end{equation} 
The distance from stationarity at time $t$ is then defined as the worst-case distance, maximizing over all possible starting states $x$:
\begin{equation}
    \label{eq:convergence-thm}
    d(t) := \sup_{x\in\mathcal{X}} \|P^t(x, \cdot) - \pi\|_{TV} 
\end{equation}
With respect to this distance, the mixing is the amount of time needed to reach a given $\varepsilon>0$ closeness to the stationary state $\pi$:
\begin{equation}
    \label{eq:mixing-time-cl}
    t_\text{mix}(\varepsilon) := \inf\left\{t \geq 0:\; d(t) \leq \varepsilon \right\}
\end{equation}
The mixing time clearly tells us how efficiently does a Markov chain prepare the target distribution $\pi$. Often we are interested in how the mixing time scales with the system size $n$ (or the state space size $|\mathcal{X}| = 2^n$). We distinguish between these cases
\begin{alignat*}{3}
    \text{(i)}\;   &\textit{Slow mixing}  &\quad& \text{with} &\quad& t_\text{mix}(\varepsilon) = \mathcal{O}(\exp(n) \log(1/\varepsilon)), \\
    \text{(ii)}\; &\textit{Rapid mixing}  &\quad& \text{with} &\quad& t_\text{mix}(\varepsilon) = \mathcal{O}(\mathrm{poly}(n) \log(1/\varepsilon)),\\
    \text{(ii)}\; &\textit{Optimal mixing}  &\quad& \text{with} &\quad& t_\text{mix}(\varepsilon) = \mathcal{O}(n \log(n) \log(1/\varepsilon)).
\end{alignat*}
Sometimes there could be also other parameters that can have a huge influence on the mixing times, like the temperature, which always have to be taken into account for Gibbs sampling. 
Broadly we can say, that high temperature systems unless there are geometric bottlenecks, or if the system is at a critical point of phase transition. In this latter case the system would have non-local correlations, that a local Markov chain like the Metropolis-Hastings would fail to resolve.

In order to bound the mixing times as the next step, we will make use of the following.
While the total variation distance is an $L^1$ distance, it can be rewritten in terms of the $L^1(\pi)$ distance $\|\cdot\|_{\pi, 1}$ induced by the $\pi$-weighted inner product, 
\begin{equation}
    \label{eq:tv-1norm}
    \begin{split}
        2\|P^t(x,\cdot) - \pi \|_{TV} &= \sum_{y\in\mathcal{X}} |P^t(x, y) - \pi(y)| \\
        &= \sum_{y} \pi(y) \left|\frac{P^t(x, y)}{\pi(y)} - \mathbf{1}\right| \\
        &= \left\|\frac{P^t(x, \cdot)}{\pi(\cdot)} - \mathbf{1}\right\|_{\pi, 1}.
    \end{split}
\end{equation}
After using Jensen's inequality, we can upper-bound the TV-distance from the stationary state $\pi$ via the $L^2(\pi)$ distance to $\pi$:
\begin{equation}
    \label{eq:tv-2norm}
    \begin{aligned}
          2\|P^t(x,\cdot) - \pi \|_{TV} &= \left\|\frac{P^t(x, \cdot)}{\pi(\cdot)} - \mathbf{1}\right\|_{\pi, 1} \leq \left\| \frac{P^t(x, \cdot)}{\pi(\cdot)} - \mathbf{1}\right\|_{\pi, 2} \\
          \quad \xRightarrow{\sup} d(t) &\leq \frac{1}{2} \sup_{x\in\mathcal{X}}\left\| \frac{P^t(x, \cdot)}{\pi(\cdot)} - \mathbf{1}\right\|_{\pi, 2} 
    \end{aligned}
\end{equation}

\smallskip
\subsection{Spectral bounds.} Bounding the mixing time is a central challenge in the analysis of Markov chains. While simulation can provide empirical estimates, rigorous analytical bounds are required for guaranteeing efficiency. One of the most fundamental approaches, relates the convergence rate to the spectral properties of the transition matrix $P$. While finding those spectral properties can be computationally difficult, this method at least provides a universal way to bound mixing times for all kinds problems.

Determining the mixing time is more subtle than simply identifying the decay rate of the second-largest eigenmode. While the eigenvalue spectrum dictates the asymptotic speed at which the system forgets its initial state, it does not account for the distance the system must traverse. To capture both effects, we must move from the natural metric of the total variation distance, to a metric compatible with spectral analysis.

As we have shown in \eqref{eq:tv-1norm} and \eqref{eq:tv-2norm}, the TV distance, is basically an $L^1(\pi)$ norm, and it can be upper bounded by the $L^2(\pi)$ norm. Working in the $L^2(\pi)$ space with the $\pi$-weighted inner norm is suitable, because the spectral theorem becomes applicable. This states, that for a self-adjoint operator, like our detailed balanced $P$, has an orthonormal basis $\{v_i\}$ and real eigenvalues $\lambda_i(P)$. (If the chain is taken to be `lazy', i.e. $P(x, x)\geq 1/2$, then $\lambda_i(P) \geq 0$.) 

In this function space, probabilistic notions turn into geometric objects. Probability distributions $\mu_t$ are represented by relative density vector $h_t = \mu_t / \pi$, while the stationary distribution $\pi$ corresponds to the constant vector $\mathbf{1}$ that stays constant under the transition: $P \mathbf{1} = \mathbf{1}$. Statistical expectations are now projections onto the axis given by $\pi$, $\mathbb{E}_\pi [f] = \langle f, \mathbf{1} \rangle_\pi$. Then the variance is the squared $L^2(\pi)$ norm of the orthogonal terms, $\text{Var}_\pi(f) = \|f - \mathbf{1}\|_{\pi, 2} ^2$. To see this, take the deviation from the equilibrium to be the vector $g_t = h_t - \mathbf{1}$. This is orthogonal to the stationary state because
\begin{equation}
    \mathbb{E}_\pi[g_t] \equiv \langle g_t, \mathbf{1} \rangle_\pi = \sum_x \left(\frac{\mu_t(x)}{\pi(x)}- \mathbf{1}\right) \pi(x) = \sum_x \mu_t(x) - \sum_x \pi(x) = 0.
\end{equation}
Thus, the deviation vector $g_t$ is spanned only by the rest of the eigenvector over its whole evolution,
\begin{equation}
    g_0 = \sum_{i=2}^{|\mathcal{X}|} c_i v_i \rightarrow P^t g_0 = g_t = \sum_{i=2}^{|\mathcal{X}|} c_i \lambda^t_i(P) v_i.
\end{equation}
The sum is dominated by the second-largest eigenvalue $\lambda_2$ and allows us to bound the decay of the distance (or variance) at time $t$ as,
\begin{equation}
    \begin{split}
        \sqrt{\text{Var}_\pi(h_t)} &\equiv \left\|g_t\right\|_{\pi, 2}  = \left\|\sum_{i=2}^{|\mathcal{X}|} c_i \lambda_i^t(P) v_i\right\|_{\pi, 2}\leq (1-\gamma(P))^t\sum_{i=2}^{|\mathcal{X}|} c_i v_i \\
        &\leq e^{-\gamma(P) t} \|g_0\|_{\pi, 2}.
    \end{split}
\end{equation}
Here, we used that $\lambda_2(P) = 1 - \gamma(P)$. This inequality shows already the structure we expected for the mixing process to have: the decay part is an exponential governed by the \textit{relaxation time}, $t_\text{rel} := 1 / \gamma(P)$, that multiplies a term quantifying the initial distance from stationarity. 

To derive a rigorous mixing time bound, we must still evaluate this initial distance for the worst-case starting configuration. If the system starts in a specific state $x_0$, the initial point mass will have a peak at $\mu_0(x_0) = 1$. Then the $L^2(\pi)$ norm of the initial deviation is:
\begin{equation}
    \begin{split}
            \left\|g_0\right\|_{\pi, 2} &= \left\|\frac{\mu_0}{\pi} - \mathbf{1}\right\|_{\pi, 2} = \sqrt{\sum_x \pi(x) \left(\frac{\mu_0}{\pi} - \mathbf{1}\right)^2} = \sqrt{\frac{1}{\pi(x_0)} -1} \\
            &\leq \frac{1}{\sqrt{\pi(x_0)}}.
    \end{split}
\end{equation}
Evidently, the worst-case means that the initial distribution $\mu_0$ concentrated all its probability mass on a state $x_0$ that has the least contribution to the stationary state: $\pi(x_0) = \pi_{\min} := \min_x \pi(x)$. Consequently, the mixing time is the required for the exponential decay to suppress this large initial factor, the \textit{burn-in} factor, down to a desired $\varepsilon$ size,
\begin{equation}
    \frac{1}{\sqrt{\pi_{\min}}} e^{-\gamma(P) t} \leq \varepsilon
\end{equation}
Rearranging, simplifying and not forgetting that mixing time was defined with the TV distance, this inequality yields the standard spectral bound on the mixing time:
\begin{equation}
\label{eq:spectral-bound}
t_\text{mix}(\varepsilon) \leq \frac{1}{\gamma(P)} \ln \left(\frac{1}{2\varepsilon \sqrt{\pi_{\min}}}\right).
\end{equation}

While the spectral bound \eqref{eq:spectral-bound} is universally applicable, it is often loose for high-dimensional systems where $\pi_{\min}$ is exponentially small. This motivates the need for sharper tools, such as the \textit{logarithmic Sobolev inequalities}, that were pioneered by Leonard Gross in \cite{gross1975logarithmic}.

\smallskip
\subsection{Logarithmic Sobolev bounds.} The weakness of the $1 / \sqrt{\pi_{\min}}$ scaling in the spectral bound stemmed from the $L^2(\pi)$ norm. Essentially, we have been tracking the convergence of the variance, $\mathrm{Var}_\pi(h) = \|h - \mathbf{1}\|_{\pi, 2}^2 = \|g\|_{\pi, 2}^2$ with $h = \mu / \pi$. But it turns out that looking at the convergence of entropy instead can lead to stronger results on the mixing time. 

The \textit{relative entropy} (or \textit{Kullback-Leibler divergence}) between the evolving distribution $\mu$ and the stationary distribution $\pi$ is defined as:
\begin{equation}
    \label{eq:relative-entropy}
    \text{Ent}_\pi(h) = \sum_{x\in\mathcal{X}} \pi(x) h(x) \ln h(x).
\end{equation}
Before we bound the relative entropy, it is constructive to explain why and when this measure is better than the variance. 

Consider the near-equilibrium regime, where the relative density takes the form $h(x) = \mathbf{1} + \varepsilon g(x)$ with $\varepsilon \ll 1$, thus $\mathbb{E}_\pi[g] = 0$. By performing a Taylor expansion of the function $f(t) = t \ln t $ from \eqref{eq:relative-entropy} around $t=1$, one finds that the entropy and variance agree to leading order:
\begin{equation} 
    \text{Ent}_\pi(\mathbf{1} + \varepsilon g(x)) = \frac{1}{2} \text{Var}_\pi(\mathbf{1} + \varepsilon g(x)) + \mathcal{O}(\varepsilon^3).
\end{equation}
In other words, when the initial distribution is close to $\pi$, the variance and the relative entropy are locally equivalent and both decay exponentially. 

Where they differ significantly is in their treatment of far-from-equilibrium states. Thus, consider now a chain starting in the worst-case initial distribution $h_0$, i.e. a point mass at the least likely state,
$$
h_0(x_0) = \frac{1}{\pi_{\min}}, \quad h_0(x\neq x_0) = 0.
$$
In this case, the variance in $L^2(\pi)$ is given by
\begin{equation}
    \mathrm{Var}_\pi(h_0) = \mathbb{E}_\pi[h_0^2] - 1 = \frac{1}{\pi_{\min}} - 1. 
\end{equation}
This scales exponentially with the system size $n$, since $\pi_{\min}\sim e^{-n}$. On the other hand, the relative entropy is now 
\begin{equation}
    \label{eq:initial-entropy}
    \mathrm{Ent}_\pi(h_0) = \pi_{\min} \left(\frac{1}{\pi_{\min}} \ln \frac{1}{\pi_{\min}}\right) = \ln \frac{1}{\pi_{\min}},
\end{equation}
which scales only linearly with the system size $n$. Of course, this is not the end of the story. To exploit this better scaling, we must prove that the entropy also contracts exponentially, which requires establishing a Log-Sobolev inequality.

Just as the spectral gap $\gamma(P)$ controls the decay of the variance, we define a new parameter, the constant $\alpha$, to control the decay of the entropy. It has been shown \cite{bobkov2006modified}, that the Markov chain satisfies a modified Log-Sobolev inequality (MLSI) with constant $\alpha > 0$ if for all positive functions $h$:
\begin{equation}
    \alpha \text{Ent}_\pi(h) \leq \mathcal{E}(h, \ln h),
\end{equation}
where $\mathcal{E}(f, g) = \langle f, (\id-P)g \rangle_\pi$ is the Dirichlet form associated with the chain. This specific form is chosen because the time derivative of the entropy is exactly the negative of the Dirichlet form on the right-hand side: $\frac{d}{dt}\text{Ent}_\pi(h_t) = -\mathcal{E}(h_t, \ln h_t)$ (in the continuous time limit). Consequently, this inequality guarantees that the relative entropy contracts exponentially at a rate $\alpha$:
\begin{equation}
    \text{Ent}_\pi(h_t) \leq e^{-\alpha t} \text{Ent}_\pi(h_0).
\end{equation}
To translate this entropy decay into a bound on the mixing time, we relate the entropy back to the TV distance using \textit{Pinsker's inequality}, which states that $2\|\mu - \pi\|_{TV}^2 \leq \text{Ent}_\pi(\mu/\pi)$.
Together with the initial entropy calculated above \eqref{eq:initial-entropy}:
$$
2 d(t)^2 \leq \text{Ent}_\pi(h_t) \leq e^{-\alpha t} \ln \frac{1}{\pi_{\min}} \leq 2\varepsilon^2.
$$
Solving for $t$ yields the Log-Sobolev bound on the mixing time:
\begin{equation}
    \label{eq:log-sobolev-bound}
    t_\text{mix}(\varepsilon) \leq \frac{1}{\alpha} \ln \left( \frac{\ln(1/\pi_{\min})}{2\varepsilon^2} \right).
\end{equation}
Compared with the spectral bound \eqref{eq:spectral-bound}, we see a significant improvement for high-dimensional systems. Recall that for a system of size $n$, the inverse stationary probability scales as $1/\pi_{\min} \sim e^{-n}$. The spectral bound depends linearly on this factor inside the logarithm, resulting in a mixing time scaling of $\mathcal{O}(n)$. In contrast, the Log-Sobolev bound takes the logarithm of this factor, resulting in a scaling of $\mathcal{O}(\ln n)$.

Intuitively, this improvement arises because the entropy is less sensitive to the spiky nature of the initial point-mass distribution than the $L^2$ norm is. The latter requires the dynamics to smooth out the sharp peak of the initial state before the bound becomes `useful'. The entropy, being logarithmic, penalizes this initial peak much less, thus having a tighter bound on mixing time. 

Unfortunately, the stronger bound comes with a significant analytical cost. While the spectral gap $\gamma(P)$ is well-defined for any reversible chain (by $1-\lambda_2$), establishing a strictly positive Log-Sobolev constant $\alpha$ is much harder. It generally requires proving geometric properties of the state space, e.g. Bakry-Émery criterion which relates to the Ricci curvature of the generator \cite{bakry2006diffusions}; or product structures (non-interacting systems) \cite{diaconis1996logarithmic}, that was extended to near-product structures (weakly-interacting systems, e.g. for high-temperatures) by \cite{stroock1992logarithmic}. But these results are neither universally present, nor easy to compute. Interestingly, we can relate the two decay rates as $2\alpha \leq \gamma(P)$ \cite{diaconis1996logarithmic}, which means the rate of entropy decay is actually slower than of the variance decay. But this factor is less consequential, than the difference in handling the far equilibrium starting points. All in all, while the Log-Sobolev inequalities are standard for proving rapid mixing in high-dimensional systems, the spectral gap remains the more universal and accessible tool for general contexts.

\smallskip
\subsection{Continuous-time Markov chains.} The natural generalizations of Markov chains, where the transitions occur at random times $t\in[0, \infty)$. The discrete nature of the chains we previously introduced, did also come with some artefacts, that will not be present for the continuous case, e.g. periodicity issues. 

In the continuous limit, the transition matrix $P$ is replaced by the classical Markov semigroup $(P_t)_{t\geq 0}$ generated by $L$. No surprise here, indeed, this is the classical analogue of the generator we introduced for open quantum system evolutions in Section \ref{sec:OQS}. Only that the generator $L$ will have slightly different but analogue properties to the quantum one. The evolution of the probability distribution $\mu(t)$ is governed by the classical Master Equation:
\begin{equation}
    \frac{\dd}{\dd t} \mu(t) = \mu(t)L
\end{equation}
The solution to this differential equation is formally given by the matrix exponential $\mu(t) = \mu(t=0) e^{t L}$. Just as the Lindbladian preserves the trace and positivity of the density matrix, the classical generator $L$ preserves the normalization and positivity of the probability vector $\mu(t)$ at all times. This requires the off diagonal entries of $L$ to be non-negative, and the row sums to be zero, $\sum_y L(x, y) = 0$.

The concept of detailed balance \eqref{eq:db-classical} gets rephrased in terms of the generator. Since detailed balance should hold at all times, then the derivatives of the probability flows must match at $t=0$ as well. The generator is defined exactly as that derivative, $L = \frac{\dd}{\dd t}P_t |_{t=0}$, and thus the detailed balance condition takes the form:
\begin{equation}
    \pi(x)L(x, y) = \pi(y) L(y, x), \quad \forall x\neq y.
\end{equation}
Note, that the diagonal is already fixed by $L(x, x) = - \sum_{y\neq x} L(x, y)$. This also implies, just as detailed balance implies global balance, that $\pi L = 0$. Equivalently to the above, in the operational framework, $L$ must be self-adjoint with respect to the $\pi$-weighted inner product:
\begin{equation}
    \langle f, Lg\rangle_\pi = \langle Lf, g\rangle_\pi.
\end{equation}
This self-adjointness of $L$ ensures that the semigroup $P_t$ is also self-adjoint, thereby preserving reversibility at all times $t\geq0$.

The convergence to the stationary state is still determined by the spectral gap, only that now the spectrum of a reversible $L$ is non-positive, and the unique stationary state $\pi$ has the eigenvalue $\lambda_1(L) = 0$, such that the spectral gap is simply $\gamma(L) = -\lambda_2(L)$. Of course, the relation between the eigenvalues of $P$ and $L$ at some time $t$ is $\lambda_i(P) = e^{t \lambda_i(L)}$. Nevertheless, there is no change about the contractive nature of the deviation from the stationary state, only that now we would replace $\gamma(P)$ with $\gamma(L)$, and have
$$\|h_t - \mathbf{1}\|_{\pi, 2} \leq e^{\lambda_2(L) t}\|h_0 - \mathbf{1}\|_{\pi, 2}.$$
Thus, the spectral bounds on mixing time \eqref{eq:mixing-time-cl}, also apply for the continuous case.

Similarly, the Log-Sobolev inequality adapts as expected to the generator setting. The Dirichlet form becomes $\mathcal{E}(f, g) = \langle f, -Lg \rangle_\pi$, and the modified Log-Sobolev inequality is now,
$$
    \alpha \text{Ent}_\pi(h) \leq \langle h, -L (\ln h) \rangle_\pi.
$$
This recovers the mixing time bound \eqref{eq:log-sobolev-bound} for the continuous case as well. 

With this formalism established, we can finally tackle the quantum setting in the next section. We will see that quantum Gibbs sampling is essentially the task of designing a quantum generator (the Lindbladian) that satisfies the \textit{quantum detailed balance} with respect to the quantum Gibbs state. For the convergence analysis of quantum Markov semigroups, we will also describe the required non-commutative generalizations of the spectral and Log-Sobolev results.

\section{Quantum Markov Semigroups}

It is now time to introduce the core subject of this work: the evolution of open quantum systems towards the thermal equilibrium. Just as we described classical mixing using generators and weighted inner products, we will formulate the quantum problem using reversible quantum Markov semigroups.

In the quantum setting, the probability distributions are given by density matrices $\rho(t)$, which in general does not commute with other density matrices. Gibbs sampling then means that we
construct a Lindbladian generator $\mathcal{L}$ of the GKSL form \eqref{eq:lindblad-generator}, such that the system state $\rho(t)$ converges to the stationary density matrix, i.e. the quantum Gibbs state,
\begin{equation}
    \rho_\beta = \frac{e^{-\beta H}}{Z}, \quad Z = \tr[e^{-\beta H}].
\end{equation}
What makes this quantum compared to the similar classical definition of  $\pi_\beta$ in \eqref{eq:cl-gibbs}, is the fact that now we have non-commuting terms in $H$ contributing to the energy $E(x)$. If $H$ has exclusively commuting terms inside (e.g. classical Ising model), then it is possible to write it as a diagonal matrix in some product basis, like the computational basis ($\sigma_Z$ basis). The Gibbs state $\rho_\beta$ would also be diagonal. Showing, that in very special cases quantum Gibbs sampling reduces to the classical problem.

The evolution is governed by the quantum Master equation $\frac{\dd}{\dd t} \rho = \mathcal{L}(\rho)$ which generates the semigroup $(\mathcal{P}_t)_{t\geq 0}$, just as we described for open quantum systems \eqref{eq:oqs-evo-with-L}. For the system to globally thermalize, $\rho_\beta$ must be the unique fixed point, $\mathcal{L}(\rho_\beta) = 0$. Note, that now the operators $\mathcal{L}$ (and $\mathcal{P}$) have been upgraded to superoperators $\mathcal{L}: \mathcal{B}(\mathcal{H})\rightarrow \mathcal{B}(\mathcal{H})$ since they act on operators (density matrices) instead of vector distributions.

However, generalizing the concept of detailed balance to the quantum case is not straightforward. Classical detailed balance was defined by the reversibility condition $\pi(x) P(x, y) = \pi(y) P(y, x)$. In the quantum case, the transition probabilities are replaced by operators, and the non-commutativity of the Hamiltonian and the states, $[\rho, H] \neq 0$ in general, means there is no unique way to `weight' the inner product by the stationary state.

\smallskip
\label{parag:qdb}
\subsection{Quantum detailed balance.} To define reversibility for a quantum process, we must adopt the operational framework we introduced classically: a generator is detailed balanced if it is self-adjoint with respect to a specific $\sigma$-weighted inner product. The non-weighted inner product on the space of bounded operators $\mathcal{B(H)}$ is the \textit{Hilbert-Schmidt inner product}:
\begin{equation}
    \label{eq:HS}
    \langle A, B \rangle := \tr[A^\dagger B].
\end{equation}
In fact, the Hilbert space we get with this inner product, can be thought of as the infinite temperature limit, where the stationary state is the maximally mixed state, $\lim_{\beta\rightarrow 0}\rho_\beta = \id_\mathcal{H} / \dim \mathcal{H}$. Since we juggle several inner products in this text, let us fix the following notation: the adjoint of the superoperator $\Phi\;:\;\mathcal{B(H)}\rightarrow\mathcal{B(H)}$ is always with respect to the Hilbert-Schmidt inner product $\langle \cdot,\cdot\rangle$, and it is denoted by $\Phi^\dagger$, e.g. 
\begin{equation}
    \langle A, \Phi(B) \rangle = \langle \Phi^\dagger(A), B \rangle.
\end{equation}
Written out explicitly for a general case we get,
\begin{equation}
    \Phi(\cdot) = \sum_j \alpha_j A_j (\cdot) B_j \quad\Rightarrow\quad \Phi^\dagger(\cdot) = \sum_j \bar{\alpha}_j A^\dagger_j (\cdot) B^\dagger_j
\end{equation}

The moment we have to deal with finite temperatures, we have to weight \eqref{eq:HS}. Due to non-commutativity there is no unique way to distribute the multiplication by $\sigma$ within the trace. Thus, we find a family of weighted inner-products:
\begin{equation}
    \label{eq:weighted-HS}
    \langle A, B \rangle_{\sigma, s} := \tr[\sigma^s A^\dagger \sigma^{1-s} B], \quad s\in[0,1].
\end{equation}
While all of these choices reduce to the classical product when matrices commute, they correspond to fundamentally different physical processes in the quantum setting.
The expression above \eqref{eq:weighted-HS} defines a continuous family of inner products parametrized by $s$, but they effectively boil down to just two distinct cases. As it turns out \cite{fagnola2007generators}, the inner products in \eqref{eq:weighted-HS} lead to coinciding QMS for all $s\in[0,1] \backslash \{\frac{1}{2}\}$. Consequently, all of these inner products are equivalent to the \textit{Gelfand-Naimark-Segal (GNS) inner product} defined by $s\in\{0, 1\}$ in \cite{alicki1976detailed}, i.e. to the following,
\begin{equation}
    \label{eq:GNS-product}
    \langle A, B\rangle_{\sigma, 1} = \tr[\sigma A^\dagger B].
\end{equation}
For $s=\frac{1}{2}$ on the other hand, we get a distinct, symmetrically weighted inner product, which was introduced in \cite{goldstein1995kms} and was named after \textit{Kubo-Martin-Schwinger (KMS)}:
\begin{equation}
\langle A, B \rangle_{\sigma, \frac{1}{2}} := \tr[\sigma^{1/2} A^\dagger \sigma^{1/2} B].
\end{equation}
A generator $\mathcal{L}$ satisfies quantum detailed balance if it is self-adjoint with respect to either of these weighted inner products, and it is called GNS or KMS detailed balanced for $s = 0, 1$ or $s=\frac{1}{2}$ respectively.
\begin{equation}
    \langle A, \mathcal{L}(B) \rangle_{\sigma, s} = \langle \mathcal{L}(A), B \rangle_{\sigma, s} \quad\forall A, B \in \mathcal{B(H)}
\end{equation}
Equivalently, if detailed balance holds we can express the adjoint of the generator as
\begin{equation}
    \label{eq:general-qdb-adjoint}
    \mathcal{L^\dagger}(\cdot) = \sigma^{s-1} \mathcal{L}(\sigma^{1-s}\, \cdot\, \sigma^{s}) \sigma^{-s}.
\end{equation}
Let us check if this is the case, then the stationary state $\sigma$ is indeed stationary. Since $\mathcal{L}$ generates unital maps in the Heisenberg picture, i.e. $\mathcal{L}^\dagger(\idm) = 0$, we can write 
\begin{equation}
    \label{eq:sigma-is-fix}
    \begin{split}
        0 &= \langle A, \mathcal{L}^\dagger(\idm) \rangle_{\sigma, s} \\
        &= \langle A, \sigma^{s-1} \mathcal{L}(\sigma) \sigma^{-s} \rangle_{\sigma, s} \\
        &= \tr[A^\dagger \mathcal{L}(\sigma)] \\
        &= \langle A, \mathcal{L}(\sigma)\rangle, \qquad\forall A\in\mathcal{B(H)}.
    \end{split}
\end{equation}
Where we used \eqref{eq:weighted-HS} and \eqref{eq:general-qdb-adjoint}, and assumed that $\sigma$ is a full-rank density matrix (e.g. the Gibbs state at finite temperatures) and are invertible. 
This confirms, that if the generator and its adjoint has the relation $\eqref{eq:general-qdb-adjoint}$ then $\sigma$ is the stationary state of the generated QMS, $\mathcal{P}_t(\sigma) = \sigma$.

An alternate but equivalent definition of quantum detailed balance that we will use later is given by its \textit{discriminant} or parent Hamiltonian, 
\begin{equation}
    \label{eq:discriminant}
    \begin{split}
        \mathcal{D}(\sigma, \mathcal{L}) &:= \sigma^{-\frac{1}{4}}\mathcal{L}(\sigma^{\frac{1}{4}} \cdot \sigma^{\frac{1}{4}})\sigma^{-\frac{1}{4}} \\
        &= \sigma^{\frac{1}{4}}\mathcal{L}^\dagger(\sigma^{-\frac{1}{4}} \cdot \sigma^{-\frac{1}{4}})\sigma^{\frac{1}{4}} \equiv \mathcal{D}^\dagger(\sigma, \mathcal{L}).
    \end{split}
\end{equation}
For the last equation we used the definition of $\mathcal{L}^\dagger$ from \eqref{eq:general-qdb-adjoint}. Thus, conjugating a Lindbladian with its unique stationary state in this manner, can be thought of as a similarity transformation that makes the Lindbladian Hermitian. 
Then breaking detailed balance can be understood as having a non-Hermitian discriminant. Or equivalently, since every operator can be decomposed into a Hermitian and anti-Hermitian part, that we have a non-vanishing anti-Hermitian part.
\begin{equation}
    \begin{split}
        \mathcal{D}(\rho, \mathcal{L}) &= \mathcal{H}(\rho, \mathcal{L}) + \mathcal{A}(\rho, \mathcal{L}) \\
        \mathcal{D}(\rho, \mathcal{L})^\dagger &= \mathcal{H}(\rho, \mathcal{L}) - \mathcal{A}(\rho, \mathcal{L}).
    \end{split}
\end{equation}
\begin{definition} \label{def:approx-db}
    \textup{(Approximate detailed balance condition)}. The Lindbladian $\mathcal{L}$ with full-ranked fixed point $\rho$ satisfies the $\varepsilon-$approximate detailed balance condition if the anti-Hermitian part $\mathcal{A}$ is small
    $$
    \frac{1}{2}\left\|\mathcal{D}(\rho, \mathcal{L}) - \mathcal{D}(\rho, \mathcal{L})^\dagger \right\|_{2\rightarrow2} = \left\|\mathcal{A}(\rho, \mathcal{L})\right\|_{2\rightarrow2} \leq \varepsilon
    $$
\end{definition}
This is one of the main quantities we will track in our numerical analyses, because it gives a clear picture on how much we break detailed balance via implementation errors.

An advantage of the transformation \eqref{eq:discriminant}, is that it keeps the spectrum the same, which makes spectral analysis possible for Lindbladians through their parent Hamiltonians. This is often useful if the jump operators are local or quasi-local. In these cases we often find an already known gapped quantum lattice model, or at least have more tools to use to bound that spectral gap. They used these techniques for GNS detailed balanced Lindbladians with commuting system Hamiltonians in \cite{kastoryano2016quantum}, and for the recent KMS Lindbladians where the parent Hamiltonians are provably frustration-free in \cite{rouze2025efficient,bergamaschi2024quantum,vsmid2025polynomial}.

At the moment it might seem that the difference between GNS and KMS detailed balances is purely mathematical, and it is merely a choice of the parameter $s$. But for the thermalizing process there are clear physical consequences of choosing one or the other. To understand the distinction, we must introduce the \textit{modular operator} $\Delta_\sigma\;:\;\mathcal{B(H)} \rightarrow \mathcal{B(H)}$ defined as:
\begin{equation}
    \Delta_\sigma(A) = \sigma A \sigma^{-1}
\end{equation}
In the classical limit where all operators commute, $\Delta_\sigma$ reduces to the identity. Now it encodes the non-commutativity structure of the algebra with respect to the steady state. Since $\rho_\beta \propto e^{-\beta H}$, the action of the modular operator is intimately linked to the (imaginary) time evolution by the system Hamiltonian $H$. It was shown in \cite{alicki1976detailed} that a generator $\mathcal{L}$ satisfies the GNS detailed balance if and only if it satisfies KMS detailed balance, \textit{and} it commutes with the modular operator:
\begin{equation}
    [\mathcal{L}, \Delta_\sigma] = 0.
\end{equation}
This additional condition implies that the dissipative dynamics generates a semigroup that commutes with the free unitary evolution of the system. In other words, the dissipation does not generate coherences between energy levels and thus respects the energy conservation of the system. Notice, that this is actually in line with the standard weak-coupling limit in which we explicitly needed the secular approximation \eqref{eq:secular-approx} (dropping non-secular or coherent terms) to derive a physical generator. The KMS detailed balance however leads to more permissive yet still physical dynamics that couple populations (diagonal terms) to coherences (off-diagonal terms). The fact that it retains non-secular terms but also physical, since we started from a physical Lindbladian and restricted it to have KMS detailed balance, means that it cannot be derived in the standard weak-coupling limit with a heat bath. On what kind of first principles would lead to KMS detailed balanced thermalization, we can refer to works where they use the \textit{singular coupling limit} \cite{alicki2007quantum}, or the \textit{thermal quantum collision / repeated interaction limit} \cite{attal2007langevin}.

\smallskip
\label{subsec:davies}
\subsection{Davies generator.} Canonically, the first thermalizing Lindbladian has been written up in the standard weak-coupling limit (with Born-Markov and secular approximations) by Davies in \cite{davies1974markovian}. This has been achieved by coupling the system to an infinite free heat bath, and the resulting GNS detailed balanced generator was such a milestone that it got named after Davies.

We have already derived most of what we need for it in \eqref{eq:lindblad-microscopic} where we derived the Lindblad generator from first principles. There the main idea was to decompose the interaction operator into eigenoperators of the system Hamiltonian, $A_\nu$, such that $[H, A_\nu] = -\nu A_\nu$. These operators characterize the transitions between the energy levels with an energy separation of $\nu$. In the Davies generator the transition rates given by $\gamma_a(\nu)$ satisfy a specific condition, called the \textit{KMS condition} (not to be confused with KMS symmetry), which relates the absorption and emission via the Boltzmann factor for any $a$ indexing the set of jump operators $\{A^a\}$:
\begin{equation}
    \label{eq:KMS-condition}
    \frac{\gamma_a(-\nu)}{\gamma_a(\nu)} = e^{\beta \nu}
\end{equation}
To remind ourselves, the Kossakowski matrix $\gamma(\nu)$ is the full Fourier transform of the reservoir correlation functions \eqref{eq:half-fourier-reservoir-corr-fn}. If we assume the environment is in the Gibbs state $\rho_E = e^{-\beta H_E} / Z_E$, then the correlation functions of the reservoir indeed satisfy the KMS condition \eqref{eq:KMS-condition}. This is a necessary condition for any thermalizing dynamics, but not sufficient. In this text we will discuss two distinct ways to thermalize a system dissipatively. If we work in the standard weak-coupling limit and drop all non-secular terms, we end up with the GNS detailed balanced Lindbladian from Davies
\begin{equation}
    \label{eq:davies-generator}
    \mathcal{L}_\text{Davies}(\rho) = -\ii[H_{LS}, \rho] +  \sum_{a}\sum_\nu \gamma_{a}(\nu)\left(A^a_\nu\, \rho (A^a_\nu)^\dagger - \frac{1}{2}\left\{(A^a_\nu)^\dagger A^a_\nu, \rho\right\}\right).
\end{equation}
We can see that due to the secular approximation, all terms with $\nu \neq \nu'$ have been dropped, while keeping all terms with $\nu=\nu'$. A unitary rotation via \eqref{eq:diagonalize-jumps-in-generator} was also used such that after diagonalization we have $a=b$ compared to \eqref{eq:lindblad-microscopic}. Note, that the only other difference to the Lindbladian derived in the microscopic picture with the Bohr decomposition, is that the transition rates $\gamma_{a}(\nu)$ now fulfil \eqref{eq:KMS-condition}. 

In a few steps we can even see how the Davies generator satisfies the GNS detailed balance. From $[H, A_\nu] = \nu A_\nu$, the following relation will be quite useful
\begin{equation}
    \label{eq:sigma-A-commutation}
    \rho_\beta A_\nu = e^{-\beta \nu} A_\nu \rho_\beta.
\end{equation}
This helps to show that since the Lamb-Shift Hamiltonian $H_{LS}$ commutes with the Gibbs state, and so does the anti-commutator term $\mathcal{R} \propto \{A_\nu^\dagger A_\nu, \rho\}$ in the Lindbladian, they are both self-adjoint with respect to the GNS inner product \eqref{eq:GNS-product}. Note, that this is only so simple because we have solely secular terms in the Davies generator. As we will later see, the insight for the KMS detailed balanced Lindbladian was to keep and make the non-secular terms in the coherent and reflective part $\mathcal{R}$ comply with each other. 

The remaining transition part needs a bit more work to prove that it is self-adjoint. But in both the GNS and KMS cases, the main requirement is for the transitions to fulfill KMS condition \eqref{eq:KMS-condition}, and to have  
the Bohr frequencies of the system have $B_H = - B_H$. The following derivation should apply for all jumps $a$ of $\{\sqrt{\gamma^a(\nu)A_\nu^a}\}$.
\begin{equation}
    \label{eq:davies-gns-db}
    \begin{alignedat}{2}
        \langle X, \mathcal{T}(Y) \rangle_{\rho_\beta, 1} &= \sum_\nu \gamma(\nu) \tr\left[\rho_\beta X^\dagger A_\nu^\dagger Y A_\nu\right] \\
        &= \sum_\nu \gamma(\nu) e^{\beta \nu} \tr\left[\rho_\beta (A_\nu X A_\nu^\dagger)^\dagger Y\right]  &\quad& (\text{cyclic trace; \eqref{eq:sigma-A-commutation}})\\
        &= \sum_\nu \gamma(-\nu) \tr\left[\rho_\beta ((A_{-\nu})^\dagger X A_{-\nu})^\dagger Y\right] &\quad& (\text{KMS condition})\\
        &= \sum_\nu \gamma(\nu) \tr\left[\rho_\beta (A_{\nu}^\dagger X A_{\nu})^\dagger Y\right] &\quad& (-\nu \rightarrow \nu)\\
        &= \langle \mathcal{T}(X), Y \rangle_{\rho_\beta, 1}
    \end{alignedat}
\end{equation}
Hence, all parts of the Davies generator are self-adjoint with respect to the GNS inner product, which implies that the Gibbs state $\rho_\beta$ is the fixed point, $\mathcal{L}_\text{Davies}(\rho_\beta) = 0.$

It is natural to question the physical ubiquity of the Davies semigroup. The secular approximation, which enforces the GNS detailed balance for the generator by discarding non-secular terms, implies that the environment seems to discriminate perfectly between the system's energy levels. But it is better to think of it as the ideal case where the coupling is so weak, and the environment is so large that the relaxation time $\tau_R\sim 1 / \lambda_2(\mathcal{L})$ is orders of magnitude longer than the system's intrinsic time $\tau_S \sim 1/ \nu_{\min}$. In this asymptotically exact description, all the non-secular terms average to zero, hence their absence in the Davies generator. The other obstacle that arises, is the fact that for a simulation \textit{we} would have to be the ones (instead of Nature) that discriminate, or resolve the energy levels of the system. An implementation that constantly has to distinguish between Bohr frequencies cannot be efficiently done for larger system sizes. With growing system size the gap between the Bohr frequencies gets exponentially small. Accordingly, the amount if simulation time needed, due to the Heisenberg's uncertainty principle $\Delta E \Delta t \geq \hbar / 2$, would exponentially increase as well, leading to an intractable computation even if it was done in a quantum computer. \todo{Refer to QPE section.}

Before we move on to introduce and work with the KMS detailed balanced Lindbladian, let us first go over the notion of mixing time and the techniques to bound it in the quantum setting. 

\smallskip
\subsection{Mixing time.} Similar to the classical setting, we need to first quantify the distance from the equilibrium before and then establish rigorous bounds on the convergence speed. The natural generalization of the total variation distance \eqref{eq:tv-distance} is the \textit{trace distance}. For any two quantum states $\rho$ and $\sigma$, it is defined via the non-commutative $L^1$ norm called the Schatten 1-norm, or trace norm $\|\cdot\|_1$
\begin{equation}
    \|\rho - \sigma\|_1 = \tr\left[ \sqrt{(\rho-\sigma)^\dagger (\rho-\sigma)} \right].
\end{equation}
Operationally, it represents the maximum probability of distinguishing two states with an optimal measurement. The mixing time is defined analogously to the classical case \eqref{eq:mixing-time-cl}. We look for the time $t$ required for the evolved state $\rho(t) = \mathcal{P}_t(\rho_0)$ to be $\varepsilon$ close to the stationary Gibbs state $\rho_\beta$ for any initial state $\rho_0$,
\begin{equation} \label{eq:mixing-time-defi}
    t_\text{mix}(\varepsilon) := \inf \left\{ t \ge 0 : \sup_{\rho_0} \|\mathcal{P}_t(\rho_0) - \rho_\beta\|_1 \le \varepsilon \right\}.
\end{equation}
Intuitively, the mixing of a quantum Markov semigroup is a more complex process than its classical counterpart. While a classical chain merely has to redistribute probability mass among the states (population or diagonal elements), a quantum process must simultaneously accomplish two tasks: it must drive the populations to the Gibbs distribution \textit{and} it must destroy all off-diagonal quantum coherences between energy levels. We will soon see what role this new decoherence timescale will play.
On top of that, the non-commutative nature of the space also creates new ways to thermalize. In the classical case, if an energy barrier stopped the convergence, a large enough kick could get the evolution moving. In the quantum case, such an energy barrier could be tunnelled through (if it is not too wide). But some transitions could also be possibly blocked if the phases do not align with the target subspace. Thus, a priori it is not always clear, that quantum mixing is faster than classical. To find out more about it, a recent paper \cite{rakovszky2024bottlenecks} discusses bottlenecks in both regimes. 

\todo{Possibly mention classical energy transport for Davies.}

\smallskip
\subsection{Quantum spectral bounds.} To bound the mixing time, we once again turn to the spectral properties of the generator. But just as before, we have to change to the (non-commutative) $L^2$ space. There, we can work with an orthonormal basis due to the spectral theorem that still applies.

We consider the Hilbert space of bounded operators $\mathcal{B(H)}$ equipped with the weighted inner product $\langle \cdot, \cdot \rangle_{\sigma, s}$ for which the generator $\mathcal{L}$ satisfies detailed balance. This means that even in the quantum setting, the generator is self-adjoint to an inner-product, and accordingly has real eigenvalues. Together with the fact that $\mathcal{L}$ generates CPTP maps, we know that all the eigenvalues lie on the non-positive real line, $0 = \lambda_1(\mathcal{L}) > \lambda_2(\mathcal{L}) \geq \dots$, and the spectral gap is clearly defined as:
\begin{equation}
    \lambda(\mathcal{L}) := - \lambda_2(\mathcal{L}).
\end{equation}
The arguments follow similarly to the classical case. The generator's only fixed point is $\rho_\beta$, as shown in $\eqref{eq:sigma-is-fix}$, and the dynamics in the subspace orthogonal to this kernel are governed by the remaining eigenvalues and eigenstates. The variance of any observable (equivalently, the distance of any state from the equilibrium) contracts exponentially at a rate given by $\lambda(\mathcal{L})$.

The trace distance used for the mixing time definition, can be bounded by the weighted $L^2(\sigma)$ norm via the non-commutative Cauchy-Schwarz inequality. Without loss of generality, we can use the KMS inner product with $s=\frac{1}{2}$, and write $\|\cdot\|_{2, \text{KMS}} := \|\cdot\|_{2, \rho_\beta, 1/2}$ for the following derivations. The evolved state $\rho_t = \mathcal{P}_t(\rho_0)$ the distance from the stationary state $\rho_\beta$ takes the form:
\begin{equation}
    \label{eq:trdist-ub-kms-norm}
    \begin{split}
        \|\rho_t - \rho_\beta\|_1  &= \|\rho_\beta^{-\frac{1}{2}}(X_t - \idm)\rho_\beta^{-\frac{1}{2}} \|_1 \\
        &= \|X_t - \idm\|_{1, \text{KMS}} \\
        &\leq \|X_t - \idm\|_{2, \text{KMS}} \\
        &\leq e^{-\lambda(\mathcal{L}) t} \|X_0 - \idm\|_{2, \text{KMS}}.
    \end{split}
\end{equation}
Here, $X_t := \rho_\beta^{-\frac{1}{2}} \rho_t \rho_\beta^{-\frac{1}{2}}$ can be interpreted as the relative density of $\rho_t$ to the Gibbs state with respect to the KMS inner product. (Admittedly, there is no unique "density" any more, as for the GNS inner product we would get a different form.) The evolution of $X_t$ is governed by the KMS self-adjoint operator: $$X_t = \rho_\beta^{-\frac{1}{2}}\mathcal{P}_t\left(\rho_\beta^{\frac{1}{2}} X_0 \rho_\beta^{\frac{1}{2}}\right) \rho_\beta^{-\frac{1}{2}}.$$
For a rigorous upper-bound we have to take the worst-case scenario for the remaining norm in \eqref{eq:trdist-ub-kms-norm}, i.e. when the initial state $\rho_0$ has a small overlap with the Gibbs state $\rho_\beta.$ An example would be a high energy initial state, that tries to thermalize to a low temperature Gibbs state, i.e. where $\beta \gg 1$. After squaring and writing out the remaining norm, it simplifies to
\begin{equation}
    \begin{split}
        \|X_0 - \idm\|_{2, \text{KMS}}^2 &= \| X_0 \|_{2, \text{KMS}} - \idm \\
        &\leq \tr\left[ \rho_0 \rho_\beta^{-1} \rho_0\right] \\
        &= \bra{\psi} \rho_\beta^{-1} \ket{\psi}.
    \end{split}
\end{equation}
Here, we first used that $X_0 = \rho_\beta^{-\frac{1}{2}} \rho_0 \rho_\beta^{-\frac{1}{2}}$, then chose a pure state for $\rho_0 = \ket{\psi}\bra{\psi}$ to have an initial state whose probability mass concentrates on a state that is the least compatible with $\rho_\beta$. This allows us to bound the norm, just as in the classical case, with the smallest eigenvalue of $\rho_\beta$:
$$
\sup_{\psi} \langle \psi | \rho_\beta^{-1} | \psi \rangle = \|\rho_\beta^{-1}\|_\infty = \frac{1}{\lambda_{\min}(\rho_\beta)}.
$$
Substituting back to \eqref{eq:trdist-ub-kms-norm} we get the upper bound on the trace distance from the fixed point,
$$
\|\rho_t - \rho_\beta\|_1 \leq e^{-\lambda(\mathcal{L})t} \frac{1}{\sqrt{\lambda_{\min}(\rho_\beta)}}.
$$
Finally, solving it for time for which the distance falls below some $\varepsilon$ gives the spectral bound on the mixing time,
\begin{equation}
    \label{eq:quantum-spectral-bound}
        t_\text{mix}(\varepsilon) \leq \frac{1}{\lambda(\mathcal{L})} \ln \left( \frac{1}{\varepsilon \sqrt{\lambda_{\min}(\rho_\beta)}} \right).
\end{equation}
A neat result, since it coincides with the classical bound \eqref{eq:spectral-bound}, only that now the relaxation time due to the spectral gap of the quantum generator $\mathcal{L}$ and the burn-in factor due to the smallest eigenvalue of the quantum state $\rho_\beta$ have quite different behaviour to their classical analogues. Note, that the result is the same even if we were to use the GNS inner product in the derivation. 

Unfortunately, the problem of handling far-equilibrium initial states for this bound is also inherited from the classical result. But worry not, there has been a vast amount of research done to use the logarithmic Sobolev techniques also in the non-commuting quantum setting. \todo{CITE from rouzé}

\smallskip
\label{sec:prelim-qlsi}
\subsection{Quantum logarithmic Sobolev bounds.} The problem with the spectral bound \eqref{eq:quantum-spectral-bound} is that the factor $1 / \sqrt{\lambda_{\min}(\rho_\beta)}$ typically scales as $e^{\beta n / 2}$ for an $n$-particle system, making the bound loose for large systems. To get a tighter bound even in these cases, we once again turn to the (quantum) relative entropy of the system, and try to find that logarithmic factor that helped us out for the classical case as well, \eqref{eq:log-sobolev-bound}.

The quantum relative entropy between the states $\rho, \sigma \in \mathcal{D(H)}$ is defined as:
\todo{write it up with Ent from rouzé, its not $\text{Ent}(\rho)$}
\begin{equation} 
    \text{Ent}_{\sigma}(\rho) := D(\rho \| \sigma) = \tr[ \rho (\ln \rho - \ln \sigma)]
\end{equation}
In analogy to the classical case, we seek a constant $\alpha > 0$ such that the entropy relative to the Gibbs state $\rho_\beta$ decays exponentially, i.e.
$$
D(\rho_t \| \rho_\beta) \leq e^{-\alpha t} D(\rho_0 \| \rho_\beta).
$$
The differential condition of this decay is the \textit{Quantum Logarithmic Sobolev Inequality} (QLSI), which is satisfied for a quantum Markov semigroup generated by $\mathcal{L}$
\begin{equation}
    \alpha D(\rho \| \rho_\beta) \leq \mathcal{E}
\end{equation}
where $\mathcal{E}_L(\rho) := -\frac{d}{dt} D(\rho_t \| \sigma)|_{t=0} = -\tr[\mathcal{L}(\rho)(\ln \rho - \ln \sigma)]$ is the quantum entropy production (or Dirichlet form for the log-relative entropy).

$$\mathcal{E_L}(X, Y) := -\langle X, \mathcal{L}(Y)\rangle_{\sigma}$$

\chapter{Quantum Circuits}
\smallskip
\section{Basics}
The smallest non-trivial system in $\mathcal{D(H)}$ is the \emph{qubit}, the quantum analogue of a classical bit. Its Hilbert space is $\mathcal{H} = \mathbb{C}^2$ and a pure state takes the form
$$
\ket{\psi} = \alpha\ket{0} + \beta\ket{1}\quad\text{with}\quad |\alpha|^2 + |\beta|^2 = 1.
$$
An $n$-qubit register lives in the tensor product $\mathcal{H}^{\otimes n} = (\mathbb{C}^2)^{\otimes n}$ of dimensions $2^n$, with computational basis $\{\ket{x}\}_{x\in\{0,1\}^n}$ where $\ket{x} = \ket{x_1}\otimes\cdots\otimes\ket{x_n}$. Mixed states $\rho$ on this register are the density operators on $\mathcal{H}^{\otimes n}$, in the same sense as introduced \hyperref[sec:quantum-states]{previously}.

A \emph{quantum gates} acting on $n$ qubits is a unitary $U \in \mathcal{B}((\mathbb{C}^2)^{\otimes n})$. Any unitary on $n$ qubits can be approximated to arbitrary accuracy by composing gates from the universal set consisting of all single-qubit unitaries together with the two-qubit \textsc{CNOT} gate \cite{bible}. The single-qubit Pauli matrices that will time to time show up in the thesis are defined as
\begin{equation}
    \label{eq:paulis}
    X = \begin{pmatrix} 0 & 1 \\ 1 & 0 \end{pmatrix}, \quad
    Y = \begin{pmatrix} 0 & -\ii \\ \ii & 0 \end{pmatrix}, \quad
    Z = \begin{pmatrix} 1 & 0 \\ 0 & -1 \end{pmatrix}.
\end{equation}
They are simultaneously Hermitian and unitary, and satisfy $X^2 = Y^2 = Z^2 = \idm$. For $\sigma_i, \sigma_j \in \{X, Y, Z\}$ they fulfil the following (anti-)commutation relations
$$
[\sigma_i, \sigma_j] = 2\ii\,\epsilon_{ijk}\,\sigma_k, \qquad \{\sigma_i, \sigma_j\} = 2\,\delta_{ij}\,\idm
$$
with $\epsilon_{ijk}$ the totally antisymmetric Levi-Civita symbol. Together with $\idm$ they span $\mathcal{B}(\mathbb{C}^2)$, and tensor products of Paulis form an orthogonal basis of $\mathcal{B}((\mathbb{C}^2)^{\otimes n})$ under the Hilbert-Schmidt inner product \eqref{eq:HS}. \todo{we could reason that this is why we like to write Hams as Paulis and ref to our Ham in numerical section}

Each Pauli generates a one-parameter family of \emph{Pauli rotations} $R_\sigma(\theta) := \ee^{-\ii\theta\sigma/2}$ for $\sigma\in\{X,Y,Z\}$ (which can also be extended to multiple qubit rotations $R_{\sigma_j\sigma_k}$). One of them is of particular interest to us, since it can encode arbitrary amplitudes, namely 
\begin{equation}
    R_Y(\theta) = \cos(\tfrac{\theta}{2})\,\idm - \ii\sin(\tfrac{\theta}{2})\,Y.
\end{equation}
Then, for a well-chosen angle $\theta$ its action is $R_Y(2\arcsin\sqrt{p})\ket{0} = \sqrt{1-p}\ket{0} + \sqrt{p}\ket{1}$, i.e. it deterministically loads a probability $p\in[0, 1]$ into the amplitude of $\ket{1}$. We will also make use of the shorthand: $Y_\theta := R_Y(2\arcsin\sqrt{\theta})$ throughout, particularly for the Boltzmann and weak measurement ancillas used in \hyperref[sec:algorithm]{Sec.~\ref*{sec:algorithm}}. 

A \emph{quantum circuit} is a horizontal-time diagram in which each wire carries one qubit (or a bundled register), boxes denote gates applied at that time slice, and a filled dot $\bullet$ on a wire indicates a control: the gate on the target wire fires conditionally on the control being $\ket{1}$ (or controlled on $\ket{0}$ if it is an open dot $\circ$). A half-filled dot \tikz[baseline=-0.5ex]{\node[halfctrl]{};} on a bundled register denotes a \emph{uniform} (SELECT-style) control, under which the target gate is applied in superposition to every computational-basis state of the control register (i.e. all possible bitstrings), each contributing its own branch to the state. All circuit figures in this thesis follow this convention and were created by the quantikz package \cite{kay2018tutorial}.

There is another key element in the quantum circuit: the read-out, which is performed by a \emph{computational-basis measurement}. This is a projective measurement \eqref{eq:post_meas_state} associated with the spectral decomposition of the $Z$-operators on each qubit. On a pure state $\ket{\psi} \in (\mathbb{C}^2)^{\otimes n}$ this returns a bitstring $x\in\{0, 1\}^n$ with probability $|\langle x|\psi\rangle|^2$ and collapses the register onto $\ket{x}$. \emph{Mid-circuit} measurements are also available on several types of quantum hardware, where the outcome can be discarded (yielding a non-unitary channel of Kraus form \eqref{eq:kraus}) or post-selected on a desired result.

The three primitives introduced so far, e.g. universal gate set, fresh ancilla qubits prepared in $\ket{0}$, and measurements, already suffice to realise \emph{any} CPTP map $\Phi: \mathcal{D}(\mathcal{H}^{\otimes n}) \to \mathcal{D}(\mathcal{H}^{\otimes n})$. By Stinespring's dilation theorem (see \hyperref[sec:OQS]{Sec.~\ref*{sec:OQS}} and \cite{bible}), every quantum channel can be lifted to a unitary on a larger Hilbert space. More concretely, there exist an $m$-qubit ancilla register $A$ and a joint unitary $U$ on $n + m$ qubits such that, 
$$
\Phi(\rho) = \tr_A\!\bigl[\,U\,\bigl(\rho \otimes \ket{0}\bra{0}^{\otimes m}\bigr)\,U^\dagger\,\bigr],
$$
with the partial trace implemented either by physically discarding the ancilla or, equivalently, by measuring it in the computational basis and forgetting the outcome. The resulting average over unread outcomes reproduces the Kraus form \eqref{eq:kraus}. Since single-qubit unitaries together with \textsc{CNOT} decomposes $U$ exactly, a circuit implementing $\Phi$ to any desired accuracy requires only the three primitives above.

For any work about quantum algorithms and their realizations we cannot forget about the different timsecales in which a quantum and a classical computer operate. The gates above are physical processes on real hardware, and their wall-clock duration is a hard constant in any complexity statement. On state-of-the-art \emph{superconducting} processors, single- and two-qubit gates take $10-100\, \text{ns}$ (e.g. IBM Heron r3 pushing median two-qubit fidelity above $99.9\%$ \cite{abughanem2025ibm}). On a \emph{trapped-ion} hardware pursued by QVLS in Hannover, the most recent fully chip-integrated two-qubit register on $^{9}\text{Be}^{+}$ runs microwave-driven Mølmer–Sørensen ($R_{XX}$) gates of duration around $1\,\text{ms}$ with a composite process fidelity of $96.6\%$ \cite{pulido2024arbitrary}. By contrast, a logic gate on a classical CMOS hardware switches in $\sim 10\,\text{ps}$. The quantum-classical wall-clock gap is therefore three to eight orders of magnitude per gate, a constant overhead that any claim of ``quantum advantage" must overcome.

\begin{figure}[ht]
  \centering
  \begin{quantikz}[row sep=0.7cm, column sep=0.45cm]
  %
  % Ancilla qubit 0 (MSB)
    \lstick{$\ket{0}$}
      & \gate[4, nwires={2}]{\textsc{Prep}_f}
      & \qw
      & \qw
      & \qw
      & \ctrl{4}
      & \gate[4, nwires={2}]{\text{QFT}}
      &
      &
      & \meter{} \\
  %
  % Ancilla qubit 1
    \lstick{$\ket{0}$}
      &
      & \qw
      & \qw
      & \ctrl{3}
      & \qw
      &
      &
      &
      & \meter{} \\
  %
  % Vdots row (invisible wire, just for the ellipsis)
    \lstick{$\vdots$} \setwiretype{n}
      &
      & \vdots
      &
      &
      &
      &
      &
      &
      & \vdots \\
  %
  % Ancilla qubit r-1 (LSB)
    \lstick{$\ket{0}$}
      &
      & \ctrl{1}
      & \qw
      & \qw
      & \qw
      &
      &
      &
      & \meter{} \\
  %
  % System register (bundled)
    \lstick{$\ket{\psi}$}
      & \qwbundle{}
      & \gate{U}
      & \ \ldots\ \qw
      & \gate{U^{2^{r-2}}}
      & \gate{U^{2^{r-1}}}
      & \qw
      & \qw
      & \qw
      & \rstick{$\ket{\psi}$} \qw
  \end{quantikz}
  \caption{Phase estimation
  circuit $\mathcal{C}_f$ with $r$ ancilla qubits and a bundled system register $\ket{\psi}$. The
  state-preparation block $\textsc{Prep}_f$ loads the time-domain window
  $\ket{f} = \sum_{\bar t\in S_{t_0}} f(\bar t)\ket{\bar t}$ on the ancilla register. (i) The uniform window $f\equiv 1/\sqrt{N}$ is achieved by Hadamard gates on each qubit: this is the standard QPE. (ii) A Gaussian $f(\bar{t}) \propto \ee^{-\bar{t}^2 \sigma^2}$ leads to a more smeared phase estimate (albeit controlled) with the advantage of requiring less time evolution.}
  \label{circ:qpe}
\end{figure}

\smallskip
\section{Phase Estimation} One of the key subroutines classical and quantum Gibbs sampling algorithms rest on is the phase (or energy) estimation. This allows us to label the system with an energy register on which a transition-rate $\gamma(\omega)$, satisfying the KMS condition \eqref{eq:KMS-condition}, can act in superposition. It was first introduced as \emph{Quantum phase estimation} (QPE) by Kitaev in \cite{kitaev1997quantum,cleve1998quantum}, and it was quickly applied to Hamiltonian eigenvalue problems \cite{abrams1999quantum}. It has underpinned not only recent developments in Gibbs sampling, but also the seminal quantum Metropolis algorithm of \cite{temme2011quantum}. As QPE is an already established textbook material, we only highlight the similarities and differences of it to the phase estimation we use later on: a one-parameter generalization of QPE, that can be used to compute \emph{Operator Fourier Transforms} (OFT).

To make this shared structure visible, see the circuit in \hyperref[circ:qpe]{Fig.~\ref*{circ:qpe}}. Let the Hamiltonian $H$ have Bohr frequencies $\nu \in B_H$, fix the time unit $t_0>0$, and an ancilla register of $r$ qubits with time grid $S_{t_0} = \{0, t_0, 2t_0, \dots, (N{-}1)t_0\}$ and a dual frequency grid $S_{\omega_0}$ of spacing $\omega_0$ satisfying $t_0\,\omega_0 = 2\pi/N$. The three stages of the circuit are
\begin{enumerate}
    \item $\textbf{Prep}_f$ on the ancilla, preparing the filter state $\ket{f} := \sum_{\bar t \in S_{t_0}} f(\bar t)\,\ket{\bar t}$ for a chosen window $f : S_{t_0} \to \mathbb{C}$ with $\sum_{\bar t}|f(\bar t)|^2 = 1$.
    \item \textbf{Controlled Hamiltonian evolution,} $\sum_{\bar t}\ket{\bar t}\!\bra{\bar t}\otimes U^{\bar t/t_0}$ with $U = \ee^{+\ii H t_0}$ (plus sign, so the readout register labels $\nu$ directly in our later applications \eqref{eq:oft}). This is achieved by the cascade of evolutions $U, U^2, \ldots, U^{2^{r-1}}$.
    \item \textbf{Quantum Fourier Transform (QFT)}, $\ket{\bar t}\mapsto\frac{1}{\sqrt N}\sum_{\bar\omega}\ee^{-\ii\bar t\bar\omega}\ket{\bar\omega}$, maps the time register to the frequency register $\ket{\bar \omega}$ on $S_{\omega_0}$.
\end{enumerate}
On an eigenstate $H\ket{\psi_\nu} = \nu\ket{\psi_\nu}$ the three stages together act as
\begin{equation} \label{eq:qpe-unified}
    \mathcal{C}_f \,\bigl(\ket{0}^{\otimes r} \otimes \ket{\psi_\nu}\bigr)
\;=\; \sum_{\bar\omega \in S_{\omega_0}} \hat f(\bar\omega - \nu)\,\ket{\bar\omega} \otimes \ket{\psi_\nu},
\end{equation}
where $\hat f(\bar\omega) := \frac{1}{\sqrt{N}} \sum_{\bar t \in S_{t_0}} \ee^{-\ii \bar\omega \bar t} f(\bar t)$ is the discrete Fourier transform of $f(\bar{t})$. Clearly, the choice of the filter function $f$ can significantly influence the result of the phase estimation. Let us consider
the following two cases.

\medskip 
\noindent\textbf{(i) Uniform $f$ -- standard QPE.} To weight each label uniformly by $f(\bar t) = 1/\sqrt{N}$, the $\textsc{Prep}_f$ becomes a simultaneous application of \emph{Hadamard} gates on each of the time register ancillas. A Hadamard gate is defined as, 
$$
\mathsf{H} \;:=\; \tfrac{1}{\sqrt{2}} \begin{pmatrix} 1 & 1 \\ 1 & -1 \end{pmatrix},
$$
and it maps each ancilla $\ket 0 \mapsto (\ket 0 + \ket 1)/\sqrt 2$. The collective result is the uniform time register state $\frac{1}{\sqrt N}\sum_{\bar t}\ket{\bar t}$. The output kernel then becomes
\begin{equation}
    \hat f(\bar\omega - \nu) \;=\; \frac{1}{N}\,\frac{1 - \ee^{-\ii(\bar\omega-\nu)Nt_0}}{1 - \ee^{-\ii(\bar\omega-\nu)t_0}},
\end{equation}
or the squared form that we mostly see in QPE literature, 
\begin{equation} \label{eq:QPE-prob-kernel}
    \bigl|\hat f(\bar\omega-\nu)\bigr|^2 \;=\; \frac{\sin^2\!\bigl((\bar\omega-\nu)Nt_0/2\bigr)}{N^2\sin^2\!\bigl((\bar\omega-\nu)t_0/2\bigr)}
\end{equation}
i.e. the Dirichlet kernel, peaked at $\bar{\omega} = \nu$ with width $\sim 1 / (Nt_0)$ and polynomially decaying \cite[Fig. 8]{chen2023quantum}. Sampling the ancilla register returns an $r$-bit approximation of the scaled eigenphase $\varphi_\nu = \nu t_0/(2\pi)$. In order to resolve it up to $n$ accurate bits with failure probability $\delta$ requires, 
\begin{equation} \label{eq:qpe-ancilla-count}
    r \;=\; n + \left\lceil \log\!\left(2 + \frac{1}{2\delta}\right) \right\rceil
\end{equation}
ancilla qubits \cite[Sec.~5.2.1]{bible}, for a total controlled Hamiltonian evolution time $\sum_{k=0}^{r-1} 2^k t_0 = O(1/\varepsilon)$ at resolution $\varepsilon = 2^{-n}$.

\medskip
\noindent\textbf{(ii) Smoothly decaying $f$.} Replacing the uniform filter by a rapidly decaying one, e.g. a discretized Gaussian \eqref{eq:gaussian-filter-t}, or a Gevrey function \cite[text]{ding2025efficient}, turns $\hat{f}$ into a smooth kernel of tuneable width, smearing each eigenvalue into a packet of the window's width ($2\sigma$ for the Gaussian \eqref{eq:smoothed-jump}).

Resolution can be similarly computed: for a Gaussian width $2\sigma$, truncating the discretized window to relative $\ell_2$-tail $\varepsilon$ requires time range $Nt_0 \gtrsim \sigma^{-1}\sqrt{2\log(1/\varepsilon)}$ and frequency spacing $\omega_0 = 2\pi/(Nt_0) \lesssim \sigma$, giving
\begin{equation} \label{eq:gaussian-ancilla-count}
    r \;\geq\; \left\lceil \log\!\frac{2\pi}{\sigma t_0} \;+\; \tfrac{1}{2}\log\ln\frac{1}{\varepsilon} \right\rceil,
\end{equation}
and total evolution time $Nt_0 = \mathcal{O}\!\bigl(\sigma^{-1}\sqrt{\log(1/\varepsilon)}\bigr)$, which is linear in the inverse resolution up to $\sqrt{\log(1/\varepsilon)}$ (Gevrey windows replace the square root by $\mathrm{polylog}(1/\varepsilon)$ \cite[text]{ding2025efficient}).

Both case (i) and (ii) obey $T\cdot\Delta\omega \gtrsim 1$ with $T = Nt_0$ the total controlled Hamiltonian evolution time and $\Delta \omega$ the width of the frequency kernel $\hat{f}$. However, they have different qualities that are worth highlighting to understand what we gain (or lose) by choosing the latter one as our subroutine.

(i) Standard QPE is optimal for the classical-output task: input the eigenstate $\ket{\psi_\nu}$, return an $r$-bit estimate $\hat{\nu}$ of the eigenphase. For this task a matching quantum query lower bound holds: any algorithm that accesses $H$ only through controlled $U = \ee^{\pm \ii H t_0}$ and returns $\varepsilon$-accurate estimate with constant success probability must use total controlled-evolution time $T = \Omega(1/\varepsilon)$ \cite{giovannetti2006quantum}. This is the \emph{Heisenberg limit} $\Delta E \cdot T \sim 1$ of quantum metrology and is a genuine quantum speed-up over the ``standard quantum limit" $T\sim1/\varepsilon^2$ obtainable from $M$ independent depth $t_0$ measurements followed by classical averaging. QPE's $U^1, U^2, \ldots, U^{2^{r-1}}$ cascade attains $T = \mathcal{O}(1 / \varepsilon)$ \cite{bible} and thus saturates this lower bound. Although its Dirichlet kernel has polynomially-decaying tails \eqref{eq:QPE-prob-kernel}, and is \emph{not} a minimizer of frequency spread in the sense of case (ii), the median estimation error still decays as $1 / T$, which is all the estimation task demands.

(ii) The KMS detailed balanced Lindbladians of \hyperref[sec:kms-lindbladian]{Sec.~\ref*{sec:kms-lindbladian}} use Gaussian filtered phase estimates: by sandwiching a jump operator $A$ between filtered phases estimations, we can decompose jump operators into their Fourier components $\hat{A}(\bar{\omega})$ (\eqref{eq:oft}) and allow us to weigh them with the KMS transition rates $\gamma(\bar{\omega})$. The key metric is no longer estimation error but the $L^2$ width of $\bar{f}$. The governing bound then, is the classical Heisenberg inequality: for any $f\in L^2(\mathbb{R})$ of unit norm \cite{folland1997uncertainty}, 
$$
\sigma_t\,\sigma_\omega \;\geq\; \tfrac{1}{2},
$$
Under the circuit identification $\sigma_t \leftrightarrow T$, $\sigma_\omega \leftrightarrow \Delta\omega$ this becomes a lower bound on controlled-evolution time that no window can evade. The Gaussian \eqref{eq:gaussian-filter-t} uniquely saturates the above equality, while the Dirichlet does not. But this is a result coming from harmonic analysis, independent of any quantum metrological argument, and is not to be confused with the Heisenberg limit of (i). Although, recent results found that the Heisenberg-limit with shallower circuits are achievable \cite{ding2023even} (with just an ancilla combined with statistical phase-estimation protocols), these cannot be used in our setup since they collapse the ancilla to a classical sample, destroying the superposition of $\ket{\bar{\omega}}$.

Since we will work with OFT-based Lindbladian generators, the inherent obstacle that arises is the classical Fourier bound, not Heisenberg limit. As the Bohr spectrum $B_H$ of a non-trivial many-body Hamiltonian is exponentially dense in $n$, an attempt to exactly resolve Bohr frequencies would require $T$ exponentially growing with $n$, and thus no choice of filter function can escape this. The smoothened generators of \hyperref[sec:kms-lindbladian]{Sec.~\ref*{sec:kms-lindbladian}} take this impossibility as the starting point: they fix a controlled resolution width $\sigma$ and arrange the dissipator and Lamb-Shift such that KMS detailed balance survives even this finite-resolution smearing.

\smallskip
\label{sec:prelim-trotter}
\section{Trotterization}
\begin{figure}[ht]
  \centering
  \resizebox{\textwidth}{!}{%
  \begin{quantikz}[
    transparent, row sep=0.6cm, column sep=0.4cm,
    /tikz/operator/.append style={inner sep=3pt, minimum height=6mm, font=\small}
  ]
    \lstick{$q_1$}
  % A-block L (active)
      & \gate[2]{R_{X_1X_2}}
      & \gate[2]{R_{Y_1Y_2}}
      & \gate[2]{R_{Z_1Z_2}}
  % B-block L (q_1 idle)
      & \qw & \qw & \qw
  % C-block (boundary, q_1 leads the [3]-wide gate; R^2 = R applied twice = rotate by 2 dt)
      & \gate[3, label style={yshift=0.4cm}]{R_{X_3X_1}^{2}}
      & \gate[3, label style={yshift=0.4cm}]{R_{Y_3Y_1}^{2}}
      & \gate[3, label style={yshift=0.4cm}]{R_{Z_3Z_1}^{2}}
  % B-block R (q_1 idle)
      & \qw & \qw & \qw
  % A-block R (active)
      & \gate[2]{R_{X_1X_2}}
      & \gate[2]{R_{Y_1Y_2}}
      & \gate[2]{R_{Z_1Z_2}}
      & \qw \\
  %
    \lstick{$q_2$}
  % A-block L (empty: consumed by q_1's gate[2])
      & & &
  % B-block L (active, q_2 leads)
      & \gate[2]{R_{X_2X_3}}
      & \gate[2]{R_{Y_2Y_3}}
      & \gate[2]{R_{Z_2Z_3}}
  % C-block: q_2 passes through unaffected
      & \dashedthrough & \dashedthrough & \dashedthrough
  % B-block R (active, q_2 leads)
      & \gate[2]{R_{X_2X_3}}
      & \gate[2]{R_{Y_2Y_3}}
      & \gate[2]{R_{Z_2Z_3}}
  % A-block R (empty: consumed by q_1's gate[2])
      & & &
      & \qw \\
  %
    \lstick{$q_3$}
  % A-block L (q_3 idle)
      & \qw & \qw & \qw
  % B-block L (empty: consumed by q_2's gate[2])
      & & &
  % C-block (empty: consumed by q_1's gate[3])
      & & &
  % B-block R (empty: consumed by q_2's gate[2])
      & & &
  % A-block R (q_3 idle)
      & \qw & \qw & \qw
      & \qw
  \end{quantikz}%
  }
  \caption{One second-order Strang Trotter step for the $n{=}3$ periodic
  Heisenberg chain. The sum of Pauli strings in $H_\mathrm{Heis}$ turn into a product of 2-qubit unitary Pauli rotations,
  $R_{\sigma_j\sigma_k}(\theta) \equiv \ee^{-\ii\theta\,\sigma_j\sigma_k/2}$, where $\theta$ encodes coupling strengths of the Hamiltonian and the required time step.
  The boundary bond $(3,1)$ closes the ring and $q_2$ passes through unaffected
  (dashed wire).}
  \label{circ:trotter-strang}
\end{figure}
Up to this point we have concealed a fact that Hamiltonian time evolutions $U(t) = \ee^{\pm \ii H t}$ often cannot be exactly implemented by a quantum computer, and they have to be approximated by a product formula $S_p(t)$ instead. This is also true for the local (i.e. each term acts on a bounded number of qubits) and non-commuting Hamiltonians we care about: there is no known shallow circuit that would avoid the use of Trotterization. Here, we will give only the necessary details that are relevant for the thesis, based on the work of Childs et al. \cite{childs2021theory}.

Let us write the Hamiltonian as a sum of $\Gamma$ non-commuting summands
$$
H \;=\; \sum_{\ell=1}^{\Gamma} H_\ell,
$$
where each $H_\ell$ is either a single Pauli string or a commuting family of such strings whose exponential $\ee^{-\ii t H_\ell}$ is easy to implement (typically because all terms in $H_\ell$ share support disjointly or pairwise commute). A $p$-th order product formula $S_p(t)$ is an ordered product of single-summand evolutions $\ee^{-\ii t H_\ell}$ that together approximates the full Hamiltonian evolution $\ee^{-\ii t H}$ up to error $\mathcal{O}(t^{p + 1})$ for small $t$ \cite{childs2021theory}. The simplest one for reference, with $p=1$ is called the Lie-Trotter formula, 
\begin{equation}
    S_1(t) \;:=\; \ee^{-\ii t H_\Gamma}\cdots \ee^{-\ii t H_1},
\end{equation}
but the one we mainly use for the numerical analysis is the Strang splitting with $p=2$:
\begin{equation}\label{eq:strang}
    S_2(t) \;:=\; \Bigl(\prod_{\ell=1}^{\Gamma} \ee^{-\ii (t/2) H_\ell}\Bigr)\Bigl(\prod_{k=\Gamma}^{1} \ee^{-\ii (t/2) H_\ell}\Bigr).
\end{equation}
The former leads to errors $\|S_1(t) - \ee^{-\ii H t}\| = \mathcal{O}(t^2)$ while the latter $\|S_2(t) - \ee^{-\ii H t}\| = \mathcal{O}(t^3)$. There is also a structural difference, namely that Strang splitting is palindromic, i.e. $S_2(t)^\dagger = S_2(-t)$ which is worth preserving to avoid additional deviations from detailed balance \hyperref[rem:palindromic]{Remark~\ref*{rem:palindromic}}. The circuit realization of $S_2$ for $n=3$ qubit, periodic Heisenberg chain $H_\mathrm{Heis}$ \eqref{eq:H_heis} is given in \hyperref[circ:trotter-strang]{Fig.~\ref*{circ:trotter-strang}}. Since each $H_k$ is Hermitian by construction, each factor $\ee^{-\ii t H_\ell}$ is unitary, so $S_p(t)$ and its iterate $S_p^{(M)}(t) := S_p(t/M)^M$ are unitary for every choice of $p$ and $M$. Notably, we see that an $n$ qubit unitary $U$ is broken down into 2-qubit gates. This is important for implementation since they can be easily transpiled to native 1-qubit and 2-qubit gates that a quantum hardware uses. To motivate it even more: transpiling a general $n$ qubit unitary to native gates in a brute-force manner is currently only possible up to $n = 9$ qubits \cite{rakyta2024highly}.

Higher-order Suzuki formulas $S_{2r}$ can be built recursively from $S_2$ by the fractal construction $S_{2r}(t) = S_{2r-2}(u_r t)^2\,S_{2r-2}((1-4u_r)t)\,S_{2r-2}(u_r t)^2$ with $u_r = 1/(4 - 4^{1/(2r-1)})$, leading to errors of $\mathcal{O}(t^{2r+1})$ \cite{childs2021theory}.

A naive Taylor truncation only controls the error through $\|H_k\|$, ignoring the commutativity of the summands. The key result of \cite{childs2021theory} is that the error of any $p$-th order formula is in fact controlled by nested commutators of the summands. For some real $t$ the error is given by \cite[Theorem 6]{childs2021theory}, 
\begin{equation} \label{eq:childs-comm-bound}
    \begin{split}
        &\bigl\|S_p(t) - \ee^{-\ii H t}\bigr\|
        \;=\; \mathcal{O}\!\bigl(\tilde\alpha_{\mathrm{comm}}\,|t|^{p+1}\bigr) \\ 
        \text{with}\qquad &\tilde\alpha_{\mathrm{comm}}
        \;:=\; \sum_{\ell_1,\dots,\ell_{p+1}=1}^{\Gamma}
        \bigl\|\,[H_{\ell_{p+1}},\,\cdots\,[H_{\ell_2},H_{\ell_1}]\cdots]\,\bigr\|,
    \end{split}
\end{equation}
where the constant $\tilde{\alpha}_\mathrm{comm}$ is universal (independent of $\Gamma$, $\|H_\ell\|$, and the ordering fixed by the particular formula). More concretely, for the second-order Strang splitting,
\begin{equation} \label{eq:strang-explicit-error}
        \bigl\|S_2(t) - \ee^{-\ii H t}\bigr\|
    \;\leq\; \sum_{\ell_1 < \ell_2}\!\left(
    \frac{t^3}{12}\bigl\|[H_{\ell_2},[H_{\ell_2},H_{\ell_1}]]\bigr\|
    \;+\; \frac{t^3}{24}\bigl\|[H_{\ell_1},[H_{\ell_1},H_{\ell_2}]]\bigr\|
    \right)
    \;+\; \mathcal{O}(t^4).
\end{equation}
For longer time evolutions the target evolution $\ee^{-\ii H t}$ is split into $M$ Trotter steps of size $\delta t = t / M $. Then, a total error of $\varepsilon$ can be achieved if $M$ satisfies, 
\begin{equation} \label{eq:trott-num-bound}
    M \;\geq\; \Omega\!\left(\frac{\tilde\alpha_{\mathrm{comm}}^{1/p}\;t^{\,1+1/p}}{\varepsilon^{\,1/p}}\right).
\end{equation}
For $p=1$ this is the $\mathcal{O}(\tilde\alpha_{\mathrm{comm}}\, t^2/\varepsilon)$ worst-case count, and for Strang ($p = 2$) the $\mathcal{O}(\sqrt{\tilde\alpha_{\mathrm{comm}}}\, t^{3/2}/\sqrt{\varepsilon})$ bound is what is used for our specialized results in the Trotter error analysis of [REF] \todo{ref to trotter}

One of the Hamiltonians we use for the numerical analysis (the same we have a sketched Trotter step for in \hyperref[circ:trotter-strang]{Fig.~\ref*{circ:trotter-strang}}) is the periodic Heisenberg chain on $n$ qubits, 
\begin{equation} \label{eq:H_heis}
    H_{\mathrm{Heis}}
\;=\; J \sum_{i=1}^{n} \bigl( X_i X_{i+1} + Y_i Y_{i+1} + Z_i Z_{i+1} \bigr)
\;+\; \sum_{i=1}^{n} \zeta_i Z_i,
\end{equation}
with coupling $J$ (if positive - antiferromagnetic, if negative - ferromagnetic) and disordering external fields $\zeta_i \in [-\zeta,\zeta]$. 

In the following, we give more details about the commutant constant $\tilde{\alpha}_\mathrm{comm}$ for the family of \emph{geometrically local} Hamiltonians -- the family that the Heisenberg model \eqref{eq:H_heis} is also part of. These Hamiltonians can be written as
$$
 H \;=\; \sum_c h_c, \qquad \|h_c\| \le h_*, 
$$
where each local term $h_c$ is supported on at most $k$ sites, with maximum overlapping supports of degree $d$, making the Hamiltonian $H$ $(k, d)-$\emph{local} with $k, d = \mathcal{O}(1)$. Such an $H$ admits a partition into $\Gamma = \mathcal{O}(1)$ commuting classes, i.e. a sum of pair-wise commuting local terms. Expanding the $(p+1)$ nested commutator $[H_{\ell_{p+1}}, \cdots [H_{\ell_2}, H_{\ell_1}]]$ term-by-term over local terms, we obtain a sum over ordered $(p+1)$-tuples of bonds $(c_1, \dot, c_{p+1})$. The only surviving terms in this sum are the ones that have \emph{connected} bonds, i.e. where each $c_{i+1}$ overlaps with at least one of $c_1, \dots, c_i$, since the disjoint terms commute with each other. Therefore, an estimating $\tilde{\alpha}_\mathrm{comm}$ for $(k,d)$-local Hamiltonians reduces to the problem of counting these connected bonds. The number of connected $(p+1)$-tuples anchored at a pivot bond $c_1$ is at most $(2d)^p$, and there are $\Theta(n)$ choices of pivot bonds, giving
\begin{equation} \label{eq:alpha-comm}
\tilde\alpha^{(p)}_{\mathrm{comm}} \;\leq\; C_{d,p}\,n\,h_*^{p+1} \;=\; \mathcal{O}(n).
\end{equation}
Here, $C_{d,p} \le \Gamma^{p+1}\,(2d)^p$ is a combinatorial constant depending only on the bond-graph degree $d$ and the order $p$. The counting argument for this is present in \cite[Sec. V B]{childs2021theory} and goes as follows. The outer commuting-class grouping fixes $\Gamma = \mathcal{O}(1)$ and selects which classes contribute (e.g. even or odd sites), while inner bounded-degree bond counting yields the $\Theta(n)\cdot(2d)^p$ bound. It is the inner bond layer that carries the geometric locality assumption. For instance, an all-to-all interaction that has $\Gamma = 2$ summands but unbounded bond degree $d$, yields $\mathcal{O}(n^3)$ commutator constant rather than $\mathcal{O}(n)$ \cite[Sec. IV C]{childs2021theory}.

The 1D (periodic) Heisenberg chain \eqref{eq:H_heis} is a canonical case for $(k,d)$-local Hamiltonian. Its terms can be grouped into a couple of commuting classes: odd and even sites in case $n$ is even (or into three classes if $n$ is odd due to the boundary.) As the interaction degree on the chain is $d=2$, and the number of commuting classes $\Gamma = 2$ (or $3$), it is possible to find some estimate for the combinatorial factor in \eqref{eq:alpha-comm}. Although, with numerical heuristics \cite{childs2021theory}, we can get that constant value, for us the scaling is already sufficient, i.e. for $p=2$: $\tilde\alpha^{(2)}_{\mathrm{comm}}(H_{\mathrm{Heis}}) = \mathcal{O}(|J|^3 n)$. The same principle applies to the periodic 1D transverse-field Ising model and the 2D Heisenberg model on a square lattice, again with $\tilde\alpha^{(2)}_{\mathrm{comm}} = \mathcal{O}(n)$ scaling and constant prefactors that change only with the bond-interaction degree $d$.

\smallskip
\section{Linear Combination of Unitaries}\label{sec:prelim-LCU}

\begin{figure}[ht]
  \centering
  \begin{quantikz}[row sep=0.75cm, column sep=0.55cm]
    \lstick{$\ket{0^b}$}
      & \qwbundle{}
      & \gate{\text{PREP}}
      & \ctrl{1}
      & \gate{\text{PREP}^\dagger}
      & \rstick{$\bra{0^b}$} \qw
      \\
    \lstick{$\ket{\psi}$}
      & \qwbundle{}
      & \qw
      & \gate{U(t)}
      & \qw
      & \rstick{$\displaystyle\frac{1}{\|f\|_1}\sum_t f(t)\,U(t)\,\ket{\psi}$} \qw
  \end{quantikz}
  \caption{LCU block encoding $U_A$
  of $A = \sum_t f(t)\,U(t)$: \textsc{Prep} writes the filter state
  $\ket{f} = \frac{1}{\sqrt{\|f\|_1}}\sum_t \sqrt{f(t)}\ket{t}$, \textsc{Select}
  applies $U(t)$ controlled on the ancilla, and \textsc{Prep}$^\dagger$
  uncomputes. Post-selecting the ancilla on $\ket{0}$ leaves
  $A\ket{\psi}/\|f\|_1$ on the system, with success probability
  $\|A\ket{\psi}\|^2/\|f\|_1^2$.}
  \label{circ:lcu}
\end{figure}

Another key subroutine we will make use of is the linear combination of unitaries (LCU), along with the transformation of these block encoded operators introduced in the subsequent section. For a comprehensive introduction, please see \cite{childs2012hamiltonian,berry2015simulating}.

Fix a unitary family $\{U(t)\}_{t\in S_t}$ indexed by a discrete grid $S_t$ and a scalar weight $f : S \to \mathbb{R}_{\geq 0}$ with finite $\ell_1$ norm $\|f\|_1 = \sum_t f(t)$. The target operator
$$
A \;=\; \sum_{t \in S} f(t)\,U(t)
$$
is in general Hermitian but not unitary if $U(-t) = U(t)^\dagger$ and $f$ is real and symmetric. Its operator norm $\|A\|$ is at most $\|f\|_1$ by the triangle inequality. The LCU circuit in \hyperref[circ:lcu]{Fig.~\ref*{circ:lcu}} realises $A$ up to a known rescaling using an ancilla register of $b = \lceil \log|S_t|\rceil$ qubits and the following three ingredients: 
\begin{enumerate}
    \item \textsc{Prep} on the ancilla, prepares the filter state $\ket{f} := \frac{1}{\sqrt{\|f\|_1}}\sum_{t \in S} \sqrt{f(t)}\,\ket{t}$
    \item \textsc{Select} picks out each $U(t)$ for a given $t$, i.e. $\sum_{t}\ket{t}\!\bra{t}\otimes U(t)$
    \item $\textsc{Prep}^\dagger$ uncomputes the ancilla register.
\end{enumerate}
Writing the composite circuit as $U_A := (\textsc{Prep}^\dagger \otimes \idm)\, \textsc{Select}\, (\textsc{Prep} \otimes \idm)$, a direct calculation shows that its $\ket{0^b}$-$\ket{0^b}$ matrix element (the top left corner) is exactly, 
\begin{equation} \label{eq:lcu}
    \bra{0^b}\,U_A\,\ket{0^b} \;=\; \frac{A}{\|f\|_1},
\end{equation}
i.e. $U_A$ is a block encoding of $A$ with subnormalisation $Z_A := \|f\|_1$ \cite{low2019hamiltonian}. The subnormalisation is the price for packing a non-unitary $A$ inside a unitary: $U_A$ acts as $A / Z_A$ on the block projector $\ket{0^b}\!\bra{0^b}$ and as an unspecified completion on its orthogonal subspace, and all complexity statements about simulation times and success probabilities scale with $Z_A$.

This approach can be extended to cases where $f$ assumes signed or complex values. There are different solutions, but the one we end up using later is the following: split $\textsc{Prep}$ into an asymmetric left/right pair $\textsc{Prep}'$ / $\textsc{Prep}$. The prior writes the signed / phased amplitude $\sum_t \mathrm{sgn}(f(t))\sqrt{|f(t)|/\|f\|_1}\,\ket{t}$ on the left, and the latter writes the magnitude-only amplitude $\sum_t \sqrt{|f(t)|/\|f\|_1}\,\ket{t}$ on the right, so that the two amplitudes multiply to $f(t)/\|f\|_1$. \hyperref[circ:U_B_a]{Fig.~\ref*{circ:U_B_a}} adopts this with $\mathrm{sgn}(\cdot)$ is a unit-modulus phase in the complex case.

Applying $U_A$ on $\ket{0^b}\otimes\ket{\psi}$ and \emph{post-selecting} the ancilla on $\ket{0^b}$ leaves, up to normalisation, the state $A \ket{\psi} / Z_A$ on the system, with success probability $\|A\ket{\psi}\|^2/Z_A^2$. When $Z_A \leq 1$, the block encoding can be taken as a deterministic implementation of $A$. In the case when $Z_A \geq 1$, however, the post-selection succeeds only with $\mathcal{O}(1 / Z_A^2)$, in which case LCU is \emph{not} deterministic. For bare applications of LCU, there are already remedies for this, e.g. a repeat-until-success procedure with $\mathcal{O}(Z_A^2)$ of repetitions, or the so-called oblivious amplitude amplifications which involves additional reflections. But in our thesis, an indeterministic LCU can only happen when we create a block encoding for a Hermitian operator $B$ and then directly transform it to a coherent evolution $\ee^{-\ii B t}$. This transformation goes hand in hand with block-encodings, which is the next subroutine we will introduce. There, it becomes clear what role the subnormalization factor $Z_A$ plays in the algorithm's complexity.

\smallskip
\section{Quantum Signal Processing} \label{sec:prelim-QSP}

\begin{figure}[ht]
  \centering
  \resizebox{1.1\textwidth}{!}{%
  \begin{quantikz}[row sep=0.7cm, column sep=0.5cm]
  %
  % ==================================================================
  % Row 1: QSP ancilla.
  % Time order in this figure (left to right = forward in time):
  %   col 3:   opening rotation L_0 = R_MW(theta_0, phi_0, lambda_0)
  %   col 4-5: representative open-controlled slot, repeated d times
  %            (j = 1, ..., d, applied first in time after L_0)
  %   col 6-7: representative closed-controlled slot, repeated d times
  %            (j = d+1, ..., 2d, applied last in time)
  % j INCREASES left-to-right in this figure (matching MW2024 Eq. 42 with
  % the index convention "j = 1 is the first layer applied after L_0").
  % ==================================================================
    \lstick{$\ket{0}_\mathrm{QSP}$}
      & \qw
      & \gate{R(\theta_0,\phi_0,\lambda_0)}
      & \octrl{1}
        \gategroup[3,steps=2,style={dashed,rounded corners,fill=blue!6,inner sep=5pt},
                   background,label style={label position=above,yshift=0.25cm}]
          {\scriptsize Repeat $d$ times}
      & \gate{R(\theta_j,\phi_j,0)}
      & \qw
      & \ctrl{1}
        \gategroup[3,steps=2,style={dashed,rounded corners,fill=black!6,inner sep=5pt},
                   background,label style={label position=above,yshift=0.25cm}]
          {\scriptsize Repeat $d$ times}
      & \gate{R(\theta_{d+j},\phi_{d+j},0)}
      & \rstick{$\bra{0}_\mathrm{QSP}$}\qw
      \\
  %
  % ==================================================================
  % Row 2: Block-encoding ancilla (bundled).
  % Multi-row gates: \gate[2]{W} for the open-ctrl block (col 4),
  %                  \gate[2]{W^\dagger} for the closed-ctrl block (col 6).
  % NO uncontrolled tail in Form B.
  % PREP/PREP^dagger live INSIDE U_A inside each W and are not drawn here.
  % ==================================================================
    \lstick{$\ket{0^b}$}
      & \qwbundle{}
      & \qw
      & \gate[2]{W}
      & \qw
      & \qw
      & \gate[2]{W^\dagger}
      & \qw
      & \rstick{$\bra{0^b}$}\qw
      \\
  %
  % ==================================================================
  % Row 3: System register (bundled).
  % Cells under the multi-row gates (col 4: W, col 6: W^dagger) are LEFT
  % EMPTY so quantikz draws those gates spanning rows 2-3.
  % ==================================================================
    \lstick{$\ket{\psi}$}
      & \qwbundle{}
      & \qw
      &
      & \qw
      & \qw
      &
      & \qw
      & \rstick{$\approx \ee^{-\ii\delta A}\ket{\psi}$}\qw
  \end{quantikz}
  }
  \caption{GQSP circuit approximating
    $\ee^{-\ii\delta A}$ as a Laurent polynomial on the post-selected $\ket{0}_\mathrm{QSP}\otimes\ket{0^b}$
    block from the walk operator $W = R_b\,U_A$. The dashed boxes are repeated $d$ times: one with $W$ the other with $W^\dagger$ (each followed by a two-parameter rotation $R(\theta_\cdot,\phi_\cdot,0)$), with indices $j = 1,\dots,d$ on the first block and $d+j$ on the second. In a gate-level implementation the
    $\textsc{Prep}/\textsc{Prep}^\dagger$ inside $U_A$ are pulled out of
    the controlled walks (adjacent $\textsc{Prep}^\dagger\textsc{Prep}=I$) to avoid unnecessary controlled applications.}
  \label{circ:gqsp}
\end{figure}

Without spoiling too much of the coming theoretical discussions, let us simply claim, that for a detailed balanced Lindbladian we need to apply a coherent evolution of the form $\ee^{-\ii\delta A}$ with respect to some Hermitian operator $A = A^\dagger$. There are certainly many ways to apply such a coherent part in a Lindbladian evolution, the one we have decided upon is conceptually and in circuit construction terms the simplest while operating at the same approximating error order as the other parts of \hyperref[alg:main]{Algorithm~\ref*{alg:main}}, namely as the weak measurement scheme, see the next section.

The LCU circuit of the previous section \hyperref[circ:lcu]{Fig.~\ref*{circ:lcu}} delivers a block encoding $U_A$ on a $b$-qubit register of some Hermitian $A$ on an $n$-qubit register, 
$$
\langle 0^b |\,U_A\,|0^b \rangle \;=\; A/\alpha, \qquad \|A/\alpha\| \leq 1,
$$
with $A = \sum_j \alpha_j U_j$, $\alpha = \sum_j \alpha_j$, preparation unitary $\textsc{Prep}\ket{0} = \sum_j \sqrt{\alpha_j/\alpha}\,\ket{j}$, selection unitary $\textsc{Select} = \sum_j \ket{j}\!\bra{j}\otimes U_j$, and $U_A = (\textsc{Prep}^\dagger\otimes I)\,\textsc{Select}\,(\textsc{Prep}\otimes I)$. 

It is often the case however, that we do not purely want to apply $A$ to the system but rather some function of it, e.g. the coherent evolution $\ee^{-\ii\delta A}$, with $A = A^\dagger$. \emph{ Quantum signal processing} (QSP) is the technique that bridges the gap: given a block encoding like the LCU, it post-selectively realises an arbitrary polynomial transformation of the encoded operator's $A$ spectrum. There are a miriad of ways to implement QSP for a thorough review see \cite{skelton2025hitchhiker}. We have decided upon the \emph{Generalized QSP} (GQSP) variant of \cite{motlagh2024generalized}, which lifts the parity restriction of the original QSP construction and works on the unit circle by using more general $\mathrm{SU}(2)$ rotations. This makes the natural class for Hamiltonian simulation targets $\ee^{-\ii\delta\alpha\cos\theta}$, the class of all complex Laurent polynomials $P$ directly accessible, as we explain below.

Define the walk operator $W$ as the unitary block encoding $U_A$ followed by a reflection $R_b$ around the all-zero state of the block encoding ancilla \cite[Cor. 8]{motlagh2024generalized}
\begin{equation} \label{qsp-walk}
    W \;:=\; R_b\,U_A, \qquad R_b \;:=\; 2\,\ket{0^b}\!\bra{0^b} - I.
\end{equation}
The primitives $\textsc{Prep}/\textsc{Select}/\textsc{Prep}^\dagger$ will always be defined as part of the unitary block encoding $U_A$ -- both \hyperref[circ:lcu]{Fig.~\ref*{circ:lcu}} and \hyperref[circ:gqsp]{Fig.~\ref*{circ:gqsp}} adopt this convention. For each eigenvector 
$\ket\lambda$ of $A$ with $A\ket\lambda = \lambda\ket\lambda$, the two-dimensional subspace containing $\ket{0^b}\ket\lambda$ is invariant under $W$. Hence, in a suitable orthonormal basis of that subspace, $W$ acts as the $\mathrm{SU}(2)$ matrix, 
$$
W_\lambda \;=\; \begin{pmatrix} \lambda/\alpha & \sqrt{1-(\lambda/\alpha)^2} \\[2pt] -\sqrt{1-(\lambda/\alpha)^2} & \lambda/\alpha \end{pmatrix}.
$$
The eigenvalues are given by $\ee^{\pm\ii\theta_\lambda}$, where $\cos\theta \;=\; \lambda/\alpha$ with $\lambda \in [-\alpha,\alpha]$ and $\theta \in [0,\pi]$. Any spectral transformation on the $\lambda$-eigenspace of $A$ thus reduces to approximating the target function at $z = \ee^{\ii\theta_\lambda}$ on the unit circle $\mathbb{T} = \{z\in\mathbb{C} : |z|=1\}$.

For an ordinary polynomial $P \in \mathbb{C}[z]$ with $\|P\|_{\infty,\mathbb{T}} := \max_{z\in\mathbb{T}}|P(z)| \leq 1$, and an integer $0 \leq k \leq \deg P$, the GQSP circuit \hyperref[circ:gqsp]{Fig.~\ref*{circ:gqsp}} applies a polynomial transformation onto the qubitized subspace by interleaving walk operators and single-qubit rotations defined by 
$$                      
  R(\theta, \phi, \lambda) = \begin{pmatrix} \ee^{\ii(\phi+\lambda)}\cos\theta &
  \ee^{\ii\phi}\sin\theta \\ \ee^{\ii\lambda}\sin\theta & -\cos\theta \end{pmatrix}                  
$$
with angles found by \cite[Alg. 1]{motlagh2024generalized}. Concretely:
\begin{enumerate}
    \item Start with an initial rotation $L_0 = R(\theta_0,\phi_0)\cdot\mathrm{diag}(\ee^{\ii\lambda_0}, 1)$ on the QSP ancilla, with $\lambda_0\in \mathbb{R}$ an absorbed phase parameter distinct from the spectral values $\lambda$.
    \item This is followed by a $\deg P$ controlled walk interleaved with the QSP ancilla rotations $R(\theta_j, \phi_j, 0)$, where $j = 1, \ldots, \deg P$.
\end{enumerate}
Importantly, the $\deg P$ controlled walks can be split into two groups of polarity in order to realise negative powers of polynomials. The first $\deg P - k$ slots ($j = 1, \ldots, \deg P - k$) as open-controlled forward walks, 
$$
\Lambda(W) \;:=\; \ket{0}\!\bra{0}_{\mathrm{QSP}}\otimes W \;+\; \ket{1}\!\bra{1}_{\mathrm{QSP}}\otimes I,
$$
and the last $k$ slots ($j = \deg P - k + 1, \ldots, \deg P$) are closed-controlled inverse walks, 
$$
\Lambda'(W) \;:=\; \ket{0}\!\bra{0}_{\mathrm{QSP}}\otimes I \;+\; \ket{1}\!\bra{1}_{\mathrm{QSP}}\otimes W^\dagger.
$$
Post-selecting both ancillas on $\ket{0}_{\mathrm{QSP}}\ket{0^b}$  leads to the operator \cite[Thm. 6]{motlagh2024generalized}, 
\begin{equation} \label{eq:qsp-poly}
    P'(W_\lambda) \;=\; W_\lambda^{-k}\,P(W_\lambda),
\end{equation}
which on the unit circle $z = \ee^{\ii\theta_\lambda} \in \mathbb{T}$  realises, for each qubitised two-dimensional subspace, the Laurent polynomial:
$$
P'(z) \;=\; z^{-k}\,P(z) \;=\; \sum_{m=0}^{\deg P} c_m\,z^{m - k} \;\in\; \mathbb{C}[z, z^{-1}].
$$
The monomials here, span the bilateral range $\{z^{-k}, z^{1-k}, \ldots, z^{\deg P - k}\}$. The two slot polarities are precisely what lifts the accessible class from ordinary into Laurent polynomials. The $k$ closed-controlled $\Lambda'(W)$ slots together contribute the negative-exponent prefactor $z^{-k}$, while the $\deg P - k$ open-controlled $\Lambda(W)$ slots, paired with the QSP-ancilla rotations $\{R(\theta_j,\phi_j,0)\}$ and the opening rotation $L_0$, build up the ordinary factor $P(z)$. The angle vector $(\theta_0, \phi_0, \lambda_0; \{(\theta_j, \phi_j)\}_{j=1}^{\deg P})$ depends only on $P$, while $k$ enters purely through the slot pattern. The two are chosen independently subject to $0 \leq k \leq \deg P$.

Letting $(P, k)$ range over the constraint set $\|P\|_{\infty,\mathbb{T}} \leq 1$, $0 \leq k \leq \deg P$, GQSP allows access to the full class
$$
\mathcal{P}_{\text{GQSP}} \;=\; \bigl\{P' \in \mathbb{C}[z, z^{-1}] \;:\; P'(z) = z^{-k}P(z),\; P \in \mathbb{C}[z],\; \|P\|_{\infty,\mathbb{T}} \leq 1,\; 0 \leq k \leq \deg P \bigr\}
$$
of Laurent polynomials whose ordinary partner is sup-norm-bounded on $\mathbb{T}$ — strictly richer than the ordinary polynomials accessible from an all-$\Lambda(W)$ circuit ($k = 0$). This Laurent class is exactly what the coherent target $\ee^{-\ii\delta A}$ requires.

The Hamiltonian simulation target can be cast in the form $\ee^{-\ii\delta\alpha\cos\theta}$, which is a Laurent function of $\ee^{\ii\theta}$, \cite[Thm. 6]{motlagh2024generalized}. Using the Jacobi-Anger expansion on this form yields, 
\begin{equation} \label{eq:jacobi-anger}
    \ee^{-\ii\delta\alpha\cos\theta} \;=\; \sum_{k=-\infty}^{\infty} (-\ii)^k\,J_k(\delta\alpha)\,\ee^{\ii k\theta}, \qquad \theta \in \mathbb{R},
\end{equation}
with the $k$-th Bessel functions $J_k$. As it turns out, truncating this series at $|k|\leq d < \infty$ can still give a good approximation to the target exponential function. We find the Laurent polynomials $L_d$ of degree $d$ in the expression:
$$
L_d(\ee^{\ii\theta}) \;:=\; \sum_{k=-d}^{d} (-\ii)^k\,J_k(\delta\alpha)\,\ee^{\ii k\theta} \;\in\; \mathbb{C}[z,z^{-1}],
$$
that have a tail bound 
$$
\bigl\|\ee^{-\ii\delta\alpha\cos\theta} - L_d\bigr\|_{\infty,\mathbb{T}} \;\leq\; 2\!\!\sum_{|k|>d}\!|J_k(\delta\alpha)| \;\leq\; \mathcal{O}\!\left(\!\left(\frac{\ee\delta\alpha}{2(d+1)}\right)^{d+1}\right),
$$
justifying the use of a truncation. The trick to efficiently realise $L_d$, is to first introduce an underlying polynomial $P$ as, 
$$
P(z) \;:=\; z^{d}\,L_{d}(z) \;=\; \sum_{m=0}^{2d} c_m\,z^m, \qquad c_m \;=\; (-\ii)^{m-d}\,J_{m-d}(\delta\alpha),
$$
of degree $2d$. Since $|z^d|=1$ on $\mathbb{T}$, $\|P\|_{\infty,\mathbb{T}} = \|L_d\|_{\infty,\mathbb{T}} \leq 1 + \mathcal{O}((\ee\delta\alpha/(2(d+1)))^{d+1})$, and a rescaling $(1-\varepsilon_\text{QSP})P$ enforces \eqref{eq:qsp-poly}. Applying \cite[Thm. 6]{motlagh2024generalized} with $\deg P = 2d$ and $k = d$ produces
$$
P'(z) \;=\; z^{-d}\,P(z) \;=\; L_d(z),
$$
so the post-selected top-left block directly realises the Jacobi–Anger Laurent truncation, giving $L_d(W_\lambda) = \ee^{-\ii\delta\alpha\cos\theta_\lambda}\,I_{2\times 2} + \mathcal{O}(\varepsilon_\text{QSP}) = \ee^{-\ii\delta\lambda}\,I_{2\times 2} + \mathcal{O}(\varepsilon_\text{QSP})$ on each qubitization subspace, and hence $\ee^{-\ii\delta A} + \mathcal{O}(\varepsilon_\text{QSP})$ on the system register.

The truncation index $d$ that suffices for $\varepsilon_\mathrm{QSP}$-accuracy on $\mathbb{T}$ is \cite[Thm. 7]{motlagh2024generalized}, 
\begin{equation} \label{eq:qsp-degree}
    d \;=\; \mathcal{O}\!\left(\delta\alpha \;+\; \frac{\log(1/\varepsilon_\text{QSP})}{\log\log(1/\varepsilon_\text{QSP})}\right).
\end{equation}
Consequently, one GQSP layer applies $d$ controlled-$W$ calls, $d$ controlled-$W^\dagger$ calls, together with $2d + 1$ single-qubit rotations on the QSP ancilla. In the regime $\delta\alpha \lesssim 1$, relevant for our coherent step, already $d=1$ suffices:
$$
L_1(\ee^{\ii\theta}) \;=\; J_0(\delta\alpha) \,-\, 2\ii\,J_1(\delta\alpha)\,\cos\theta.
$$
This can be realised by one controlled-$W$, one controlled-$W^\dagger$, and three rotations $(R_0, R_1, R_2)$. Inserting $J_0(z) = 1 + \mathcal{O}(z^2)$ and $2J_1(z) = z + \mathcal{O}(z^3)$ yields
$$
L_1(\ee^{\ii\theta}) \;=\; 1 \,-\, \ii\,\delta\alpha\,\cos\theta \,+\, \mathcal{O}((\delta\alpha)^2) \;=\; \ee^{-\ii\delta\alpha\cos\theta} \,+\, \mathcal{O}((\delta\alpha)^2).
$$
Larger subnormalization $\alpha$ or a larger time step $\delta$ at fixed accuracy could push $d \geq 2$ \eqref{eq:qsp-degree}, and the cost grows linearly with $\delta \alpha$. 

Given the polynomial coefficients $\{c_m\}_{m=0}^{2d}$, the angles $\{(\theta_0,\phi_0,\lambda_0)\}\cup\{(\theta_j,\phi_j)\}_{j=1}^{2d}$ are classically precomputed for the quantum circuit. First step is to construct a complementary polynomial $Q\in\mathbb{C}[z]$ with $\deg Q = 2d$ and on $\mathbb{T}$ satisfying
$$
|P(z)|^2 + |Q(z)|^2 \;=\; 1.
$$
The existence of $Q$ for a given $P$ is proven by \cite[Thm. 4]{motlagh2024generalized}. But we deviate from their proposed way, and use the contour-integral / Fast Fourier Transform construction of \cite{berntson2025complementary}. They do not use optimization to solve the problem, moreover, they can evaluate $Q$ with explicit error bounds up front, and work in classical time $\tilde{O}(d)$. Once we have $Q$, using the recursive extraction of \cite[Alg. 1]{motlagh2024generalized} results in the $2d + 1$ rotation triplets from $(P, Q)$ in $\mathcal{O}(d^2)$ arithmetic operations. Both stages are performed once, before any quantum computation, for each value $\delta \alpha$. The same angles are then reused at every quantum call of the layer. 

\smallskip
\section{Weak measurement} \label{sec:prelim-weak-meas}

\begin{figure}[H]
  \centering
  \begin{quantikz}[row sep=0.85cm, column sep=0.55cm]
    \lstick{$\ket{0}_\delta$}
      & \qw
      & \qw
      & \qw
      & \gate{Y_\delta}
      & \octrl{1}
      & \meter{}
      \\
    \lstick{$\ket{0^b}$}
      & \qwbundle{}
      & \gate[2]{U_L}
      & \qw
      & \octrl{-1}
      & \gate[2]{U_L^{\dagger}}
      & \rstick{$\bra{0^b}$} \qw
      \\
    \lstick{$\rho$}
      & \qwbundle{}
      &
      & \qw
      & \qw
      &
      & \rstick{$\approx\ \rho + \tfrac{\delta}{Z_L^{2}}\,\mathcal{D}_L(\rho) + \mathcal{O}(\delta^{2})$} \qw
  \end{quantikz}
  \caption{Weak measurement as a Lindblad dissipator, with
  $\bra{0^b} U_L \ket{0^b} = L/Z_L$ a block encoding of the jump
  $L$. The $Y_\delta$ rotation is controlled on the block register
  $\ket{0^b}$, so that it fires only on the post-selectable branch of
  $U_L$. Post-selecting $b$ on $\ket{0^b}$ and measuring $q_{\delta}$
  realises one Euler step of $\mathcal{D}_L(\rho) = L\rho L^{\dagger}
  - \tfrac{1}{2}\{L^{\dagger}L,\rho\}$ at rate
  $\gamma = \delta/(Z_L^{2}\,\dd t)$.}
  \label{circ:weak-measurement}
\end{figure}

There is no one-size-fits-all quantum simulation of Lindbladian evolutions. Based on the Lindbladian's specifics, one method could work while another one not. The \emph{weak measurement} scheme we introduce here, established by \cite[Thm. III.1]{chen2023quantum}, fits indeed best to the quantum algorithm in the centre of this thesis, \hyperref[alg:main]{Alg.\ref*{alg:main}}.
A Lindbladian dissipator for a jump operator $L$ is given by $\mathcal{D}_L(\rho) := L\rho L^\dagger - \tfrac{1}{2}\{L^\dagger L,\rho\}$. It is a combination of a transition term and a reflection term, that together leads to a trace-preserving map $\rho \mapsto \ee^{t\mathcal{D}_L}\rho$. A unitary block-encoding $U_L$ like the LCU \hyperref[circ:lcu]{Fig.~\ref*{circ:lcu}} could already act as the transition term $L\rho L^\dagger$: post-select on the block-encoding ancillas as $|0^b\rangle$ to have $U_L$ imprint $L\ket{\psi} / Z_L$ on the system (with $\|L/Z_L\|\leq 1$). Clearly, the same technique could not deliver the remaining reflection part of the Lindbladian. Chen et al. solves this in \cite{chen2023quantum} by weaving a weak measurement ancilla $\ket{0}_\delta$ between $U_L$ and a conditional $U_L^\dagger$ in such a way that post-selecting on all ancillas to be zero implements an Euler step of $\ee^{\delta\mathcal{D}_L/Z_L^2}$ on $\rho$ up to $\mathcal{O}(\delta^2)$ (\hyperref[circ:weak-measurement]{Fig.~\ref*{circ:weak-measurement}}).

The circuit walkthrough after initialising the three registers in $\ket{0}_\delta \otimes \ket{0^b} \otimes \rho$ is the following:
\begin{enumerate}[(i)]
    \item \emph{Block encoding.} $U_L$ acts on $(b, S)$ s.t. on the $\ket{0^b}$ block of the state imprints the amplitude $L/Z_L$ on $S$, leaving an orthogonal artefact on $(b, S)$:
    $$
    |\phi^\perp\rangle := \left(\idm - |0^b\rangle\!\langle0^b|\!\otimes\idm\right)\,U_L\,|0^b\rangle\!\ket\psi.
    $$
    \item \emph{Ancilla rotation.} $Y_\delta = R_Y(2\arcsin\sqrt\delta)$ is controlled on the block encoding being successful, i.e. $\ket{0^b}$, and leads to
    \begin{equation}
        \begin{split}
            Y_\delta\,\left(\ket{0}_\delta\otimes|0^b\rangle\right) &\;=\; \left(\sqrt{1-\delta}\,\ket{0}_\delta + \sqrt{\delta}\,\ket{1}_\delta\right)\otimes |0^b\rangle, \\
            Y_\delta\,\left(\ket{0}_\delta\otimes |\phi^\perp\rangle\right) &\;=\; \ket{0}_\delta\otimes|\phi^\perp\rangle.
        \end{split}
    \end{equation}
    In other words, the weak measurement ancilla acquires amplitude $\sqrt{\delta}$ on $\ket{1}_\delta$ precisely where $U_L$ succeeded, and carries $\ket{0}_\delta$ unchanged on the leaked $\ket{\phi^\perp}$ branch.
    \item \emph{Controlled undo.} The open-controlled $U_L^\dagger$ fires only on the $\ket{0}_\delta$ branch. On the $\ket{0^b}$-block of $U_L\ket{0^b}\ket\psi$ this uncomputes to $\ket{0^b}\ket\psi$. The artefact part $\ket{\phi^\perp}$ survives without any change to it. On the $\ket{1}_\delta$ branch $U_L^\dagger$ is absent, thus $\ket{0^b}\otimes(L/Z_L)\ket\psi$ stands.
    \item \emph{Measure and post-select.} Finally, measure all ancillas, and post-select only on the block encoding ones as $\ket{0^b}$.
\end{enumerate}

We can understand what is happening in this construction by Kraus operators \eqref{eq:kraus}: in the \emph{no-jump} branch  ($\ket{0}_\delta$ outcome, $\ket{0^b}$ post-selected), the uncomputation $U_L^\dagger U_L = \idm$ acts on the block, while the $\ket{\phi^\perp}$ artefact has $\ket{0^b}$ post-selection failure probability $\|L\ket\psi\|^2/Z_L^2$. In the \emph{jump} branch ($\ket{1}_\delta$ outcome, $\ket{0^b}$ post-selected) the full amplitude $L\ket\psi/Z_L$ survives, weighted by the $\sqrt{\delta}$ of the weak rotation. The results for a given Lindblad operator $L$ derived by \cite{chen2023quantum} can be summarized as the Kraus jumps
\begin{equation}
    M_0 \;=\; \sqrt{1-\delta}\,\idm \;-\; \tfrac{1}{2}\,\tfrac{\delta}{Z_L^2}\,L^\dagger L \;+\; \mathcal{O}(\delta^2), \qquad M_1 \;=\; \sqrt{\delta}\;\frac{L}{Z_L}.
\end{equation}
By tracing out the ancillas (summing over the branches) and expanding in $\delta$,
\begin{equation} \label{eq:weak-meas-euler}
    \rho \;\mapsto\; M_0\,\rho\,M_0^\dagger + M_1\,\rho\,M_1^\dagger \;=\; \rho \;+\; \frac{\delta}{Z_L^{2}}\Bigl(L\rho L^\dagger - \tfrac{1}{2}\{L^\dagger L,\rho\}\Bigr) \;+\; \mathcal{O}(\delta^{2}).
\end{equation}
which is precisely one Euler step of $\dot\rho = (1/Z_L^2)\,\mathcal{D}_L(\rho)$ up to $\mathcal{O}(\delta^2)$ errors. The post-selection on $b$ succeeds with probability $1 - \delta\,\|L\ket\psi\|^2/Z_L^2 + \mathcal{O}(\delta^2) = 1 - \mathcal{O}(\delta)$, i.e. unlike the post-selection loss ($1/Z_L^2$) of LCU \eqref{eq:lcu}, here there is no need for oblivious amplitude amplificiation or any other trick to boost the success probability. The smallness of $\delta$ is enough to make this into a deterministic mechanism. Or put differently, even the artefact blocks contribute to the Lindbladian correctly -- as the reflection part -- up to $\mathcal{O}(\delta^2)$ errors. The extension to multiple jumps is straight-forward, which is precisely what we can see in \cite{chen2023quantum,chen2023efficient}, and also in our slightly modified variant \hyperref[alg:diss]{Alg.~\ref*{alg:diss}}.

The Euler-step accuracy of \eqref{eq:weak-meas-euler} is only first order in $\delta$ so reaching a diamond-norm error $\varepsilon$ over the total time $T$ costs $T/\delta = \mathcal{O}(T^2 / \varepsilon)$ calls to $U_L$, i.e. polynomial in $1 / \varepsilon$. One might therefore ask why we do not instead pick a higher-order Lindblad simulator, such as the one proposed by \cite{li2022simulating} that uses Duhamel's principle. This could simulate a generator $\mathcal{L}(\rho) = -\ii[H,\rho] + \sum_{j=1}^m \mathcal{D}_{L_j}(\rho)$ with $\mathcal{O}(\tau\,\mathrm{polylog}(t/\varepsilon))$ block-encoding queries to $U_H$ and to $U_{L_j}$ and $\mathcal{O}(m\tau \mathrm{polylog}(t / \varepsilon))$ additional 1- and 2-qubit gates, where $\tau := t (\alpha_0 + \tfrac{1}{2}\sum_{j=1}^m \alpha_j^2)$. Accordingly, that higher-order simulators are meant to take in a \emph{finite} list $\{L_j\}_{j=1}^m$ of jump operators, each separately block encoded by $U_{L_j}$, leading to a cost linear in $m$. This breaks down in case the generator is of the form \eqref{eq:ckg-L}
$$
\mathcal{L}(\rho) \;=\;\int \dd\omega\;\gamma(\omega)\,\mathcal{D}_{L_j(\omega)}(\rho),
$$
with the continuously parametrized jump operators $L_j(\omega)$, that is used later to estimate the Bohr-frequencies jump components for the detailed balanced Davies generator \eqref{eq:davies-generator}. Even after a quadrature discretization of $\int \dd \omega$ to $\omega_0\sum_{\bar{\omega} \in \Omega}$, $m$ is at least the size of the frequency grid $\Omega$, and each grid point would need its own block encoding rather than being addressed coherently through the $\Omega$ register of the phase estimation \hyperref[circ:qpe]{Fig.~\ref*{circ:qpe}}. The weak measurement primitive is exactly what one uses when the integral cannot be separated from the block encodings. A single $U_L$ call already block-encodes the full sum of jumps $\{L_j(\omega)\}_{\omega\in\Omega}$. Then, one Euler step of \eqref{eq:weak-meas-euler} realises one step of the \emph{whole} Lindbladian per call, with cost that scales with the resolution of the phase estimations rather than with the number of effective jumps.

Undeniably, there is a structural difference between the two Lindbladian simulations. The higher-order series scheme buys $\mathrm{polylog}(1/\varepsilon)$ accurate channel implementation via LCUs, but only when the jump set is finite, and preferably not a large set. In this sense, the continuously parametrized jumps $\{L_j(\omega)\}_{\omega\in\Omega}$ simply does not compress into this scheme. An alternative KMS detialed balanced Lindbladian construction \cite{ding2025efficient} that came after the Chen et al. results, manages to confine the jumps to a finite set. While that allows them to utilize the higher-order simulators of \cite{li2022simulating}, their construction introduces more cost at other parts of the algorithm, making it less clear which solution is better or worse in general. We discuss this in more details in [REF] \todo{IF you write the comparison section, ref to it.}



\part{Gibbs Sampling on a Quantum Computer}
%\chapter{Quantum Metropolis Sampling}
%We have now come to the point where we can turn our attention to recent developments in quantum Gibbs sampling. The moment we mentioned Gibbs sampling in the classical setting, we also learnt the recipe that researchers used for the past decades in all kinds of scenarios. So much so, that the \hyperref[parag:MH]{Metropolis-Hastings algorithm} is now a ubiquitous word of the sampling jargon, and it became one of the most important algorithms of the 20th century \cite{dongarra2000guest}. 
%While we have no doubt the classical Metropolis algorithm will continue to advance and have even more successes, it naturally has scaling limitations, just as most algorithms do. The dimensional explosion of the state space in quantum systems and similarly in large classical systems can pose a problem for even the state-of-the-art implementations of the algorithm. \textcolor{purple}{[cite]} Now that quantum computers have been on the horizon for some years, it could be beneficial to explore its capabilities for (Gibbs) sampling. Then one of the natural questions to ask is,
%$$\textit{Can we generalize the Metropolis algorithm to the quantum setting?}$$
%The answer has only been recently given in a foundational work \cite{temme2011quantum}, that established the notion of the \textit{quantum Metropolis algorithm.} 

%If we recall the steps of the classical algorithm, it becomes quickly evident why the quantum generalization kept the community waiting. We start in some initial state and then propose a new state: if the energy of the new configuration is lower accept it, otherwise reject it with probability given by the (local) Metropolis rule, $\exp(\beta(E_\text{old} - E_\text{new}))$. Classically, we can compute the energy of the new state, and if it has to be rejected, then we throw it away and use a copy of the previous state. In quantum mechanics however, we cannot copy, or \textit{clone} states \cite{wootters1982single}. This could pose a problem, because some operations in a quantum computer, like the energy measurement of states, are irreversible. One way to get around this, is to disturb the system as weakly as possible when we measure, such that the rejected state is still in some relation with the initial state. If some overlap between them survives, then using a Marriott-Watrous amplification scheme \cite{marriott2005quantum} could rewind back iteratively to the initial state. Crucially, this procedure together with the rest of the quantum subroutines, cannot break the detailed balanced condition, or else there is no guarantee for sampling from the Gibbs distribution. 

\chapter{Dissipative Quantum Gibbs Sampling}

In the previous chapter we have seen and analysed the very first breakthrough algorithm in quantum Gibbs sampling. The quantum Metropolis algorithm has kicked off a myriad of research, trying to better understand it, use it as a subroutine for other problems \todo{cites}, and also to possibly circumvent its main drawbacks, namely, the costly rewinds and the lack of exact detailed balance. There is no clear-cut, perfect solution to either of these problems. 

One way to get rid of the rewinds, is to work with a continuous-time Markov sampling instead of a discrete one. In this case, we never have to worry about a subpar jump and its rewind, but rather worry about the rate of the jumps. As we will see in this chapter, we are constrained to use low rates for the jumps, or else the approximations would get out of hand. Nevertheless, this might be an improvement to the Metropolis algorithm, since the rewind itself required a weak measurement, i.e. a weakly disturbed system, as well. 

Concerning the question of detailed balance, it must be stated that it is not a priori clear how and if at all $exact$ detailed balance is necessary for Gibbs sampling \textit{in practice}. Of course, if we have detailed balance we can work with well established techniques for reversible Markov semigroups; it gives guarantees that the sampling eventually reach the Gibbs state and thus helps with rigorous mixing time bound proofs. However, implementing an exact detailed balanced algorithm on a computer could easily break that nice, yet fragile symmetry. Even the algorithm we will introduce here, will break its proven detailed balance the moment we put it in quantum circuit form. The hope is that once there is an exact detailed balance framework established, it can be used as a blueprint for follow-up algorithms that break the symmetry in a controlled and rigorous way. 

The goal for this chapter is to (i) introduce a recently proposed approximate and exact detailed balanced Lindbladians and their implementations from Chen et al. \cite{chen2023quantum,chen2023efficient}; (ii) Analyse it numerically to understand their $real$ complexity, e.g. mixing times for given systems, gate complexities, temperature dependencies; (iii) Work out how best can we implement them on a quantum device.

\smallskip
\subsection{Approximate GNS Lindbladian.} After learning that the Davies generator has a proven GNS detailed balance in \hyperref[subsec:davies]{Section \ref*{subsec:davies}}, it is clear that the crux is not only about having detailed balance, but also about having an efficient implementation for it. The main obstacle with the Davies generator was that it required perfect energy estimates for the jumps. For small systems this might even be possible, and we could distinguish each and every Bohr frequency component of a given jump, then weigh each of them  with the prescribed transition rates (that satisfy KMS condition). But as we increase the system size and consider more interesting systems, the Bohr frequencies get exponentially close to each other. This poses an inherent metrology problem for implementations, because of the energy-time uncertainty principle. Which means that the more precise an energy value needs to be estimated, the longer Hamiltonian time evolution is necessary to do in the quantum computer (e.g. with QPE), thus leading to exponential costs.

A natural question to ask is, what happens with detailed balance if we allow a controlled error around the Bohr frequencies we need to estimate. As it turns out from \cite{chen2023quantum}, these smoothened Davies generators still \textit{approximately} satisfy detailed balance. To see why, let us start from the Davies generator, which fulfils exact GNS detailed balance \eqref{eq:davies-gns-db}. There, it was useful to work with the Heisenberg evolution of the jumps $A$:
$$
    A(t) = \ee^{\ii Ht} A \ee^{-\ii Ht} = \sum_{\nu\in B_H} A_\nu \ee^{\ii \nu t},
$$
where the Bohr frequencies $\nu\in B_H$ are all the energy differences in the Hamiltonian's spectrum, $\text{Spec}(H)$. The key step introduced in the Chen et al. papers is to use a filter function $f: \mathbb{R}\rightarrow \mathbb{C}$ and define the \textit{Operator Fourier Transform} (OFT) of a jump $A$ as
\begin{equation}
    \label{eq:oft}
    \hat{A}(\omega) = \frac{1}{\sqrt{2\pi}} \int_{-\infty}^{\infty} f(t) A(t) \ee^{-\ii \omega t} \dd t = \sum_{\nu \in B_H} A_\nu \hat{f}(\omega - \nu),
\end{equation}
with the Fourier transform of $f(t)$,
$$
    \hat{f}(\omega) = \frac{1}{\sqrt{2\pi}} \int_{-\infty}^{\infty} f(t) \ee^{-\ii \omega t} \dd t.
$$
As a reminder: in a quantum computer we always have to work in the time domain, and the Fourier transform is to the tool that brings us to the energy domain. In order to implement the Davies generator, one would need a Fourier transform that maps exactly to $A_\nu$. That can be also understood as an OFT, but with $f(t) = 1$ and thus $\hat{f}(\omega) = \sqrt{2\pi} \delta(\omega)$:
\begin{equation}
    \hat{A}_{f=1}(\omega) = \sqrt{2\pi}\sum_{\nu \in B_H} A_\nu \delta(\omega - \nu).
\end{equation}
This can then only be exactly achieved if the Fourier transform is computed over the full, infinitely large time domain. Hence, the idea here is to choose an $f(t)$ such that it has some decaying tail that effectively truncates the Fourier transforms to some more manageable subdomains. From \eqref{eq:oft} we can see the effect of the non-trivial filter function: the OFT leads to $\hat{A}(\omega)$ instead of $A_\nu$, which includes neighbouring energy contributions centred around $A_\nu$. What happens with the Davies generator and its GNS detailed balance if we swap out the jump operators, to their smoothened out form? In \cite{chen2023quantum} it was shown, that it leads to an approximately GNS detailed balance Lindbladian. 

There is some freedom in how we choose our filter function, but it necessarily has to be well-behaving $L^1$ integrable distribution. Clearly, for a not too costly implementation this function $f(t)$ has to decay rapidly in the time domain as $|t|\rightarrow \infty$, and its Fourier integral needs to be efficiently amenable to quadrature techniques. Moreover, its Fourier transform $\hat{f}(\omega)$ has to be smooth too.

One possible choice that fulfils all the above requirements is, to use a Gaussian filter function,
\begin{equation}
    \label{eq:gaussian-filter-t}
    f(t) := \left(\frac{2 \sigma^2}{\pi}\right)^{\frac{1}{4}}  \exp(-\sigma^2 t^2) \quad\text{with}\quad \int_{-\infty}^{\infty} |f(t)|^2 \dd t = 1.
\end{equation}
The width of the Gaussian $1 / (2 \sigma)$ is a tuneable parameter, that determines how strongly we smoothen the jumps and thus the generator, and it is directly affecting the cost of the quantum algorithm. Gaussians have the nice property, that their Fourier transforms are also Gaussians, but with inverse width $2\sigma$,
\begin{equation} \label{eq:gaussian-filter-w}
    \hat{f}(\omega) = \left(2\pi \sigma^2 \right)^{-\frac{1}{4}} \exp\left(-\frac{\omega^2}{4\sigma^2}\right).
\end{equation}
Then in this case, the smooth jump operators that the Fourier transforms are mapping to, can be understood as Gaussian distributions centred around each Bohr frequency $\nu$ with an uncertainty of $2\sigma$,
\begin{equation}
    \label{eq:smoothed-jump}
    \hat{A}(\omega) = \left(2\pi \sigma^2 \right)^{-\frac{1}{4}} \sum_{\nu \in B_H} \exp\left(- \frac{(\omega - \nu)^2}{4\sigma^2}\right)A_\nu.
\end{equation}
\todo{figure for Gaussian smoothing around $\nu$}
Note, that a Gaussian filter is not the only function that could be used here. Gevrey functions also fulfil the necessary requirements and have been successfully used for KMS detailed balanced Lindbladians in \cite{ding2025efficient} that we will look at later. Until then let us keep ourselves to Gaussians. 

The Davies generator with the smoothened jump operators \eqref{eq:smoothed-jump} we will call after the authors, \textit{Chen-Kastoryano-Brandão-Gilyén} or \textit{CKBG Lindbladian}, and it takes the following form in the Schrödinger picture,
\begin{equation}
    \label{eq:L-ckbg}
    \mathcal{L}_\text{CKBG}(\cdot) = \sum_{a} \int_{-\infty}^{\infty} \gamma(\omega)\left(\hat{A}^a(\omega)(\cdot)\hat{A}(\omega)^\dagger -\frac{1}{2}\left\{\hat{A}^a(\omega)^\dagger \hat{A}^a(\omega), \cdot \right\}\right)\dd\omega .
\end{equation}

The sum runs over the set of jumps, $\{A^a\}_{a\in\mathcal{A}}$, that can be chosen in any way that ensures ergodicity. Usually, with a varied enough set (breaking any symmetries that the Hilbert space might have) where each jump also has its adjoint included works well,
\begin{equation}
    \label{eq:ergodic-set-jumps}
    \{A^a\}_{a\in\mathcal{A}} = \{(A^a)^\dagger\}_{a\in\mathcal{A}} \quad\text{and}\quad \left\|\sum_a A^{a\dagger} A^a\right\| \leq 1.
\end{equation}
Here, the norm will be used to normalize the jumps and its purpose is (i) to make block encodings (i.e. implementation) possible \todo{cite to block encoding} and (ii) it also helps in analysis since it fixes a variable, namely the "strength" of the Lindbladian for different number of jump operators.
The naive set of single-site Pauli jumps work already quite well, and since they are unitary we can avoid using some block encodings and excess ancillas in the circuits. (In this case, the normalization factor in \eqref{eq:ergodic-set-jumps} becomes the total number of jumps $|\mathcal{A}|$.) Even if we choose a single-site jump as the core, the Kraus operators of the Lindbladian, $\{L_j\}_{j\in\mathcal{J}} = \{\sqrt{\gamma(\omega)}\hat{A}(\omega)\}$ will not be constraint to just that one site. We can see this from the OFT construction \eqref{eq:oft}, where it becomes clear, that as the Pauli operators evolve over time in the Heisenberg picture, its support will also spread out spatially. We can estimate this size by the Lieb-Robinson bound to see that for the small system sizes that we can numerically simulate, these jumps spread out over the whole lattice. See the Appendix for more details [APPENDIX] \todo{appendix LR bound}.

The transition rates $\gamma(\omega)$ can be any bounded function, $\gamma(\omega) \in [0, 1]$, that satisfy the KMS condition \eqref{eq:KMS-condition}, which now has the form
\begin{equation}
    \label{eq:KMS-condition-2}
    \frac{\gamma(\omega)}{\gamma(-\omega)} = \ee^{-\beta \omega}.
\end{equation}
Let us first choose a tuneable Gaussian for these transition rates,
\begin{equation} \label{eq:gaussian-transition}
    \gamma(\omega) = \exp\left( -\frac{(\omega + \omega_\gamma)^2}{2\sigma_\gamma^2} \right),
\end{equation}
where the centre of the Gaussian is at $\omega_\gamma$ and its width is given by $\sqrt{2}\sigma_\gamma.$ In other words, most of the jumps are suppressed in this case and only the ones that lead to an energy difference with $\omega \approx \omega_\gamma$ contribute meaningfully. The KMS condition \eqref{eq:KMS-condition-2} for this Gaussian $\gamma(\omega)$ actually fixes the following relation between the parameters $\beta, \omega_\gamma, \sigma_\gamma$
\begin{equation} \label{eq:beta-approx-gns-gaussian}
    \beta := \frac{2\omega_\gamma}{\sigma_\gamma^2}
\end{equation}

In [FIGURE] \todo{add figure of gammas} you can see how constraining it is compared to transitions that we will consider later on. But before we could do that, we need to see what kind of detailed balance, if any, we have with these Gaussian smoothed jump operators \eqref{eq:oft}. 

In its current form \eqref{eq:L-ckbg} has been shown to satisfy \textit{approximate GNS detailed balance} in \cite{chen2023quantum}. Quantitatively, we have understood the severity of breaking detailed balance as the norm of the anti-Hermitian part $\|\mathcal{A}(\rho, \mathcal{L})\|_{2\rightarrow 2}$ of the generator's discriminant $\mathcal{D}(\rho, \mathcal{L})$, see \hyperref[def:approx-db]{Definition \ref*{def:approx-db}}. Thus, if we can choose algorithm parameters such that the norm of the anti-Hermitian part is below some $\varepsilon$, we can also bound how well does the fixed point $\rho_\text{fix}(\mathcal{L})$ approximate the target state $\rho$.

\begin{corollary}
    \textup{(Fixed point accuracy \cite{chen2023quantum})} If a Lindbladian $\mathcal{L}$ satisfies $\varepsilon$-approximate $\rho$-detailed balance condition, then its fixed point $\rho_\text{fix}(\mathcal{L})$ deviates from $\rho$ by at most $$\|\rho_\text{fix}(\mathcal{L}) - \rho\|_1 \leq 20 t_\text{mix}(\mathcal{L})\,\varepsilon .$$
\end{corollary}
Chen et al. provides the proof of this, which is indeed less than obvious. The core thought is that perturbations to the parent Hamiltonian $\mathcal{H}$ with $\mathcal{A}$, leads to only small changes for the eigenvectors with separated eigenvalues. In fact, one can show that the top eigenvalue of $\mathcal{H}$ will not deviate further from $0$ (its detailed balanced value) than $$|\lambda_1(\mathcal{H})| \leq \|\mathcal{A}\|_{2\rightarrow 2}.$$ We can combine all of this with the CKBG Lindbladian and find that it is approximately GNS detailed balanced.
\begin{theorem}
    \textup{(Approximate detailed balance for $\mathcal{L}_\text{CKBG}$  \cite{chen2023quantum})}. Any Lindbladian $\mathcal{L}_\textup{CKBG}$ of the form \eqref{eq:L-ckbg} with jumps \eqref{eq:ergodic-set-jumps}, satisfying the KMS condition \eqref{eq:KMS-condition-2}, and having Gaussian filter functions \eqref{eq:gaussian-filter-t} has an approximate detailed balance. Hence, its fixed point $\rho_\textup{fix}(\mathcal{L_\textup{CKBG}})$ is approximately the Gibbs state $\rho_\beta$
    $$
    \left\|\rho_\textup{fix}(\mathcal{L}_\textup{CKBG}) - \rho_\beta\right\|_1 = \mathcal{O}\left(\sigma \beta \sqrt{\log\frac{1}{\sigma \beta}}\,t_\textup{mix}(\mathcal{L}_\textup{CKBG})\right).
    $$
\end{theorem}
The main takeaway of these results is that if some error appears on the generator level, it will get amplified with time, and can have significant impact on Gibbs sampling results.
This is even worse for cases (usually the interesting ones) where fast or rapid mixing has not been proven, and thus can very well make $t_\text{mix}(\mathcal{L})$ exponentially scaling with system sizes and with inverse temperature $\beta$.

One of the main culprits leading to deviations from the Gibbs state is the presence of non-secular terms in the CKBG Lindbladian \eqref{eq:L-ckbg}. For the \hyperref[subsec:davies]{Davies generator} we had to apply the secular approximation in order to get a correct, physical generator. But here, with the smoothened jump operators \eqref{eq:oft} we have reintroduced some off-diagonal entries in the Kossakowski matrix. To better see this, let us derive its analytical form by computing the integral over the energies in \eqref{eq:L-ckbg} for the transition part of the Lindbladian.
\begin{equation} \label{eq:gauss-transition-part-of-lindbladian}
    \begin{split}
        \mathcal{T} = &\int_{-\infty}^\infty \gamma(\omega) \hat{A}(\omega) (\cdot) \hat{A}(\omega)^\dagger \dd \omega \\
        &= \frac{1}{\sigma \sqrt{2\pi}} \sum_{\nu_1, \nu_2 \in B_H} \int_{-\infty}^{\infty} \ee^{-\frac{(\omega + \omega_\gamma)^2}{2\sigma_\gamma^2}}\ee^{-\frac{(\omega - \nu_1)^2}{4\sigma^2}}\ee^{-\frac{(\omega - \nu_2)^2}{4\sigma^2}}A_{\nu_1}(\cdot)A_{\nu_2}^\dagger \dd \omega \\
        &=: \sum_{\nu_1, \nu_2 \in B_H} \alpha_{\nu_1, \nu_2} A_{\nu_1}(\cdot)A_{\nu_2}^\dagger \dd \omega
    \end{split}
\end{equation}
with the Kossakowski matrix,
\begin{equation} \label{eq:kossakowski-gaussian}
    \alpha_{\nu_1,\nu_2} = \frac{\sigma_\gamma}{\sqrt{\sigma^2 + \sigma_\gamma^2}} \exp\left(-\frac{(\nu_1 + \nu_2 + 2\omega_\gamma)^2}{8(\sigma^2 + \sigma_\gamma^2)}\right)\exp\left(-\frac{(\nu_1 - \nu_2)^2}{8\sigma^2}\right).
\end{equation}
There are a few things we can read off from this expression: (i) There are non-secular terms present $\nu_1 \neq \nu_2$ due to the Gaussian smearing around the $\nu$'s; (ii) It is a positive semi-definite matrix (and Hermitian), which makes the generator $\mathcal{L}$ to be a Lindbladian generator; (iii) 
The Gaussian filters are not only breaking the secular structure but also raising the temperature. To see this latter point, let us reconsider what we are aiming to approximate: the GNS detailed balance that was physically derived by retaining diagonal terms, and drops all off-diagonal ones via the secular approximation. Setting $\nu_1 = \nu_2 = \nu$ in \eqref{eq:kossakowski-gaussian} yields, 
\begin{equation}
    \alpha_{\nu, \nu} = \frac{\sigma_\gamma}{\sqrt{\sigma^2 + \sigma_\gamma^2}} \exp\left(-\frac{(\nu + \omega_\gamma)^2}{2(\sigma^2 + \sigma_\gamma^2)}\right).
\end{equation}
These diagonal entries have to satisfy the KMS condition,
\begin{equation} \label{eq:kossakowski-diag-condition}
    \frac{\alpha_{\nu, \nu}}{\alpha_{-\nu, -\nu}} = \exp\left(-\frac{2\nu\omega_\gamma}{\sigma^2 + \sigma_\gamma^2}\right) =: e^{-\beta_\text{eff} \nu}
\end{equation}
with 
\begin{equation} \label{eq:beta-eff}
    \beta_\text{eff} = \frac{2 \omega_\gamma}{\sigma^2 + \sigma_\gamma^2} < \beta = \frac{2\omega_\gamma}{\sigma_\gamma^2}.
\end{equation}
In other words, Gibbs sampling with the $\mathcal{L}_\text{CKBG}$ from \eqref{eq:L-ckbg} and with $\sigma > 0$, effectively leads to some higher temperature Gibbs state. Though because of the additional non-secular errors, the error structure is not as clear-cut as just "raised temperatures". It can now also be clearly seen, that by taking the $\sigma \rightarrow 0$ limit, all errors to the GNS detailed balance vanish as well: $\beta_\text{eff} \rightarrow \beta$ and as the Gaussian filters become delta peaks in this limit, they also eliminate any off-diagonal elements and only $\nu_1 = \nu_2$ remains. 

\label{sec:kms-lindbladian}
\smallskip
\subsection{KMS Lindbladian} While the previous approximate GNS detailed balanced Lindbladian was a result on its own, its improvement is what got researchers (with myself included) thrilled about quantum Gibbs sampling. In \cite{chen2023efficient} Chen et al. have managed to come up with a structure for the Lindbladian where they do not fight the above-mentioned errors, but rather lean into them and fix them all at once with an additional coherent term. The result is an exact KMS detailed balanced Lindbladian instead of an approximation of the GNS detailed balance and the Davies generator with high costs. In order to better see how terms in the Lindbladian interact with each other to get KMS detailed balance, it will be useful to write the Lindbladian in this general form:
\begin{equation}
    \mathcal{L} = \mathcal{T} - \frac{1}{2}\mathcal{R} - \ii \mathcal{B}
\end{equation}
with the transition $\mathcal{T}$, decay $\mathcal{R}$ and coherent $\mathcal{B}$ parts, 
\begin{alignat}{2} \label{eq:ckg-L-parts}
\mathcal{T}(\cdot) 
  &= \sum_a\sum_{\nu_1,\nu_2} \alpha_{\nu_1,\nu_2} A^a_{\nu_1}(\cdot)(A^a_{\nu_2})^\dagger
  &\quad&\text{with}\quad
  \mathcal{T}^\dagger(\cdot) = \sum_a\sum_{\nu_1,\nu_2} \bar\alpha_{\nu_1,\nu_2}(A^a_{\nu_1})^\dagger(\cdot)A^a_{\nu_2},
  \\
\mathcal{R}(\cdot) 
  &= \sum_a\sum_{\nu_1,\nu_2} \alpha_{\nu_1,\nu_2}\{(A^a_{\nu_2})^\dagger A^a_{\nu_1},\,\cdot\,\}
  &\quad&\text{with}\quad
  \mathcal{R}^\dagger = \mathcal{R},
  \\
\mathcal{B}(\cdot) 
  &= \sum_a[B_a,\,\cdot\,]
  &\quad&\text{with}\quad
  (-i\mathcal{B})^\dagger = i\mathcal{B}.
\end{alignat}
For $\mathcal{R}^\dagger$ we implicitly took the Kossakowski matrix to be positive semidefinite and thus Hermitian, $\alpha_{\nu_1, \nu_2} = \bar\alpha_{\nu_2, \nu_1}$. Though, this needs to be checked for new constructions, without it, we would not have a GKSL Lindbladian.

Other than the additional coherent term $\mathcal{B}$, the rest of the setup is unchanged: we take the same Gaussian filter functions for the jump operators \eqref{eq:oft}; we keep the Gaussian transition weights $\gamma(\omega)$ \eqref{eq:gaussian-transition}; Consequently, we have the exact same Kossakowski matrix $\alpha_{\nu_1, \nu_2}$ too \eqref{eq:kossakowski-gaussian} with one small but pivotal difference, the temperature. We set the temperature such that the diagonal entries of $\alpha_{\nu_1,\nu_2}$ satisfy the KMS condition, i.e. as $\beta := \beta_\text{eff}$. This instantly makes $\gamma(\omega)$ itself break the original KMS condition, however, the goal here is to satisfy the condition \textit{with} the Gaussian filters together. From the relation \eqref{eq:beta-eff}, it is clear that the Gaussian parameters and the temperature are intertwined. In the KMS case, there is a convenient way to choose them to keep the Kossakowski matrix necessarily skew-symmetric, by 
\begin{equation}
    \sigma = \omega_\gamma = \sigma_\gamma = \frac{1}{\beta}.
\end{equation}
Nevertheless, we will keep the following expressions general, because it is not a priori clear how best to choose these parameters for a faster Gibbs state preparation. It is best to keep in mind, that since the mixing time and these Gaussian parameters are all temperature dependent, they can affect one another. This is one of the questions we have tackled in our numerical analysis, see \todo{ref to results: mixing time dependence on more or less off-diagonal terms.}

Previously the off-diagonal terms $\nu_1 \neq \nu_2$ were either eliminated via secular approximations (Davies) or were the reason for deviation from GNS detailed balance (CKBG Lindbladian). Now these terms are part of the design. But that also means we must have a condition on how these off-diagonal terms relate to each other just like what we had for the diagonal part in \eqref{eq:kossakowski-diag-condition}, in order to attain KMS detailed balanced.
In the following we will see parts of the Lindbladian \eqref{eq:ckg-L-parts} behave differently in this KMS space. In fact, while the transition part $\mathcal{T}$ keeps its structure under Gibbs conjugations, the decay $\mathcal{R}$ and coherent $\mathcal{B}$ parts mix together.

Let us start with the transition part: for it to satisfy the KMS detailed balance condition, it needs to be self-adjoint with respect to the KMS inner product, or equivalently: 
\begin{equation} \label{eq:transition-self-adjoint}
    \mathcal{T}^\dagger = \Gamma_\rho^{-1}\, \circ \, \mathcal{T}\, \circ\, \Gamma_\rho^{1},
\end{equation}
with $\Gamma_\rho(\cdot) := \rho_\beta^{\frac{1}{2}} (\cdot) \rho_\beta^{\frac{1}{2}}$. The equality indeed holds if each jump $A^a$ has its Hermitian conjugate in the jumps set $\mathcal{A}$, and most importantly if the Kossakowski matrix has a specific skew-symmetry:
\begin{equation} \label{eq:skew-symmetry}
    \alpha_{\nu_1, \nu_2} = \alpha_{-\nu_2, -\nu_1} e^{-\beta(\nu_1 + \nu_2) / 2},
\end{equation}
which can even be thought of as the off-diagonal extension of the diagonal condition we found previously \eqref{eq:kossakowski-diag-condition}.
This skew-symmetry has to be satisfied for any given combination of filter function $f(t)$, acceptance rates $\gamma(\omega)$ and $\beta$. For our current Gaussian case \eqref{eq:kossakowski-gaussian}, it was shown that the inverse temperature has to take its already introduced form, $\beta = \beta_\text{eff}$ from \eqref{eq:beta-eff}. 

Now let us look at the remaining parts of the Lindbladian, $\mathcal{R}$ and $\mathcal{B}$. As a quick check, we can conjugate $\mathcal{R}$ with the Gibbs state to see what kinds of terms we get: 
\begin{equation} \label{eq:decay-gibbs-conjugation}
    \Gamma_\rho^{-1}\, \circ \, \mathcal{R}\, \circ\, \Gamma_\rho^{1}(\cdot) =\sum_{\nu_1, \nu_2} \left( e^{\beta(\nu_1 - \nu_2)/2}R_{\nu_1, \nu_2}(\cdot) + (\cdot) R_{\nu_1, \nu_2} e^{-\beta(\nu_1 - \nu_2)/2}\right),
\end{equation}
with $R_{\nu_1, \nu_2} := \sum_a \alpha_{\nu_1, \nu_2} (A^a_{\nu_2})^\dagger A^a_{\nu_1}$. This clearly shows, that the two terms pick up a different factor in the anti-commutator of $\mathcal{R}$ and thus no simple condition on $\alpha_{\nu_1, \nu_2}$ like the skew-symmetry, can lead to KMS detailed balance alone. However, the coherent term $\mathcal{B}$ leads to similar terms under the KMS conjugation as the decay part, 
\begin{equation} \label{eq:coherent-gibbs-conjugation}
    \Gamma_\rho^{-1}\circ\mathcal{B}\circ\Gamma_\rho(\cdot) = \sum_\nu \left(e^{+\beta\nu/2}\,B_\nu\, (\cdot) \;-\; (\cdot)\,B_\nu e^{-\beta\nu/2}\right),
\end{equation}
with the usual Bohr frequency decomposition $B = \sum_\nu B_\nu$. Then, one can dream up a form for $B$ such that all the detailed balance violating terms in the decay part in \eqref{eq:decay-gibbs-conjugation} cancel (while the skew-symmetry \eqref{eq:skew-symmetry} for the Kossakowski matrix holds). Which is exactly what Chen et al. did in \cite{chen2023efficient}.

Then, the key equation that has to hold is
\begin{equation}
    \left(-\frac{1}{2}\mathcal{R} + \ii \mathcal{B}\right) = \Gamma_\rho^{-1}\circ\left(-\frac{1}{2}\mathcal{R} - \ii \mathcal{B}\right)\circ\Gamma_\rho.
\end{equation}
With the Hilbert-Schmidt adjoint on the LHS and the Gibbs conjugation on the RHS. In other words, we are looking for an expression of $\mathcal{B}$ such that $\mathcal{R}$ and $\mathcal{B}$ together are self-adjoint with respect to the KMS inner product. If we write all these terms out using \eqref{eq:decay-gibbs-conjugation} and \eqref{eq:coherent-gibbs-conjugation}, and identify $\nu = \nu_1 - \nu_2$, then we can match the prefactors of the same type of terms on both hands. For instance, the left-multiplication prefactors give: 
\begin{equation}
    -\frac{1}{2}R_{\nu} + \ii B_\nu = -\frac{1}{2} \ee^{\beta\nu / 2}R_{\nu} - \ii \ee^{\beta\nu / 2} B_\nu
\end{equation}
which rearranges to
\begin{equation}
    \ii B_\nu\left(1 + \ee^{\beta\nu / 2}\right) = -\frac{1}{2}\left(\ee^{\beta\nu / 2} - 1\right) R_{\nu}.
\end{equation}
This determines $B_\nu$ uniquely at each Bohr frequency,
\begin{equation}
    B_\nu = \frac{1}{2\ii}\,\frac{1 - \ee^{\beta\nu /2}}{1 + \ee^{\beta\nu /2}}\;R_\nu = \frac{1}{2\ii}\tanh\!\left(-\frac{\beta\nu}{4}\right)\,R_\nu,
\end{equation}
where we used the identity: $\frac{1 - \ee^x}{1 + \ee^x} = \tanh(-x / 2).$ The prefactors for the right-multiplications lead to the same result. Restoring the indices and summing up over them leads to the full coherent term
\begin{equation} \label{eq:B-bohr-domain}
    B = \frac{1}{2\ii}\sum_a\sum_{\nu_1,\nu_2}\tanh\!\left(-\frac{\beta(\nu_1-\nu_2)}{4}\right)\;\alpha_{\nu_1\nu_2}\;(A^a_{\nu_2})^\dagger A^a_{\nu_1}.
\end{equation}
Concluding this section, we have successfully found a unique and Hermitian $B$ which makes the decay and coherent parts of the Lindbladian, $\mathcal{R}$ and $\mathcal{B}$ KMS detailed balanced. Since $\alpha_{\nu_1, \nu_2}$ has been kept skew-symmetric, the translation part $\mathcal{T}$ is also correct. Combining everything together, we get a KMS detailed balanced Lindbladian, that we will call from now on as the CKG Lindbladian,
\begin{equation}\label{eq:ckg-L}
    \mathcal{L}_{\text{CKG}}(\cdot) = -\ii[B, \cdot\,] + \sum_a \int_{-\infty}^\infty \gamma(\omega) \left(\hat{A}^a(\omega)(\cdot)\hat{A}(\omega)^\dagger -\frac{1}{2}\left\{\hat{A}^a(\omega)^\dagger \hat{A}^a(\omega), \cdot \right\}\right)\dd\omega .
\end{equation}

\smallskip
\subsection{Designing transition weights}
\todo{plot for gammas}
Only allowing transitions with energies concentrated as a Gaussian around some $\omega_\gamma$ can be very restricting, which inevitably leads to slow thermalization. The question is, \textit{can we come up with more accepting rates $\gamma(\omega)$ while keeping the KMS detailed balance of the CKG Lindbladians?} The answer is yes, by finding a clear connection between the condition on the transition weights $\gamma(\omega)$ and the necessary skew-symmetry of the Kossakowski matrix $\alpha_{\nu_1, \nu_2}$.
\begin{lemma} \label{lem:shifted-kms-condition}
    \textup{(Shifted KMS condition for detailed balance, \cite{ramkumar2024mixing})} If $\tilde{\gamma}(\omega)$ satisfies the KMS condition $\tilde{\gamma}(\omega) = \tilde{\gamma}(-\omega)e^{-\beta \omega},$ then the Kossakowski matrix 
    $$
    \alpha_{\nu_1, \nu_2} = \int_{-\infty}^\infty \gamma(\omega) \hat{f}(\omega - \nu_1) \hat{f}(\omega - \nu_2) \dd \omega \quad\textup{with}\quad \gamma(\omega) = \tilde{\gamma}\left(\omega + \frac{\beta \sigma^2}{2}\right),
    $$
    will be skew-symmetric by \eqref{eq:skew-symmetry}. Here, $\hat{f}$ is any filter with $|\hat{f}(\omega)|^2 = |\hat{f}(-\omega)|^2$, like the usual Gaussian filter from \eqref{eq:gaussian-filter-w}.
\end{lemma}
\noindent The proof from the authors follows by direct verification of the self-adjointness condition \eqref{eq:transition-self-adjoint} for $\mathcal{T}$. 

Hence, we can take any transition weight function $\tilde{\gamma}(\omega)$ that satisfies the KMS condition, shift it to $\gamma(\omega)$ and find a KMS detailed balanced CKG Lindbladian. This is quite interesting, since we can import canonically well-working transition weights that has been used for Gibbs sampling in the past. One such function, is the \textit{Metropolis function}, which takes the following form after the prescribed shift from \hyperref[lem:shifted-kms-condition]{Lemma~\ref*{lem:shifted-kms-condition}}: 
\begin{equation} \label{eq:shifted-metro-func}
    \gamma_\text{M}(\omega)= \exp\left(-\beta \max\left(\omega + \frac{\beta\sigma^2}{2}, 0\right)\right).
\end{equation}

From Fig. \todo{fig} it becomes apparent why the Metropolis function would lead to better mixing compared to the previous Gaussian case. It accepts almost all transitions that lower the energy with probability 1, and only suppresses energy increasing jumps. Since the Metropolis function is the centrepiece of the Metropolis-Hastings algorithm \cite{metropolis1953equation, hastings1970monte} which is already decades old, the community had time to rigorously analyse why it works so well. 
Peskun showed in \cite{peskun1973optimum} that among all reversible chains with the same proposal and target, the one with the pointwise maximal off-diagonal transition rates (between states $i \neq j$), has the smallest asymptotic variance for any estimator. The (unshifted) Metropolis acceptance is the unique pointwise maximum satisfying detailed balance with $\|\gamma\|_\infty\leq 1$. Then, Tierney extends these results in \cite{tierney1998note} by finding that the spectral gap among these reversible chains, the spectral gap is also the largest for the Metropolis transition, hence it is the fastest mixing of them all.

These are strong results, but they apply only for classical Markov chains and as it turns out, they do not generally hold in the quantum case. It is an open question if there exists a quantum variant of the above results. But at the moment it is not clear if and how the off-diagonal elements of a Kossakowski matrix $\alpha_{\nu_1, \nu_2}$ improve upon the mixing times of quantum Markov semigroups. There are recent numerical evidences in the general KMS detailed balanced framework \cite{chen2024boosting,hahn2025efficient} suggesting that there is not one choice that fits all problems. But also shows a gap in analytical results.

We can however look at just the CKG Lindbladian case and ask, for given Gaussian filter functions \eqref{eq:gaussian-filter-w}, would the Metropolis transition rate $\gamma_M$ lead to the fastest mixing?  

\begin{proposition} \label{prop:optim-metro}
    \textup{(Optimality of Metropolis weights for CKG Lindbladians)} Assume that the CKG Lindbladian $\mathcal{L}_\textup{CKG}^{\gamma_M}$ with transition weight $\gamma_M$ \eqref{eq:shifted-metro-func} is primitive. For any $\gamma: \mathbb{R} \rightarrow [0, \infty)$ that satisfies the shifted KMS condition \hyperref[lem:shifted-kms-condition]{Lemma~\ref*{lem:shifted-kms-condition}} and $\|\gamma\|_\infty \leq 1$: 
    $$\textup{Gap}\left(\mathcal{L}_\textup{CKG}^{\gamma_M}\right) \geq \textup{Gap}\left(\mathcal{L}_\textup{CKG}^{\gamma}\right) \geq 0.$$
\end{proposition}

\begin{proof}
    Define $\Delta := \gamma_M - \gamma$. We claim $\Delta \geq 0$. For the piece of the function where $\omega \geq -\beta\sigma^2 / 2$, the shifted KMS condition gives
    $$
    \gamma(\omega) = \gamma(-\omega - \beta\sigma^2)e^{-\beta(\omega + \beta\sigma^2/2)} \leq e^{-\beta(\omega + \beta\sigma^2 / 2)} \equiv \gamma_M(\omega),
    $$
    using $\|\gamma\|_\infty \leq 1$. For the other piece, $\omega < -\beta\sigma^2/2$, we directly have $\gamma(\omega) \leq 1 = \gamma_M(\omega)$. This argument is analogous to Peskun \cite{peskun1973optimum}. The CKG shifted KMS condition is a linear identity in $\gamma$: since both $\gamma_M, \gamma$ satisfy it, so does $\Delta$. Thus, $\mathcal{L}_\textup{CKG}^\Delta$ is a valid GKLS generator which satisfies KMS detailed balance.

    Since the map $\gamma\rightarrow\mathcal{L}^\gamma$ is linear: for any non-negative $\gamma_1, \gamma_2$, both satisfying the CKG shifted KMS condition, and any $c_1, c_2 \geq 0$, we have $\mathcal{L}_\textup{CKG}^{c_1\gamma_1 + c_2\gamma_2} = c_1 \mathcal{L}_\textup{CKG}^{\gamma_1} + c_2 \mathcal{L}_\textup{CKG}^{\gamma_2}$. The reason behind this is that every construction in the CKG Lindbladian involves an integral linear in $\alpha_{\nu_1, \nu_2}$ and correspondingly linear in $\gamma$. Then we can write: $$\mathcal{L}_\textup{CKG}^{\gamma_M} = \mathcal{L}_\textup{CKG}^\gamma + \mathcal{L}_\textup{CKG}^\Delta, 
    \quad\textup{or}\quad
     \left(\mathcal{L}_\textup{CKG}^{\gamma_M}\right)^\dagger = \left(\mathcal{L}_\textup{CKG}^\gamma\right)^\dagger + \left(\mathcal{L}_\textup{CKG}^\Delta\right)^\dagger.$$
    Recall that the spectral gap is defined as
    $\textup{Gap}\left(\mathcal{L}^\dagger\right) = \inf_{X} \mathcal{E}(X, X) / \|X\|^2_\textup{KMS}$ with $X\neq 0$ and $\tr(\rho_\beta X) = 0$, and with the Dirichlet form $\mathcal{E}^\gamma(X, X) := -\langle X, (\mathcal{L}^\gamma)^\dagger(X) \rangle_\textup{KMS}$. See also [REF] \todo{refer to preliminaries} and note that the Dirichlet form is defined in the Heisenberg picture, hence the $\mathcal{L}^\dagger$. Consequently, the proof boils down to showing that $$\mathcal{E}^\Delta(X, X) \geq 0 \quad\forall X \in \mathcal{B(H)},$$
    which is true since: (i) the Schrödinger and Heisenberg operators share the same spectra; (ii) $\mathcal{L}_\textup{CKG}^\Delta$ is a GKLS generator and thus has $\textup{Re}(\lambda) \leq 0$ for all eigenvalues $\lambda$; (iii) $\mathcal{L}_\textup{CKG}^\Delta$ has KMS detailed balance, i.e. it is self-adjoint with respect to the KMS inner-product which makes all eigenvalues on this space real. Combining all of this leads to $\langle X, (\mathcal{L}^\Delta)^\dagger(X) \rangle_\textup{KMS} \leq 0$ and finally to the above statement.
    
    The last step is to relate the Dirichlet forms and the spectral gaps. For all $X$ with $\tr(\rho_\beta X) = 0$, $X \neq 0$:
    $$  
    \mathcal{E}^{\gamma_M}(X, X) = \mathcal{E}^{\gamma}(X, X) + \mathcal{E}^{\Delta}(X, X) \geq \mathcal{E}^{\gamma}(X, X).
    $$
    Since $\|X\|^2_\textup{KMS}$ is independent of $\gamma$, dividing with it and taking the infimum yields:
    $$
    \textup{Gap}\left(\mathcal{L}_\textup{CKG}^{\gamma_M}\right) \geq \textup{Gap}\left(\mathcal{L}_\textup{CKG}^{\gamma}\right)
    $$
    The primitivity assumption ensures $\textup{Gap}\left(\mathcal{L}_\textup{CKG}^{\gamma_M}\right) > 0$.
\end{proof}

Although, we cannot say anything definite about optimal Kossakowski matrices, with the above proposition, we have at least the motivation to numerically analyse CKG Lindbladians with the Metropolis function $\gamma_M$ or (risking foreshadowing here) with its smoothened variants.

There is however an interesting technique that is worth mentioning, as it can help us design transition functions and Kossakowski matrices.
In \cite{chen2023efficient} the authors found this shifted Metropolis function $\gamma_M$ by different means. The idea is to take a linear combination of detailed balanced Lindbladians in order to speed up the mixing. Due to linearity, a linear combination of skew-symmetric Kossakowski matrices also respects the symmetry,
\begin{equation}\label{eq:lin-comb-alpha}
    \alpha_{\nu_1,\nu_2}^{(s, g)} := \int_{\frac{\beta\sigma^2}{2}(1 + s)}^\infty g(x) \alpha_{\nu_1,\nu_2}^{(\omega_\gamma, \sigma_\gamma)} \dd x \quad\text{with}\quad (\omega_\gamma(x), \sigma_\gamma(x)) = \left(x, \sqrt{\frac{2x}{\beta} - \sigma^2}\right).
\end{equation}
Where $g \in \ell_1(\mathbb{R})$ is necessary to keep the resulting Kossakowski matrix Hermitian. This is a linear combination of the skew-symmetric Gaussian-case Kossakowski matrices \eqref{eq:kossakowski-gaussian} by sweeping with the centre of the transition Gaussians $\omega_\gamma$. Hence, we find a transition that admits a larger set of energy differences.

By choosing the $g$ as the inverse $\ell_1$ weight of the Gaussians, $g(x) = \frac{1}{\sqrt{2\pi}\sigma_\gamma}$ and $s = 0$, Chen et al. managed to find the shifted Metropolis function \eqref{eq:shifted-metro-func}, after applying the linear combination to the Gaussian transition weight $\gamma_G(\omega)$ with the same $\omega_\gamma(x), \sigma_\gamma(x)$ parameters as in \eqref{eq:lin-comb-alpha},
\begin{equation} \label{eq:gamma_M-from-lin-comb-of-G}
    \gamma_M(\omega) = \int_{\beta\sigma^2 / 2}^\infty g(x)\gamma_G(\omega) \dd x.
\end{equation}
But while this leads to better mixing than the Gaussian case, it has some problems the moment we want to implement it. 

The source of all these problems is the fact that though $\gamma_M$ is Lipschitz continuous, it has a kink at $\omega = -\beta\sigma^2 / 2$ (and the unshifted variant at $\omega = 0$), see \todo{plot}. To be clear up front, the problems that this fact leads to are not dominating the asymptotic error / cost scalings of the algorithm, but they seem to be unnecessary costs, that for implementations prior fault-tolerance can be important. 
\begin{itemize}
    \item On one side this would lead to extra quantum circuit synthesis costs, since in the quantum algorithm at one point we have to apply a rotation onto an ancilla qubits based on $\gamma(\omega)$ \todo{refer to q. alg and circuit}. There are multiple ways to do this, but one is to use QSP, which has to approximate $\gamma_M$ via some polynomials. Due to the kink however we would need to have higher degree polynomial for approximations, and thus a higher gate cost \todo{refer to something possibly}. Chen et al. notes \cite[footnote 33.]{chen2023quantum} that we could do this also piecewise, if we check the energy register via a comparator circuit block (to decide between left and right of the kink) and then apply QSP on smoother pieces of functions. But the costs are higher if the kink is not exactly at $\omega = 0$, since the comparator circuit then has to check more bits in the register, not just the sign bit. 
    \item The other problem is more important however: since there is an integral over $\omega$ in the CKBG and CKG Lindbladians, we would have to approximate this with a quadrature sum in the quantum computer. A classical result of Fourier analysis \cite[Chapter 2.]{stein2011fourier} states that if a function $f$ is $k$ times continuously differentiable with $f^{(k)}$, its Fourier transform decays as $\hat{f}(t) = \mathcal{O}(|t|^{-(k+1)})$. The Metropolis $\gamma_M$ has a discontinuous first derivative at the kink, and thus its Fourier transform decays only as $\mathcal{O}(1 / t)$. Of course the convolution with the Gaussian filters with width $ \sigma^{-1}$ smoothen this kink and improves the scaling, but the algebraic prefactor persists $\mathcal{O}(\frac{e^{-\sigma^2 t^2}}{\sigma^2 t^2})$. Hence, this algebraic tail requires a larger truncation range and a finer grid (more estimating ancilla qubits) to achieve the same precision as an analytic transition weight.
    \item Finally, the existence of the kink is tightly connected to having a constant piece of the Metropolis function that equates to 1 for $\omega \leq \beta\sigma^2 / 2$. When passing to the time-domain representation of the Lindbladian (the one that runs in the quantum computer) via a Fourier transform, this constant piece produces a $\delta(t)$ singularity in the integral kernel. In practice this singularity has to be tamed by some regularization technique which incurs more costs as well. \todo{refer to Metropolis tame and maybe our taming}
\end{itemize}

But we can have our cake and eat it too, almost that is. We can have close to optimal transitions while not stumbling into all the additional implementation costs. The solution is to smoothen the kink of the Metropolis function $\gamma_M$, which eliminates the first two previously mentioned problems. After a long battle of trying to improve on \cite{chen2023quantum}'s way of regularizing the singularity in the time domain, it seems that the incurring additional costs cannot be made better. The remnants of the battle can be found in [REF to APPENDIX] \todo{appendix where parameter a is back in the lin comb.}.

To designing a better transition function $\gamma$, we will make use of the linear combination that was introduced in \eqref{eq:lin-comb-alpha}. The key is to choose $s \neq 0$ to shift the lower limit of the integration away from the kink, in a tuneable manner, i.e. the larger $s$ is the stronger we smoothen. Explicitly writing out $\gamma_G$ and $g(x)$ yields,
\begin{equation} \label{eq:smooth-metro-int-defi}
    \gamma_M^{(s)}(\omega) = \int_{\frac{\beta\sigma^2}{2}(1 + s)}^\infty \frac{1}{\sqrt{2\pi \sigma_\gamma^2(x)}} \exp\left(\frac{-(\omega + x)^2}{2 \sigma_\gamma^2(x)}\right) \dd x.
\end{equation}
The resulting the transition function cleanly separates into the Metropolis function $\gamma_M$ with an additional factor dependent on $s, \beta, \sigma$:
\begin{equation} \label{eq:smooth-metro}
    \gamma_M^{(s)}(\omega) = \gamma_M^{(s=0)}(\omega) \times \frac{1}{2}\left[ \erfc\left(z_-\right) + e^{\beta |\tilde{\omega}|} \erfc\left(z_+\right) \right]
\end{equation}
with $\tilde{\omega} = \omega + \beta\sigma^2 / 2$ and 
$$
z_\pm := \frac{\beta \sigma \sqrt{s}}{2\sqrt{2}} \pm \frac{|\tilde{\omega}|}{\sigma\sqrt{2s}} 
$$
\todo{check with numerics, if these were the correct factors here.}
Of course, there is a drawback here, since \hyperref[prop:optim-metro]{Proposition~\ref*{prop:optim-metro}} we know that any other choice of transition function $\gamma$ could lead to a smaller spectral gap. More concretely in our case, all $\gamma^{(s)}$ are pointwise smaller or equal to $\gamma_M$ since the right multiplying factor in \eqref{eq:smooth-metro} is $\leq 1$ for $s > 0$. But in practice we only need a small $s$, which is enough to tame quadrature errors (and circuit costs), while keeping almost the same spectral gap of the CKG Lindbladian, see the numerical analysis [REF] \todo{ref}. The parameter $s$ creates a family of Metropolis-like transitions, with $s\rightarrow 0^+$ recovering the Metropolis function $\gamma_M$ (the square bracket with the error functions equals 2 in this limit). Also, we get the Glauber-like transition weights if we choose $s = 2$, which has been found by Chen et al. in \cite{chen2023quantum}. But that is a stark smoothening, which leads to noticeably worse mixing times in the simulations.

Let us state rigorously in what sense the smooth Metropolis $\gamma^{(s > 0)}_M$ improves upon its kinky version $\gamma^{(s = 0)}_M$. First, we have to recall the definition of Gevrey functions \cite{adwan2017global}.
\begin{definition}
    \textup{(Gevrey functions).} Let $\Omega \subseteq \mathbb{R}$. A $C^\infty$ function $h: \Omega \to \mathbb{C}$ is a Gevrey function of order $\kappa \geq 0$ if there exist constants $C_1, C_2 > 0$ such that
    \[
    \|h^{(n)}\|_{L^\infty(\Omega)} \leq C_1\, C_2^n\, n^{n \kappa} \quad \text{for all } n \geq 0,
    \]
    For the given constants $C_1, C_2, \kappa$, the set of Gevrey functions is denoted by $\mathcal{G}_{C_1, C_2}^\kappa(\Omega)$.
\end{definition}
Being in a Gevrey set tells us something about the regularity of a function. For instance, with order $\kappa = 0$ the set of real analytic functions is recovered. Gevrey order with finite $\kappa$ provides a quantitative interpolation: integrals over $\mathbb{R}$ admit quadrature errors decaying as $\exp(-c N^{1 / \kappa})$ under optimal truncation of the Euler-Maclaurin expansion, where $N$ is the number of quadrature nodes \cite{trefethen2014exponentially}.
Conversely, if the integrand has only finite regularity, for instance due to a kink, the quadrature errors converge only at a polynomial rate $\mathcal{O}(N^{-k})$ determined by the order of smoothness $k$ \cite{trefethen2014exponentially}.

In the following proposition we will denote the integrand in \eqref{eq:smooth-metro-int-defi} as $F(\omega, x)$, i.e. $\gamma_M^{(s)}(\omega) = \int_{x_0}^\infty F(\omega, x)\,dx$, with the lower limit $x_0 := \beta \sigma^2 (1 + s) / 2$.
\begin{proposition} \label{prop:smooth-metro-gevrey}
    \textup{(Smooth Metropolis is Gevrey-1/2).} For $s > 0, \beta > 0$, the smooth Metropolis weight $\gamma^{(s)}_M$ defined by \eqref{eq:smooth-metro-int-defi} satisfies:
    \begin{enumerate}[(i)]
        \item For all $n\geq 1$ and $\omega \in \mathbb{R}$:
        \[
        \left|\left(\gamma_M^{(s)}\right)^{(n)}(\omega)\right| \leq \frac{1}{2\pi}\left(\frac{2\beta^2}{s}\right)^{n/2} \Gamma\!\left(\frac{n}{2}\right),
        \]
        with $f^{(n)}$ denoting the $n$-th derivative, and $\Gamma(z) = \int_{0}^{\infty} t^{z - 1} \ee^{-t} \dd t$ is the Euler gamma function.
        \item Consequently, $\gamma_M^{(s)} \in \mathcal{G}^{1/2}_{C_1, C_2}(\mathbb{R})$ with 
        \[
        C_1 = 1, \qquad C_2 = \frac{1}{\sigma\sqrt{\ee\,s}}.
        \]
        %\item For $s = 0$, the Metropolis weight $\gamma_M^{(0)}(\omega) = e^{-\beta\max(\omega + 1/(2\beta),\, 0)}$ is Lipschitz but not $C^1$, hence not in any Gevrey class.
    \end{enumerate}
\end{proposition}
\begin{proof}
    For a fixed $x \geq x_0$, the function $\omega \mapsto F(\omega, x)$ is a Gaussian with mean $-x$ and variance $\sigma_\gamma^2(x) \geq s \sigma^2 / 2 > 0$. Its Fourier transform in $\omega$ is
    \[
    \mathcal{F}[F(\cdot, x)](\xi) = \frac{1}{\sqrt{2\pi}}\,\ee^{\ii x\xi}\, \ee^{-\sigma_\gamma^2(x)\, \xi^2/2}.
    \]
    Substituting $\sigma_\gamma^2(x) = 2x / \beta  - \sigma^2$:
    \[
    \mathcal{F}[F(\cdot, x)](\xi) = \frac{1}{\sqrt{2\pi}}\,\ee^{\ii x\xi - (x/\beta - \sigma^2/2)\,\xi^2} = \frac{1}{\sqrt{2\pi}}\,\ee^{\sigma^2\xi^2/2}\, \ee^{x(\ii\xi - \xi^2/\beta)}.
    \]
    Then for, $n \geq 1$, differentiation under the integral is justified by the locally uniform integrability of $|\partial_\omega^n F(\omega, x)|$ over $x \in [x_0, \infty)$: the Gaussian decay $F(\omega, x) \lesssim e^{-\beta x/4}$ provides the domination, and $\sigma_\gamma > 0$ prevents singularities. Therefore:
    \[
    \left(\gamma_M^{(s)}\right)^{(n)}(\omega) = \int_{x_0}^\infty \frac{\partial^n F}{\partial \omega^n}(\omega, x)\, \dd x.
    \]
    Taking the Fourier transform of both sides in $\omega$ and applying Fubini's theorem:
    We compute the integral for the linear combination. For real $\xi \neq 0$
    \begin{equation*}
        \begin{split}
            \mathcal{F}\left[\left(\gamma_M^{(s)}\right)^{(n)}\right](\xi) &= \frac{(\ii\xi)^n}{\sqrt{2\pi}} \int_{x_0}^\infty \ee^{\sigma^2\xi^2/2}\, \ee^{x(\ii\xi - \xi^2/\beta)}\, \dd x. \\
            &= \frac{(i\xi)^n \cdot \beta}{\sqrt{2\pi}} \cdot \frac{e^{\sigma^2\xi^2/2 + x_0(i\xi - \xi^2/\beta)}}{\xi(\xi - i\beta)}.
        \end{split}
    \end{equation*}
    Where for computing the integral, $\mathrm{Re}(\xi^2/\beta - i\xi) = \xi^2/\beta > 0$, thus we do not run into singularities.
    Substituting $x_0$ and doing some arithmetics we arrive at
    \[
    \mathcal{F}[(\gamma_M^{(s)})^{(n)}](\xi) = \frac{(i\xi)^{n-1} \cdot i\beta}{\sqrt{2\pi}} \cdot \frac{e^{-s\,\sigma^2\xi^2/2\; +\; i\beta\sigma^2(1+s)\xi/2}}{\xi - i\beta},
    \]
    with modulus:
    \[
    \left|\mathcal{F}[(\gamma_M^{(s)})^{(n)}](\xi)\right| = \frac{\beta\, |\xi|^{n-1}}{\sqrt{2\pi}\,\sqrt{\xi^2 + \beta^2}}\; e^{-s\,\sigma^2\xi^2/2}.
    \]
    Since $\mathcal{F}[(\gamma_M^{(s)})^{(n)}] \in L^1(\mathbb{R})$ for $n \geq 1$, we can now apply the Fourier inversion formula and bound the derivatives of the smooth Metropolis weights,
    \[
    \left|(\gamma_M^{(s)})^{(n)}(\omega)\right| \leq \frac{1}{\sqrt{2\pi}} \int_{\mathbb{R}} \left|\mathcal{F}[(\gamma_M^{(s)})^{(n)}](\xi)\right| d\xi = \frac{\beta}{2\pi} \int_{\mathbb{R}} \frac{|\xi|^{n-1}\, e^{-s\,\sigma^2\xi^2/2}}{\sqrt{\xi^2 + \beta^2}}\, d\xi.
    \]
    Using $\sqrt{\xi^2 + \beta^2} \geq \beta$ for all $\xi\in\mathbb{R}$ (which is a loose bound, but a tighter solution still would lead to the same Gevrey order):
    \[
    \left|(\gamma_M^{(s)})^{(n)}(\omega)\right| \leq \frac{1}{2\pi} \int_{\mathbb{R}} |\xi|^{n-1}\, e^{-s\,\sigma^2\xi^2/2}\, d\xi = \frac{1}{\pi} \int_0^\infty \xi^{n-1}\, e^{-s\,\sigma^2\xi^2/2}\, d\xi.
    \]
    This Gaussian integral can be evaluated with $\alpha := s \sigma^2 / 2$ and the identity $\int_0^\infty \xi^{n-1} e^{-\alpha \xi^2}\, d\xi = \frac{1}{2}\,\alpha^{-n/2}\,\Gamma(n/2)$:
    \[
    \left|\left(\gamma_M^{(s)}\right)^{(n)}(\omega)\right| \leq \frac{1}{2\pi}\left(\frac{2}{s\,\sigma^2}\right)^{n/2} \Gamma\!\left(\frac{n}{2}\right),
    \] 
    proving the part $(i)$ of the proposition.
    The last step is to show that 
    \[
    \frac{1}{2\pi}(2/(s\sigma^2))^{n/2}\,\Gamma(n/2) \leq C_1\, C_2^n\, n^{n/2}
    \]
    for suitable constants $C_1, C_2$. Let us handle the different cases of $n$ separately: 
    \begin{itemize}
        \item $n = 0$: $\|\gamma_M^{(s)}\|_\infty \leq 1$ since $\gamma_M^{(s)} \leq \gamma_M^{(s=0)} \leq 1$.
        \item $n = 1$: Substituting $\Gamma(1 / 2) = \sqrt{\pi}$ into (i) gives $\| (\gamma_M^{(s)})'\|_\infty \leq \frac{1}{\sqrt{2\pi} \sigma s}$ 
        \item $n \geq 2$: in these cases the Stirling upper bound gives $\Gamma(z) \leq \sqrt{2\pi} (z / \ee)^z$. Applying this with $z = n /2 \geq 1$ to (i) gives, 
        \[
        \| (\gamma_M^{(s)})^{(n)}\|_\infty \leq \frac{1}{\sqrt{2\pi}}\left(\frac{1}{\ee s \sigma^2}\right)^{n / 2} n^{n / 2}
        \]
    \end{itemize}
    From which we find that $C_1 = 1$ and $C_2 = \frac{1}{\sigma \sqrt{\ee s}}$ (where we could absorb the prefactor of $C_2$, since $1 / \sqrt{2\pi} < C_1$).
\end{proof}
\noindent An important remark here is that the Metropolis function $\gamma_M^{(s = 0)}$, has a left derivative at the kink $\omega + \beta\sigma^2 / 2$ equal to $0$, and a right derivative equal to $-\beta$, which makes it only $C^0$ but not $C^1$, and in particular it is not a Gevrey function of any order. This qualitative difference between transition rates will become important for our quadrature error discussions [REF] \todo{ref}.

\smallskip
\subsection{Improved Kossakowski matrices.}
For bookkeeping and later analysis, it is worth to write down what form the Kossakowski matrix takes if instead of the Gaussian transition $\gamma_G$ we use the smoothened Metropolis functions $\gamma_M^{(s)}$. Similarly to \eqref{eq:kossakowski-gaussian}, integrating over $\omega$ with the kernel consisting of $\gamma_M^{(s)}$ and the Gaussian filters $\hat{f}(\omega - \nu)$, we find, 
\begin{equation} \label{eq:kossakowski-smooth-metro}
    \alpha_{\nu_1, \nu_2}^{(s)} = \frac{1}{2}e^{-\frac{\nu_-^2}{8\sigma^2}}e^{-\frac{\beta}{4}(\nu_+ +|\nu_+|)}\left[\text{erfc}(\zeta_-) + e^{\frac{\beta}{2}|\nu_+|}\text{erfc}(\zeta_+)\right]
\end{equation}
with $\nu_\pm := \nu_1 \pm  \nu_2$ and 
$$
\zeta_\pm := \frac{\beta \sigma \sqrt{1 + s}}{\sqrt{8}} \pm \frac{|\nu_+|}{\sigma\sqrt{8(1 + s)}}.
$$
A few things to note here:
\begin{enumerate}[(i)]
    \item $\alpha^{(s)}_{\nu_1,\nu_2}$ is positive semidefinite, since $\gamma_M^{(s)}\geq 0$ and $\hat{f}\geq 0$, guaranteeing GKSL form for the Lindbladian. It also satisfies the skew-symmetry \eqref{eq:skew-symmetry}, i.e. leads to a KMS detailed balanced Lindbladian as it was promised via linear combinations of detailed balanced Lindbladians.
    \item In the limit of $s\rightarrow 1$ we again recover the Metropolis Kossakowski matrix that Chen et al. found \cite[Appendix B]{chen2023quantum}, exactly how we previously recovered $\gamma_M$ from $\gamma_M^{(s)}$.
    \item In the $\sigma\rightarrow 0$ limit consistently finds the diagonal Kossakowski matrix, i.e. the Davies generator, since the off-diagonal suppression $e^{-\nu_-^2 / (8\sigma^2)}$ forces all terms with $\nu_1 \neq \nu_2$ to vanish.
    \item Due to the factorizing Gaussian filter functions, $\alpha_{\nu_1,\nu_2}^{(s)}$ can be separated into to two parts, that each depend on either $\nu_+$ or $\nu_-$. This will be important for later, when we derive time-domain functions, since it allows for a double Fourier transformation. Note, that this is also true for the case with the Gaussian transition weight \eqref{eq:kossakowski-gaussian}.
    \item Another implication of the factorizing Gaussian filters is a simple shift relation between the Kossakowski matrices for the GNS and KMS detailed balanced Lindbladians (i.e. between CKBG and CKG Lindbladians), analogous to \hyperref[lem:shifted-kms-condition]{Lemma~\ref*{lem:shifted-kms-condition}}:
\end{enumerate}

\begin{corollary} \label{cor:kossakowski-shift}
    \textup{(Kossakowski shift between CKG and CKBG).}
    Let $\tilde{\gamma}(\omega) = \tilde{\gamma}(-\omega)\,e^{-\beta\omega}$, and let
    $\hat{f}$ be the Gaussian filter with width~$\sigma$ from
    \eqref{eq:gaussian-filter-w}.
    Denote by $\alpha^{\textup{KMS}}_{\nu_1,\nu_2}$ the Kossakowski matrix
    of the CKG Lindbladian, which uses the shifted transition weight
    $\gamma(\omega) = \tilde{\gamma}(\omega + \beta\sigma^2/2)$ from
    \hyperref[lem:shifted-kms-condition]{Lemma~\ref*{lem:shifted-kms-condition}}:
    %
    $$
      \alpha^{\textup{KMS}}_{\nu_1,\nu_2}
      \;=\; \int_{-\infty}^{\infty}
        \tilde{\gamma}\!\bigl(\omega + \tfrac{\beta\sigma^2}{2}\bigr)\,
        \hat{f}(\omega - \nu_1)\,\hat{f}(\omega - \nu_2)\,\dd\omega.
    $$
    %
    Then the Kossakowski matrix of the CKBG Lindbladian, which uses
    $\tilde{\gamma}$ directly,
    %
    $$
      \alpha^{\textup{GNS}}_{\nu_1,\nu_2}
      \;=\; \int_{-\infty}^{\infty}
        \tilde{\gamma}(\omega)\,
        \hat{f}(\omega - \nu_1)\,\hat{f}(\omega - \nu_2)\,\dd\omega,
    $$
    %
    satisfies
    %
    \begin{equation} \label{eq:kossakowski-shift}
      \alpha^{\textup{GNS}}_{\nu_1,\nu_2}
      \;=\;
      \alpha^{\textup{KMS}}_{\,\nu_1 - \beta\sigma^2/2,\;\,\nu_2 - \beta\sigma^2/2}.
    \end{equation}
    %
    Equivalently, with $\nu_{\pm} := \nu_1 \pm \nu_2$, the off-diagonal
    suppression factor $e^{-\nu_-^2/(8\sigma^2)}$ is invariant under the
    shift, while every occurrence of $\nu_+$ in
    $\alpha^{\textup{KMS}}$ is replaced by $\nu_+ - \beta\sigma^2$.
\end{corollary}

\begin{proof}
    Substitute $u = \omega + \beta\sigma^2/2$ in the defining integral of
    $\alpha^{\textup{KMS}}_{\nu_1 - \beta\sigma^2/2,\;\nu_2 - \beta\sigma^2/2}$:
    %
    \begin{align*}
      \alpha^{\textup{KMS}}_{\nu_1 - \beta\sigma^2/2,\;\nu_2 - \beta\sigma^2/2}
      &= \int_{-\infty}^{\infty}
        \tilde{\gamma}\!\bigl(\omega + \tfrac{\beta\sigma^2}{2}\bigr)\,
        \hat{f}\!\bigl(\omega - \nu_1 + \tfrac{\beta\sigma^2}{2}\bigr)\,
        \hat{f}\!\bigl(\omega - \nu_2 + \tfrac{\beta\sigma^2}{2}\bigr)\,
        \dd\omega \\[4pt]
      &= \int_{-\infty}^{\infty}
        \tilde{\gamma}(u)\,
        \hat{f}(u - \nu_1)\,\hat{f}(u - \nu_2)\,\dd u
      \;=\; \alpha^{\textup{GNS}}_{\nu_1,\nu_2}.
    \end{align*}
    %
    For the $\nu_\pm$ decomposition, the Gaussian filter gives
    %
    $$
      \hat{f}(\omega - \nu_1)\,\hat{f}(\omega - \nu_2)
      \;\propto\;
      \exp\!\Bigl(-\frac{(\omega - \nu_+/2)^2}{2\sigma^2}\Bigr)\,
      \exp\!\Bigl(-\frac{\nu_-^2}{8\sigma^2}\Bigr).
    $$
    %
    The shift $(\nu_1, \nu_2) \mapsto
    (\nu_1 - \beta\sigma^2/2,\;\nu_2 - \beta\sigma^2/2)$ sends
    $\nu_- \mapsto \nu_-$ and $\nu_+ \mapsto \nu_+ - \beta\sigma^2$,
    leaving the $\nu_-$-dependent factor unchanged, and shifting $\nu_+$ as prescribed.
\end{proof}
\noindent Though it is a simple corollary to \hyperref[lem:shifted-kms-condition]{Lemma~\ref*{lem:shifted-kms-condition}}, we will make use of it in our numerical analysis, where we look at the mixing performances of the approximately GNS detailed balanced CKBG Lindbladian [REF to num analysis of GNS] \todo{ref}.

\smallskip
\subsection{Time, time, time.} \todo{add norms and LCU subnormalizations here.}
As we have seen, the energy-domain representation is natural for analysing detailed balance, but if we want to run any algorithm in a quantum computer, we have to understand and analyse the given algorithm also in the time-domain. We have already seen one of the core subroutines in the time-domain that \cite{chen2023efficient,chen2023quantum} used, the Operator Fourier Transform, \eqref{eq:oft} and its circuit form [REF to Quantum Circuits chapter] \todo{ref}. With this the dissipative part of the Lindbladians is more or less covered, though how one implements the corresponding open quantum system evolution, $\ee^{t\mathcal{L}}(\rho)$, will be left for the next paragraph.

Then, we are left to work out what form the coherent part $G$ \eqref{eq:B-bohr-domain} takes in the time-domain. The factorization property that we have noted in the previous chapter is the key to derive the time-domain representation. For any transition weight $\gamma$ (e.g. $\gamma_G, \gamma_M^{(s)}$) and with the Gaussian filters \eqref{eq:gaussian-filter-w} the Kossakowski matrix splits as
\begin{equation}
    \frac{1}{2\pi}\alpha_{\nu_1,\nu_2} = \hat{f}_+(\nu_+)\,\hat{f}_-(\nu_-), \qquad \nu_{\pm} := \nu_1 \pm \nu_2,
\end{equation}
where the $\hat{f}_-$ factor encodes the off-diagonal suppression and the coherent-term structure ($\tanh$ piece), while $\hat{f}_+$ carries all dependence on the transition weight $\gamma$. By \cite[Corollary~A.1]{chen2023efficient}, the two-dimensional Fourier transform over the Bohr frequencies $(\nu_1,\nu_2)$ then decouples into two independent one-dimensional transforms, yielding a nested double integral in the time domain. 

After executing the two independent Fourier transforms on $\hat{f}_\pm(\nu_\pm)$, we find that the coherent term $G$ takes the following general form in the time-domain \cite[Corollary III.1]{chen2023efficient}:
\begin{equation} \label{eq:coherent-time-domain}
    B = \sum_{a \in \mathcal{A}} \int_{-\infty}^{\infty} b_-(t)\, e^{-iHt / \sigma} \left[\int_{-\infty}^{\infty} b_+(t')\, A^{a\dagger}(\beta t')\, A^a(-\beta t')\,\dd t' \right] e^{iHt/\sigma}\,\dd t
\end{equation}
Here $b_-(t)$ comes from $\hat{f}_-(\nu_-)$, and it is universal across all $\gamma$ choices. It takes the following form:
\begin{equation} \label{eq:b_minus}
    b_-(t) = \frac{2\sqrt{\pi}}{\beta \sigma}e^{\beta^2\sigma^2/8}\left[\frac{1}{\cosh(\frac{2\pi t}{\beta \sigma})}\ast_t \sin(-\beta \sigma t) e^{-2t^2}\right] \quad\text{with}\quad \|b_-\|_1 \leq \frac{\pi}{\sqrt{2}}e^{\beta^2 \sigma^2 / 8},
\end{equation}
where $\ast_t$ denotes the convolution over the variable $t$, 
$$
    (f \ast_t g)(t) := \int_{-\infty}^{\infty} f(s) g(t - s)\dd s.
$$
Importantly, $b_-(t)$ is real function that has a support over $t = \mathcal{O}(\beta \sigma)$ due to the width of the $1 / \cosh$ kernel. The time evolution of the outer integral is $e^{\pm \beta H t / \sigma}$, which implies Hamiltonian simulation time of $\mathcal{O}(\beta)$.

Since $b_+$ depends on $\gamma$, we will have to look at the Gaussian $\gamma_G$ and the smooth Metropolis $\gamma_M^{(s)}$ separately. The Gaussian case is definitely the simpler one out of the two. It takes the following form:
\begin{equation} \label{eq:b_plus-gauss}
    b_+(t) = \frac{2 \sigma_\gamma}{\pi\sqrt{\pi}}e^{-4t^2\beta^2 \omega_\gamma - 2it\beta\omega_\gamma} \quad\text{with}\quad \|b_+\|_1 = \frac{\beta}{\omega_\gamma} \sqrt{\frac{\sigma_\gamma}{2\pi}}.
\end{equation}
This is a rapidly decaying function with width $\mathcal{O}(1 / \sqrt{\beta \omega_\gamma}) = \mathcal{O}(1)$, if we use that $\omega_\gamma$ has a natural unit of $1 / \beta$. Hence, the inner integral also needs an $\mathcal{O}(\beta)$ amount of Hamiltonian simulation time in this case.

The Metropolis case is slightly more complicated. But this should not surprise us after the foreshadowing in the discussions of \eqref{eq:gamma_M-from-lin-comb-of-G}. Since in the energy-domain the Metropolis function has a constant piece for most of the negative energies $\omega < -\beta\sigma^2 / 2$, the Fourier transform produces a $\delta(t)$ contribution in the time domain. In the smooth case $s \neq 0$, $\gamma_M^{(s)}$ is only asymptotically reaching that flat $\gamma = 1$ portion. Nevertheless, the near singularity behaviour also persists in these smoother cases. Despite the divergent $\ell_1$ norm of $b_+^{(s)}$ the operator $B_M$ is actually bounded. The key idea is that the $1/t$ singularity, when sandwiched between operator valued functions $A^\dagger(t)A(t)$, can be absorbed into $H \textup{sinc}(2Ht)$ terms that are bounded at $t=0$, \cite[Proposition B.2]{chen2023efficient}. Though mathematically this is sound, and implementing it would also be asymptotically efficient with the Quantum Singular Value Transformation (QSVT), but would also be cumbersome. Moreover, it would be a significant change to the LCU quantum circuit [ref]\todo{ref to quantum circuits section} that we use for the Gaussian case coherent term $B_G$. 
For LCU circuits we care about the subnormalization of the operator we want to implement, as it determines the quantum circuit cost. For the exact Metropolis term $B_M$ it bounded by \cite[Proposition B.2]{chen2023efficient}
\begin{equation}
    \|B^{(s=0)}\| \equiv\|B_M\| \leq \left\|\sum_a A^{a\dagger}A^a\right\| \mathcal{O}\left(\log\left(\beta \| H \|\right)\right)
\end{equation}
Instead of implementing the Metropolis coherent term exactly, we can regularize $b_+^{(s)}$ around $t=0$ by replacing the singularity near $|t|\leq\eta$ with a smooth cutoff \cite[Proposition B.1]{chen2023efficient}, which leads to the following expression:
\begin{equation} \label{eq:b_plus-s-eta}
    b_+^{(s, \eta)}(t) = \frac{1}{2\sqrt{2}\pi^2}\frac{\ee^{-\sigma^2\beta^2 t(2t + \ii)(1 + s)} + \mathbb{I}(|t|\leq\eta)\ii(2t + \ii)}{t(2t + \ii)}, 
\end{equation}
with the $\ell_1$ norm, 
$$
\|b_+^{(s, \eta)}\|_1 < \frac{2}{5 \pi^{3/2}} - \frac{\ln(1 + s)}{2\sqrt{2}\pi^2}+ \frac{1}{\sqrt{2}\pi^2}\ln(1 / \eta).
$$

\todo{write somewhere: B M as A'A additional term in it}
Consistency check: Chen's norm bound is recovered for $s=0$, albeit with the corrected prefactors. There was a slight prefactor discrepancy in their paper, the corrected norm is smaller by a factor of $1 / \pi^{3 / 2}$, which improves their implementation by a constant factor. For the $s$-family the log-term is negative for $s > 0$, the norm is strictly smaller than the Metropolis case with $s = 0$, which also helps by a $\ln(1 + s)$ decrease for the LCU implementations. 

In \eqref{eq:b_plus-s-eta} it also becomes clear that smoothing with $s$ indeed does not shift the poles of the expression above $(t = 0, -\ii / 2)$, and thus the regularization is still necessary. 
In \cite[Corollary III.2]{chen2023efficient} they have shown that the error between the regularized coherent term with $B^{(s = 0)}_\eta$ and the non-regularized one $B^{(s = 0)}$ is bounded. Their result also adapts to the smooth Metropolis variants since the error is purely due to the regularized region $|t| \leq \eta$ where the smoothening plays no role, 
\begin{equation}
    \|B^{(s)} - B_\eta^{(s)} \| \leq \left\|\sum_a A^{a\dagger}A^a\right\| \min\left(\frac{\eta\beta\|H\|}{\sqrt{2}\pi}, \mathcal{O}\left(\left(\eta \beta \|H\|\right)^3\right)\right).
\end{equation}
To achieve an $\varepsilon$ accuracy in this approximation, we would have to set
\begin{equation}
    \eta = \frac{\varepsilon}{\beta \|H\| \|\sum_a A^{a\dagger}A^a\|},
\end{equation}
which in turn leads to the following subnormalization for the LCU implementation
\begin{equation} \label{eq:subnorm-eta-B}
    \|b_-\|_1 \|b_+^{(s,\eta)}\|_1 = \mathcal{O}\left(\log\frac{\beta \|H\| \|\sum_a A^{a\dagger}A^a\| }{\varepsilon}\right).
\end{equation}
Combining it all, we end up with a Hamiltonian simulation time required for the coherent term $B^{(s, \eta)}$ of $\mathcal{O}(\beta \log(1 / \varepsilon))$. This is a $\log$ factor more than what we have seen for the Gaussian case

We could try to use different regularization techniques, like the $\ii \varepsilon$ prescription from quantum field theory, i.e. to shift the pole away from the real time values to the imaginary ones, such that the integral of the time-domain would not run into it. In our case this would materialize by an exponentially decaying factor multiplying the weights of the linear combination, i.e. $\ee^{- \varepsilon \beta x} g(x)$. But while it does make both poles imaginary, avoiding Chen's regularization and its additional $\log$ costs, it does it at the price of an $e^{-\varepsilon \beta / 2}$ reduction in the transition weight near the kink, which degrades the spectral gap. As it turns out, this degradation would lead to higher costs asymptotically, than the $\eta$ regularization \eqref{eq:subnorm-eta-B}. Hence, our regularization has been tucked away in the appendix as a failed but hopefully constructive attempt \todo{ref}. For the main text, we will stick to the above form for $b_+^{(s)}$ \eqref{eq:b_plus-s-eta}.

\label{sec:quad-errors-diss}
\smallskip
\subsection{Quadrature errors for $\mathcal{L}_\text{diss}$.}
The next step towards implementation is the realization that every integral expression in the CKG Lindbladian \eqref{eq:ckg-L}, both the coherent term $B$ and the dissipative terms, have to be approximated via quadrature sums in the quantum computer. This introduces two types of errors that are interconnected with each other:
\begin{enumerate}[(i)]
    \item \emph{Truncation errors} occur due to the restriction of infinite integrals to a finite window $t\in[-T/2, T/2)$ and $\omega \in [-W/2, W/2)$. These errors are controlled by the decaying tail of the respective integrands.
    \item \emph{Discretization / Aliasing errors} occur since the integrals are replaced by a sum on the finite grids. These are controlled by the smoothness of the integrand in the given truncated domain.
\end{enumerate}
We will tackle these errors for the dissipative and coherent parts separately. First let us cast the CKG Lindbladian \eqref{eq:ckg-L} into its discretized form. The dissipative part becomes,
\begin{equation} \label{eq:discr-dissiative-L}
    \mathcal{\bar{L}}_\text{diss}(\cdot) = \sum_a \sum_{\bar{\omega} \in S_{\omega_0}} \gamma(\bar{\omega}) \hat{A}(\bar{\omega})(\cdot)\hat{A}(\bar{\omega})^\dagger - \frac{1}{2}\left\{\hat{A}(\bar{\omega})^\dagger \hat{A}(\bar{\omega}), \cdot\right\},
\end{equation}
with the discretized OFT, 
\begin{equation} \label{eq:oft-discr}
    \hat{A}(\bar{\omega}) = \frac{1}{\sqrt{N}}\sum_{\bar{t}\in S_{t_0}} \bar{f}(\bar{t})\ee^{-\ii \bar{\omega} \bar{t}} A(\bar{t}) =  \frac{1}{\sqrt{N}}\sum_{\bar{t}\in S_{t_0}} \bar{f}(\bar{t})\ee^{-\ii \bar{\omega} \bar{t}} \ee^{\ii H \bar{t}} A \ee^{-\ii H \bar{t}}.
\end{equation}
Here, we merged the Riemann-sum prefactor $t_0$ into the \emph{discretized} filter function $\bar{f}(t) := \sqrt{t_0} f(t)$. This lets us work with the natural Discrete Fourier Transform prefactor $1 / \sqrt{N}$, with $N$ the number of grid points. (It is also the reason why there is no explicit $\omega_0$ in \eqref{eq:discr-dissiative-L}.)
The infinite integrals are now replaced by discretized sums, where we will denote the discrete values by $\bar{t}, \bar{\omega}$ whose sums run over
\begin{equation*}
    \begin{split}
    S^{\lceil N \rfloor} &:= \left\{-\lceil(N-1)/2\rceil, \dots, -1, 0, 1, \dots, \lfloor (N - 1) / 2 \rfloor\right\} \\
    S^{\lceil N \rfloor}_{\omega_0} &:= \omega_0 S^{\lceil N \rfloor}, \quad S^{\lceil N \rfloor}_{t_0} := t_0 S^{\lceil N \rfloor},
\end{split}
\end{equation*}
with $\omega_0t_0 = \frac{2\pi}{N}$, where $N = 2^r$ is determined by the number of estimating qubits $r$ in the quantum circuit. Hence, the resolution can be tuned by choosing a larger $r$. But there is an interplay between time- and energy-domains due to the Fourier duality. Since we have integrals in both domains, the ideal scenario is when the integrands and their Fourier transforms are both well-behaving, i.e. smooth and rapidly decaying, functions.

In order to quantify the quadrature errors for the dissipative part of the Lindbladian (but also for the coherent part), we follow the strategy of \cite[Appendix C]{chen2023quantum}. All parts of the algorithm rely on the OFT, thus it is essential to understand how it behaves under discretization / truncation \eqref{eq:oft-discr}. The central difficulty here can be immediately revealed if we look at its decomposition $\hat{A}(\omega) = \sum_{\nu_\in B(H)} \hat{f}(\omega - \nu) A_\nu$. A sum over the Bohr frequencies $B(H)$ means, that a quadrature error in $\hat{f}$ would aggregate over all pairs of Bohr frequencies $|B(H)|^2$ -- a potentially exponentially large number. 

Chen et al. circumvent this problem by introducing a \emph{rounding Hamiltonian} $\bar{H}$ whose eigenvalues are the result of rounding each eigenvalue of $H$ to the nearest integer multiple of some $\eta >0$. Then $\|H - \bar{H}\| \leq \eta$ and $|B(\bar{H})| \leq 4 \|H\| / \eta + 1$. To be clear, this is only a proof technique, neither the quantum circuit nor our numerical simulations ever use $\bar{H}$; both work with the true (but rescaled to $\|H\| = 0.5$ for the QFT) Hamiltonian $H$.

Then, the total error between the continuous $\mathcal{L}_\text{diss}$ and the discretized $\mathcal{\bar{L}}_\text{diss}$ is bounded by the following triangle inequality,
\begin{equation} \label{eq:main-lindbladian-tri-ineq-discr}
    \begin{split}
    \|\mathcal{L}_{(f,H)} &- \bar{\mathcal{L}}_{(f,H)}\|_{1 \to 1}
\;\leq\; \\
&\underbrace{\|\mathcal{L}_{(f,H)} - \mathcal{L}_{(f,\bar{H})}\|_{1 \to 1}}_{\text{(i) continuous: }H \to \bar{H}}
+\;
\underbrace{\|\mathcal{L}_{(f,\bar{H})} - \bar{\mathcal{L}}_{(f,\bar{H})}\|_{1 \to 1}}_{\text{(ii) continuous} \leftrightarrow \text{discrete for }\bar{H}}
+\;
\underbrace{\|\bar{\mathcal{L}}_{(f,\bar{H})} - \bar{\mathcal{L}}_{(f,H)}\|_{1 \to 1}}_{\text{(iii) discrete: }\bar{H} \to H}.
\end{split}
\end{equation}
Terms (i) and (iii) are perturbation bounds comparing two Lindbladians that use the same grid but different Hamiltonians. They are controlled by \cite[Corollary A.1, A.2]{chen2023quantum},
\begin{equation}
    \text{(i), (iii)} \;\leq\; 
8\big(\|f - f_T\|_2 + T\|H - \bar{H}\|\big) 
\;\leq\; 8\big(\|f - f_T\|_2 + T\eta\big),
\end{equation}
where $f_T(t) := f(t) \mathbf{1}(t \in [-T/2, T/2))$ captures the truncation error. These terms are independent of the choice of $\gamma$, and depend only on the Gaussian filter tail and the rounding resolution.

The term (ii) is where the smoothness of $\gamma$ enters. By \cite[Lemma C.1]{chen2023quantum}, the operator error decomposes into $|B(\bar{H})|^2$ - many scalar integrals of the form
\begin{equation}
    \begin{split}
        &\left|\alpha_{\nu_1, \nu_2} - \bar{\alpha}_{\nu_1, \nu_2}\right| = \\
        &\left|\int_{-\infty}^{\infty} \gamma(\omega)\,\hat{f}^*(\omega - \nu_1)\,\hat{f}(\omega - \nu_2)\,\mathrm{d}\omega
\;-\;
\sum_{\bar{\omega} \in S_{\omega_0}} 
\tilde{\gamma}(\bar{\omega})\,
\bar{\mathcal{F}}\!\big[\bar{f}_T \cdot e^{i\nu_1\bar{t}}\big]^*\!(\bar{\omega})\;
\bar{\mathcal{F}}\!\big[\bar{f}_T \cdot e^{i\nu_2\bar{t}}\big](\bar{\omega})
        \right|,
    \end{split}
\end{equation}
with the discrete Fourier transform $\mathcal{\bar{F}}$.
In other words, the problem reduces to how well we can estimate the exact Kossakowski matrix $\alpha_{\nu_1, \nu_2}$ by quadrature sums. The norm of this error can be split further into three parts via another triangle inequality: frequency-tail truncation due to the $W$ window, i.e. $\gamma_W(\omega) := \gamma(\omega)\mathbf{1}(\omega \in [-W/2, W/2))$ in both the continuous and discretized settings; and the Riemann-sum error for the truncated integral. The convergence of the Riemann sum is governed by the regularity of the integrand 
\begin{equation}\label{eq:g-integrand}
    g_{\nu_1, \nu_2}(\omega) = \gamma(\omega)\hat{f}^*(\omega - \nu_2) \hat{f}(\omega - \nu_1).
\end{equation}
Since $\hat{f}$ is our smooth Gaussian filter, all comes down to $\gamma$ and its smoothness. 

We now go through each choice for the transition weight $\gamma$, clarifying what regularity each possesses and what quadrature error scaling it implies for their respective dissipative Lindbladians. In each case, we require the total discretization error to satisfy the following bound,
\begin{equation}
    \|\mathcal{L}^{(f,H)}_\text{diss} - \bar{\mathcal{L}}^{(f,H)}_\text{diss}\|_{1\to 1} \leq \varepsilon.
\end{equation}

\begin{enumerate}[(a)]
    \item \textbf{Gaussian weight $\gamma_G$}
    
    For the Gaussian weight, the integrand \eqref{eq:g-integrand} in (ii) is a product of Gaussians. Since the Gaussian is Gevrey-1/2 \cite[Proposition B.1]{hoepfner2019global} and self-dual under the Fourier transform, the Riemann sum converges exponentially in the number of grid points $N$. Likewise, the Gaussian filter tail error $\|f - f_T\|_2$ decays as $\mathcal{O}(e^{-T^2/2\sigma^2})$, and so does the frequency-domain tail \cite[Proposition A.9]{chen2023quantum}. Then, for the dissipative part of the Lindbladian to achieve an $\varepsilon$ precision, it requires 
    $$
    r = \mathcal{O}\!\left(\log\!\left(|A|(\|H\|{+}1)(\sigma{+}1/\sigma)(\beta{+}1)\right) + \log^2(1/\varepsilon)\right)
    $$
    estimating qubits, i.e. a polylogarithmic in $1/\varepsilon$ \cite[Corollary C.2]{chen2023quantum}. The dependence on $\log(1/\varepsilon)$ arises from the interplay between the truncation window $T = \Theta(\sigma \sqrt{\log(1/\varepsilon)})$, the frequency bandwidth $W = \Theta(\sigma^{-1}\sqrt{\log(1/\varepsilon)} + 2\|H\|)$, and the Fourier label condition on $\omega_0, t_0$. 
    
   That $\log(1/\varepsilon)$ is optimal follows from the classical uncertainty principles. Any filter achieving $\varepsilon$ truncation error on both time and frequency domains simultaneously, the product $T\cdot W$ scales at least as $\log(1 \varepsilon)$. That the Gaussian uniquely saturates this uncertainty bound is a classical result, see \cite{folland1997uncertainty}.
    
    \item \textbf{Metropolis weight $\gamma_M^{(s=0)}$}
    
    For the Metropolis weight $\gamma_M$ (with $s = 0$), the Lipschitz kink at $\omega = -\beta\sigma^2 / 2$, with Lipschitz constant $\beta$, limits $\gamma_M$ to $C^{0, 1}$ (in Hölder's notation), i.e. Lipschitz but not $C^1$. 
    
    By \cite[Proposition C.1]{chen2023quantum}, the Lipschitz constant $\beta$ enters the grid spacing constraint as
    $$
        \omega_0 = \mathcal{O}\!\left(\frac{\varepsilon}{\beta\, T\, |B(\bar{H})|^2}\right),
    $$
    which forces the number of estimating qubits to scale polynomially in $1 / varepsilon$ \cite[Corollary C.2]{chen2023quantum}: 
    $$
        r = \mathcal{O}\left(\log\tfrac{|A|(\|H\|+1)(\sigma + 1/\sigma)(\beta+1)}{\varepsilon}\right).
    $$
    Thus, the kink causes a qualitative degradation: polynomial many estimating qubits (rather than polylogarithmic) are needed to achieve the same precision.
    
    To see the mechanism concretely: after splitting the integral at $\omega_\text{kink} = -\beta\sigma^2/2$ and applying the Euler-Maclaurin formula on each half, the trapezoidal error for the $(\nu_1, \nu_2)$ Kossakowski entry is 
    $$
       \frac{\omega_0^2}{12}\,|\Delta g'| = \frac{\beta\,\omega_0^2}{12}\,|\hat{f}(\omega_{\mathrm{kink}} - \nu_1)|\hat{f}(\omega_{\mathrm{kink}} - \nu_2)|, 
    $$
    Summing over rounded Bohr-frequency pairs via \cite[Lemma C.1]{chen2023quantum} and using that $\sum_\nu \|A_\nu\|^2 \leq 1$ together with $|\hat{f}| \leq \|\hat{f}\|_\infty$ this aggregates to a total quadrature error of $\mathcal{O}(\beta \omega_0^2)$. Every other error contributions in the triangle inequalities (e.g. Gaussian aliasing, time truncation, spectral leakage from off-grid Bohr frequencies) otherwise converges exponentially in $1 / \omega_0$, which makes this algebraic quadrature error scaling stick out.

    \item \textbf{Metropolis weight $\gamma_M^{(s > 0)}$}
    
    In \hyperref[prop:smooth-metro-gevrey]{Proposition~\ref*{prop:smooth-metro-gevrey}}, we have shown that the smooth Metropolis function with $s > 0$ belongs to the Gevrey-1/2 class, the same regularity class as the Gaussian transition weights $\gamma_G$, albeit with different Gevrey constants $C_1, C_2$.
    
    An important property of Gevrey functions, is that their product is also Gevrey \cite[Lemma 26]{ding2025efficient}. Consequently, the Kossakowski integrand $g_{\nu_1, \nu_2}$ as a product of Gevrey-1/2 functions -- smooth Metropolis with $C_{1,2}^{(\gamma)}$ and Gaussian filters $\hat{f}$ with $C_{1,2}^{(\hat{f})}$ -- satisfies
    \begin{equation} \label{eq:gevrey-product-integrand}
        \begin{split}
            \hat{g}_{\nu_1,\nu_2} \in G^{1/2}_{C_1^{(\hat{g})},\, C_2^{(\hat{g})}} \qquad\text{with}& \\
            C_2^{(\hat{g})} = C_2^{(\gamma)} + 2\,C_2^{(\hat{f})} = \frac{1}{\sigma\sqrt{\ee\,s}} + \frac{\sqrt{2}}{\sigma \sqrt{\ee}} \quad\text{and}&\quad C_1^{(\hat{g})} = C_1^{(\gamma)}\left(C_1^{(\hat{f})}\right)^2.
        \end{split}
    \end{equation}
    Where $C_1^{(\hat{g})}$ is just the squared normalization of the Gaussian filters $\hat{f}$, since $C_1^{(\gamma)} = 1$.

    The trapezoidal quadrature error for this integrand is governed by the Poisson summation formula \cite[Chapter 5]{trefethen2014exponentially}, specifically for $g \in C^\infty(\mathbb{R}) \cap L^1(\mathbb{R})$ with all derivatives in $L^1(\mathbb{R})$ is
    \begin{equation}
        E_{\nu_1,\nu_2} := \omega_0 \sum_{n \in \mathbb{Z}} \hat{g}_{\nu_1,\nu_2}(n\omega_0) - \int_{\mathbb{R}} \hat{g}_{\nu_1,\nu_2}(\omega)\,\dd\omega = \sum_{k \neq 0} g_{\nu_1,\nu_2}\!\left(\frac{2\pi k}{\omega_0}\right).
    \end{equation}
    Bounding the Fourier transform $|g(\xi)|$ via $2K$-fold integration by parts gives $|g(\xi)| \leq |\xi|^{-2K}\,\|\hat{g}^{(2K)}\|_{L^1}$, hence
    \begin{equation} \label{eq:remainder-g-discr-1}
        |E_{\nu_1,\nu_2}| \leq 2\sum_{k=1}^{\infty}\left(\frac{\omega_0}{2\pi k}\right)^{2K}\|\hat{g}^{(2K)}\|_{L^1} \leq 4\left(\frac{\omega_0}{2\pi}\right)^{2K}\|\hat{g}^{(2K)}\|_{L^1},
    \end{equation}
    where we used $\sum_{k=1}^\infty k^{-2K} \leq \zeta(2) \leq 2$ for $K \geq 1$. 

    But in its current form with $L^1$ we cannot use Gevrey properties yet which are formulated in $L^\infty$. Since the Gaussian filter factor confine the integrand $\hat{g}_{\nu_1, \nu_2}$ to an effective support of width $W_\text{eff} = \mathcal{O}(\sigma)$ (which is $K$ independent), we have $\|\hat{g}^{(2K)}\|_{L^1} \leq W_\mathrm{eff}\,\|\hat{g}^{(2K)}\|_{L^\infty}$. The Gevrey-1/2 bound from $g \in \mathcal{G}^{1/2}_{C_1^{(\hat{g})}, C_2^{(\hat{g})}}$ gives 
    $$
       \|\hat{g}^{(2K)}\|_{L^\infty} \leq C_1^{(\hat{g})}\,(C_2^{(\hat{g})})^{2K}\,(2K)^{K}. 
    $$
    Combining it with \eqref{eq:remainder-g-discr-1}, 
    \begin{equation}
        |E_{\nu_1,\nu_2}| \leq 4\,W_\mathrm{eff}\,C_1^{(\hat{g})}\left(\frac{C_2^{(\hat{g})}\,\omega_0}{2\pi}\right)^{2K}(2K)^K.
    \end{equation}
    Then, we can optimize in $K$ by the standard saddle-point method by first setting 
    $$
        \Phi(K) := 2K\ln\!\left(\frac{C_2^{(g)}\omega_0}{2\pi}\right) + K\ln(2K),
    $$
    and then using the condition $\Phi'(K^*) = 0$ leads to 
    $$
        2\ln\!\left(\frac{C_2^{(g)}\omega_0}{2\pi}\right) + \ln(2K^*) + 1 = 0 \qquad\Longrightarrow\qquad 2K^* = \frac{1}{e}\!\left(\frac{2\pi}{C_2^{(g)}\omega_0}\right)^{\!2}.
    $$
    Inserting $K^*$ back we find that the optimum evaluates to
    $$
        \Phi(K^*) = -K^* = -\frac{2\pi^2}{e\,(C_2^{(g)}\omega_0)^2}.
    $$
    Hence, the per $(\nu_1, \nu_2)$ entry Riemann-sum error satisfies:
    \begin{equation}\label{eq:per-entry-error-smooth-metro}
    |E_{\nu_1,\nu_2}| = \mathcal{O}\!\left(\exp\!\left(-\frac{2\pi^2}{e\,(C_2^{(g)}\,\omega_0)^2}\right)\right).
    \end{equation}

    In the regime $s \ll 1$ (close to the Metropolis limit), the transition weight contribution $C_2^{(\gamma)} = 1/(\sigma\sqrt{es})$ dominates the filter contribution $2C_2^{(\hat{f})} = \sqrt{2}/(\sigma\sqrt{e})$, since $C_2^{(\gamma)}/2C_2^{(\hat{f})} = 1/(\sqrt{2s}) \gg 1$. In this regime $C_2^{(g)} \approx C_2^{(\gamma)}$ and the bound simplifies to
    \begin{equation}\label{eq:smooth-metro-dominant-regime}
        |E_{\nu_1,\nu_2}| = \mathcal{O}\!\left(\exp\!\left(-\frac{2\pi^2\,s\,\sigma^2}{\omega_0^2}\right)\right),
    \end{equation}
    where the factor $e$ from $C_2^{(\gamma)} = 1/(\sigma\sqrt{es})$ cancels the $1/e$ from the saddle-point optimization. 

    For the general regime where both $C_2^{(\gamma)}$ and $C_2^{(\hat{f})}$ contribute comparably, the full additive constant $C_2^{(g)} = C_2^{(\gamma)} + 2C_2^{(\hat{f})}$ must be retained. However, in general we have found that the smooth Metropolis requires the same polylogarithmic estimating qubit scaling as the Gaussian case (a), 
    \begin{equation}
        r = \mathcal{O}\!\left(\log\!\left(|A|(\|H\|{+}1)(\sigma{+}1/\sigma)(\beta{+}1)\right) + \frac{1}{s}\log^2(1/\varepsilon)\right).
    \end{equation}

    \begin{table}[ht]
    \centering
    \begin{tabular}{lll}
    \toprule
    Transition weight & Regularity & Quadrature error in $\omega_0$ \\
    \midrule
    Gaussian $\gamma_G$ & Gevrey-$1/2$ & $\mathcal{O}(e^{-c_G/\omega_0^2})$ — Gaussian \\
    Smooth Metropolis $\gamma_M^{(s)}$, $s > 0$ & Gevrey-$1/2$ & $\mathcal{O}(e^{-c_s/\omega_0^2})$ — Gaussian ($c_s \leq c_G$) \\
    Metropolis $\gamma_M$ ($s = 0$) & Lipschitz, not $C^1$ & $\mathcal{O}(\beta\,\omega_0^2)$ — algebraic \\
    \bottomrule
    \end{tabular}
    \caption{Convergence hierarchy of the trapezoidal quadrature error for the three transition weights.}
    \label{tab:quadrature-hierarchy}
    \end{table}
\end{enumerate}

The above discussions can be summarized in a table of the convergence hierarchy \hyperref[tab:quadrature-hierarchy]{Table~\ref*{tab:quadrature-hierarchy}}. The Gaussian and smooth Metropolis fall in the same qualitative convergence class. However, the Gaussian case with $\gamma_G$ has an optimal constant $c_G$ because the Kossakowski integrand $\hat{g}_{\nu_1, \nu_2} = \gamma_G \cdot |\hat{f}|^2$ is itself a Gaussian, whose Poisson summation aliases are exactly Gaussian and therefore optimally localized. For the smooth Metropolis, the integrand is note a pure Gaussian, so $c_s$ is generally smaller. A qualitative difference can be found between the kinky Metropolis $\gamma_M$ with $s=0$ and its smooth variant: the latter converges super-algebraically (restoring polylogarithmic qubit scaling), while the former is limited to an algebraic convergence due to its derivative discontinuity. Thus, for the dissipative Lindbladian, the smooth Metropolis $\gamma_M^{(s)}$ with $s > 0$ needs less amount of estimating qubits in general, while maintaining the broad Metropolis-like transition window, which keeps the spectral gap (and thus mixing time) favourable.

% Remark (comparison with DLL24). Ding, Li, and Lin [DLL24, Proposition 16] establish an analogous exponential quadrature rate $\exp(-c'((M{-}1)\tau)^{1/s}/(C\,e))$ for their Lindbladian, using compactly supported Gevrey filters and the Poisson summation formula for exact alias cancellation [DLL24, Lemma 31]. Their proof does not use Euler–Maclaurin optimization; instead, compact support gives exact cancellation of all aliases within the support bandwidth, so only the Paley–Wiener tail of the Fourier-domain kernel contributes [DLL24, Lemma 29]. Our setting differs: the Gaussian filter $\hat{f}$ is not compactly supported, so the DLL24 alias-cancellation technique yields no asymptotic gain over the CKBG23 rounding approach for our case ($a = 0$). However, both routes yield the same qualitative conclusion — Gevrey-$1/2$ regularity implies super-algebraic (Gaussian-type) convergence in the grid spacing.

The following remark has to be made on the frequency register size $N$. In the algorithm from Chen et al. the DFT gird parameters $N, \omega_0, t_0$ are set by: (i) the frequency range $N\omega_0/2 \geq 2\|H\| + \mathcal{O}(\sigma^{-1}\sqrt{\log(1/\varepsilon)})$, (ii) the Gaussian truncation in time $Nt_0/2 \geq \mathcal{O}(\sigma\sqrt{\log(1/\varepsilon)})$, and (iii) the DFT aliasing condition \cite[Proposition A.7]{chen2023quantum}. None of these conditions involve $\gamma$. This is to say, that for analytically prescribed (fine) grids where all error sources are already exponentially small, the number of frequency-register qubits $\lceil\log_2 N\rceil$ is the same for all three transition weights. However, in practice quadrature errors can be slightly below the algorithmic errors e.g. prescribed by $\delta$ in the weak measurement scheme, and yet achieve good results. If we operate on such a coarser grid (larger $\omega_0$, fewer qubits), then the exponential terms remain small while $\mathcal{O}(\beta\omega_0^2)$ dominates for the Metropolis case, making the kink the actual bottleneck. Smoothing with $s > 0$ removes this bottleneck entirely.

\todo{we might needs different structure than this paragraph / subsection..}
\label{sec:quad-errors-B}
\smallskip
\subsection{Quadrature errors for $B$.}
The remaining part in the CKG Lindbladian \eqref{eq:ckg-L} that has to be discretized is the coherent term $B$, which takes the following form
\begin{equation} \label{eq:B-discr}
    \bar{B} = \sum_a t_0\sum_{\bar{t}\in S_{t_0}} b_-(\bar{t}) \ee^{-\ii H \bar{t} / \sigma} \left[t_0'\sum_{\bar{t}' \in S_{t_0'}} b_+(\bar{t}') A^{a\dagger}(\beta \bar{t}') A^a(-\beta \bar{t}')\right]\ee^{\ii H \bar{t} / \sigma}.
\end{equation}
Note, that for the coherent term we keep the Riemann-sum factor explicit (since here there is no Fourier factors present). However, in quantum computer, the time registers have to then encode e.g. $\sqrt{t_0 b_-(\bar{t})} / \sqrt{\alpha_-}$, where $\alpha_- = t_0 \sum_{\bar{t}} |b_-(\bar{t})| \approx \|b_-\|_1$ is the subnormlaization of $b_-$.
Due to the nested integral form (which arose from the factorizing Gaussian filters), the quadrature error $\|B - \bar{B}\|$ decomposes into independent contributions from the outer ($b_-$) and inner ($b_+$) integrals via a triangle inequality. Concretely, define the continuous inner operator,
\begin{equation}
    I_+^a := \int_{-\infty}^{\infty} b_+(t')\, A^{a\dagger}(\beta t')\, A^a(-\beta t')\,\dd t',
\end{equation}
whose discretized counterpart is $\bar{I}_+^a$, and its outer conjugated operator $O_a(t) := \ee^{-\ii H t / \sigma} I_+^a \ee^{\ii H t}$. Then, the quadrature error for the full nested integral can be written as
\begin{equation} \label{eq:B-error-factorize}
    \begin{split}
        &\|B - \bar{B}\| \;\leq\; \\
    &\underbrace{\sum_a\left\|\int_\mathbb{R} b_-(t)\,O_a(t)\,dt - t_0\sum_{\bar{t}\in S_{t_0}} b_-(\bar{t})\,O_a(\bar{t})\right\|}_{\text{outer quadrature error (with exact } I_+^a\text{)}}
    \;+\;
    \underbrace{\|b_-\|_1 \cdot \sum_a \|I_+^a - \bar{I}_+^a\|.}_{\text{inner quadrature error, weighted by } \|b_-\|_1}
    \end{split}
\end{equation}
Although the errors factorize, a coupling between them arises nevertheless: 
the outer Riemann-sum error involves $\|I^a_+\| \leq \|b_+\|_1 \|A^a\|^2 \leq \|b_+\|_1$ through the norm of the operator being conjugated via $O_a$. Thus, the outer grid requirement inherits the subnormalization from the inner integral. 

We proceed by analysing the integrals one by one and then combine the results.

\medskip
\noindent\textbf{Outer integral ($b_-$).}
The kernel $b_-(t)$ is universal across all transition weight choices $\gamma$ and takes the form \eqref{eq:b_minus}. It is a convolution of two entire functions: $1/\cosh(2\pi t / (\beta\sigma))$, which decays exponentially as $\sim 2\ee^{-2\pi |t| / (\beta\sigma)}$, and $\sin(-\beta\sigma t)\,\ee^{-2t^2}$, which decays as a Gaussian. The convolution inherits the slower (exponential) decay  rate, giving and effective width $\mathcal{O}(\beta\sigma)$ as noted in \cite[footnote 13]{chen2023efficient}.

The outer integral involves Hamiltonian evolution $\ee^{\pm \ii H \bar{t} / \sigma}$, requiring simulation time of $\mathcal{O}(|t| / \sigma)$. Truncating to $|t| \leq T_-  / 2$ introduces a tail error that decays exponentially:
\begin{equation}
    \text{(outer truncation error)} = \mathcal{O}\!\left(\ee^{-c_-\, T_- / (\beta\sigma)}\right),
\end{equation}
for a constant $c_- > 0$ that depends on the $1/\cosh$ decay rate. Setting $T_- = \Theta(\beta\sigma\log(1/\varepsilon))$ controls this error to $\varepsilon$, corresponding to a Hamiltonian simulation time of $\mathcal{O}(\beta\log(1/\varepsilon))$ in the outer integral.

For the Riemann-sum error, we apply the operator-valued discretization bound from \cite[Lemma III.1]{chen2023efficient}:
\begin{equation}\label{eq:riemann-sum-operator}
    \left\|\int_{-T/2}^{T/2} O(t)\,f(t)\,\dd t - \sum_{\bar{t}\in S_{t_0}} O(\bar{t})\,f(\bar{t})\,t_0\right\| \;\leq\; \frac{T^2}{2|S_{t_0}|}\!\left(\|[O, H_\text{eff}]\|\sup_t|f(t)| + \|O\|\sup_t|f'(t)|\right),
\end{equation}
with $O(t) = \ee^{-\ii H_\text{eff} t} \,O\, \ee^{\ii H_\text{eff} t}$. Specifically, for the outer integral we have $f = b_-$, $H_\text{eff} = H / \sigma$ and $O = \bar{I}^a_+$, i.e. the discretized inner block. So the outer operator-norm error $\delta_-^a$ for each $a$ is:
\begin{equation}
    \delta^a_- \;\leq\; \frac{T_-^2}{2N_-}\left( \left\|[\bar{I}^a_+, H / \sigma]\right\|\|b_-\|_\infty + \|\bar{I}^a_+ \|\; \|b_-'\|_\infty\right).
\end{equation}
Using the commutator bound $\left\|[\bar{I}^a_+, H / \sigma]\right\|\leq 2 \|H\|\; \|\bar{I}^a_+\| / \sigma$, and the norm bound of $\|\bar{I}^a_+\|$ we find the full discretization error for the outer integral after summing over $a$: 
\begin{equation}
    \delta_- \leq \frac{T_-^2 |\mathcal{A}| \|b_+\|_1}{2 N_-} \left(\frac{2\|H\|}{\sigma}\|b_-\|_\infty + \|b_-'\|_\infty\right).
\end{equation}
Setting $\delta \leq \varepsilon$ with $T_- = \Theta(\beta\sigma\log(1/\varepsilon))$:
\begin{equation} \label{eq:r_-}
    r_- = \mathcal{O}\!\left(\log\!\left(\frac{\beta\sigma\|H\|\,|\mathcal{A}|\,\|b_+\|_1 \log(1/\varepsilon)}{\varepsilon}\right)\right).
\end{equation}
For the Gaussian transition weight $\|b_+\|_1 = \mathcal{O}(1)$, and thus the inner integral does not lead to additional costs. However, for the Metropolis-like transitions (both smooth and non-smooth) we have $\|b_+\|_1 = \mathcal{O}(\log(\beta\|H\|/\varepsilon))$, contributing as an additive $\log\log(\beta\|H\|/\varepsilon)$ term to $r_-$ (which is indeed not the end of the world).

\medskip
\noindent\textbf{Inner integral ($b_+$).}
Before discussing the different cases, we first have to establish the Riemann-sum bound for the inner integral \eqref{eq:B-discr}. Since 
$$
A^{a\dagger}(\beta t')\,A^a(-\beta t') = e^{i\beta H t'}\,A^{a\dagger}\,e^{-2i\beta H t'}\,A^a\,e^{i\beta H t'}
$$
is not of the Heisenberg conjugation form we cannot directly apply \cite[Lemma III.1]{chen2023efficient} like before. Instead, we use the general scalar Riemann-sum bound
\begin{equation}
    \begin{split}
        &\left\|\int_{-T_+/2}^{T_+/2} b_+(t')A^{a\dagger}(\beta t')\,A^a(-\beta t')\,dt' - \sum_{S_{t_0'}} b_+(\bar{t}')A^{a\dagger}(\beta \bar{t}')\,A^a(-\beta \bar{t}')\,t_0'\right\| \\
        &\leq \frac{T_+^2}{2N_+}\sup_{t'}\left\|\left(b_+(t')A^{a\dagger}(\beta t')\,A^a(-\beta t')\right)'(t')\right\|.
    \end{split}
\end{equation}
Where the norm of the derivative can be upper-bounded by the derivative of the parts. Specifically if we use that $\|A^{a\dagger}(\beta t')\,A^a(-\beta t')\| \leq \|A^a\|^2 \leq 1$ and that via the Leibniz rule we get
\begin{equation}
    \left\|\left(A^{a\dagger}(\beta t')\,A^a(-\beta t')\right)'(t')\right\| \leq 4\beta\|H\|\,\|A^a\|^2,
\end{equation}
then combining everything leads to the bound, 
\begin{equation}
    \left\|\left(b_+(t')A^{a\dagger}(\beta t')\,A^a(-\beta t')\right)'(t')\right\| \leq \|A^a\|^2\left(|b_+'(t')| + 4\beta\|H\|\,|b_+(t')|\right).
\end{equation}
This yields a per-$a$ discretization error $\delta_+^a$:
\begin{equation}
    \delta_+^a \leq \frac{T_+^2\,\|A^a\|^2}{2N_+}\sup_{t'}\!\left(|b_+'(t')| + 4\beta\|H\|\,|b_+(t')|\right).
\end{equation}
Summing over $a$:
\begin{equation} \label{eq:inner-discr-error}
    \delta_+ \leq \frac{T_+^2\,|A|}{2N_+}\sup_{t'}\!\left(|b_+'(t')| + 4\beta\|H\|\,|b_+(t')|\right).
\end{equation}
\begin{enumerate}[(a)]
    \item \textbf{Gaussian case}
    
    For the Gaussian weight $\gamma_G$, the time-domain kernel is given by a Gaussian \eqref{eq:b_plus-gauss} (surprise!) with width $\mathcal{O}(1)$. Hence, $\|b_+^{(G)}\|_\infty, \|(b_+^{(G)})'\|_\infty = \mathcal{O}(1)$.

    Correspondingly, the Gaussian decay gives 
    $$
    \text{(inner truncation)} = \mathcal{O}\!\left(e^{-c_G\, T_+^2}\right),
    $$
    so $T_+ = \Theta(\sqrt{\log(1/\varepsilon)})$ suffices. The Hamiltonian simulation time for the inner integral is $\beta T_+ = \Theta(\beta\sqrt{\log(1/\varepsilon)})$.

    Inserting $\|b_+^{(G)}\|_\infty, \|(b_+^{(G)})'\|_\infty$ into the general discretization error bound above \eqref{eq:inner-discr-error}:
    \begin{equation}
        \delta_+^{(G)} \leq \frac{T_+^2\,|\mathcal{A}|}{2N_+}\,\mathcal{O}(\beta\|H\|).
    \end{equation}
    Setting $\delta^{(G)}_+ \leq \varepsilon$, 
    \begin{equation} \label{eq:r_+_gaussian}
        r_+^{(G)} = \mathcal{O}\!\left(\log\!\left(\frac{\beta\|H\|\,|\mathcal{A}|\,\log(1/\varepsilon)}{\varepsilon}\right)\right).
    \end{equation}
    \item \textbf{Metropolis cases $(s \geq 0)$}
    
    Unlike for the dissipative case $\mathcal{\bar{L}}_\text{diss}$, here we find no qualitative difference between the smooth or kinky Metropolis variants, since both of them require the same $\eta$-regularization in the time-domain, which becomes the dominant factor. 

    However, for completeness and also to ground all of our later numerical results via analytic forms, let us derive the necessary quantities here.

    For $|t'|$ large enough that the indicator $\mathbb{I}(|t'|\leq\eta)$ is inactive, the kernel $b_+^{(s,\eta)}(t')$ is dominated by:
    $$
    |b_+^{(s,\eta)}(t')| \leq \frac{1}{2\sqrt{2}\pi^2}\,\frac{e^{-2\sigma^2\beta^2(1+s)\,t'^2}}{|t'|\,|2t'+i|} \quad\text{for } |t'| > \eta.
    $$
    The Gaussian decay rate $2\sigma^2\beta^2(1+s)$ increases with $s$, narrowing the effective support:
    \begin{equation} \label{eq:T_+_metro}
        T_+ = \Theta\!\left(\frac{\sqrt{\log(1/\varepsilon)}}{\sigma\beta\sqrt{1+s}}\right).
    \end{equation}

    Then, the required Hamiltonian simulation time for the inner integral in \eqref{eq:B-discr} is $\beta T_+$.

    The indicator function $\mathbb{I}(|t'|\leq\eta)$ in $b_+^{(s,\eta)}$ introduces jump discontinuities at $|t'| = \eta$. Since the Riemann-sum bound requires smoothness of the integrand, we split the inner integral at $t' = \pm\eta$ into three pieces:

    \textit{Central piece $[-\eta, \eta]$}: On the interval regularization removes the $1 / t'$ singularity by design, making $b_+^{(s,\eta)}$ bounded near the origin. Both $b_+^{(s,\eta)}$ and its derivatives are $\mathcal{O}(1)$ on $[-\eta,\eta]$. With a uniform grid of step size $t_0' = T_+/N_+$, the number of grid points in $[-\eta,\eta]$ is $\sim 2\eta N_+/T_+$, and the Riemann-sum error from this piece is
    $$
    \mathcal{O}\!\left(\frac{(2\eta)^2}{2\cdot (2\eta N_+/T_+)}\cdot \mathcal{O}(\beta\|H\|)\right) = \mathcal{O}\!\left(\frac{\eta T_+\beta\|H\|}{N_+}\right),
    $$
    which is a factor $\sim \eta/T_+$ of the outer-piece error and therefore negligible.

    \textit{Outer pieces $[-T_+/2, -\eta] \cup [\eta, T_+/2]$}: On these intervals $b_+^{(s,\eta)}(t')$ is smooth, with the Gaussian envelope providing rapid decay. The kernel has the scaling $|b_+^{(s,\eta)}(\eta^+)| = \mathcal{O}(1/\eta)$, while its derivative satisfies $\|(b_+^{(s,\eta)})'\|_{L^\infty(|t'|>\eta)} = \mathcal{O}(1/\eta^2)$ near $t' = \eta$ with Gaussian decay for larger $|t'|$. Applying the general discretization bound \eqref{eq:inner-discr-error} to the outer pieces the dominant contribution (the derivative) gives:
    \begin{equation}
        \delta_+^a \leq \frac{T_+^2\,\|A^a\|^2}{2N_+}\,\mathcal{O}\!\left(\frac{1}{\eta^2}\right).
    \end{equation}
    Summing over $a$, substituting $T_+$ and setting $\eta = \varepsilon/(\beta\|H\|\|\sum_a A^{a\dagger}A^a\|)$ leads to the scaling of the necessary estimating qubits for the inner integral of $B$:
    \begin{equation} \label{eq:r_+_metro}
        r_+^{(s)} = \mathcal{O}\!\left(\log\!\left(\frac{\|H\|\left\|\sum_a A^{a\dagger}A^a\right\| |\mathcal{A}|\,\log(1/\varepsilon)}{\sigma(1+s)\,\varepsilon}\right)\right).
    \end{equation}
    This is a polylogarithmic in $1 / \varepsilon$ for all $s \geq 0$. The $s$-dependence appears only as a mild $\log(1 + s)$ reduction inside the algorithm.
\end{enumerate}
As we will see later in the algorithm section \todo{ref}, to implement the block encoding of the coherent term $B$, we need two time registers to compute the nested integrals in \eqref{eq:B-discr}. Accordingly, the total number of estimating qubits is the sum of the above, separately derived expressions, $r = r_- + r_+$. At this discretization / estimating-qubit level, the outer $r_-$ and inner $r_+$ integrals contribute essentially the same way. Where it will matter more that the Metropolis cases had to be regularized with $\eta$ is in their implementation via the LCU method. There the complexity depends on the subnormalizations for which $\|b^{(s)}_+\|_1$ indeed leads to additional costs. \todo{ref}

% Comment on Ding's Poission technique that works for smooth variants but not for time singular Metropolis functions

\smallskip
\subsection{Hamiltonian Trotterization.} 

\todo{could add a comment about childs constants for Heis or Isin}
\todo{Add that discussion into the Trotter prelims}
The CKG Lindbladian \eqref{eq:ckg-L} is abundant with exact Hamiltonian time evolutions $\ee^{\pm \ii H t}$. Any implementation on a quantum computer must approximate these using some product formula when $H = \sum_{k = 1}^K H_k$ has non-commuting terms. With this subsection we have managed to reach the deepest point of the dark, error-analysis-cave, where we quantify the incurring errors due to the replacement of $e^{\pm iH\bar{t}}$ in the discretized OFT and coherent term with a product formula $S_p^{(M)}(\bar{t}) := S_p(\bar{t}/M)^M$ of order $p$ with $M$ Trotter steps. 

In \hyperref[sec:prelim-trotter]{Section~\ref*{sec:prelim-trotter}} the general Trotterization strategy and the approximation errors they lead to have been introduced using \cite{childs2021theory} as the main reference. Here, we make use of those results, namely that by \cite[Theorem 6]{childs2021theory} a $p$-th order product formula $S_p^{(M)}(t)$ for $\ee^{\ii Ht}$ satisfies the commutator-scaling bound:
\begin{equation} \label{eq:childs-trotter-bound}
    \bigl\|S_p^{(M)}(t) - e^{iHt}\bigr\|
    = \mathcal{O}\!\left(\frac{\tilde{\alpha}_{\mathrm{comm}}\, |t|^{p+1}}{M^p}\right),
\end{equation}
where
$$
\tilde{\alpha}_{\mathrm{comm}}
:= \sum_{k_1, k_2, \ldots, k_{p+1} = 1}^{K}
\bigl\|\bigl[H_{k_{p+1}}, \cdots [H_{k_2}, H_{k_1}]\cdots\bigr]\bigr\|
$$
is the commutator-scaling constant. With $M$ Trotter steps of size $\delta t = t / M$, the total error is $\|S_p^{(M)}(t) - e^{iHt}\| \leq \varepsilon_p(t, M)$ with $\varepsilon_p(t, M) = O(\tilde{\alpha}_{\mathrm{comm}}\, |t|^{p+1}/M^p)$. 

We use the following notation throughout: $\hat{A}(\bar{\omega})$ is the discretized OFT operator with \emph{exact} Hamiltonian evolution; $\tilde{A}(\bar{\omega})$ is the same object with every $\ee^{\pm \ii H\bar{t}}$ replaced by $S_p(\pm \bar{t})$. Similarly, $\mathcal{\bar{L}}_\text{diss}$ is built from $\hat{A}(\bar{\omega})$, while $\mathcal{\tilde{L}}_\text{diss}$ is its Trotterized counterpart built from $\tilde{A}(\bar{\omega})$.

First, we tackle the dissipative part of the Lindbladian, i.e. the Trotter errors propagating from the OFTs \eqref{eq:oft-discr} within. The Trotterized version replaces $\ee^{\ii H\bar{t}} A_a \ee^{-\ii H\bar{t}}$ with $S_p^{(M)}(\bar{t})\, A_a\, S_p^{(M)}(\bar{t})^\dagger$:
\begin{equation} \label{eq:trotter-oft}
    \tilde{A}_a(\bar{\omega})
    = \frac{1}{\sqrt{N}} \sum_{\bar{t} \in S_{t_0}} \ee^{-\ii\bar{\omega}\bar{t}}\, \bar{f}(\bar{t})\, S_p^{(M)}(\bar{t})\, A_a\, S_p^{(M)}(\bar{t})^\dagger.
\end{equation}

\begin{proposition} \label{prop:trotter-diss}
    \textup{(Trotter error for $\mathcal{L}_\textup{diss}$).} Let $H = \sum_{k=1}^K H_k$ with non-commuting terms, and let $\{A_a\}_{a \in \mathcal{A}}$ satisfy the normalization $\|\sum_a A_a^\dagger A_a\| \leq 1$. Let $\bar{f} : S_{t_0} \to \mathbb{R}$ be some discretized filter function with $\|\bar{f}\|_2^2 := \sum_{\bar{t} \in S_{t_0}} |\bar{f}(\bar{t})|^2 = 1$, and let the transition weight $\gamma$ satisfy $\|\gamma\|_\infty \leq 1$. With the Trotter product formula defined by \eqref{eq:childs-trotter-bound} and the corresponding Trotterized OFT \eqref{eq:trotter-oft} the total incurring Trotter error for the $\mathcal{\bar{L}}_\textup{diss}$ can be bounded as:
    $$
    \bigl\|\bar{\mathcal{L}}_{\mathrm{diss}} - \tilde{\mathcal{L}}_{\mathrm{diss}}\bigr\|_{1 \to 1}
    \leq 8\,\sqrt{\sum_{\bar{t} \in S_{t_0}} |\bar{f}(\bar{t})|^2\, \varepsilon_p^2(\bar{t}, M)},
    $$
    where $\varepsilon_p(\bar{t}, M) := \|S_p^{(M)}(\bar{t}) - \ee^{\ii H\bar{t}}\|$ is the product-formula error at time $\bar{t}$.

    For the second-order Strang splitting ($p = 2$), define the explicit commutator-scaling constant
    $$
    \tilde{\alpha}_{\mathrm{comm}}^{(2)}
    := \sum_{k_1 < k_2}
    \left(
      \frac{1}{12}\bigl\|[H_{k_2},[H_{k_2},H_{k_1}]]\bigr\|
    + \frac{1}{24}\bigl\|[H_{k_1},[H_{k_1},H_{k_2}]]\bigr\|
    \right),
    $$
    which is the leading-order prefactor from the tight second-order analysis of \cite[Proposition 10]{childs2021theory}. Then at leading order in $1 / M$:
    $$
    \bigl\|\bar{\mathcal{L}}_{\mathrm{diss}}
          - \tilde{\mathcal{L}}_{\mathrm{diss}}\bigr\|_{1 \to 1}
    \;\leq\;
    \frac{\sqrt{15}\;\tilde{\alpha}_{\mathrm{comm}}^{(2)}}
         {M^2\,\sigma^3}
    \;+\; \mathcal{O}\!\left(\frac{1}{M^3\,\sigma^4}\right).
    $$
\end{proposition}
\begin{proof}
    \noindent\textbf{(i) Bound $\bigl\|\bar{\mathcal{L}}_{\mathrm{diss}} - \tilde{\mathcal{L}}_{\mathrm{diss}}\bigr\|_{1 \to 1}$.}
    Define the Stinespring isometries with the discretized OFT \eqref{eq:oft-discr} and Trotterized OFT \eqref{eq:trotter-oft}, 
    $$
    \bar{V} := \sum_{a,\bar{\omega}} \sqrt{\gamma(\bar{\omega})}\, |\bar{\omega}, a\rangle\langle \bar{0}|\, \hat{A}_a(\bar{\omega}), \qquad
    \tilde{V} := \sum_{a,\bar{\omega}} \sqrt{\gamma(\bar{\omega})}\, |\bar{\omega}, a\rangle\langle \bar{0}|\, \tilde{A}_a(\bar{\omega}).
    $$
    Both $\mathcal{\bar{L}}_\text{diss}$ and $\mathcal{\tilde{L}}_\text{diss}$ admit the Stinespring factorization 
    $$
    \mathcal{L}[{\rho}] = \tr_{\mathrm{anc}}(V \rho\, V^\dagger) - \frac{1}{2}\tr_{\bar{0}}\{V^\dagger V, \rho\}
    $$
    Writing $\delta V := \bar{V} - \tilde{V}$ and expanding both the transition and decay terms, a triangle inequality on the $1 \to 1$ norm gives (following the proof technique of \cite[Lemma A.1]{chen2023quantum}):
    $$
    \bigl\|\bar{\mathcal{L}}_{\mathrm{diss}} - \tilde{\mathcal{L}}_{\mathrm{diss}}\bigr\|_{1 \to 1}
    \leq \|\delta V\|\bigl(\|\bar{V}\| + \|\tilde{V}\|\bigr) + \|\delta V\|\bigl(\|\bar{V}\| + \|\tilde{V}\|\bigr) = 2\,\|\delta V\|\bigl(\|\bar{V}\| + \|\tilde{V}\|\bigr).
    $$
    Both $\bar{V}$ and $\tilde{V}$ are sub-isometries with $\|V\| \leq 1$: this follows from the operator Parseval identity (\cite[Proposition A.1]{chen2023quantum}), $\|\gamma\|_\infty \leq 1$, $\|\bar{f}\|_2 = 1$, and the CKG normalization $\|\sum_a A_a^\dagger A_a\| \leq 1$. (For $\tilde{V}$, the same bound holds since $S_p^{(M)}(\bar{t})$ is unitary, so the Parseval identity applies identically with $S_p^{(M)}$ in place of $U$.) Therefore,
    $$
    \bigl\|\bar{\mathcal{L}}_{\mathrm{diss}} - \tilde{\mathcal{L}}_{\mathrm{diss}}\bigr\|_{1 \to 1}
    \leq 4\,\|\bar{V} - \tilde{V}\|.
    $$
    Note, that the key thought here is to use the perturbation techniques introduced in Chen et al. for the error analysis of discretization, with the understanding that Trotterization is just another kind of perturbation in this picture.

    \medskip
    \noindent\textbf{(ii) Bound $\|\bar{V} - \tilde{V}\|$}. We claim that,
    $$
    \|\bar{V} - \tilde{V}\|^2 \leq \sum_{\bar{t} \in S_{t_0}} |\bar{f}(\bar{t})|^2 \, \left\|\sum_a \delta C_a(\bar{t})^\dagger\, \delta C_a(\bar{t})\right\|,
    $$
    where
    $$
    \delta C_a(\bar{t}) := U(\bar{t})^\dagger A_a\, U(\bar{t}) - S_p^{(M)}(\bar{t})^\dagger A_a\, S_p^{(M)}(\bar{t})
    $$
    is there per-time-point Heisenberg-picture error due to Trotterization.

    By linearity, we can write,
    $$
    \bar{V} - \tilde{V} = \sum_{a,\bar{\omega}} \sqrt{\gamma(\bar{\omega})}\, |\bar{\omega}, a\rangle\langle \bar{0}|\, \delta A_a(\bar{\omega}),
    $$
    with $\delta A_a(\bar{\omega}) := \hat{A}_a(\bar{\omega}) - \tilde{A}_a(\bar{\omega}) = \frac{1}{\sqrt{N}} \sum_{\bar{t}} e^{-i\bar{\omega}\bar{t}}\, \bar{f}(\bar{t})\, \delta C_a(\bar{t})$
    Using the orthonormality of the ancillas $\langle \bar{\omega}', a' | \bar{\omega}, a\rangle = \delta_{aa'}\delta_{\bar{\omega}\bar{\omega}'}$, we find
    $$
    (\bar{V} - \tilde{V})^\dagger(\bar{V} - \tilde{V}) = \sum_{a,\bar{\omega}} \gamma(\bar{\omega})\, \delta A_a(\bar{\omega})^\dagger\, \delta A_a(\bar{\omega}) \preceq \sum_{a,\bar{\omega}} \delta A_a(\bar{\omega})^\dagger\, \delta A_a(\bar{\omega}),
    $$
    In the last step we used $\gamma(\omega) \leq 1$. Then we can apply Parseval's identity from \cite[Proposition A.1]{chen2023quantum},
    $$
    \sum_{\bar{\omega}} \delta A_a(\bar{\omega})^\dagger\, \delta A_a(\bar{\omega}) = \sum_{\bar{t}} |\bar{f}(\bar{t})|^2\, \delta C_a(\bar{t})^\dagger\, \delta C_a(\bar{t}).
    $$
    Inserting this back we get, 
    $$
    (\bar{V} - \tilde{V})^\dagger(\bar{V} - \tilde{V}) \preceq \sum_{\bar{t}} |\bar{f}(\bar{t})|^2 \sum_a \delta C_a(\bar{t})^\dagger\, \delta C_a(\bar{t}).
    $$
    Each summand $|\bar{f}(\bar{t})|^2 \sum_a \delta C_a(\bar{t})^\dagger\, \delta C_a(\bar{t})$ is positive semidefinite, and after using triangle inequality, we find the claimed bound on $\|\bar{V} - \tilde{V}\|^2$. 

    \medskip
    \noindent\textbf{(iii) Bound $\left\|\sum_a \delta C_a(\bar{t})^\dagger\, \delta C_a(\bar{t})\right\|$}. We claim that,
    $$
    \Bigl\|\sum_a \delta C_a(\bar{t})^\dagger\, \delta C_a(\bar{t})\Bigr\| \leq 4\,\varepsilon_p^2(\bar{t}, M).
    $$
    Define $\Delta(\bar{t}) := U(\bar{t}) - S_p^{(M)}(\bar{t})$, and its norm (the Trotter error) satisfies $\|\Delta(\bar{t})\| \leq \varepsilon_p(\bar{t}, M)$. Decomposing $\delta C_a(\bar{t})$:
    $$
    \begin{aligned}
        \delta C_a(\bar{t}) &= U^\dagger A_a\, U - (S_p^{(M)})^\dagger A_a\, S_p^{(M)} \\
    &= (S_p^{(M)} + \Delta)^\dagger A_a\, U - (S_p^{(M)})^\dagger A_a (U - \Delta) \\
    &= \Delta^\dagger A_a\, U + (S_p^{(M)})^\dagger A_a\, \Delta \\
    &=: X_a + Y_a,
    \end{aligned}
    $$
    Since $(X_a - Y_a)^\dagger(X_a - Y_a) \succeq 0$ implies $X_a^\dagger Y_a + Y_a^\dagger X_a \preceq X_a^\dagger X_a + Y_a^\dagger Y_a$, we get
    $$
    \delta C_a^\dagger\, \delta C_a = (X_a + Y_a)^\dagger(X_a + Y_a) \preceq 2\bigl(X_a^\dagger X_a + Y_a^\dagger Y_a\bigr).
    $$
    We can bound $\sum_a X_a^\dagger X_a$ by using that $\Delta\Delta^\dagger \preceq \|\Delta\|^2\, \mathbf{1} \preceq \varepsilon_p^2\, \mathbf{1}$, that $U$ is unitary and $\|\sum_a A_a^\dagger A_a\| \leq 1$,
    $$
    \sum_a X_a^\dagger X_a = \sum_a U^\dagger A_a^\dagger\, \Delta\Delta^\dagger\, A_a\, U
    \preceq \varepsilon_p^2 \sum_a U^\dagger A_a^\dagger A_a\, U
    \preceq \varepsilon_p^2\, \mathbf{1}.
    $$
    The result is the same for $\sum_a Y_a^\dagger Y_a$ but by using that $S_p^{(M)}$ is unitary.
    Combining them gives, 
    $$
    \sum_a \delta C_a^\dagger\, \delta C_a \preceq 2\bigl(\varepsilon_p^2\, \mathbf{1} + \varepsilon_p^2\, \mathbf{1}\bigr) = 4\,\varepsilon_p^2\, \mathbf{1}.
    $$
    Taking the operator norm finally proves the upper-bound on $\left\|\sum_a \delta C_a(\bar{t})^\dagger\, \delta C_a(\bar{t})\right\|$.

    \medskip
    \noindent\textbf{(iv) Assemble the cake.} Putting steps (i)-(iii) together, we find that the Trotter error on the full superoperator $\mathcal{\bar{L}}_\text{diss}$ is bounded by
    \begin{align*}
        \bigl\|\bar{\mathcal{L}}_{\mathrm{diss}} - \tilde{\mathcal{L}}_{\mathrm{diss}}\bigr\|_{1 \to 1}
    &\leq 4\,\|\bar{V} - \tilde{V}\| \tag{i}\\
    &\leq 8 \,\sqrt{\sum_{\bar{t}} |\bar{f}(\bar{t})|^2\, \left\|\sum_a \delta C_a(\bar{t})^\dagger\, \delta C_a(\bar{t})\right\|} \tag{ii}\\
    &\leq 8 \,\sqrt{\sum_{\bar{t}} |\bar{f}(\bar{t})|^2\, \varepsilon_p^2(\bar{t}, M)}. \tag{iii}
    \end{align*}

    \medskip
    \noindent\textbf{(v) Strang splitting and Gaussian filters.} In our numerical analysis we will use Strang splitting, i.e. Trotterization of order $p = 2$. In this case, the error bound specializes as follows.
    \begin{align*}
    \bigl\|S_2(\delta t) - \ee^{\ii H \delta t}\bigr\|
    &\leq \sum_{k_1 < k_2}
    \left(
      \frac{1}{12}\bigl\|[H_{k_2},[H_{k_2},H_{k_1}]]\bigr\|
    + \frac{1}{24}\bigl\|[H_{k_1},[H_{k_1},H_{k_2}]]\bigr\|
    \right)
    + \mathcal{O}\bigl(|\delta t|^4\bigr) \\
    &= \tilde{\alpha}_{\mathrm{comm}}^{(2)}\,|\delta t|^3
    + \mathcal{O}\bigl(|\delta t|^4\bigr),
    \end{align*}
    where the leading-order coefficients $1/12$ and $1/24$ are tight, and the $\mathcal{O}(|\delta t|^4)$ correction arises from the multi-term bootstrapping of \cite[Proposition~9]{childs2021theory}.
    With $M$ Trotter steps of size $\delta t = \bar{t}/M$ and a triangle inequality
    over the $M$ steps:
    $$
    \varepsilon_2(\bar{t}, M)
    \leq \tilde{\alpha}_{\mathrm{comm}}^{(2)}\,
        \frac{|\bar{t}|^3}{M^2}
    + \mathcal{O}\!\left(\frac{|\bar{t}|^4}{M^3}\right).
    $$
    Substituting the leading term into the assembled bound and evaluating the sixth moment of $|\bar{f}|^2$, where in the CKG case the filter is the Gaussian $f(t) = \ee^{-\sigma^2 t^2}\big/\bigl(\pi/(2\sigma^2)\bigr)^{1/4}$ with squared density $|f(t)|^2 \propto \ee^{-2\sigma^2 t^2}$, whose even moments evaluate to
    $$
    \int_{-\infty}^{\infty} |f(t)|^2\,t^{2k}\,\dd t
    = \frac{(2k-1)!!}{(2\sigma)^{2k}}.
    $$
    The discrete sum $\sum_{\bar{t}} |\bar{f}(\bar{t})|^2\,|\bar{t}|^6$
    agrees with its continuous counterpart up to an exponentially small
    correction (by the quadrature analysis of the preceding subsection).
    Setting $k = 3$ yields:
    $$
    \bigl\|\bar{\mathcal{L}}_{\mathrm{diss}}
        - \tilde{\mathcal{L}}_{\mathrm{diss}}\bigr\|_{1 \to 1}
    \leq \frac{8\,\tilde{\alpha}_{\mathrm{comm}}^{(2)}}{M^2}\;
        \sqrt{\frac{15}{64\,\sigma^6}}
    + \mathcal{O}\!\left(\frac{1}{M^3\,\sigma^4}\right)
    = \frac{\sqrt{15}\;\tilde{\alpha}_{\mathrm{comm}}^{(2)}}
        {M^2\,\sigma^3}
    + \mathcal{O}\!\left(\frac{1}{M^3\,\sigma^4}\right).
    $$
    The $\mathcal{O}(1/(M^3\sigma^4))$ correction arises from the $\mathcal{O}(|\bar{t}|^4/M^3)$ sub-leading Trotter error propagated through the eighth moment of $|\bar{f}|^2$.
\end{proof}
\noindent The main takeaway here is that the Trotter errors from the OFT propagate to the superoperator level in an expected manner, but with a crucial realization: Since the natural unit of the Gaussian filter width is the temperature $\sigma \sim 1 / \beta$, that means the lower the temperature is the more the  number of Trotter steps $M$ have to compensate to control the errors. Not even higher-orders $ p>2$ can help here, as this relation between $M$ and $\sigma$ is independent of the order. This was to be expected since the lower the temperature, the better the resolution in the energy-domain has to be, and thus the longer the respective time-evolutions are in the time-domain. Nevertheless, this has not been written out explicitly for a Trotter-domain before. 

\medskip
With the Trotter errors for the dissipative part of the Lindbladian $\mathcal{L}_\text{diss}$ fully understood, it remains to see how Trotterization affects the coherent term $B$. Replacing all exact Hamiltonian time evolutions in \eqref{eq:B-discr} with the Trotterized evolution $S_p^{(M)}(-\bar{t})$ leads to
\begin{equation} \label{eq:B-trotter}
    \begin{split}
         \tilde{B} &= \sum_a t_0\sum_{\bar{t}} b_-(\bar{t}) S_p^{(M)}(-\bar{t} / \sigma) \left[t_0'\sum_{\bar{t}'} b_+(\bar{t}') \tilde{A}_{a\dagger}(\beta \bar{t}') \tilde{A}_a(-\beta \bar{t}')\right]S_p^{(M)}(\bar{t} / \sigma) \\
         &=: \sum_{a} t_0 \sum_{\bar{t} \in S_{t_0}} b_-(\bar{t})\, S_p^{(M)}(-\bar{t}/\sigma)\, \tilde{O}_a\, S_p^{(M)}(\bar{t}/\sigma),
    \end{split}
\end{equation}
with $\tilde{O}_a$ denoting the inner Trotterized integral, and the Trotterized Heisenberg evolution of order $p$ and steps $M$, 
$$
\tilde{A}_a(t) := S_p^{(M)}(t) A_a S_p^{(M)}(-t).
$$
In the following we will use $\|b_-\|_{1} := t_0 \sum_{\bar{t}} |b_-(\bar{t})|$ and $\|b_+\|_{1} := t_0' \sum_{\bar{t}'} |b_+(\bar{t}')|$ for the discrete $\ell_1$ norms (with quadrature weights).

\todo{ref to impl}
There are a few ways to implement the evolution via the CKG Lindbladian, especially in terms of how we add the coherent evolution via $B$, but in all cases the core error boils down to controlling $\|\bar{B} - \tilde{B}\|$.
\begin{proposition} \label{prop:trotter-B}
    \textup{(Trotter error for the coherent term $B$).} Let $H = \sum_{k=1}^K H_k$ with non-commuting terms, let the jump operators satisfy $\bigl\|\sum_a A_a^\dagger A_a\bigr\| \leq 1$, and let $b_\pm$ be the time-domain weight functions of the coherent term with finite discrete norms $\|b_-\|_{1}, \|b_+\|_{1} < \infty$. Let $S_p^{(M)}(t)$ be a $p$-th order product formula with $M$ Trotter steps satisfying the commutator-scaling bound $\varepsilon_p(t, M) = \mathcal{O}(\tilde{\alpha}_{\mathrm{comm}}\,|t|^{p+1}/M^p)$.

    Then the operator-norm Trotter error of the coherent term satisfies, 
    $$
    \|\bar{B} - \tilde{B}\|
    \;\leq\;
    2\,\|b_+\|_{1}\; t_0 \sum_{\bar{t}} |b_-(\bar{t})|\, \varepsilon_p\!\bigl(|\bar{t}|/\sigma,\, M\bigr)
    \;+\;
    4\,\|b_-\|_{1}\; t_0' \sum_{\bar{t}'} |b_+(\bar{t}')|\, \varepsilon_p\!\bigl(\beta|\bar{t}'|,\, M\bigr).
    $$
    Here, the first term captures the outer Trotter error (replacing the outer exact time evolution $\ee^{\pm \ii H \bar{t} / \sigma}$), while the second captures the inner Trotterization error (replacing $\ee^{\pm \ii H\beta\bar{t}'}$ inside $A_a(\beta \bar{t}')$).

    For the special case of Strang splitting ($p = 2$), substituting $\varepsilon_2(t, M) \leq \tilde{\alpha}_{\mathrm{comm}}^{(2)}\,|t|^3/M^2 + \mathcal{O}(|t|^4/M^3)$ gives the following result at leading order:
    $$
    \|\bar{B} - \tilde{B}\|
    \;\leq\;
    \frac{\tilde{\alpha}_{\mathrm{comm}}^{(2)}}{M^2}
    \left(
    \frac{2\,\|b_+\|_1}{\sigma^3}\, \mu_3(|b_-|)
    \;+\;
    4\,\beta^3\,\|b_-\|_1\, \mu_3(|b_+|)
    \right)
    + \mathcal{O}\!\left(\frac{\beta^4}{M^3}\right),
    $$
    where $\mu_3(|b_\pm|) := t_0^{(\prime)} \sum |b_\pm(\bar{t})|\,|\bar{t}|^3$ denotes the third absolute moment of the respective weight function (with $t_0$ or $t_0'$).
\end{proposition}
\begin{proof}
    \hfill\break
    \noindent\textbf{(i) Decompose the Trotter error.}
    For each outer grid point $\bar{t}$ and jump operator $a$, define:
    $$
    \bar{Q}_a(\bar{t}) := U(-\bar{t}/\sigma)\, \bar{O}_a\, U(\bar{t}/\sigma), \qquad
    \tilde{Q}_a(\bar{t}) := S(-\bar{t}/\sigma)\, \tilde{O}_a\, S(\bar{t}/\sigma).
    $$
    Then, to bound an error between these two expressions, follow these steps in order: use the three-term telescoping identity for products of the form $PQR - P'Q'R'$, simplify, sum over $a$, use that $\|U\| = 1 = \|S\|$. At this point we arrive at
    $$
    \Bigl\|\sum_a \bigl[\bar{Q}_a(\bar{t}) - \tilde{Q}_a(\bar{t})\bigr]\Bigr\|
    \leq
    \varepsilon_p(|\bar{t}|/\sigma, M)\, \Bigl\|\sum_a \bar{O}_a\Bigr\|
    + \Bigl\|\sum_a \bigl[\bar{O}_a - \tilde{O}_a\bigr]\Bigr\|
    + \Bigl\|\sum_a \tilde{O}_a\Bigr\|\, \varepsilon_p(|\bar{t}|/\sigma, M),
    $$

    \noindent\textbf{(ii) Bound $\Bigl\|\sum_a \bar{O}_a\Bigr\|$ and $\Bigl\|\sum_a \tilde{O}_a\Bigr\|$.} The inner operator $O_a$ has the following general form: $A_a^\dagger W A_a$ for some unitary $W$ in both the Trotterized case \eqref{eq:B-trotter} and discretized causes \eqref{eq:B-discr}. Since, $-\mathbf{1} \preceq W \preceq \mathbf{1}$, we also have $-A_a^\dagger A_a \preceq A^\dagger_a W A_a \preceq A_a^\dagger A_a$. Summing over $a$ and taking norm then leads to:
    $$
    \Bigl\|\sum_a A_a^\dagger\, W\, A_a\Bigr\| \leq \Bigl\|\sum_a A_a^\dagger A_a\Bigr\| \leq 1,
    $$
    Therefore:
    $$
    \Bigl\|\sum_a \bar{O}_a\Bigr\| \leq t_0' \sum_{\bar{t}'} |b_+(\bar{t}')| = \|b_+\|_1.
    $$
    The identical bound holds for $\bigl\|\sum_a \tilde{O}_a\bigr\| \leq \|b_+\|_1$ since $S$ is also unitary.

    \noindent\textbf{(ii) Bound $\Bigl\|\sum_a \bigl[\bar{O}_a - \tilde{O}_a\bigr]\Bigr\|$.} We have
    $$
    \sum_a [\bar{O}_a - \tilde{O}_a] = t_0' \sum_{\bar{t}'} b_+(\bar{t}') \sum_a \Bigl[\hat{A}_a^\dagger(\beta\bar{t}')\, \hat{A}_a(-\beta\bar{t}') - \tilde{A}_a^\dagger(\beta\bar{t}')\, \tilde{A}_a(-\beta\bar{t}')\Bigr],
    $$
    so it suffices to bound $\bigl\|\sum_a \bigl[\hat{A}_a^\dagger(\beta\bar{t}')\hat{A}_a(-\beta\bar{t}') - \tilde{A}_a^\dagger(\beta\bar{t}')\tilde{A}_a(-\beta\bar{t}')\bigr]\bigr\|$ for each $\bar{t}'$. Define $\delta_a(t) := \hat{A}_a(t) - \tilde{A}_a(t)$ for the per-operator Heisenberg-picture Trotter error. Then,
    $$
    \sum_a \Bigl[\hat{A}_a^\dagger(\beta\bar{t}')\hat{A}_a(-\beta\bar{t}') - \tilde{A}_a^\dagger(\beta\bar{t}')\tilde{A}_a(-\beta\bar{t}')\Bigr]
    = \sum_a \delta_a(\beta\bar{t}')^\dagger\, \hat{A}_a(-\beta\bar{t}') + \sum_a \tilde{A}_a^\dagger(\beta\bar{t}')\, \delta_a(-\beta\bar{t}').
    $$
    We bound each sum using the operator Cauchy-Schwarz inequality, which in general can be written for operators $\{X_a\}, \{Y_a\}$ as
    $$
    \Bigl\|\sum_a X_a Y_a\Bigr\|^2 \leq \Bigl\|\sum_a X_a X_a^\dagger\Bigr\| \cdot \Bigl\|\sum_a Y_a^\dagger Y_a\Bigr\|.
    $$
    The first sum with $X_a = \delta_a(\beta\bar{t}')^\dagger$ and $Y_a = \hat{A}_a(-\beta\bar{t}')$ is bounded by two corresponding factors. The left one is
    $$
    \bigl\|\sum_a \delta_a(\beta\bar{t}')^\dagger\, \delta_a(\beta\bar{t}')\bigr\| \leq 4\,\varepsilon_p^2(\beta|\bar{t}'|, M),
    $$
    which was shown in step (iii) of \hyperref[prop:trotter-diss]{Proposition~\ref*{prop:trotter-diss}}. The right factor for the bound is 
    $$
    \bigl\|\sum_a \hat{A}_a(-\beta\bar{t}')^\dagger\, \hat{A}_a(-\beta\bar{t}')\bigr\| = \bigl\|\sum_a A_a^\dagger A_a\bigr\| \leq 1.
    $$
    Therefore, $\bigl\|\sum_a \delta_a(\beta\bar{t}')^\dagger\, \hat{A}_a(-\beta\bar{t}')\bigr\| \leq 2\,\varepsilon_p(\beta|\bar{t}'|, M).$ An identical bound holds for the second sum as well. Combining them via triangle inequality leads to,
    $$
    \Bigl\|\sum_a [\bar{O}_a - \tilde{O}_a]\Bigr\|
    \leq 4\, t_0' \sum_{\bar{t}'} |b_+(\bar{t}')|\, \varepsilon_p(\beta|\bar{t}'|, M).
    $$
    \noindent\textbf{(iv) Assemble another cake.} From steps (i)-(iii):
    $$
    \Bigl\|\sum_a [\bar{Q}_a(\bar{t}) - \tilde{Q}_a(\bar{t})]\Bigr\|
    \leq 2\,\|b_+\|_1\, \varepsilon_p(|\bar{t}|/\sigma, M)
    + 4\, t_0' \sum_{\bar{t}'} |b_+(\bar{t}')|\, \varepsilon_p(\beta|\bar{t}'|, M).
    $$
    Taking the outer $b_-$-weighted sum of $\|\|\bar{B} - \tilde{B}\|\|$, then inserting the above bound in yields our result
    \begin{align*}
        \|\bar{B} - \tilde{B}\|
    &= \Bigl\|t_0 \sum_{\bar{t}} b_-(\bar{t})\, \sum_a [\bar{Q}_a(\bar{t}) - \tilde{Q}_a(\bar{t})]\Bigr\|
    \leq t_0 \sum_{\bar{t}} |b_-(\bar{t})| \cdot \Bigl\|\sum_a [\bar{Q}_a(\bar{t}) - \tilde{Q}_a(\bar{t})]\Bigr\| \\
    & \leq 2\,\|b_+\|_1\, t_0 \sum_{\bar{t}} |b_-(\bar{t})|\, \varepsilon_p(|\bar{t}|/\sigma, M)
    + 4\,\|b_-\|_1\, t_0' \sum_{\bar{t}'} |b_+(\bar{t}')|\, \varepsilon_p(\beta|\bar{t}'|, M).
    \end{align*}    
    \noindent\textbf{(v) Strang splitting specialization.} For $p = 2$, the leading-order Trotter error is $\varepsilon_2(t, M) \leq \tilde{\alpha}_{\mathrm{comm}}^{(2)}|t|^3/M^2$. Define the third absolute moment of a weight function $w$ on grid $S_{\tau}$:
    $$
    \mu_3(|w|) := \tau \sum_{\bar{t} \in S_\tau} |w(\bar{t})|\, |\bar{t}|^3.
    $$
    Then the outer term becomes, $t_0 \sum_{\bar{t}} |b_-(\bar{t})|\, \varepsilon_2(|\bar{t}|/\sigma, M) \leq \frac{\tilde{\alpha}_{\mathrm{comm}}^{(2)}}{M^2\, \sigma^3}\, \mu_3(|b_-|)$, while the inner term, $t_0' \sum_{\bar{t}'} |b_+(\bar{t}')|\, \varepsilon_2(\beta|\bar{t}'|, M) \leq \frac{\tilde{\alpha}_{\mathrm{comm}}^{(2)}\, \beta^3}{M^2}\, \mu_3(|b_+|)$. Hence, up to leading order we find
    $$
    \|\bar{B} - \tilde{B}\|
    \leq \frac{\tilde{\alpha}_{\mathrm{comm}}^{(2)}}{M^2}
    \left(
    \frac{2\,\|b_+\|_1}{\sigma^3}\, \mu_3(|b_-|)
    + 4\,\beta^3\, \|b_-\|_1\, \mu_3(|b_+|)
    \right)
    + \mathcal{O}\!\left(\frac{\beta^4}{M^3}\right).
    $$
\end{proof}
The two terms present in the Trotter error bound of $\|\bar{B} - \tilde{B}\|$ arise from the nested outer and inner time-sums in $B$ \eqref{eq:B-discr}. Qualitatively they contribute similarly. The $\log$ scaling of $\|b^{(s, \eta)}_+\|_1$ \eqref{eq:b_plus-s-eta} that appears for Metropolis-like transitions (due to the $\eta$-regularization) is the only difference between the errors. Note that at leading-order, the third absolute moments $\mu_3(|b_\pm|)$ are $\mathcal{O}(1)$ for both Gaussian and (smooth) Metropolis-like weights: the reason is that the $|t|^3$ factor within renders the $1 / t$ singularity integrable. Thus, the moment decays dominantly as a Gaussian.

In principle, we could compose the two Trotter error results of \hyperref[prop:trotter-diss]{Proposition~\ref*{prop:trotter-diss}} and \hyperref[prop:trotter-B]{Proposition~\ref*{prop:trotter-B}} into one complete generator level error bound by a simple triangle inequality. But this gives no additional insight: in our implementation we Trotter split the generator and apply the coherent evolution by $B$ separate from the dissipative part $\mathcal{L}_\text{diss}$. Moreover, in implementation we have the flexibility of choosing different number of Trotter steps for these two parts since to control Trotter errors on OFT's might be simpler than controlling it on some less regular functions $b_\pm$ in $B$.

It is worth clarifying which structural properties of the CKG Lindbladian survive discretization and Trotterization. The three levels of approximation are: 
\begin{itemize}
    \item \textbf{Continuous} $\mathcal{L}$: exact KMS DB by construction  \hyperref[sec:kms-lindbladian]{KMS Lindbladian} based on \cite[Theorem I.1]{chen2023efficient}.
    \item \textbf{Discretized} $\mathcal{\bar{L}}$: approximate KMS DB. Quadrature errors break detailed balance. The size of the violation is controlled by the analysis in \hyperref[sec:quad-errors-diss]{Quadrature errors for $\mathcal{L}_\text{diss}$} and \hyperref[sec:quad-errors-B]{Quadrature errors for $B$}.
    \item \textbf{Trotterized} $\mathcal{\tilde{L}}$: further approximates KMS DB, with additional Trotter errors from \hyperref[prop:trotter-diss]{Proposition~\ref*{prop:trotter-diss}} and \hyperref[prop:trotter-B]{Proposition~\ref*{prop:trotter-B}} on top of the quadrature errors.
\end{itemize}
The exact KMS DB of the continuous Lindbladian relies on several ingredients acting in unison: the shifted KMS condition on the transition weight $\gamma$, the unique construction of the coherent term $B$ cancelling the anti-Hermitian residual of the quantum discriminant, the properties of the OFT construction. Not all of these are relevant at the discretized level, but one structural identity does remain exact here -- and the question is whether the imminent Trotterization preserves it. We cover this in the following remark.

\begin{remark} \label{rem:palindromic}
    \textup{(Palindromic product formulas).} The identity that survives discretization and can be made to survive Trotterization as well, is the \emph{OFT adjoint identity}:
    $$
    \hat{A}_{a'}(-\bar{\omega})^\dagger = \hat{A}_a(\bar{\omega}), \qquad \text{with } A_{a'} = A_a^\dagger \in \mathcal{A}.
    $$
    At the discretized level with exact Hamiltonian evolution, this identity holds exactly because the proof only requires that, (i) $\bar{f}(\bar{t}) \in \mathbb{R}$, (ii) $\bar{f}(-\bar{t}) = \bar{f}(\bar{t})$ (the used Gaussian filter is even), and (iii) $U(\bar{t})^\dagger = U(-\bar{t})$. All three hold for the discretized OFT \eqref{eq:oft-discr}. This identity ensures that the discriminant $\mathcal{D}_\beta$ from \eqref{eq:discriminant} satisfies $\mathcal{D}_\beta = \mathcal{D}_\beta^\dagger$, i.e. its anti-Hermitian part $\mathcal{A}(\rho_\beta, \mathcal{\bar{L}})$
    -- the measure of DB violation -- from \hyperref[def:approx-db]{Definition~\ref*{def:approx-db}} vanishes.

    At the Trotterized level, input (iii) must be replaced by $S_p^{(M)}(\bar{t})^\dagger = S_p^{(M)}(-\bar{t})$, which holds if and only if the product formula is \emph{palindromic} (or adjoint-symmetric).
\end{remark}
\begin{proof}
    A product formula $S_p(\delta t) = \prod_{j=1}^{J} e^{-ic_j H_{\pi(j)} \delta t}$ is palindromic if $c_{J+1-j} = c_j$ and $\pi(J+1-j) = \pi(j)$ for all $j$. For such a formula:
    $$
    S_p(\delta t)^\dagger = \prod_{j=J}^{1} e^{ic_j H_{\pi(j)} \delta t} = \prod_{j=1}^{J} e^{ic_{J+1-j} H_{\pi(J+1-j)} \delta t} = \prod_{j=1}^{J} e^{ic_j H_{\pi(j)}\delta t} = S_p(-\delta t),
    $$
    where the second step is the palindromic reindexing. This naturally extends to $M$ steps $S_p^{(M)}(\bar{t})^\dagger = S_p^{(M)}(-\bar{t})$. Strang splitting (and all other even Trotter formulas) satisfies this by construction.
    Then, we can explicitly check for $\tilde{A}_a(\bar{\omega})$ from \eqref{eq:trotter-oft}:
    \begin{align*}
        \tilde{A}_a(\bar{\omega})^\dagger &= \frac{1}{\sqrt{N}} \sum_{\bar{t} \in S_{t_0}} e^{i\bar{\omega}\bar{t}}\, \bar{f}(\bar{t})\, S_p^{(M)}(\bar{t})\, A_a^\dagger\, S_p^{(M)}(\bar{t})^\dagger \\
        &= \frac{1}{\sqrt{N}} \sum_{\bar{t}} e^{-i\bar{\omega}\bar{t}}\, \bar{f}(\bar{t})\, S_p^{(M)}(\bar{t})^\dagger\, A_{a'}\, S_p^{(M)}(\bar{t}),
    \end{align*}
    where we sent $\bar{t} \to -\bar{t}$ on the symmetric grid $S_{t_0}$, applied $\bar{f}(-\bar{t}) = \bar{f}(\bar{t})$ and the palindromic identity $S_p^{(M)}(-\bar{t}) = S_p^{(M)}(\bar{t})^\dagger$. We can also compute $\tilde{A}_{a'}(-\bar{\omega})$ separately, 
    \begin{align*}
        \tilde{A}_{a'}(-\bar{\omega}) &= \frac{1}{\sqrt{N}} \sum_{\bar{t}} e^{i\bar{\omega}\bar{t}}\, \bar{f}(\bar{t})\, S_p^{(M)}(\bar{t})\, A_{a'}\, S_p^{(M)}(\bar{t})^\dagger \\
        &= \frac{1}{\sqrt{N}} \sum_{\bar{t}} e^{-i\bar{\omega}\bar{t}}\, \bar{f}(\bar{t})\, S_p^{(M)}(\bar{t})^\dagger\, A_{a'}\, S_p^{(M)}(\bar{t}),
    \end{align*}
    and see that the two expressions coincide, $\tilde{A}_a(\bar{\omega})^\dagger = \tilde{A}_{a'}(-\bar{\omega})$.
\end{proof}
The anti-Hermitian part $\mathcal{A}(\rho_\beta, \mathcal{\bar{L}})$ of the quantum discriminant receives contributions from all sources of error: quadrature, Trotter, and possibly from non-palindromic product formulas (e.g. Lie-Trotter) violating the OFT adjoint identity. A palindromic formula eliminates the last channel entirely, leaving the Trotter contribution to the DB violation coming solely from the failure of $S_p^{(M)}(\bar{t})$ to commute with $\rho_\beta^{1/2}$ (defined with the exact Hamiltonian $H$).
For a non-palindromic formula (necessarily odd-order) formula of order $p$, the violation can be quantified with the commutator constant $\tilde{\alpha}_\text{comm}$ from \cite{childs2021theory} as 
$$
\|S_p^{(M)}(\bar{t})^\dagger - S_p^{(M)}(-\bar{t})\| = \mathcal{O}\left(\frac{\tilde{\alpha}_\text{comm} |\bar{t}|^{p + 1}}{M^p}\right), 
$$
so the violation enters at the same order as the Trotter error itself. In other words, the asymptotic scaling of the total incurring Trotter errors $\mathcal{O}(1 / M^p)$ is the same for both palindromic and non-palindromic formulas, but the constant prefactor is strictly smaller in the former case.

For commuting Hamiltonians, the product formula is exact for any $M$ and there is no Trotterization error. But the CKG framework was designed to handle non-commuting Hamiltonians. The commuting case is more naturally handles by the Davies generator directly \cite[text]{kastoryano2016quantum,bardet2023rapid}, for which the CKG Gaussian-smearing and coherent term $B$ are unnecessary.

\smallskip
\subsection{Generator splitting.}
In the implementation we use two distinct decomposition of the CKG generator \eqref{eq:ckg-L}. 

(i) First, the full generator is decomposed into per-jump contributions $\mathcal{L}_a$,
$$
\mathcal{L} = \sum_a \mathcal{L}_a.
$$
There are two ways to do this, each of them preserves the Gibbs state as the fixed point. We can either sweep over $A_a \in \mathcal{A}$ or randomly sample the jumps $A_a$ with probability $p_a$, where each contribution is rescaled by $1 / p_a$ such that after many samples we find the full generator $\mathcal{L}$ on average. Asymptotically both strategies give the same complexity, but there is some qualitative differences between the two. While the sequential Lie-Trotter splitting breaks KMS DB in general, at leading order it does so by an antisymmetric term which does not change the spectral gap. On the other hand, the random sampling is essentially a convex combination of KMS DB Lindbladians which leads to an KMS DB evolution on average [REF] \todo{ref}. Even though this sounds advantageous, it also means that the leading order errors do perturb the spectral gap. Since these errors are of the same order as errors from other sources in the algorithm, it made no significant difference which strategy we chose for our numerical simulations. But it is hard to argue against the better error scaling for an implementation, so we will mainly discuss the sequential Lie-Trotter splitting. (For other theoretical works, where implementation via discretization, Trotterization, etc. is not relevant, then keeping KMS DB might be more appealing since it keeps the quantum discriminant Hermitian.)

(ii) Second, for each fixed $a$, the generator $\mathcal{L}_a$ is further split into a coherent and a dissipative part, each implemented as their own, separate primitive:
$$
 \mathcal L_a = \mathcal L_{a,\mathrm{coh}} + \mathcal L_{a,\mathrm{diss}},
    \qquad
    \mathcal L_{a,\mathrm{coh}}(\cdot) := - i [B_a,\cdot].
$$
These two splittings are qualitatively different from each other: while the jump-wise splitting preserves the Gibbs state as the fixed point, the coherent-dissipative splitting generally does not. Consequently, the former introduces errors in the spectrum that can slow the mixing time, while the latter produces a genuine fixed-point bias. In either case, we analyse the incurring errors here and find a suiting way to control them for our purposes.

It is important to keep in mind that for the implementation cost, the relevant baseline error comes from the weak measurement simulation of the dissipative part $\mathcal{L}_\mathrm{diss}$. This primitive is only first-order accurate in the step size, i.e. for a step $\delta$ it has a local error $\mathcal{O}(\delta^2)$, and over the total simulated time $t_\text{mix}$ the accumulated error is $\mathcal{O}(t_\text{mix} \delta)$. Accordingly, all the generator decomposition we introduce in (i) and (ii) should respect this scaling as well.

\medskip
\noindent\textbf{(i) Jump-wise splitting.} Each of the Lindbladians $\mathcal{L}_a$ individually satisfies exact KMS DB for $\rho_\beta$, i.e. they are self-adjoint with respect to the KMS inner product ($\mathcal{L}_a^* = \mathcal{L}_a$). Consequently, the Gibbs state is the fixed point of each generator $\mathcal{L}_a(\rho_\beta) = 0$, and hence also of each semigroup channel $e^{\delta\mathcal{L}_a}(\rho_\beta) = \rho_\beta$. Note that, the kernel of an individual $\mathcal{L}_a$ is generically larger than $\mathrm{Span}\{\rho_\beta\}$. Primitivity is a property of the full generator $\mathcal{L} = \sum_a \mathcal{L}_a$ only, not of each summand.

The simplest (yet sufficient) way to split the generator sequentially, is to use the first-order Lie-Trotter splitting:
\begin{equation}
    \Phi_{\mathcal{A}} := \ee^{\delta \mathcal{L}_{M_\mathcal{A}}} \circ \cdots \circ \ee^{\delta \mathcal{L}_2} \circ \ee^{\delta \mathcal{L}_1},
\end{equation}
with $M_\mathcal{A} := |\mathcal{A}|$ denoting the number of jump operators. By the standard first-order product formula bound, applied to the superoperators $\mathcal{L}_a$, we can write  
\begin{equation} \label{eq:gen-split-approx-error}
    \left\|\Phi_{\mathcal{A}} - e^{\delta \mathcal{L}}\right\|_{1\to 1} \leq \frac{\delta^2}{2} \left\|\sum_{a < b} [\mathcal{L}_a, \mathcal{L}_b]\right\|_{1\to 1} + O(\delta^3).
\end{equation}
Under the normalization $\|\sum_a A_a^\dagger A_a\| \leq 1$, each per-jump Lindbladian has norm
\begin{equation}
    \|\mathcal{L}_a\|_{1 \to 1} = O\!\left(\frac{\log(\beta\|H\|)}{M_\mathcal{A}}\right),
\end{equation}
where the $\log(\beta\|H\|)$ factor originates from the Metropolis-like coherent term $B_a$ [REF within thesis]\todo{ref}. The dissipative part alone satisfies $\|\mathcal{L}_{a,\mathrm{diss}}\|_{1\to 1} \leq 4\|A_a^\dagger A_a\|$, so the coherent contribution dominates. Applying the triangle inequality to the commutator sum in \eqref{eq:gen-split-approx-error}:
\begin{equation}
    \left\|\sum_{a<b} [\mathcal{L}_a, \mathcal{L}_b]\right\|_{1\to 1} 
    \leq \binom{M_\mathcal{A}}{2} \cdot 2 \max_{a,b} \|\mathcal{L}_a\|_{1\to 1}\|\mathcal{L}_b\|_{1\to 1} 
    = O\!\left(\log^2(\beta\|H\|)\right),
\end{equation}
for which the $M_\mathcal{A}$ dependence cancelled. The per-step approximation error is therefore
\begin{equation} \label{eq:gen-split-explicit-bound}
    \left\|\Phi_{\mathcal{A}} - e^{\delta \mathcal{L}}\right\|_{1\to 1} = O\!\left(\delta^2 \log^2(\beta\|H\|)\right),
\end{equation}
independently of $M_\mathcal{A}$. For Gaussian transition weights, where the $\eta$-regularization of the coherent term is unnecessary, the $\log^2$ factor is absent and the bound simplifies to $O(\delta^2)$.

There is some structure to this error however, that is worth exploring.
Since each $\ee^{\delta \mathcal{L}_a}$ is KMS DB, and the KMS adjoint reverses the composition, $({\Phi_1 \circ \Phi_2})^* = \Phi_2^* \circ \Phi_1^*$, we obtain: 
\begin{equation}
    \Phi_{\mathcal{A}}^* = e^{\delta \mathcal{L}_1} \circ e^{\delta \mathcal{L}_2} \circ \cdots \circ e^{\delta \mathcal{L}_{M_\mathcal{A}}} \neq \Phi_{\mathcal{A}}.
\end{equation}
This is the same kind of DB violation as what we have seen for non-palindromic Trotter products for the Hamiltonian time evolutions in the previous section: the choice of ordering in the product formula breaks time-reversal symmetry. The DB violation is quantified by the difference between the forward and reverse orderings. By Baker-Campbell-Hausdorff (BCH), the effective generators of the two orderings differ in the sign of the leading commutator term, 
$$
\delta \sum_a \mathcal{L}_a \pm \frac{\delta^2}{2}\sum_{a < b}[\mathcal{L}_a, \mathcal{L}_b] + O(\delta^3),
$$
yielding the DB violation at leading order, 
\begin{equation}
    \left\|\Phi_{\mathcal{A}} - \Phi_{\mathcal{A}}^*\right\|_{1\to 1} = O\!\left(\delta^2 \left\|\sum_{a<b}[\mathcal{L}_a, \mathcal{L}_b]\right\|_{1\to 1}\right),
\end{equation}
which is of the same order as the approximation error \eqref{eq:gen-split-approx-error}. Crucially, even though the sequential product breaks detailed balance, it still has the Gibbs state as its fixed point exactly: since each factor fixes $\rho_\beta$, so does $\Phi_{\mathcal{A}}(\rho_\beta) = \rho_\beta$. Moreover, the DB violation has a definite algebraic structure: the commutator of any two KMS self-adjoint generators is KMS anti-self-adjoint: 
$$
[\mathcal{L}_a, \mathcal{L}_b]^* = [\mathcal{L}_b^*, \mathcal{L}_a^*] = -[\mathcal{L}_a, \mathcal{L}_b].
$$
Hence, the leading BCH correction involving the commutator is a purely antisymmetric perturbation of the generator. This has a direct spectral consequence. Since $\mathcal{L}$ is KMS self-adjoint, its eigenvalues in the KMS inner product are real. By standard perturbation theory, an anti-self-adjoint term can only lead to a purely imaginary perturbation to the spectrum $\{\lambda_i\}$ of $\mathcal{L}$ or written out:
\begin{equation} \label{eq:gen-split-imaginary-delta2-errors}
    \Delta\lambda_k = \frac{\delta}{2}\langle \psi_k, \sum_{a<b}[\mathcal{L}_a, \mathcal{L}_b](\psi_k) \rangle_{\mathrm{KMS}} \in \ii\mathbb{R},
\end{equation}
where $\{\psi_k\}$ are the KMS-orthonormal eigenvectors of $\mathcal{L}$. Therefore, the jump-wise generator splitting is less harmful than a generic first-order product formula: for the $\delta$ step channel it lets the spectral gap of $\Phi_\mathcal{A}$ agree with that of $\ee^{\delta \mathcal{L}}$ up to $\mathcal{O}(\delta^3)$.

The sequential product $\Phi_\text{seq}$ leading to this specific structure for the error terms, raises a question, namely: could this anti-KMS-adjoint error term accelerate the convergence to the Gibbs state? The only reason we would even ask such a question is because this is precisely the structure that has been shown to accelerate mixing in the continuous generator setting. For classical diffusions, \cite{hwang2005accelerating} showed that such an antisymmetric addition can only improve the spectral gap. A similar fact for the quantum case was established by \cite{li2025speeding}, showing that for a DB QMS this could mean an up to quadratic speed-up in mixing. But none of these papers are one-to-one applicable to our case, nor did the simulations show any noticeable speed-up. It seems that turning these pesky errors into an advantage keeps eluding us.
\todo{ref, we could check this numerically.}

As a final note to the jump-wise splitting: higher order product formulas would make no meaningful difference here, especially if the weak measurement scheme within $\mathcal{L}_a$ already introduces errors of $\mathcal{O}(\delta^2)$ per step.

\medskip
\noindent\textbf{(ii) Coherent--Dissipative}
Splitting the coherent term from the dissipative one and implementing their respective evolutions separately from each other, is a qualitatively different process compared to the jump-wise splitting. The coherent part is chosen precisely such that the sum of the coherent and dissipative parts are exactly KMS DB. Unlike in the previous case (i), now the separate terms do not satisfy KMS DB, thus the emerging errors already shift the fixed point at leading orders. To see this let us write out the first-order splitting:
\begin{equation}
    \ee^{\delta \mathcal L_{a,\mathrm{coh}}}\ee^{\delta \mathcal L_{a,\mathrm{diss}}}
    = \exp\!\left(
        \delta \mathcal L_a
        + \frac{\delta^2}{2}[\mathcal L_{a,\mathrm{coh}},\mathcal L_{a,\mathrm{diss}}]
        + O(\delta^3)
    \right),
\end{equation}
so that the effective generator is 
\begin{equation} \label{eq:gen-split-coh-diss}
    \mathcal L_{a,\mathrm{eff}}
    = \mathcal L_a
    + \frac{\delta}{2}[\mathcal L_{a,\mathrm{coh}},\mathcal L_{a,\mathrm{diss}}]
    + O(\delta^2).
\end{equation}
Since neither $L_{a,\mathrm{coh}}$ nor $L_{a,\mathrm{diss}}$ satisfy KMS DB, their commutator does not become KMS anti-self-adjoint. Therefore, the fixed point is now shifted, and the leading perturbation now has an impact even on the real part of the spectrum. 

The already used and convenient way to express this is via the fixed-point perturbation bound of \cite[Lemma II.1]{chen2023quantum} \todo{ref to where we used this} 
\begin{equation}
        \|\rho_{\mathrm{fix}}(\mathcal L_1)-\rho_{\mathrm{fix}}(\mathcal L_2)\|_1
    \le 4\, \|\mathcal L_1-\mathcal L_2\|_{1\to1}\, t_{\mathrm{mix}}(\mathcal L_1).
\end{equation}
Applying this with $\mathcal L_1=\mathcal L_a$ and $\mathcal L_2=\mathcal L_{a,\mathrm{eff}}$ yields a fixed-point shift of order $\mathcal O\left(t_\text{mix}\, \delta\right)$. If the substeps (coherent evolution via $B$ and weak measurement scheme for $\mathcal{L}_\text{diss}$) were implemented exactly, then a Strang splitting would improve this BCH error to second order in the effective generator:
\begin{equation}
    e^{\frac{\delta}{2}\mathcal L_{a,\mathrm{coh}}}
    e^{\delta \mathcal L_{a,\mathrm{diss}}}
    e^{\frac{\delta}{2}\mathcal L_{a,\mathrm{coh}}}
    = \exp\!\left(\delta \mathcal L_a + O(\delta^3)\right),
\end{equation}
which would reduce the induced fixed-point bias to $\mathcal O(t_{\mathrm{mix}}\delta^2)$ over the full evolution. However, this improvement is not useful in our implementation, since any more accuracy beyond $\mathcal O\left(t_\text{mix}\, \delta\right)$ is washed away by the errors due to the weak measurement subroutine. 

For this reason, higher-order coherent-dissipative product formulas do not improve the asymptotic end-to-end scaling in our setting. If the weak measurement scheme were to be replaced by some more accurate variant, then we could always adjust all the above product formulas to a higher order variant to comply with the baseline error. \eqref{eq:relative-entropy}

% Use Trotter splitting also for coherent and dissipative parts. t_mix scaling becomes poly with this compared to linear for Chen's, but doesn't need QSVT, doesn't need a ton of repetitions due to subnorm, coherent ancillas are independent and incoherent from dissipative ancillas. It is a valid realization with NISQ in mind, even if asymptotically it's worse.

\smallskip
\label{sec:algorithm}
\subsection{Algorithm.}

\begin{algorithm}[H]
\caption{\textsc{CKG-GibbsSampler}}
\label{alg:main}
\begin{algorithmic}[1]
\small
\Require
  System Hamiltonian $H = \sum_{k=1}^{K} H_k$ on $n$ qubits;
  jump set $\{A_a\}_{a\in\mathcal{A}}$ with $\bigl\lVert\sum_a A_a^\dagger A_a\bigr\rVert \le 1$;
  inverse temperature $\beta$;
  Gaussian filter width $\sigma$ for filter function $f$ \eqref{eq:gaussian-filter-t};
  transition weight $\gamma$ (e.g.\ $\gamma_M^{(s)}$ from~\eqref{eq:smooth-metro});
  functions $b_\pm$ for the coherent term \eqref{eq:b_plus-s-eta} and \eqref{eq:b_minus};
  step size $\delta$;
  target trace-distance accuracy $\varepsilon$
\Ensure System register $S$ in state
  $\rho \approx \rho_\beta = \ee^{-\beta H}/\mathrm{tr}(\ee^{-\beta H})$
\Statex
\State $L \gets \bigl\lceil t_{\mathrm{mix}}(\mathcal{L})\,
       \log(2/\varepsilon)\,/\,\delta \bigr\rceil$
       \Comment{total Lindbladian steps}
\State Initialise $S$ in an arbitrary state $\rho_0$
\For{$\ell = 1, \ldots, L$}
  \For{$a = 1, \ldots, M_{\mathcal{A}}$}
       \Comment{sequential sweep over jumps}
    \State \Call{CoherentStep}{$S, a, \delta, b_\pm$}
           \Comment{$\ee^{-\mathrm{i}\delta B_a}$, Alg.~\ref{alg:coh}}
    \State \Call{DissipativeStep}{$S, a, \delta, f, \gamma$}
           \Comment{$\ee^{\delta\mathcal{L}_{a,\mathrm{diss}}}$, Alg.~\ref{alg:diss}}
  \EndFor
\EndFor
\State \Return $S$
\end{algorithmic}
\end{algorithm}


\begin{figure}[ht]
  \centering
  \resizebox{\textwidth}{!}{%
  \begin{quantikz}[row sep=0.6cm, column sep=0.5cm]
  %
  % Cell-by-cell layout (15 cells per row):
  %   1: \lstick          (left label)
  %   2: wire decl        (\qw or \qwbundle{r})
  %   3: R_0 / qw         (coherent: first QSP rotation)
  %   4: octrl / U_{B_a}  (Lambda(U_{B_a}))
  %   5: octrl / R_T      (Lambda(R_T))
  %   6: R_k / qw         (coherent: inner-loop QSP rotation)
  %   7: meter / qw       (terminate coherent ancilla)
  %   8: setwiretype{n}   (transition cell A: hide outgoing wire)
  %   9: new lstick       (transition cell B: re-init wire as dissipative ancilla)
  %  10: U_{diss} / qw    (dissipative forward)
  %  11: Y_delta / octrl  (weak measurement rotation)
  %  12: U_{diss}^dag     (dissipative reverse, 0-ctrl by q_delta)
  %  13: meter / rstick   (terminate dissipative ancilla / system output)
  %
  % Meter and setwiretype{n} are placed in SEPARATE cells so that the
  % wire leading into the meter stays visible (otherwise quantikz2 applies
  % {n} retroactively to the incoming wire and breaks the meter).
  %
  % ==================================================================
  % Row 1:  q_QSP  (coherent, cells 3-7)  -->  q_delta  (dissipative, cells 9-13)
  % ==================================================================
    \lstick{$\ket{0}_{\mathrm{QSP}}$}
      & \qw
      & \gate{R(\theta_0,\phi_0,\lambda_0)}
      & \octrl{1}
      \gategroup[4,steps=3,style={dashed,rounded corners,fill=blue!6,inner sep=5pt},
                 background,label style={label position=above,yshift=0.30cm,anchor=south}]
        {\scriptsize \textsc{CoherentStep}: repeat $d$ times with angles $(\theta_k,\phi_k)\Rightarrow\ee^{-\ii\delta B_a}$}
      & \octrl{1}
      & \gate{R(\theta_k,\phi_k,0)}
      & \meter{}
      & \setwiretype{n}
      & \lstick{$\ket{0}_\delta$} \setwiretype{q} \wireoverride{n}
      & \qw
      & \qw
      \gategroup[4,steps=4,style={dashed,rounded corners,fill=red!6,inner sep=5pt},
                 background,label style={label position=below,yshift=-0.40cm,anchor=north}]
        {\scriptsize \textsc{DissipativeStep}: weak measurement simulation of $\ee^{\delta\mathcal{L}_{a,\mathrm{diss}}}$}
      & \gate{Y_\delta}
      & \octrl{1}
      & \meter{}
      \\
  %
  % ==================================================================
  % Row 2:  T_+  -->  q_gamma  (holds the multi-row gate labels in both phases:
  %         U_{B_a}/R_T on coherent, U_{diss}/U_{diss}^dag on dissipative.
  %         q_gamma sits directly under q_delta so the weak measurement
  %         control is the minimal \octrl{-1}.)
  % ==================================================================
    \lstick{$\ket{0}_{T_+}$}
      & \qwbundle{r_+}
      & \qw
      & \gate[3]{U_{B_a}}
      & \gate[2]{R_T}
      & \qw
      & \meter{}
      & \setwiretype{n}
      & \lstick{$\ket{0}_\gamma$} \setwiretype{q} \wireoverride{n}
      & \qw
      & \gate[3]{U_{\mathrm{diss}}}
      & \octrl{-1}
      & \gate[3]{U_{\mathrm{diss}}^\dagger}
      & \meter{}
      \\
  %
  % ==================================================================
  % Row 3:  T_-  -->  Omega  (bundled throughout; cells 4, 5, 10, 12 empty as
  %         they are consumed by the multi-row gates sitting on row 2 above.)
  % ==================================================================
    \lstick{$\ket{0}_{T_-}$}
      & \qwbundle{r_-}
      & \qw
      &
      &
      & \qw
      & \meter{}
      & \setwiretype{n}
      & \lstick{$\ket{0}_\Omega$} \setwiretype{q} \wireoverride{n}
      & \qwbundle{r}
      &
      & \qw
      &
      & \meter{}
      \\
  %
  % ==================================================================
  % Row 4:  S  (system register, active throughout with no wire-type change)
  % ==================================================================
    \lstick{$\rho$}
      & \qwbundle{n}
      & \qw
      &
      & \qw
      & \qw
      & \qw
      & \qw
      & \qw
      & \qw
      &
      & \qw
      &
      & \qw
      & \rstick{$\ee^{\delta\mathcal{L}_a}(\rho)+\mathcal{O}(\delta^2)$}\qw
  \end{quantikz}%
  }
  \caption{Algorithm~\ref*{alg:main}: one $\delta$-step of the CKG Lindbladian
  $\ee^{\delta\mathcal{L}_a} = \ee^{-\ii\delta B_a}\cdot
  \ee^{\delta\mathcal{L}_{a,\mathrm{diss}}}+\mathcal{O}(\delta^2)$ for a fixed
  jump $A_a$, split into a GQSP \emph{coherent block} (blue, $d$ repetitions
  of $\Lambda(U_{B_a}),\Lambda(R_T),R(\theta_k,\phi_k,0)$) and a
  weak measurement \emph{dissipative block} (red).  Mid-circuit measurements
  mark ancilla reuse: $q_{\mathrm{QSP}}, T_\pm \!\to q_\delta, q_\gamma, \Omega$.}
  \label{circ:algorithm-delta-step}
\end{figure}

It is time to assemble the subroutines discussed in the preceding subsections into a complete quantum algorithm for Gibbs state preparation via the CKG Lindbladian \eqref{eq:ckg-L}. The resulting procedure, Algorithm~\eqref{alg:main}, simulates the continuous-time Markov semigroup $\ee^{t \mathcal{L}}$ by repeated application of small $\delta$-step channels that approximate $\ee^{\delta \mathcal{L}_a}$ for each jump $a\in\mathcal{A}_a$. Its structure draws on three key ingredients: the \emph{weak measurement simulation scheme} of \cite{chen2023quantum}, which implements the dissipative channels $\ee^{\delta \mathcal{L}_{a, \text{diss}}} + \mathcal{O}(\delta^2)$; The \emph{coherent correction term} $B_a$ \eqref{eq:coherent-time-domain} introduced in \cite{chen2023efficient} (or its Trotterized version \eqref{eq:B-trotter}), that makes the CKG generator $\mathcal{L}$ exactly KMS detailed balanced; and \emph{Generalized Quantum Signal Processing} (GQSP) from \cite{motlagh2024generalized}, which provides a near-optimal method to realize the unitary $\ee^{-\ii \delta B_a}$ up to $\mathcal{O}(\delta^2)$ errors from the block encoding of $B_a$.
The end result is a quantum channel that, after $L$ iterations of the outer loop, produces a state $\rho$ that is at most $\varepsilon$ trace distance from the Gibbs state $\rho_\beta$.

We now give a few comments on Algorithm~\eqref{alg:main}. 
\begin{itemize}
    \item The system register $S$ can be initalised in any arbitrary state $\rho_0$, due to the primitivity of $\mathcal{L}$: $\rho_\beta$ is the unique fixed point and the semigroup is ergodic, thus any initial state converges to $\rho_\beta$ at the same asymptotic rate. In practice however, it is ideal to choose a state that has a reasonable overlap with $\rho_\beta$ (and is not too costly to prepare), e.g. the maximally mixed state $\mathbf{1} / d$, to eliminate part of the initial sampling steps, related to the factor $1/\sqrt{\lambda_{\min}(\rho_\beta)}$.
    \item From the mixing time definition \eqref{eq:mixing-time-defi} and the spectral bound \eqref{eq:quantum-spectral-bound}, after evolving some initial state $\rho_0$ for a total Lindbladian time $t = t_{\mathrm{mix}}(\mathcal{L})\,\log(2/\varepsilon)$ we are guaranteed to have $\|e^{t\mathcal{L}}(\rho_0) - \rho_\beta\|_1 \leq \varepsilon$.
    \item For each $\delta$ step of the outer-loop we perform a sequential sweep over all $M_\mathcal{A} = |\mathcal{A}|$ jumps. This is a Lie-Trotter splitting of $e^{\delta\mathcal{L}}$ to $\Phi_\mathcal{A} := e^{\delta\mathcal{L}_{M_\mathcal{A}}} \circ \cdots \circ e^{\delta\mathcal{L}_1}$, introducing approximation errors of $\mathcal{O}(\delta^2)$ \eqref{eq:gen-split-approx-error} but perturbing the real part of the spectrum only at $\mathcal{O}(\delta^3)$ \eqref{eq:gen-split-imaginary-delta2-errors}.
    \item In the inner-loop for a fixed $a$ we split the generator further as $\mathcal{L}_a = \mathcal{L}_{a,\mathrm{coh}} + \mathcal{L}_{a,\mathrm{diss}}$, introducing an error to the channel $\ee^{\delta \mathcal{L}_a}$ of order $\mathcal{O}(\delta^2)$ order \eqref{eq:gen-split-coh-diss}.
    \item \textsc{CoherentStep} is the subalgorithm that implements the unitary $\ee^{-\ii \delta B_a} + \mathcal{O}(\delta^2)$ via GQSP, see Alg.~\eqref{alg:coh}.
    \item \textsc{DissipativeStep} is the subalgorithm that implements the dissipation $e^{\delta\mathcal{L}_{a,\mathrm{diss}}}(\rho) + \mathcal{O}(\delta^2)$ via the weak measurement scheme, see Alg.~\eqref{alg:diss}.
\end{itemize}
\medskip 

%% ---------------------------------------------------------------
\setcounter{algorithm}{0}
\renewcommand{\thealgorithm}{\arabic{algorithm}a}
\begin{algorithm}[!htbp]
\caption{\textsc{CoherentStep}: Hamiltonian simulation of
         $e^{-\mathrm{i}\delta B_a}$ via block encoding and QSP}
\label{alg:coh}
\begin{algorithmic}[1]
\small
\Require System register $S$, jump index $a$, step size $\delta$
\Ensure  $S$ evolved by $e^{-\mathrm{i}\delta B_a}$
  up to QSP and Trotter precision of $\mathcal{O}(\delta^2)$
\Statex
\State \textbf{Ancillas:}
  outer time register $T_-$ ($r_-$ qubits),
  inner time register $T_+$ ($r_+$ qubits),
  QSP signal qubit $q_{\mathrm{QSP}}$
\Statex
\Statex\textit{--- Block encoding $U_{B_a}$ of
  $B_a/\alpha$~\eqref{eq:B-discr}, with
  $\alpha := \lVert b_-\rVert_1\,\lVert b_+\rVert_1$ ---}
\Statex\textit{Forward pass:}
\State \textsc{StatePrep}$(T_-)$:
  $|0\rangle \to |b_-\rangle :=
   \displaystyle\sum_{\bar{t}}
   \sqrt{\frac{t_0\,|b_-(\bar{t})|}{\alpha_-}}\;
   \mathrm{sgn}\bigl(b_-(\bar{t})\bigr)\,|\bar{t}\rangle$
  \Comment{outer kernel~\eqref{eq:b_minus}; $\alpha_- = \lVert b_-\rVert_1$}
\State \textsc{Ctrl-Trott}$(T_-, S)$: apply
  $S_2^{(M_-)}(-\bar{t}/\sigma)$ on $S$, controlled on $T_-$
  \Comment{outer: $\approx e^{-\mathrm{i}H\bar{t}/\sigma}$}
\State \textsc{StatePrep}$(T_+)$:
  $|0\rangle \to |b_+\rangle :=
   \displaystyle\sum_{\bar{t}'}
   \sqrt{\frac{t_0'\,|b_+(\bar{t}')|}{\alpha_+}}\;
   \mathrm{sgn}\bigl(b_+(\bar{t}')\bigr)\,|\bar{t}'\rangle$
  \Comment{inner kernel~\eqref{eq:b_plus-s-eta}; $\alpha_+ = \lVert b_+\rVert_1$}
\Statex
\Statex \textit{Inner Heisenberg evolution}
  $$\tilde{A}_a^\dagger(\beta\bar{t}')\,\tilde{A}_a(-\beta\bar{t}')
  = S_2^{(M_+)}(\beta\bar{t}')\,A_a^\dagger\,
    S_2^{(M_+)}(-2\beta\bar{t}')\,A_a\,
    S_2^{(M_+)}(\beta\bar{t}'),$$
  \textit{controlled on $T_+$}:
\State \textsc{Ctrl-Trott}$(T_+, S)$: apply
  $S_2^{(M_+)}(+\beta\bar{t}')$ on $S$
  \Comment{$\approx e^{+\mathrm{i}\beta H\bar{t}'}$}
\State Apply $A_a^\dagger$ on $S$
\State \textsc{Ctrl-Trott}$(T_+, S)$: apply
  $S_2^{(M_+)}(-2\beta\bar{t}')$ on $S$
  \Comment{$\approx e^{-2\mathrm{i}\beta H\bar{t}'}$}
\State Apply $A_a$ on $S$
\State \textsc{Ctrl-Trott}$(T_+, S)$: apply
  $S_2^{(M_+)}(+\beta\bar{t}')$ on $S$
  \Comment{$\approx e^{+\mathrm{i}\beta H\bar{t}'}$}
\Statex
\State \textsc{StatePrep}$^\dagger(T_+)$:
  $|b_+\rangle \to |0\rangle$
  \Comment{reflect / uncompute inner register}
\State \textsc{Ctrl-Trott}$(T_-, S)$: apply
  $S_2^{(M_-)}(+\bar{t}/\sigma)$ on $S$
  \Comment{undo outer evolution}
\State \textsc{StatePrep}$^\dagger(T_-)$:
  $|b_-\rangle \to |0\rangle$
  \Comment{reflect / uncompute outer register}
\Statex \Comment{Block encoding:
  $\langle 0|_{T_-}\!\langle 0|_{T_+}\;
   U_{B_a}\;
   |0\rangle_{T_-}\!|0\rangle_{T_+}
   \;=\; B_a\,/\,\alpha$
  acting on $S$}
\Statex
\Statex \textit{--- QSP: implement $e^{-\mathrm{i}\delta B_a}$
  from the block encoding ---}
\State Apply QSP~\cite{motlagh2024generalized} with precomputed rotation angles     
    $\{(\theta_k,\phi_k)\}_{k=0}^{d}$ for the Laurent polynomial $P(e^{i\theta}) \approx e^{-\mathrm{i}\delta\alpha\,\sin\theta}$
\Statex \Comment{Degree                                        
    $d = 1$ calls to $U_{B_a}$}
\State Measure and reset ancillas $q_{\mathrm{QSP}}$, $T_\pm$
\end{algorithmic}
\end{algorithm}
\renewcommand{\thealgorithm}{\arabic{algorithm}}
%% ---------------------------------------------------------------
\begin{figure}[ht]
  \centering
  \resizebox{1.05\textwidth}{!}{%
  \begin{quantikz}[row sep=0.9cm, column sep=0.42cm]
  % Row 1:  T_-  (outer time register; controls cells 4 and 12)
  % ==================================================================
    \lstick{$\ket{0}_{T_-}$}
      & \qwbundle{r_-}
      & \gate{\text{PREP}'_-}
      & \ctrl[style=halfctrl]{2}
      & \qw
      & \qw
      & \qw
      & \qw
      & \qw
      & \qw
      & \qw
      & \ctrl[style=halfctrl]{2}
      & \gate{\text{PREP}_-^\dagger}
      & \rstick{$\bra{0}_{T_-}$}\qw
      \\
  %
  % ==================================================================
  % Row 2:  T_+  (inner time register; controls cells 6, 8, 10;
  %              prepared in cell 5, uncomputed in cell 11)
  % ==================================================================
    \lstick{$\ket{0}_{T_+}$}
      & \qwbundle{r_+}
      & \qw
      & \qw
      & \gate{\text{PREP}'_+}
      \gategroup[2,steps=7,style={dashed,rounded corners,fill=blue!6,inner sep=5pt},
                 background,label style={label position=below,yshift=-0.35cm,anchor=north}]
        {\scriptsize \textsc{Inner sum}:
         $\sum_{\bar{t}'} b_+(\bar{t}')\tilde{A}_a^\dagger(\beta\bar t')\,\tilde{A}_a(-\beta\bar t')$
         controlled on $T_+$}
      & \ctrl[style=halfctrl]{1}
      & \qw
      & \ctrl[style=halfctrl]{1}
      & \qw
      & \ctrl[style=halfctrl]{1}
      & \gate{\text{PREP}_+^\dagger}
      & \qw
      & \qw
      & \rstick{$\bra{0}_{T_+}$}\qw
      \\
  %
  % ==================================================================
  % Row 3:  S  (system register; target of every controlled evolution
  %             and of the two jump applications A_a^dag, A_a)
  % ==================================================================
    \lstick{$\ket{\psi}$}
      & \qwbundle{n}
      & \qw
      & \gate{S_2\left(\frac{\bar t}{\sigma}\right)}
      & \qw
      & \gate{S_2(\beta\bar t')}
      & \gate{A_a}
      & \gate{S_2(-2\beta\bar t')}
      & \gate{A_a^\dagger}
      & \gate{S_2(\beta\bar t')}
      & \qw
      & \gate{S_2\left(-\frac{\bar t}{\sigma}\right)}
      & \qw
      & \rstick{$\tilde{B}_a\ket{\psi} / \alpha$}\qw
  \end{quantikz}%
  }
  \caption{Block encoding $U_{B_a}$ of the Trotterized coherent
  Metropolis term $\tilde{B}_a/\alpha$ with $\alpha=\lVert b_-\rVert_1\lVert b_+\rVert_1$
  (Alg.~\ref*{alg:coh}, \eqref{eq:B-block-encoding}), on outer/inner time
  registers $T_\mp$. Half-filled controls denote uniform (SELECT) control
  over the bundled time registers.
  $\text{PREP}'_\pm$ prepares the signed state
  $\sum_{\bar t}[b_\pm(\bar t)/\sqrt{|b_\pm(\bar t)|}]\ket{\bar t}/\sqrt{\lVert b_\pm\rVert_1}$
  and $\text{PREP}_\pm$ the magnitude state
  $\sum_{\bar t}\sqrt{|b_\pm(\bar t)|/\lVert b_\pm\rVert_1}\ket{\bar t}$,
  so that amplitudes multiply to the signed kernel $b_\pm(\bar t)$. Controlled
  evolutions are Trotterized with the second-order Strang splitting $S_2$ with $M_\pm$ Trotter steps respectively.}
  \label{circ:U_B_a}
\end{figure}

\medskip
\noindent \textsc{CoherentStep:} \textbf{Coherent evolution of $B$ with QSP.} The \textsc{CoherentStep}, \hyperref[alg:coh]{Algorithm~\ref*{alg:coh}}, implements the unitary evolution $\ee^{-\ii \delta B_a}$ generated by the per-jump Hermitian term $B_a$. This is the component that distinguishes the CKG algorithm from the CKBG one: without it (and without the shifted transitions weights) the algorithm reduces to the CKBG Gibbs sampler [ref to text \todo{ref}], which satisfies only approximate GNS detailed balance. Recall from the \hyperref[sec:kms-lindbladian]{KMS Lindbladian~\ref*{sec:kms-lindbladian}} subsection that the coherent term $B$ was derived \eqref{eq:B-bohr-domain} to cancel the emergent anti-Hermitian part of the quantum discriminant, upgrading the Lindbladian to exact KMS detailed balance. Its per-jump contribution $B_a$ enters the coherent-dissipative splitting \eqref{eq:gen-split-coh-diss} as the generator $\mathcal{L}_{a,\mathrm{coh}}(\cdot) = -i[B_a, \cdot\,]$, with the first-order splitting error of $\mathcal{O}(\delta^2)$ already matched by the weak measurement error of the dissipative step. Of course, discretization and Trotterization errors analysed in the preceding subsections will re-introduce controlled violations of the detailed balance, but these are precisely the errors we have learned to bound.

The subalgorithm breaks down to two steps: first we construct a block encoding $U_B$, satisfying
\begin{equation} \label{eq:B-block-encoding}
    \langle 0|_{T_-}\langle 0|_{T_+}\,U_{B_a}\,|0\rangle_{T_-}|0\rangle_{T_+} = B_a/\alpha,
\end{equation}
where $\alpha = \|b_-\|_1\,\|b_+\|_1$ is the subnormalisation factor of $B_a$. Then, we use GQSP \cite{motlagh2024generalized} (in most cases with degree $d = 1$) to implement $\ee^{-\ii \delta B_a}$ from $U_{B_a}$. In the following, we will give a few more details about each of the steps.

The block encoding of the coherent term \eqref{eq:B-block-encoding} containing the nested time integrals of \eqref{eq:coherent-time-domain} is done by a nested LCU. [Ref to prelims \todo{ref}]. This has been established in the CKG paper \cite[Proposition III.1]{chen2023efficient}, where they also prove its asymptotic scaling. The full nested double-integral structure with two time registers has been left implicit by them. Here we unpack it completely.

The starting point is the nested integral form of $B_a$ \eqref{eq:coherent-time-domain} which is discretized by nested Riemann-sums \eqref{eq:B-discr} with prefactors $t_0, t_0'$. In the LCU framework \cite{childs2012hamiltonian,berry2015simulating} this manifests as a state preparation / controlled-evolution / unpreparation on two separate ancilla registers: an outer time register $T_-$ of $r_-$ qubits, and an inner time register $T_+$ of $r_+$ qubits. 

Steps 2 and 4 of \hyperref[alg:coh]{Algorithm~\ref*{alg:coh}} prepare these registers by
\begin{equation}
    \begin{split}
        |0\rangle_{T_-} &\;\longrightarrow\; |b_-\rangle := \sum_{\bar t}\, \sqrt{\frac{t_0\,|b_-(\bar t)|}{\alpha_-}}\;\mathrm{sgn}\!\big(b_-(\bar t)\big)\,|\bar t\rangle, \qquad \alpha_- = \|b_-\|_1, \\
        |0\rangle_{T_+} &\;\longrightarrow\; |b_+\rangle := \sum_{\bar t'}\, \sqrt{\frac{t_0'\,|b_+(\bar t')|}{\alpha_+}}\;\mathrm{sgn}\!\big(b_+(\bar t')\big)\,|\bar t'\rangle, \qquad \alpha_+ = \|b_+\|_1.
    \end{split}
\end{equation}
Here, the absolute values go into the amplitudes (ensuring normalization $\langle b_\pm | b_\pm \rangle = 1$) and the signs go into the phases. For the outer kernel $b_-$ \eqref{eq:b_minus}, which is real-valued, $\mathrm{sgn}(b_-(\bar t)) \in \{-1,+1\}$ is the usual sign function. For the inner kernel $\mathrm{sgn}(b_+(\bar t'))$ \eqref{eq:b_plus-s-eta}, which is complex-valued in general, the notation $\mathrm{sgn}(b_+(\bar t'))$ should be understood as the unit-modulus phase $b_+(\bar t')/|b_+(\bar t')|$. The same convention applies to the Gaussian case \eqref{eq:b_plus-gauss}. 

The state preparation costs for $\ket{b_\pm}$ depend on the register sizes and the regularity of the kernels. For actual costs it is best to look at our numerical analysis [REF, \todo{ref}], but we can already say something qualitative here based on our quadrature error results for $r_-$ \eqref{eq:r_-}, and $r_+$ in the Gaussian case \eqref{eq:r_+_gaussian}, in the Metropolis case \eqref{eq:r_+_metro}. For our parameter regimes (moderate temperatures $\beta \leq 20$ and system sizes $n \leq 12$) the register sizes $r_- = 5-8$ and $r_+ = 10-15$ are already sufficient for our purposes to block encode $B_a$. There are several ways to prepare a quantum state, and as usual, the question boils down to space-time trade-offs: either with ancillas and fewer gates, or no ancillas and more gates. Since we are aiming to specify details for an implementation on a partially fault-tolerant quantum hardware which would likely not have an abundance of qubits, the ideal state preparation is one that uses no ancillas and its gate complexity is subdominant to the other parts of the Gibbs preparation. One such algorithm is to use uniformly controlled rotations \cite{mottonen2004transformation,shende2008cnot,plesch2011quantum} at cost $\mathcal{O}(2^r)$ gates. At $r \leq 12$, this is at most a few thousand gates, definitely below the cost of implementing controlled Hamiltonian evolutions. But for lower temperatures, or higher precisions for Gibbs sampling (smaller $\delta$) we would likely find ourselves in the regimes where the exponential scaling of the state preparation becomes dominant. In this case, we would suggest using the QSP-based method from \cite{mcardle2022quantum} that achieves $\mathcal{O}(r \cdot \log(1 / \varepsilon))$ gate complexity at the cost of 3-4 ancilla qubits by exploiting the (possibly piecewise) smoothness of $b_\pm$ via polynomial approximation. [REF \todo{ref to prelims}]

After preparing the time registers, we proceed to apply the controlled Hamiltonian evolutions. As written in \hyperref[alg:coh]{Algorithm~\ref*{alg:coh}}, all of these time evolutions are implemented by second order Trotterization [Ref to prelims \todo{ref}]: $S_2^{(M_-)}(\pm \bar{t}/\sigma)$ for the outer-, and $S_2^{(M_+)}(\pm \beta\bar{t})$ for the inner time sum (with $M_\pm$ many Trotter steps). For the desired quadrature errors, we need to retain contributions to the time sums up to $T_- = \Theta(\beta\sigma\log(1/\varepsilon))$ and $T_+ = \Theta(\sqrt{\log(1/\varepsilon)}/(\sigma\beta\sqrt{1+s}))$ (in the Metropolis case \eqref{eq:T_+_metro}). Hence, the maximum Hamiltonian simulation times are $\mathcal{O}(\beta\log(1/\varepsilon))$ and $\mathcal{O}(\sqrt{\log(1/\varepsilon)}/(\sigma\sqrt{1+s}))$ respectively. Note, that we decided for the Strang splitting opposed to the Lie-Trotter splitting because: (i) for these time scales its second order nature actually leads to less gate complexity, and (ii) it also avoids some additional errors to the detailed balance since it is palindromic $S_p(t)^\dagger = S_p(-t)$, see \hyperref[rem:palindromic]{Remark~\ref*{rem:palindromic}}.

The subnormalisation $\alpha = \|b_-\|_1\,\|b_+\|_1$ is a critical parameter: it enters both the QSP degree (see below) and the success probability of the block encoding. Its value depends on the choice of transition weight $\gamma$. More concretely, $b_-$ is universal for all cases and its norm is upper bounded, $\|b_-\|_1 \leq (\pi/\sqrt{2})\,e^{\beta^2\sigma^2/8}$ \eqref{eq:b_minus}. While this is ideal, and so is $b_+$ in the Gaussian case with $\|b_+\|_1 = (\beta/\omega_\gamma)\sqrt{\sigma_\gamma/(2\pi)}$ \eqref{eq:b_plus-gauss}, the same cannot be said about $b_+$ in the Metropolis cases \eqref{eq:b_plus-s-eta}. There, due to the singularity in the time-domain, we have an additional factor, and have
$$
\alpha = \|b_-\|_1\,\|b_+^{(s,\eta)}\|_1 = \mathcal{O}\!\left(\log\frac{\beta\|H\|\,\|\textstyle\sum_a A^{a\dagger}A^a\|}{\varepsilon}\right).
$$
How this effectively enters the total simulation time for the coherent evolution, we will see in the following QSP section.

\medskip
Given the block encoding $U_{B_a}$, with $\langle 0|U_{B_a}|0\rangle = B_a/\alpha$, we wish to implement the unitary $\ee^{-\ii\delta B_a}$. Since $U_{B_a}$ is a unitary that block-encodes a Hermitian operator $B_a/\alpha$ with $\|B_a/\alpha\| \leq 1$, its eigenvalues come in conjugate pairs $\ee^{\pm \ii\theta}$ where $\sin\theta = \lambda$ for each eigenvalue $\lambda$ of $B_a/\alpha$.  Thus, implementing $\ee^{-\ii\delta B_a}$ on the eigenspaces of $U_{B_a}$ reduces to finding a Laurent polynomial $P(\ee^{\ii\theta}) \approx \ee^{-\ii\delta\alpha\sin\theta}$ on the unit circle. We use GQSP of \cite{motlagh2024generalized} for this task. The basics were already covered in [REF to prelims \todo{ref}]. Here we only clarify the specific features to our application.

The degree of the Laurent polynomial required for $\varepsilon_\text{QSP}$-approximation of $\ee^{-\ii\delta\alpha\sin\theta}$ is given by \cite[Theorem 7.]{motlagh2024generalized} as
$$
d = \mathcal{O}\!\left(\delta\,\alpha + \frac{\log(1/\varepsilon_\mathrm{QSP})}{\log\log(1/\varepsilon_\mathrm{QSP})}\right),
$$
which in turn means this amount of applications of $U_{B_a}$. Indeed, $\delta$ in our cases is quite small, so much so, that even for lower temperatures $\beta = 20$ (for which $\alpha$ becomes larger) we can effectively use $d = 1$ and still have better precisions than what is set as a baseline by the weak measurement scheme. Of course, for larger systems moderate temperatures would manifest for larger $\beta$ than for smaller systems. Meaning, that though for the humble system sizes that we can classically simulate it seems that degree $d = 1$ is sufficient for the QSP, the same would not hold for larger systems. There, the irregularity of the Metropolis case function $b_+^{(s, \eta)}$ would eventually require higher QSP degrees and with it longer simulation times.

\setcounter{algorithm}{0}
\renewcommand{\thealgorithm}{\arabic{algorithm}b}
\begin{algorithm}[!htbp]
\caption{\textsc{DissipativeStep}: weak measurement simulation of
         $e^{\delta\mathcal{L}_{a,\mathrm{diss}}}$}
\label{alg:diss}
\begin{algorithmic}[1]
\small
\Require System register $S$, jump index $a$, step size $\delta$, filter function $f$, transition weight $\gamma$
\Ensure  $S$ evolved by
  $e^{\delta\mathcal{L}_{a,\mathrm{diss}}}(\rho) + \mathcal{O}(\delta^2)$
\Statex
\State \textbf{Ancillas:}
  frequency register $\Omega$ ($r$ qubits),
  Boltzmann qubit $q_\gamma$,
  weak measurement qubit $q_\delta$
  \Comment{all initialised to $|0\rangle$}
\Statex
\Statex \hfill \textit{--- Forward: Trotterized OFT~\eqref{eq:trotter-oft} ---}
\State \textsc{StatePrep}$(\Omega)$:
  $|0\rangle^{\otimes r} \;\longrightarrow\;
   |f\rangle := \sum_{\bar{t}\in S_{t_0}^{[N]}}
   \bar{f}(\bar{t})\,|\bar{t}\rangle$
   \Comment{Gaussian filter amplitudes~\eqref{eq:gaussian-filter-t}}
\State \textsc{Ctrl-Trott}$^{-}(\Omega, S)$:
  $|\bar{t}\rangle\,|\psi\rangle \;\longrightarrow\;
   |\bar{t}\rangle\; S_p^{(M)}(\bar{t})^\dagger \,|\psi\rangle$
   \Comment{$\approx e^{-\mathrm{i}H\bar{t}}$}
\State Apply $A_a$ on $S$
   \Comment{jump operator (e.g.\ single-site Pauli)}
\State \textsc{Ctrl-Trott}$^{+}(\Omega, S)$:
  $|\bar{t}\rangle\,|\psi'\rangle \;\longrightarrow\;
   |\bar{t}\rangle\; S_p^{(M)}(\bar{t})\,|\psi'\rangle$
   \Comment{$\approx e^{+\mathrm{i}H\bar{t}}$}
\State \textsc{QFT}$(\Omega)$
  \Comment{time $\to$ energy: register now encodes $|\bar{\omega}\rangle$}
\Statex \Comment{Joint state:
  $\displaystyle\sum_{\bar\omega}
   |\bar\omega\rangle \otimes
   \tilde{A}_a(\bar\omega)\,|\psi\rangle$}
\Statex
\Statex \textit{--- Boltzmann filter ---}
\State Controlled on $|\bar\omega\rangle$, rotate $q_\gamma$:\;
  $R_Y\!\Bigl(2\arcsin\!\sqrt{1-\gamma(\bar\omega)}\Bigr)$
  \Comment{see [REF to text]}
\Statex \Comment{$|0\rangle_{q_\gamma} \to
   \sqrt{\gamma(\bar\omega)}\,|0\rangle
   + \sqrt{1 - \gamma(\bar\omega)}\,|1\rangle$}
\Statex
\Statex \textit{--- Weak measurement ---}
\State Controlled on $q_\gamma = |0\rangle$, rotate $q_\delta$:\;
  $R_Y\!\bigl(2\arcsin\!\sqrt{\delta}\bigr)$
  \Comment{step-size acceptance}
\Statex \Comment{In the $q_\gamma\!=\!|0\rangle$ branch:
  $|0\rangle_{q_\delta} \to
   \sqrt{1-\delta}\,|0\rangle + \sqrt{\delta}\,|1\rangle$;
  \quad net jump probability $= \gamma(\bar\omega)\,\delta$}
\Statex
\Statex \textit{--- Reverse: uncompute $U$ (controlled on $q_\delta = |0\rangle$) ---}
\Statex \Comment{No-jump branch ($q_\delta\!=\!|0\rangle$): full uncomputation;
  jump branch ($q_\delta\!=\!|1\rangle$): $\Omega, q_\gamma, S$ remain entangled}
\State Undo Boltzmann rotation on $q_\gamma$
  \Comment{undo step~7}
\State \textsc{QFT}$^\dagger(\Omega)$
  \Comment{undo step~6}
\State \textsc{Ctrl-Trott}$^{-}(\Omega, S)$
  \Comment{undo step~5}
\State Apply $A_a^\dagger$ on $S$
  \Comment{undo step~4; for Pauli jumps $A_a^\dagger = A_a$}
\State \textsc{Ctrl-Trott}$^{+}(\Omega, S)$
  \Comment{undo step~3}
\State \textsc{StatePrep}$^\dagger(\Omega)$:
  $|f\rangle \to |0\rangle^{\otimes r}$
  \Comment{undo step~2}
\Statex
\Statex \textit{--- Discard ancillas ---}
\State Measure and reset $q_\delta$, $q_\gamma$, and $\Omega$
  \Comment{discard outcomes}
\Statex \Comment{Net channel on $S$:
  $\;\rho \;\mapsto\; e^{\delta\mathcal{L}_{a,\mathrm{diss}}}(\rho)
   + \mathcal{O}(\delta^2)$}
\end{algorithmic}
\end{algorithm}
\renewcommand{\thealgorithm}{\arabic{algorithm}}
\setcounter{algorithm}{1}
%% ---------------------------------------------------------------
\begin{figure}[ht]
  \centering
  \resizebox{1.05\textwidth}{!}{%
  \begin{quantikz}[row sep=0.9cm, column sep=0.42cm]
  % ==================================================================
    \lstick{$\ket{0}_\gamma$}
      & \qw
      & \qw
      & \qw
      & \qw
      & \qw
      & \qw
      & \gate{Y_{1-\gamma}}
      & \rstick{$\bra{0}_\gamma$}\qw
      \\
  %
  % ==================================================================
  % Row 2:  Omega  (time/frequency register, middle wire;
  %                  controls cells 4, 6 (downward to S) and cell 8 (upward to q_gamma),
  %                  prepared in cell 3, QFT'd in cell 7)
  % ==================================================================
    \lstick{$\ket{0}_\Omega$}
      & \qwbundle{r}
      & \gate{\text{PREP}}
      \gategroup[2,steps=5,style={dashed,rounded corners,fill=red!6,inner sep=5pt},
                 background,label style={label position=below,yshift=-0.40cm,anchor=north}]
        {\scriptsize $U_\mathrm{OFT}$}
      & \ctrl[style=halfctrl]{1}
      & \qw
      & \ctrl[style=halfctrl]{1}
      & \gate{\text{QFT}}
      & \ctrl[style=halfctrl]{-1}
      & \rstick{$\ket{\bar\omega}_\Omega$}\qw
      \\
  %
  % ==================================================================
  % Row 3:  S  (system register, bottom wire; target of the two controlled
  %             Trotters and of the bare jump A_a in between)
  % ==================================================================
    \lstick{$\ket{\psi}_S$}
      & \qwbundle{n}
      & \qw
      & \gate{S_2(-\bar t)}
      & \gate{A_a}
      & \gate{S_2(\bar t)}
      & \qw
      & \qw
      & \rstick{$\sqrt{\gamma(\bar{\omega})}\tilde A_a(\bar\omega)\ket{\psi}$}\qw
  \end{quantikz}%
  }
  \caption{Dissipative unitary $U_{\mathrm{diss}}$ (Alg.~\ref*{alg:diss}) on
  $q_\gamma\otimes\Omega\otimes S$ with PREP that prepares the time register as the Gaussian filter $\ket{\bar{f}(\bar{t})}_\Omega$, and with Boltzmann rotation
  $Y_{1-\gamma(\bar{\omega})}$.
  The red dashed block is the OFT,
  $\bra{\bar\omega}_\Omega U_{\mathrm{OFT}}\ket{0}_\Omega = \tilde A_a(\bar\omega)$.}
  \label{circ:u-diss}
\end{figure}

\medskip
\noindent\textsc{DissipativeStep:} \textbf{Weak measurement simulation.}
The \textsc{DissipativeStep}, described in \hyperref[alg:diss]{Algorithm~\ref*{alg:diss}}, implements the dissipative channel $\ee^{\delta\mathcal{L}_{a,\mathrm{diss}}}(\rho) + \mathcal{O}(\delta^2)$ for a single jump $A_a$. This part is inherited from \cite[Theorem III.1]{chen2023quantum} up to some small changes that we will highlight here. 

The forward pass (steps 2--6) that encodes the Trotterized OFT $\tilde{A}_a(\bar{\omega})$ \eqref{eq:trotter-oft} is denoted by the unitary $U_\mathrm{diss}$ that acts on the joint system $\Omega \otimes S$. Concretely, 
\begin{equation}
    U_{\mathrm{diss}} := \mathrm{QFT}(\Omega) \cdot \mathrm{Ctrl\text{-}Trott}^+(\Omega, S) \cdot A_a \cdot \mathrm{Ctrl\text{-}Trott}^-(\Omega, S) \cdot \mathrm{StatePrep}(\Omega),
\end{equation}
where $\mathrm{StatePrep}(\Omega)$ prepares the Gaussian filter state $\ket{f}$ on the time/frequency register, and the Heisenberg evolutions are implemented by Trotterized evolutions $S_2^{(M)}(\bar t)$. The net effect is then, 
\begin{equation}
    |0\rangle_\Omega |\psi\rangle_S \mapsto U_{\mathrm{diss}}\,|0\rangle_\Omega\,|\psi\rangle_S = \sum_{\bar\omega} |\bar\omega\rangle_\Omega \otimes \tilde{A}_a(\bar\omega)\,|\psi\rangle_S.
\end{equation}
The unitary $U_\mathrm{diss}$ is the dissipative analogue of the block encoding of the coherent part $U_{B_a}$, with a small difference: it is a unitary on the full $\Omega \otimes S$ space, and projecting onto a given $\ket{\bar{\omega}}$ selects the corresponding filtered jump component $\tilde{A}_a(\bar{\omega})$, i.e. $\langle \bar\omega|U_{\mathrm{diss}}|0\rangle = \tilde{A}_a(\bar\omega)$. 

\todo{this maybe at the end?}
Our construction is a simplified version of the one presented in \cite{chen2023quantum}. We split the generator jump-wise and implement dissipation of just $A_a$ in a given $\delta$ step, and use only single-site Pauli operators that are unitary ($A^2_a = \mathbf{1}$) and Hermitian ($A_a^\dagger = A_a$) jumps. Hence, we  do not need ancillas to block encode the sum of jumps, or to block encode non-unitary jumps either. A caveat here is that we would not apply the sum of jumps $\sum_a A_a$ at once, which has an effect on the mixing of the effective generator, but these errors are controlled to relevant orders \eqref{eq:gen-split-approx-error}, \eqref{eq:gen-split-imaginary-delta2-errors}. The other caveat, is that we certainly do not know what is a \emph{good} set of jump operators, ones that would mix the system quicker than other sets. However, after trying out different jump sets and seeing no real difference for our setup, we ended up sticking with the naive choice of single Pauli jumps. Then, the only ancillas we need are the ones for the time/frequency register $\Omega$, the Boltzmann qubit $q_\gamma$, and the weak measurement qubit $q_\delta$, a total of $r+2$ ancilla qubits.

Given the energy-resolved state after $U_\mathrm{diss}$, the remaining steps are the following:
\begin{enumerate}
    \item \emph{Boltzmann filter} (step 7): controlled on the energy register $\ket{\bar{\omega}}$, rotate the Boltzmann ancilla qubit $q_\gamma$ by $R_Y(2\arcsin\sqrt{1 - \gamma(\bar\omega)})$. This maps 
    $$|0\rangle_{\gamma} \to \sqrt{\gamma(\bar\omega)}\,|0\rangle_\gamma + \sqrt{1-\gamma(\bar\omega)}\,|1\rangle_\gamma,
    $$
    encoding the transition weight $\gamma$ into the amplited of the $\ket{0}_\gamma$ branch.
    \item \emph{Weak measurement} (step 8): controlled on $q_\gamma = \ket{0}_\gamma$ rotate the weak measurement ancilla qubit $q_\delta$ by $R_Y(2\arcsin\sqrt\delta)$. In the acceptance branch,
    $$
    |0\rangle_{\delta} \to \sqrt{1-\delta}\,|0\rangle_\delta + \sqrt\delta\,|1\rangle_\delta. 
    $$
    The net probability of the "jump" outcome ($q_\gamma = \ket{1}_\gamma$) is $\gamma(\bar{\omega}) \cdot \delta$, i.e. first-order in $\delta$ which is required by the weak measurement approximation $e^{\delta\mathcal{L}_{a,\mathrm{diss}}} = \mathrm{id} + \delta\mathcal{L}_{a,\mathrm{diss}} + \mathcal{O}(\delta^2)$. 
\end{enumerate}
The small jump probability is precisely what makes this a \emph{weak} measurement: the system state is only weakly disturbed per step, such that the accumulated effect of $L \cdot M_\mathcal{A}$ steps reproduces the dissipative semigroup to the desired accuracy.  The $\mathcal{O}(\delta^2)$ local error per step is the baseline that all other error sources (Trotterisation, quadrature, generator splitting) are tuned to match.

A comment has to be made regarding the Boltzmann rotation, as it could also lead to additional circuit costs. It requires a circuit that, controlled on the $r$-qubit energy register $\ket{\bar{\omega}}$ rotates $q_\gamma$ by the angle $\theta(\bar\omega) = 2\arcsin\sqrt{1 - \gamma(\bar\omega)}$. Crucially, this is a \emph{function-controlled} rotation: the classical function $\gamma$ must be compiled into the quantum circuit and its cost depends on the regularity of $\gamma$ and on the register size $r$.

For large $r$ (say $r \geq 15$) the standard approach is to use QSP/QSVT \cite{gilyen2019quantum} to implement a polynomial approximation of $\sqrt{\gamma(\bar{\omega})}$ as a function of the block-encoded normalised energy $\bar\omega / (\omega_0 N/2)$. The required polynomial degree $d_\gamma$ depends on which transition weight $\gamma$ we pick:
\begin{itemize}
    \item For the Gaussian $\gamma_\mathrm{G}$ (analytic) or the smooth Metropolis $\gamma_\mathrm{M}^{(s)}$ with $s>0$ (Gevrey-1/2, \hyperref[prop:smooth-metro-gevrey]{Propositon~\ref*{prop:smooth-metro-gevrey}}): $d_\gamma = \mathcal{O}(\mathrm{polylog}(1/\varepsilon_\gamma))$, since Gevrey-1/2 functions admit super-algebraically converging polynomial approximations.
    \item For the kinky Metropolis $\gamma_\mathrm{M}$ ($s = 0$, only Lipschitz at $\omega = -\beta\sigma^2 / 2$): $d_\gamma = \mathcal{O}(\mathrm{poly}(1/\varepsilon_\gamma))$, since Lipschitz functions require polynomially many terms. This can indeed be circumvented by splitting the domain piecewise (via a comparator circuit block on $\ket{\bar{\omega}}$) as \cite{chen2023quantum} notes, but that is still some additional circuit cost compared to the smooth Metropolis family $s>0$.
\end{itemize}

For moderate $r$ (around $r \leq 12)$, one can bypass QSP and just use uniformly controlled $R_Y$ rotations \cite{mottonen2004transformation,shende2008cnot}, the same we proposed for $\ket{b_\pm}$ in the \textsc{CoherentStep}. Although the number of CNOT gates in this case is exponentially growing $\mathcal{O}(2^r)$, at these moderate regimes, the cost is still subdominant to the other parts of the circuit. This is the approach that we use in practice. 

In either case the Boltzmann rotation is not a bottleneck: for smooth $\gamma$ and large $r$, QSP gives polylogarithmic overhead. For moderate $r$, the brute-force decomposition is cheap. 

\todo{maybe need better structuring with subsections...}
After the weak measurement, we apply the reverse pass $U_{\mathrm{diss}}^\dagger$ controlled on $q_\delta = |0\rangle$. Then after measuring and resetting all ancillas, it has been shown in \cite[Theorem III.1]{chen2023quantum} that the resulting state is 
$$
\rho \mapsto \rho + \delta\,\mathcal{L}_{a,\mathrm{diss}}(\rho) + \mathcal{O}(\delta^2),
$$
which is the first-order expansion of the dissipative semigroup channel we want to implement. 

Compared to the CoherentStep, the DissipativeStep is structurally simpler: it uses a single time/frequency register $\Omega$ (versus the nested $T_-, T_+$ of the coherent term), involves no QSP, and its Hamiltonian simulation time per step is $\mathcal{O}(\sigma\sqrt{\log(1/\varepsilon)})$ (set by the Gaussian filter truncation $T = \Theta(\sigma\sqrt{\log(1/\varepsilon)})$ and the controlled Trotter evolutions). 

\todo{add circuit figures}

\smallskip
\subsection{Parent Hamiltonian.} Beyond its role as a design principle for the algorithm, KMS detailed balance carries another analytical payoff that we have already touched upon: the similarity transformation leads to a \emph{Hermitian} quantum discriminant $\mathcal{D}(\rho_\beta, \mathcal{L})$ \eqref{eq:discriminant}. This transformation preserves the spectrum of the Lindbladian $\mathcal{L}$, which also keeps it negative semi-definite. Hence, to define a PSD parent Hamiltonian, we flip the signs, 
\begin{equation} \label{eq:parent-ham-defi}
    \mathcal{H}_\beta \;:=\; -\mathcal{D}(\rho_\beta, \mathcal{L}) \;\geq\; 0.
\end{equation}
Its ground state is the purified Gibbs state (since $\mathcal{H}_\beta$ acts on the double Hilbert space $\mathcal{H}\otimes\mathcal{H}$),
\begin{equation}
    |\sqrt{\rho_\beta}\rangle\rangle := \frac{1}{\sqrt{Z}}\sum_i e^{-\beta E_i/2}|\psi_i\rangle
   \otimes |\psi_i^*\rangle \in \mathcal{H}\otimes\mathcal{H},
\end{equation}
where $\{|\psi_i\rangle\}$ are the energy eigenstates and $|\psi_i^*\rangle$ denotes the entry-wise complex conjugate. As the spectrum is preserved, the spectral gap of the parent Hamiltonian equals that of the Lindbladian, $\lambda_\text{gap}(\mathcal{H}_\beta) = \lambda_\text{gap}(\mathcal{L})$. Consequently, the question of mixing time, reduces to bounding the spectral gap of a Hermitian PSD operator, based on the quantum spectral bound \eqref{eq:quantum-spectral-bound},
\begin{equation}\label{eq:parent-ham-gap-mixing}
    t_\text{mix}(\varepsilon) \;\leq\; \frac{1}{\lambda_\text{gap}(\mathcal{H}_\beta)} \ln \left( \frac{1}{\varepsilon\sqrt{\lambda_{\min}(\rho_\beta)}} \right).
\end{equation}
By reformulating the problem as a Hermitian spectral gap, we can tap into a well-developed research area of Hamiltonian spectral theory, e.g. perturbation theory for gapped self-adjoint operators \cite{kato1966perturbation}, detectability lemma \cite{aharonov2009detectability,anshu2016simple}, or specialized variational (min-max) principles \cite{bhatia2013matrix}.

But having a Hermitian parent Hamiltonian for detailed balanced generators is nothing new [REF]. But the parent Hamiltonians of CKG Lindbladians were also shown to have an additional structural property \cite[Proposition I.1]{chen2023efficient}: they are \emph{frustration-free}. Intuitively, this can be understood by the fact that the CKG Lindbladian decomposes as $\mathcal{L} = \sum_{a\in\mathcal{A}} \mathcal{L}_a$, where each per-jump term $\mathcal{L}_a$ individually satisfies the KMS DB condition $\mathcal{L}_a(\rho_\beta) = 0$. The discriminant transformation carries this over to the parent Hamiltonian as well,
\begin{equation}
    \mathcal{H}_\beta = \sum_{a\in\mathcal{A}} \mathcal{H}_\beta^a, \qquad \mathcal{H}_\beta^a \geq 0, \qquad \mathcal{H}_\beta^a\,|\sqrt{\rho_\beta}\rangle\!\rangle = 0 \quad\text{for each}\;\; a\in\mathcal{A}.
\end{equation}
In other words, every local term $\mathcal{H}_\beta^a$ is PSD and has (only) the purified Gibbs state in its kernel. For lattice Hamiltonians with local jumps $A_a$, the parent Hamiltonian even inherits the quasi-locality of the Lindbladian, making each $\mathcal{H}_\beta^a$ to act on a patch of radius $\tilde{\mathcal{O}}(\beta)$ of radius around the support of $A_a$. Frustration-freeness has the following three concrete analytical uses:
\begin{enumerate}
    \item \emph{Spectral gap stability.} In \cite{rouze2025efficient} Rouzé et al. proved that for 1D lattice Hamiltonians, the CKG parent Hamiltonian satisfies local topological quantum order at $\beta = 0$ where the gap is explicitly computable. Using the Michalakis-Zwolak stability theorem for frustration-free gapped Hamiltonians \cite{michalakis2013stability}, they showed that the gap persists at all sufficiently small inverse temperatures $\beta$, establishing the first unconditional rapid mixing results for a KMS DB Lindbladian with non-commuting $H$. This was recently strengthened by \cite{bergamaschi2024quantum}, who showed that 1D spin chains mix rapidly at \emph{all} finite temperatures ($\beta < \infty$), not just in the high-temperature regime. Without frustration-freeness, the stability theorem simply does not apply.
    \item \emph{Detectability lemma and gap amplification.} Since the ground state projector factorizes into local projectors (one per kernel of $\mathcal{H}_\beta^a$), the detectability lemma of \cite{anshu2016simple} provides a way to amplify the spectral gap via coarse-graining: grouping the local terms into $\mathcal{O}(1)$ layers of commuting projectors yields a product of projections whose spectral gap lower-bounds that of $\mathcal{H}_\beta$ up to a constant factor. This is also a key ingredient in the proof of \cite{rouze2025efficient}.

    In an even more recent paper \cite{fang2026quantum} the authors took this one step further, using the detectability lemma \cite{aharonov2009detectability,anshu2016simple}. They could circumvent the Lindbladian simulation and prepare the Gibbs state from the parent Hamiltonian. Their construction yields an $\mathcal{O}(|\mathcal{A}|)$-factor gate saving over simulation-based methods in general, and when the parent Hamiltonian is strictly local (e.g. for commuting $H$), a quadratic improvement in the spectral-gap dependence via QSVT. 

    \item \emph{Gap-to-MLSI upgrade and rapid mixing.} As we have discussed in the preliminaries [REF] the spectral bound \eqref{eq:parent-ham-gap-mixing} could in principle be upgraded to a tighter MLSI bound (albeit the MLSI constant has only been found to a few special cases so far). The frustration-free decomposition of $\mathcal{H}_\beta$ can help with this conversion, which is exactly what the authors used in \cite{kochanowski2025rapid}: they proved such a gap-to-MLSI equivalence for 1D Davies dynamics (GNS DB) of commuting Hamiltonians at all positive temperatures, strengthening earlier 1D rapid-mixing results. Indeed, the analogous statement in the non-commuting KMS DB setting remains open.
\end{enumerate}

Beyond proving mixing time bounds, the Hermitian and frustration-free structure of $\mathcal{H}_\beta$ also allows for some numerical techniques that could help with classical simulations as well. 

Without any special structure the Lindblad dynamics can already be simulated without having to retain the full eigenspace: in \cite{minganti2022arnoldi} the authors apply Arnoldi iteration directly to the non-Hermitian Liouvillian $\mathcal{L}$ of dimension $d^2$, maintaining a full upper-Hessenberg matrix and reorthogonalising against all $K$ Krylov basis vectors accumulated basis vectors at every step. This procedure has a cost of $\mathcal{O}(K d^2)$. Non-normality of $\mathcal{L}$ however further complicates the relationship between its spectral gap and the iteration count. 

\emph{Hermeticity}, obtained by the discriminant transformation $\mathcal{L} \mapsto \mathcal{H}_\beta$ reduces Arnoldi to the Lánczos iteration. The Hessenberg matrix becomes tridiagonal, and the iterations involve only the last two basis vectors, reducing the per-step cost to $\mathcal{O}(d^2)$. Moreover, the convergence of the Ritz values is now controlled directly by the spectral gap $\lambda_\text{gap}(\mathcal{H}_\beta)$ and the overlap of the starting vector with the kernel. 

\emph{Frustration-freeness} adds a structural layer on top due to the detectability lemma \cite{aharonov2009detectability,anshu2016simple}. Because the purified Gibbs state lies in the common kernel of the PSD local terms $\{\mathcal{H}_\beta^a\}$, the lemma turns this collection into an explicit operator (as the product of the corresponding local kernel projectors) whose repeated action drives any state towards $|\rho_\beta \rangle\rangle$ at a rate set by $\lambda_\text{gap}(\mathcal{H}_\beta)$. As we have already noted, on the quantum side this has been recently used by \cite{fang2026quantum}. On the classical side, the same lemma can be used in tensor network methods, where it is the analytical backbone of the polynomial-time ground state algorithms for 1D gapped local Hamiltonians \cite{landau2015polynomial, arad2017rigorous}.

\newpage

% \include{results}

\part{Review}
\part[QuantumFurnace.jl]{\textmd{\textsc{QuantumFurnace.jl}}}

% --- BIBLIOGRAPHY ---
\printbibliography[title={References}, heading=bibintoc]

\begin{appendices}
    \chapter{Quantum this quantum that, la la la}

\end{appendices}

\end{document}